% concatenated from arxiv_template.tex, macros/command.tex, macros/symbol.tex, sections/introduction.tex, sections/problem_settings.tex, sections/main_results.tex, sections/hyper_contractivity.tex, sections/experiments.tex, appendix.tex, appendix/notations.tex, sections/discussion.tex, sections/sequence_model.tex, appendix/sketch_of_proof_of_main_theorem.tex, appendix/variances.tex, appendix/bias.tex, appendix/learning_curve.tex, appendix/learn_2.tex, appendix/learn_2_complex.tex, appendix/auxiliary_for_learning_curve.tex, appendix/kernel_approx_sphere.tex, appendix/kernel_approx_general.tex, appendix/proof_other_results.tex, appendix/verify_assump_filter.tex, appendix/auxiliary_lemma.tex, appendix/further_experiments.tex
\documentclass{article}


\usepackage{PRIMEarxiv}

%%% Change the font size for readability when writing
\makeatletter
\renewcommand{\normalsize}{\@setfontsize\normalsize{11.5pt}{13.2pt}}
\normalsize
\makeatother
%%% 


\usepackage[utf8]{inputenc} % allow utf-8 input
\usepackage[T1]{fontenc}    % use 8-bit T1 fonts    
\usepackage{url}            % simple URL typesetting
\usepackage{booktabs}       % professional-quality tables
\RequirePackage{amsthm,amsmath,amsfonts,amssymb}     % blackboard math symbols
\usepackage{nicefrac}       % compact symbols for 1/2, etc.
\usepackage{microtype}      % microtypography
\usepackage{lipsum}
\usepackage{fancyhdr}       % header

\usepackage{physics,mathtools,mathrsfs}

%%% Change the citation for easy check
% \RequirePackage[numbers]{natbib}
\RequirePackage[authoryear]{natbib}%% uncomment this for author-year citations

% \defcitealias{zhang2023optimality_2}{Zhang et al. (2024a)}
% \defcitealias{zhang2024optimal}{Zhang et al. (2024b)}

%\RequirePackage[colorlinks,citecolor=blue,urlcolor=blue]{hyperref}%% uncomment this for coloring bibliography citations and linked URLs
%\RequirePackage{graphicx}%% uncomment this for including figures
\RequirePackage[colorlinks,citecolor=blue,urlcolor=blue]{hyperref}
\RequirePackage{graphicx}
\usepackage{subfigure}
\usepackage{caption}
\usepackage{subcaption}  % For sub-figures
\usepackage{makecell}
\usepackage{multirow}

%%%%%%%%%%%%%%%%%%%%%%%%%%%%%%%%%%%%%%%%%%%%%%
\theoremstyle{plain}
\newtheorem{axiom}{Axiom}
\newtheorem{claim}[axiom]{Claim}
\newtheorem{theorem}{Theorem}[section]
\newtheorem{lemma}[theorem]{Lemma}
\newtheorem{corollary}[theorem]{Corollary}
\newtheorem{proposition}[theorem]{Proposition}
\newtheorem{example}{Example}
\newtheorem{claimh}{Hypothesis}
%%%%%%%%%%%%%%%%%%%%%%%%%%%%%%%%%%%%%%%%%%%%%%
\theoremstyle{definition}
\newtheorem{assumption}{Assumption} % for better visualization

%%%%%%%%%%%%%%%%%%%%%%%%%%%%%%%%%%%%%%%%%%%%%%
\theoremstyle{remark}
\newtheorem{definition}[theorem]{Definition}
\newtheorem*{fact}{Fact}

\newtheorem{remark}[theorem]{Remark}
%%%%%%%%%%%%%%%%%%%%%%%%%%%%%%%%%%%%%%%%%%%%%%
%% Please put your definitions here:        %%








\newcommand{\nbullet}[1][$\bullet$]{\vspace*{2mm}\noindent \hspace{1em}{\bf#1}\quad }


\newcommand{\qbullte}[1][$\circ$]{\vspace*{2mm}\noindent \hspace{1em} {\bf #1}\quad }


\renewcommand{\qed}{\hfill$\blacksquare$}
\newcommand{\qqed}{\hfill$\square$}

\newcommand{\huaH}{\mathcal{H}}
\newcommand{\huaM}{\mathcal{M}}
\newcommand{\huaX}{\mathcal{X}}
\newcommand{\rR}{\mathbb{R}}
\newcommand{\kK}{\mathbb{K}}


% \newcommand*{\bm}{\boldsymbol}

\newcommand*{\dx}{\mathrm{d}}
\newcommand*{\vect}{\operatorname{vec}}
\newcommand*{\unif}{\operatorname{unif}}
\newcommand*{\tran}{^{\mathsf{T}}}
\newcommand*{\NTK}{\mathtt{NT}}
\newcommand*{\NN}{\mathtt{NN}}
\newcommand*{\LI}{\mathtt{LI}}
\newcommand*{\Prob}{\mathbf{P}}
\newcommand*{\Expc}{\mathbf{E}}
\newcommand*{\comp}{^{\mathsf{c}}}


\newcommand*{\inner}{\mathtt{in}}
\newcommand*{\fnt}{\mathtt{FNT}}
\newcommand*{\poprisk}{\mathcal{L}}
\newcommand*{\emprisk}{\hat{\mathcal{L}}_{n}}
\newcommand*{\excrisk}{\mathcal{E}}




\newcommand{\bbA}{\mathbb{A}}
\newcommand{\bbB}{\mathbb{B}}
\newcommand{\bbC}{\mathbb{C}}
\newcommand{\bbD}{\mathbb{D}}
\newcommand{\bbE}{\mathbb{E}}
\newcommand{\bbF}{\mathbb{F}}
\newcommand{\bbG}{\mathbb{G}}
\newcommand{\bbH}{\mathbb{H}}
\newcommand{\bbI}{\mathbb{I}}
\newcommand{\bbJ}{\mathbb{J}}
\newcommand{\bbK}{\mathbb{K}}
\newcommand{\bbL}{\mathbb{L}}
\newcommand{\bbM}{\mathbb{M}}
\newcommand{\bbN}{\mathbb{N}}
\newcommand{\bbO}{\mathbb{O}}
\newcommand{\bbP}{\mathbb{P}}
\newcommand{\bbQ}{\mathbb{Q}}
\newcommand{\bbR}{\mathbb{R}}
\newcommand{\bbS}{\mathbb{S}}
\newcommand{\bbT}{\mathbb{T}}
\newcommand{\bbU}{\mathbb{U}}
\newcommand{\bbV}{\mathbb{V}}
\newcommand{\bbW}{\mathbb{W}}
\newcommand{\bbX}{\mathbb{X}}
\newcommand{\bbY}{\mathbb{Y}}
\newcommand{\bbZ}{\mathbb{Z}}


\newcommand{\calA}{\mathcal{A}}
\newcommand{\calB}{\mathcal{B}}
\newcommand{\calC}{\mathcal{C}}
\newcommand{\calD}{\mathcal{D}}
\newcommand{\calE}{\mathcal{E}}
\newcommand{\calF}{\mathcal{F}}
\newcommand{\calG}{\mathcal{G}}
\newcommand{\calH}{\mathcal{H}}
\newcommand{\calI}{\mathcal{I}}
\newcommand{\calJ}{\mathcal{J}}
\newcommand{\calK}{\mathcal{K}}
\newcommand{\calL}{\mathcal{L}}
\newcommand{\calM}{\mathcal{M}}
\newcommand{\calN}{\mathcal{N}}
\newcommand{\calO}{\mathcal{O}}
\newcommand{\calP}{\mathcal{P}}
\newcommand{\calQ}{\mathcal{Q}}
\newcommand{\calR}{\mathcal{R}}
\newcommand{\calS}{\mathcal{S}}
\newcommand{\calT}{\mathcal{T}}
\newcommand{\calU}{\mathcal{U}}
\newcommand{\calV}{\mathcal{V}}
\newcommand{\calW}{\mathcal{W}}
\newcommand{\calX}{\mathcal{X}}
\newcommand{\calY}{\mathcal{Y}}
\newcommand{\calZ}{\mathcal{Z}}


\newcommand{\scrA}{\mathscr{A}}
\newcommand{\scrB}{\mathscr{B}}
\newcommand{\scrC}{\mathscr{C}}
\newcommand{\scrD}{\mathscr{D}}
\newcommand{\scrE}{\mathscr{E}}
\newcommand{\scrF}{\mathscr{F}}
\newcommand{\scrG}{\mathscr{G}}
\newcommand{\scrH}{\mathscr{H}}
\newcommand{\scrI}{\mathscr{I}}
\newcommand{\scrJ}{\mathscr{J}}
\newcommand{\scrK}{\mathscr{K}}
\newcommand{\scrL}{\mathscr{L}}
\newcommand{\scrM}{\mathscr{M}}
\newcommand{\scrN}{\mathscr{N}}
\newcommand{\scrO}{\mathscr{O}}
\newcommand{\scrP}{\mathscr{P}}
\newcommand{\scrQ}{\mathscr{Q}}
\newcommand{\scrR}{\mathscr{R}}
\newcommand{\scrS}{\mathscr{S}}
\newcommand{\scrT}{\mathscr{T}}
\newcommand{\scrU}{\mathscr{U}}
\newcommand{\scrV}{\mathscr{V}}
\newcommand{\scrW}{\mathscr{W}}
\newcommand{\scrX}{\mathscr{X}}
\newcommand{\scrY}{\mathscr{Y}}
\newcommand{\scrZ}{\mathscr{Z}}


\newcommand{\itA}{\mathit{A}}
\newcommand{\itB}{\mathit{B}}
\newcommand{\itC}{\mathit{C}}
\newcommand{\itD}{\mathit{D}}
\newcommand{\itE}{\mathit{E}}
\newcommand{\itF}{\mathit{F}}
\newcommand{\itG}{\mathit{G}}
\newcommand{\itH}{\mathit{H}}
\newcommand{\itI}{\mathit{I}}
\newcommand{\itJ}{\mathit{J}}
\newcommand{\itK}{\mathit{K}}
\newcommand{\itL}{\mathit{L}}
\newcommand{\itM}{\mathit{M}}
\newcommand{\itN}{\mathit{N}}
\newcommand{\itO}{\mathit{O}}
\newcommand{\itP}{\mathit{P}}
\newcommand{\itQ}{\mathit{Q}}
\newcommand{\itR}{\mathit{R}}
\newcommand{\itS}{\mathit{S}}
\newcommand{\itT}{\mathit{T}}
\newcommand{\itU}{\mathit{U}}
\newcommand{\itV}{\mathit{V}}
\newcommand{\itW}{\mathit{W}}
\newcommand{\itX}{\mathit{X}}
\newcommand{\itY}{\mathit{Y}}
\newcommand{\itZ}{\mathit{Z}}


\newcommand{\frakA}{\mathfrak{A}}
\newcommand{\frakB}{\mathfrak{B}}
\newcommand{\frakC}{\mathfrak{C}}
\newcommand{\frakD}{\mathfrak{D}}
\newcommand{\frakE}{\mathfrak{E}}
\newcommand{\frakF}{\mathfrak{F}}
\newcommand{\frakG}{\mathfrak{G}}
\newcommand{\frakH}{\mathfrak{H}}
\newcommand{\frakI}{\mathfrak{I}}
\newcommand{\frakJ}{\mathfrak{J}}
\newcommand{\frakK}{\mathfrak{K}}
\newcommand{\frakL}{\mathfrak{L}}
\newcommand{\frakM}{\mathfrak{M}}
\newcommand{\frakN}{\mathfrak{N}}
\newcommand{\frakO}{\mathfrak{O}}
\newcommand{\frakP}{\mathfrak{P}}
\newcommand{\frakQ}{\mathfrak{Q}}
\newcommand{\frakR}{\mathfrak{R}}
\newcommand{\frakS}{\mathfrak{S}}
\newcommand{\frakT}{\mathfrak{T}}
\newcommand{\frakU}{\mathfrak{U}}
\newcommand{\frakV}{\mathfrak{V}}
\newcommand{\frakW}{\mathfrak{W}}
\newcommand{\frakX}{\mathfrak{X}}
\newcommand{\frakY}{\mathfrak{Y}}
\newcommand{\frakZ}{\mathfrak{Z}}


\newcommand{\bba}{\mathbf{a}}
\newcommand{\bbb}{\mathbf{b}}
\newcommand{\bbc}{\mathbf{c}}
\newcommand{\bbd}{\mathbf{d}}
\newcommand{\bbe}{\mathbf{e}}
\newcommand{\bbf}{\mathbf{f}}
\newcommand{\bbg}{\mathbf{g}}
\newcommand{\bbh}{\mathbf{h}}
\newcommand{\bbi}{\mathbf{i}}
\newcommand{\bbj}{\mathbf{j}}
\newcommand{\bbk}{\mathbf{k}}
\newcommand{\bbl}{\mathbf{l}}
\newcommand{\bbm}{\mathbf{m}}
\newcommand{\bbn}{\mathbf{n}}
\newcommand{\bbo}{\mathbf{o}}
\newcommand{\bbp}{\mathbf{p}}
\newcommand{\bbq}{\mathbf{q}}
\newcommand{\bbr}{\mathbf{r}}
\newcommand{\bbs}{\mathbf{s}}
\newcommand{\bbt}{\mathbf{t}}
\newcommand{\bbu}{\mathbf{u}}
\newcommand{\bbv}{\mathbf{v}}
\newcommand{\bbw}{\mathbf{w}}
\newcommand{\bbx}{\mathbf{x}}
\newcommand{\bby}{\mathbf{y}}
\newcommand{\bbz}{\mathbf{z}}


\newcommand{\reg}{\varphi_\lambda}
\newcommand{\rem}{\psi_\lambda}
\newcommand{\tildereg}{\varphi_{\tilde{\lambda}}}
\newcommand{\tilderem}{\psi_{\tilde{\lambda}}}
\newcommand{\ep}{\varepsilon}
\newcommand{\R}{\mathbb{R}}
\newcommand{\E}{\mathbb{E}}
\newcommand{\K}{\mathbb{K}}
\newcommand{\mr}{\mathrm}
\newcommand{\hf}{\frac{1}{2}}
\providecommand{\ang}[1]{\left\langle{#1}\right\rangle}
\newcommand{\fstar}{f_{\rho}^*}
\newcommand{\flam}{f_{\lambda}}
\newcommand{\gtl}{\tilde{g}_X}

\newcommand*{\gf}{\mathtt{GF}}
\newcommand*{\krr}{\mathtt{KRR}}

\newcommand{\green}{\color{green}}
% \newcommand{\green}{}

\usepackage{xcolor}
\definecolor{myblue}{HTML}{1f77b4}
\definecolor{mygreen}{HTML}{2ca02c}
\definecolor{myorange}{HTML}{ff7f0e}

%%%%%%%%% For commenting in draft
\usepackage{comment}
\newcommand{\dhmargin}[2]{{\color{cyan}#1}\marginpar{\color{red}\raggedright\footnotesize [DH]:#2}}%comment on side
\newcommand{\dhcom}[2]{{\color{cyan}#1}{\color{red}[#2]} }%comment on side
\newcommand{\dhnote}[1]{{\color{red}[#1]} } %note inline
\newcommand{\dhadd}[1]{{\color{red}[Add: #1]} } %add
\newcommand{\dhrev}[2]{{\color{cyan}[Old:]\textit{#1}} {\color{magenta}[New]#2}} %direct rewrite
\newcommand{\dhdel}[1]{{\color{cyan}[Delete:]\textit{#1}} } %delte
%%%%%%%%%%%%


%Header
\pagestyle{fancy}
\thispagestyle{empty}
\rhead{ \textit{ }} 

% Update your Headers here
\fancyhead[LO]{Learning curves in large dimensions}
% \fancyhead[RE]{Firstauthor and Secondauthor} % Firstauthor et al. if more than 2 - must use \documentclass[twoside]{article}



  
\title{
    Learning Curves and Benign Overfitting of Spectral Algorithms in Large Dimensions
}


\author{
  Weihao Lu\thanks{Department of Statistics and Data Science, National University of Singapore. Email: \texttt{weihaolu@nus.edu.sg}} 
  \And
  Qian Lin\thanks{Department of Statistics and Data Science, Tsinghua University.  Email: \texttt{qianlin@tsinghua.edu.cn}} 
  \And
  Yingcun Xia\thanks{Department of Statistics and Data Science, National University of Singapore. Email: \texttt{yingcun.xia@nus.edu.sg}}
  \And
  Dongming Huang\thanks{Department of Statistics and Data Science, National University of Singapore. (Corresponding author). \\ Email: \texttt{huang.dongming@nus.edu.sg}} 
}





\begin{document}
\maketitle


\begin{abstract}
    Existing large-dimensional theory for spectral algorithms resolves either the optimally tuned point or the interpolation limit, but leaves the under-regularized regime unexplored. We study the learning curve and benign overfitting of spectral algorithms in the large-dimensional setting where the sample size and dimension are of comparable order, i.e., $n \asymp d^{\gamma}$ for some $\gamma>0$. We first consider inner-product kernels on the sphere $\mathbb{S}^{d-1}$ and establish a sharp asymptotic characterization of the excess risk across the full regularization path under various source conditions $s \geq 0$, where $s$ measures the relative smoothness of the regression function. Our results reveal that the learning curve is not simply U-shaped but instead consists of three distinct regimes: over-regularized, under-regularized, and interpolation regimes. This characterization allows us to fully capture the benign overfitting phenomenon, demonstrating that benign overfitting arises consistently across both the under-regularized and interpolation regimes whenever $s$ is positive but no larger than a critical threshold. We further show that, in the sufficiently regularized regime, the kernel learning curve is recovered by an associated sequence model. Finally, we extend the learning-curve analysis to large-dimensional KRR for a class of kernels on general domains in $\mathbb{R}^d$ whose low-degree eigenspaces satisfy spectral-scaling and hyper-contractivity conditions.
\end{abstract}


\noindent {\bf Keywords: }
Spectral algorithms, learning curves, high dimension, benign overfitting.


% \clearpage





\section{Introduction}\label{sec_intro}


Nonparametric regression studies the estimation of an unknown function $f_{\star}:\mathbb{R}^d\to\mathbb{R}$ from $n$ i.i.d. samples $\{(x_i,y_i)\}_{i=1}^n$, where $y_i = f_{\star}(x_i) + \epsilon_i$ and $\epsilon_i$ represents noise. 
The performance of an estimator $\hat{f}$ is typically measured by the \textit{excess risk} 
\[
\|\hat f-f_{\star}\|_{L^2}^2:=\mathbb{E}_x[(\hat f(x)-f_{\star}(x))^2].
\]
Among the most widely studied estimators for this problem are spectral algorithms, including kernel ridge regression (KRR) and kernel gradient flow (KGF) \citep{Caponnetto2006OptimalRF,caponnetto2007optimal}. 
These methods are attractive both for their theoretical tractability and for their strong empirical performance across a range of kernel regression problems, and they are typically analyzed under source conditions of the form $f_{\star}\in[\mathcal H]^s$, where $[\mathcal H]^s$ denotes the interpolation space associated with the reproducing kernel Hilbert space $\mathcal H$ and $s>0$ measures the smoothness of the regression function; see Section~\ref{subsec:interpolation_space}. 
% todo: check if the subsection contains all info

In the fixed-dimensional regime, the performance of spectral algorithms is well understood. Under polynomial eigenvalue decay of the kernel operator, properly regularized spectral methods such as KRR and KGF can achieve minimax-optimal excess-risk rates for $0<s\leq 2$ \citep{Caponnetto2006OptimalRF, caponnetto2007optimal, raskutti2014early, Yao2007OnES, Lin_Optimal_2020, zhang2023optimality}. 
For smoother signals, however, this picture changes: KRR exhibits the \emph{saturation effect} and fails to attain the optimal rate \citep{bauer2007_RegularizationAlgorithms, Gerfo_spectral_2008, dicker2017kernel, li2022saturation}, whereas KGF can still remain rate-optimal. At the interpolation limit, several works also show that kernel interpolation can be inconsistent \citep{pmlr-v99-rakhlin19a,beaglehole2022kernel,buchholz2022_KernelInterpolation,li2023kernel}. 
These results provide the fixed-dimensional benchmark against which the large-dimensional regime should be compared. 



Recent work has renewed interest in spectral algorithms in the large-dimensional regime, where the input dimension $d$ grows simultaneously with the sample size $n$ as $n\asymp d^{\gamma}$ for some $\gamma>0$, and the input is distributed on or near the sphere $\mathbb{S}^{d-1}$. 
This regime is also closely connected to the modern benign-overfitting literature and to the lazy-training view of wide neural networks, which links gradient-based training to kernel methods \citep{bartlett2020benign,Jacot_NTK_2018,Arora_on_2019,Hu_Regularization_2021}. 
%; see Section~XXX for a review. 
In this regime, the behavior of spectral algorithms can differ substantially from the fixed-dimensional benchmark. 
Existing results have characterized the spectra of empirical inner-product kernels, identified the polynomial-approximation barrier for KRR and KGF, and established sharp asymptotic properties such as minimax optimality, saturation, and Pinsker-type risk bounds \citep{Liang_Just_2019,ghorbani2021linearized,mei2022generalization,lu2023optimal,zhang2024optimal,lu2024saturation,lu2024pinsker}. 
%; see Section~XXX for a review. 
At the interpolation limit, several recent works further showed that kernel interpolation can still generalize and can even achieve minimax-optimal rates under suitable relationships between the smoothness parameter $s$ and the dimensional scaling exponent $\gamma$. 
Such a phenomenon is known as \emph{benign overfitting} in large-dimensional kernel interpolation, which challenges the traditional view of the bias-variance trade-off. 
% We defer a more detailed discussion of this literature to Section \ref{sec_related_work}. 


Despite these advances, existing theory for large-dimensional spectral algorithms is confined to two extremes: 
\begin{enumerate}
    \item the optimally tuned case, where the regularization parameter $\lambda$ is chosen to minimize the excess risk; 
    \item the perfectly interpolated case, where $\lambda=0$ and the spectral algorithm degenerates to kernel interpolation.
\end{enumerate}
Much less is understood about the intermediate { under-regularized} regime, where $\lambda$ is positive but too small to yield optimal performance. 
This intermediate regime is particularly important from a practical perspective, because computational limitations often yield solutions in the { under-regularized} regime rather than exact interpolation. 
For example, \cite{zhang2016understanding} trained neural networks using gradient descent, and stopped training at a large but finite time, which corresponds to a small but nonzero $\lambda$. 
The benign overfitting phenomenon they observed likely occurs in the { under-regularized} regime rather than in the exact interpolation limit ($\lambda=0$).
Therefore, a fundamental question remains unresolved: 
\begin{itemize}
    \item []
   \textit{ Is benign overfitting a phenomenon that is unique to exact interpolation (i.e., $\lambda=0$), or does it persist throughout the { under-regularized} regime? }
\end{itemize}



% However, these results do not fully address whether benign overfitting can persist into the overfitting regime. 
% Importantly, studying benign overfitting of spectral algorithms in the overfitting regime is closely aligned with a common practical setting, namely early stopping, as considered in seminal works on benign overfitting. 
% For example, \cite{zhang2016understanding} trained the neural networks with gradient descent, and stopped training at a finite time $t=1/\lambda \geq T$, which corresponds to a small but nonzero $\lambda$.
% This suggests that the benign overfitting they observed likely occurs in the overfitting regime.


A natural way to address this question is through the \emph{learning curve} of a spectral algorithm, namely the dependence of its excess risk on the regularization level $\lambda$. 
Studying the learning curve does more than identify the optimal choice of $\lambda$; it reveals the entire transition from underfitting to overfitting and then to interpolation, and therefore provides the right framework for locating where benign overfitting begins. 
While learning curves of spectral algorithms are relatively well understood in fixed-dimensional settings, the corresponding large-dimensional theory remains fragmentary and is available mainly for special cases of KRR and interpolation.
% In fixed-dimensional settings, learning curves of spectral methods are relatively well-understood and are typically governed by a classical bias--variance trade-off. 
% whereas in large-dimensional settings, existing results are limited and mostly concern special estimators or special regimes; 
% see Section~XXX for a review. 
Determining these large-dimensional learning curves is therefore a necessary step toward a complete understanding of the generalization behavior of spectral algorithms.



In this paper, we study the learning curves of large-dimensional spectral algorithms in the spherical inner-product-kernel setting, and we extend the analysis for KRR to a structured class of general-domain kernels.
Prior results treat either optimal tuning or exact interpolation, whereas we determine the full regularization path and its consequences for benign overfitting. 
This fills a major gap in the large-dimensional theory of spectral algorithms. 


Our main contributions are as follows.
\begin{itemize}
    \item For inner-product kernels on $\mathbb{S}^{d-1}$, we derive a sharp asymptotic characterization of the excess risk of large-dimensional spectral algorithms for regression functions in $[\mathcal H]^s$; see Theorem~\ref{theorem_learn_curve}. 
    As $\lambda$ decreases, the learning curve exhibits three distinct regimes, namely {over-regularized}, {under-regularized}, and interpolation.
    
    % The resulting learning curve is not simply U-shaped; instead, as $\lambda$ decreases, it exhibits three distinct regimes, namely {over-regularized}, {under-regularized}, and interpolation.

    \item We then identify the benign-overfitting region for large-dimensional spectral algorithms; see Theorem~\ref{coroll_benign_overfit}. 
    In particular, benign overfitting is not confined to interpolation; instead, it persists throughout the under-regularized regime and at interpolation if and only if
    \[
    0<s\leq \Gamma(\gamma),
    \]
    where $\Gamma(\gamma)$ is defined in \eqref{eqn_threshold_s}.

    \item 

    We further show that, in the sufficiently regularized regime $\lambda\gg n^{-1}$, the learning curve of kernel regression is recovered by an associated sequence model; see Proposition~\ref{prop_learning_curve_of_sequence}. This comparison isolates the mechanism behind the kernel learning curve in a simpler and more tractable setting.

    \item We show that the spherical KRR learning-curve result extends to a structured class of general-domain kernels satisfying the spectral-scaling and hyper-contractivity assumptions; see Theorem~\ref{theorem_learn_curve_general}.
   %  \dhrev{
   %  Finally, we extend the analysis beyond spherical inner-product kernels. 
   %  Under a spectral-scaling assumption and a hyper-contractivity assumption, we establish the learning curves for large-dimensional KRR on general domains in $\mathbb{R}^d$; see Theorem~\ref{theorem_learn_curve_general}. 
   % This resulting rate on general domains has the same qualitative form as in the spherical KRR case.
   % }{
   % We further show that the spherical KRR learning-curve result extends to a structured class of general-domain kernels satisfying the spectral-scaling and hyper-contractivity assumptions; see Theorem~\ref{theorem_learn_curve_general}. 
   % }
    
    
    % This shows that the identified large-dimensional qualitative phenomena persist beyond the spherical setting.
\end{itemize}



The remainder of the paper is organized as follows. Section~\ref{sec_settings} introduces the kernel model, interpolation spaces, and spectral algorithms used throughout the paper. Section~\ref{sec_main_results} states the main learning-curve and benign-overfitting results for inner-product kernels on the sphere. 
%Section~\ref{sec_related_work} compares our results with the existing literature and discusses the sequence-model perspective. 
Section~\ref{sec_general} extends the analysis to KRR on general domains, and Section~\ref{sec_experiments} reports numerical experiments.


\subsection{ Literature review}
Recent research has renewed its focus on spectral algorithms in large-dimensional settings, where the sample size $n$ scales with the data dimension $d$ as $n \asymp d^{\gamma}$ for some $\gamma > 0$, and the input distributions are defined on (or near) the sphere  $\mathbb{S}^{d-1}$ (e.g., the standard normal distribution).
For instance, \cite{Karoui_spectrum_2010} investigated the spectral properties of empirical inner-product kernels when the inputs are drawn from normal distributions with $\gamma = 1$. 
A series of studies \citep{Liang_Just_2019, liang2020multiple, Ghorbani_When_2021, ghorbani2021linearized, mei2021learning, mei2022generalization, misiakiewicz_learning_2021, aerni2023strong, barzilai2023generalization, Ghosh_three_2021} further analyzed the phenomena of the polynomial approximation barrier for KRR and KGF in the regime $\gamma>0$ and $s=0$. 
However, when \( s > 0 \), these results are not precise enough to provide an exact convergence rate of the excess risk.

Other works \citep{lu2023optimal, zhang2024optimal, lu2024saturation} addressed the minimax rate-optimality and saturation effect of spectral algorithms, while \cite{lu2024pinsker} established the Pinsker bound (the precise asymptotic minimax risk bound) for kernel regression problems. 
Nevertheless, these results do not apply to the under‑regularized or interpolation regimes.
In addition, some studies have focused on kernel interpolation algorithms to explore the phenomenon of ``benign overfitting'' in large-dimensional settings. For example,  \cite{Liang_Just_2019} demonstrated that kernel interpolation can generalize, and further research \citep{liang2020multiple, aerni2023strong, barzilai2023generalization, zhang2024phase} determined the convergence rate of the excess risk of kernel interpolation and showed that kernel interpolation is minimax rate-optimal under certain relationships between $s$ and $\gamma$.
It is worth noting that these results either provide only upper bounds on the excess risk, or are difficult to extend to the over‑regularized and under‑regularized regimes, as well as to more general spectral algorithms.


In the fixed-dimensional regime, several studies \citep{Bordelon_Spectrum_2020, Cui2021GeneralizationER, jin2021learning, li2023asymptotic, cheng2024comprehensive} demonstrated an exact bias-variance trade-off, yielding U-shaped learning curves for the excess risk of KRR. 
Moreover, \cite{li2024generalization} emphasized the importance of studying the learning curves of KGF and related spectral algorithms, particularly because neural networks are typically trained via gradient descent. By developing an analytic functional framework, \cite{li2024generalization} derived U-shaped excess risk curves for fixed-dimensional analytic spectral algorithms.
Unfortunately, their method depends on integral operator concentration inequalities that demand a lower bound on $\lambda$; in large dimensions, these inequalities are valid only in the under‑regularized regime, which prevents us from extending their approach to obtain complete learning curves.




% In contrast, s
Studies on the learning curves of large-dimensional spectral algorithms remain relatively scarce and mainly focus on KRR. 
For instance, 
\cite{xiao2022precise, mei2022generalization} characterized the behavior of the learning curve of KRR when $s=0$, 
\cite{barzilai2023generalization} provided upper bounds on the learning curve of kernel interpolation, 
\cite{aerni2023strong} characterized the behavior of the learning curve of KRR for certain types of regression functions,
and \cite{cheng2022dimension, misiakiewicz2024non} derived tractable approximations for prediction risk for ridge regression with random design and for KRR under specific conditions. 
However, when $s>0$, the results in \cite{cheng2022dimension} and \cite{misiakiewicz2024non} rely on assumptions that become increasingly restrictive as $d$ grows.


A comprehensive comparison between our results and the existing literature is provided in Appendix B.
Moreover, a detailed discussion specifically addressing the most closely related works \citep{mei2022generalization, misiakiewicz2024non, zhang2024phase} is provided in Remark \ref{remark_compare_with_mei_mis} and after Proposition B.3.



% \dhnote{
% To avoid the most possible confusion, briefly compare with Montanari's group and direct referees to appendix. 
% Only add the ones that are not covered in above discussion and in Remark 4.3. 
% }
%  We provide a detailed comparison between our results and the existing literature in
% Appendix B.
% %defer a more detailed discussion of above literature to 
% % \ref{sec_related_work}.
% {\color{red}
% [The existing work that is most likely to be mentioned by the reviewer to ask us to show our novelty compared to them is \cite{mei2022generalization, misiakiewicz2024non, zhang2024phase}. We have covered them in the above discussion, and in Remark 4.3 and the appendix. How about adding the following paragraph?]
% }








\iffalse
\subsection{Relation between neural networks and kernel methods}
Recent years have witnessed significant advances in neural network methodologies, especially in large-dimensional applications such as computer vision \citep{he2016deep, krizhevsky2017imagenet} and natural language processing \citep{Devlin_BERT_2019}. 
A notable distinction between neural networks and traditional nonparametric models is the phenomenon of ``benign overfitting'' \citep{bartlett2020benign, frei2022benign, mallinar2022benign, sanyal2020benign}, where neural networks, even when trained until they overfit, can achieve strong generalization performance \citep{zhang2016understanding, zhang2021understanding, belkin2018understand, belkin2019does, belkin2021fit, nakkiran2021deep}.
To explain this behavior, researchers have investigated the ``lazy regime'' observed in wide neural networks, where the network parameters remain close to their initialization values during training \citep{arora2019fine, Du_gradient_2019_b, Du_gradient_2019_a, li2018learning}. 
Subsequent work \citep{Jacot_NTK_2018, Arora_on_2019, Hu_Regularization_2021, Namjoon_Non_2022, jianfa2022generalization, li2023statistical} rigorously demonstrated that, as the width tends to infinity, the time-varying neural network kernel converges uniformly to the time-invariant neural tangent kernel (NTK). Consequently, the excess risk of kernel gradient flow with the NTK closely approximates that of neural networks trained in the lazy regime.
This connection motivates recent resurgence of interest in large-dimensional kernel regression. 

\subsection{Spectral methods in large dimensions}
Recent research has renewed its focus on spectral algorithms in large-dimensional settings, where the sample size $n$ scales with the data dimension $d$ as $n \asymp d^{\gamma}$ for some $\gamma > 0$, and the input distributions are defined on (or near) the sphere  $\mathbb{S}^{d-1}$ (e.g., the standard normal distribution).
For instance, \cite{Karoui_spectrum_2010} investigated the spectral properties of empirical inner-product kernels when the inputs are drawn from normal distributions with $\gamma = 1$. A series of studies \citep{Liang_Just_2019, liang2020multiple, Ghorbani_When_2021, ghorbani2021linearized, mei2021learning, mei2022generalization, misiakiewicz_learning_2021, aerni2023strong, barzilai2023generalization, Ghosh_three_2021} further analyzed the phenomena of the polynomial approximation barrier for KRR and KGF in the regime $\gamma>0$ and $s=0$. Other works \citep{lu2023optimal, zhang2024optimal, lu2024saturation} addressed the minimax rate-optimality and saturation effect of spectral algorithms, while \cite{lu2024pinsker} established the Pinsker bound (the precise asymptotic minimax risk bound) for kernel regression problems. 
In addition, some studies have focused on kernel interpolation algorithms to explore the phenomenon of ``benign overfitting'' in large-dimensional settings. For example,  \cite{Liang_Just_2019} demonstrated that kernel interpolation can generalize, and further research \citep{liang2020multiple, aerni2023strong, barzilai2023generalization, zhang2024phase} determined the convergence rate of the excess risk of kernel interpolation and showed that kernel interpolation is minimax rate-optimal under certain relationships between $s$ and $\gamma$.

\subsection{Learning curves for spectral methods}

In the fixed-dimensional regime, several studies \citep{Bordelon_Spectrum_2020, Cui2021GeneralizationER, jin2021learning, li2023asymptotic, cheng2024comprehensive} demonstrated an exact bias-variance trade-off, yielding U-shaped learning curves for the excess risk of KRR. 
Moreover, \cite{li2024generalization} emphasized the importance of studying the learning curves of KGF and related spectral algorithms, particularly because neural networks are typically trained via gradient descent. By developing an analytic functional framework, \cite{li2024generalization} derived U-shaped excess risk curves for fixed-dimensional analytic spectral algorithms.
In contrast, studies on the learning curves of large-dimensional spectral algorithms remain relatively scarce and mainly focus on KRR. 
For instance, 
\cite{xiao2022precise} and \cite{ mei2022generalization} characterized the behavior of the learning curve of KRR when $s=0$, 
\cite{barzilai2023generalization} provided upper bounds on the learning curve of kernel interpolation, 
\cite{aerni2023strong} characterized the behavior of the learning curve of KRR for certain types of regression functions,
and \cite{cheng2022dimension} and \cite{misiakiewicz2024non} derived tractable approximations for prediction risk for ridge regression with random design and for KRR under specific conditions.

\iffalse % dh 2026-03-26
\newpage 
\color{blue}
Nonparametric regression aims to estimate an unknown function $f_{\star}: \mathbb{R}^d \to \mathbb{R}$ from $n$ i.i.d. samples $\{(x_i, y_i)\}_{i=1}^n$, where $y_i = f_{\star}(x_i) + \epsilon_i$ and $\epsilon_i$ represents noise. 
The performance of an estimator $\hat{f}$ is typically measured by the \textit{excess risk} $\|\hat{f} - f_{\star}\|_{L^2}^2 := \mathbb{E}_x[(\hat{f}(x) - f_{\star}(x) )^2]$. 
Spectral algorithms, such as kernel ridge regression (KRR) and kernel gradient flow (KGF), have been extensively studied as powerful tools for this task \citep{Caponnetto2006OptimalRF, caponnetto2007optimal}. 
In particular, when $f_{\star}$ belongs to the reproducing kernel Hilbert space (RKHS) $\mathcal{H}$ \citep{Caponnetto2006OptimalRF, caponnetto2007optimal, Yao2007OnES, raskutti2014early, Liang_Just_2019} or its \emph{interpolation spaces} $[\mathcal{H}]^s$ with smoothness parameter $s > 0$ (see Section~\ref{subsec:interpolation_space} and \citet*{fischer2020_SobolevNorm, zhang2023optimality, zhang2023optimality_2,zhang2024optimal}), these methods have demonstrated excellent theoretical guarantees and empirical performance.


Early theoretical studies focused on fixed-dimensional settings (i.e., constant $d$), assuming that the eigenvalues of the kernel operator decay polynomially. Specifically, they assumed there exist constants $0 < \mathfrak{c} \le \mathfrak{C} < \infty$ and $\beta > 1$ such that $\mathfrak{c}\, j^{-\beta} \le \lambda_j \le \mathfrak{C}\, j^{-\beta}$ for all $j \ge 1$.
Under this assumption, when $f_{\star} \in [\mathcal{H}]^{s}$ with $0 < s \leq 2$, the minimax rate of the excess risk is lower bounded by $n^{-{\beta}/({\beta+1})}$, and appropriately regularized spectral algorithms (e.g., KRR and KGF) can achieve this rate \citep{Caponnetto2006OptimalRF, caponnetto2007optimal, raskutti2014early, Yao2007OnES, Lin_Optimal_2020, zhang2023optimality}. 
In contrast, when the regression function is smoother, i.e., when $f_{\star} \in [\mathcal{H}]^{s}$ with $s > 2$, the \emph{saturation effect} occurs. In this regime, KRR achieves an excess risk of order $n^{-{2\beta}/({2\beta+1})}$, which fails to reach the information-theoretic lower bound $n^{-{s\beta}/({s\beta+1})}$ \citep{bauer2007_RegularizationAlgorithms, Gerfo_spectral_2008, dicker2017kernel, li2022saturation}. Meanwhile, KGF remains minimax rate-optimal \citep{Yao2007OnES, Lin_Optimal_2020}. 
Moreover, several studies discussed the inconsistency of kernel interpolation \citep{pmlr-v99-rakhlin19a, beaglehole2022kernel, buchholz2022_KernelInterpolation, li2023kernel}. 

Recent years have witnessed significant advances in neural network methodologies, especially in large-dimensional applications such as computer vision \citep{he2016deep, krizhevsky2017imagenet} and natural language processing \citep{Devlin_BERT_2019}. 
A notable distinction between neural networks and traditional nonparametric models is the phenomenon of ``benign overfitting'' \citep{bartlett2020benign, frei2022benign, mallinar2022benign, sanyal2020benign}, where neural networks, even when trained until they overfit, can achieve strong generalization performance \citep{zhang2016understanding, zhang2021understanding, belkin2018understand, belkin2019does, belkin2021fit, nakkiran2021deep}.
To explain this behavior, researchers have investigated the ``lazy regime'' observed in wide neural networks, where the network parameters remain close to their initialization values during training \citep{arora2019fine, Du_gradient_2019_b, Du_gradient_2019_a, li2018learning}. 
Subsequent work \citep{Jacot_NTK_2018, Arora_on_2019, Hu_Regularization_2021, Namjoon_Non_2022, jianfa2022generalization, li2023statistical} rigorously demonstrated that, as the width tends to infinity, the time-varying neural network kernel converges uniformly to the time-invariant neural tangent kernel (NTK). Consequently, the excess risk of kernel gradient flow with the NTK closely approximates that of neural networks trained in the lazy regime.


Motivated by these developments, recent research has renewed its focus on spectral algorithms in large-dimensional settings, where the sample size $n$ scales with the data dimension $d$ as $n \asymp d^{\gamma}$ for some $\gamma > 0$, and the input distributions are defined on (or near) the sphere  $\mathbb{S}^{d-1}$ (e.g., the standard normal distribution).
For instance, \cite{Karoui_spectrum_2010} investigated the spectral properties of empirical inner-product kernels when the inputs are drawn from normal distributions with $\gamma = 1$. A series of studies \citep{Liang_Just_2019, liang2020multiple, Ghorbani_When_2021, ghorbani2021linearized, mei2021learning, mei2022generalization, misiakiewicz_learning_2021, aerni2023strong, barzilai2023generalization, Ghosh_three_2021} further analyzed the phenomena of the polynomial approximation barrier for KRR and KGF in the regime $\gamma>0$ and $s=0$. Other works \citep{lu2023optimal, zhang2024optimal, lu2024saturation} addressed the minimax rate-optimality and saturation effect of spectral algorithms, while \cite{lu2024pinsker} established the Pinsker bound (the precise asymptotic minimax risk bound) for kernel regression problems. 
In addition, some studies have focused on kernel interpolation algorithms to explore the phenomenon of ``benign overfitting'' in large-dimensional settings. For example,  \cite{Liang_Just_2019} demonstrated that kernel interpolation can generalize, and further research \citep{liang2020multiple, aerni2023strong, barzilai2023generalization, zhang2024phase} determined the convergence rate of the excess risk of kernel interpolation and showed that kernel interpolation is minimax rate-optimal under certain relationships between $s$ and $\gamma$.
We will provide a detailed discussion of recent advance in large-dimensional spectral algorithms in Section \ref{sec_related_work}.



Benign overfitting in large-dimensional kernel interpolation challenges the traditional view of the bias-variance trade-off. 
However, existing theoretical analyses of large-dimensional spectral algorithms are largely confined to two extremes: 
(i) the optimally tuned case, where the regularization parameter $\lambda$ is chosen to minimize the excess risk;
(ii) the perfectly interpolated case, where $\lambda=0$ and the spectral algorithm degenerates to kernel interpolation.
This leaves a critical gap in understanding the under-regularized regime, where $\lambda$ is positive but too small for the solution to attain optimal performance.
Importantly, studying benign overfitting of spectral algorithms in the under-regularized regime is closely aligned with a common practical setting, namely early stopping, as considered in seminal works on benign overfitting. 
For example, \cite{zhang2016understanding} trained the neural networks with gradient descent, and stopped training at a finite time $t=1/\lambda \geq T$, which corresponds to a small but nonzero $\lambda$.
This suggests that the benign overfitting phenomenon they observed likely occurs in the under-regularized regime rather than in the perfect interpolation limit ($\lambda=0$).
Therefore, a fundamental question remains unresolved: does benign overfitting extend to the under-regularized regime, or is it a phenomenon exclusive to the interpolation limit ($\lambda=0$)?







%On one hand, classical theory predicts a U-shaped learning curve in noisy settings, suggesting that interpolation (i.e., $\lambda = 0$) should perform much worse than methods with optimal regularization.On the other hand, recent findings show that kernel interpolation can generalize well, even achieving minimax optimal rates.
%These contrasting observations indicate that analyses focusing solely on optimally tuned algorithms or interpolation are not enough to fully explain benign overfitting and related phenomena in spectral algorithms.


To bridge this gap and deepen the understanding of the generalization behavior of spectral algorithms, recent research has focused on characterizing the \emph{learning curves} of spectral algorithms, which precisely describe the orders of the excess risks as the regularization parameters vary. 
In the fixed-dimensional regime, several studies \citep{Bordelon_Spectrum_2020, Cui2021GeneralizationER, jin2021learning, li2023asymptotic, cheng2024comprehensive} demonstrated an exact bias-variance trade-off, yielding U-shaped learning curves for the excess risk of KRR. 
Moreover, \cite{li2024generalization} emphasized the importance of studying the learning curves of KGF and related spectral algorithms, particularly because neural networks are typically trained via gradient descent. By developing an analytic functional framework, \cite{li2024generalization} derived U-shaped excess risk curves for fixed-dimensional analytic spectral algorithms.
In contrast, studies on the learning curves of large-dimensional spectral algorithms remain relatively scarce and mainly focus on KRR. 
For instance, 
\cite{xiao2022precise} and \cite{ mei2022generalization} characterized the behavior of the learning curve of KRR when $s=0$, 
\cite{barzilai2023generalization} provided upper bounds on the learning curve of kernel interpolation, 
\cite{aerni2023strong} characterized the behavior of the learning curve of KRR for certain types of regression functions,
and \cite{cheng2022dimension} and \cite{misiakiewicz2024non} derived tractable approximations for prediction risk for ridge regression with random design and for KRR under specific conditions.



Motivated by the recent resurgence of interest in large-dimensional kernel regression and the benign overfitting phenomenon, we explore the learning curves of spectral algorithms and the benign overfitting phenomenon of kernel ridge regularization in large dimensions.  
More precisely,
\begin{itemize}
    \item We first focus on inner-product kernels on the sphere $\mathbb{S}^{d-1}$. 
    In Theorem \ref{theorem_learn_curve}, we characterize the learning curves of large-dimensional spectral algorithms for regression functions belonging to the interpolation space $[\calH]^{s}$ defined in Section \ref{subsec:interpolation_space}. 
    This result implies that the learning curve of a large-dimensional spectral algorithm is not simply U-shape; instead, it consists of three distinct regimes as $\lambda$ decreases: the over-regularized regime, the under-regularized regime, and the interpolation regime.

    \item In Theorem~\ref{coroll_benign_overfit}, we fully capture the benign overfitting phenomenon of large-dimensional spectral algorithms. 
    Consequently, when $f_{\star}$ is less smooth, benign overfitting can occur in both the under-regularized and interpolation regimes. 
    Moreover, we show that benign overfitting behaves consistently across both the under-regularized and interpolation regimes: it occurs if and only if $0<s \leq \Gamma(\gamma)$, where $\Gamma(\gamma)$ is the threshold defined in (\ref{eqn_threshold_s}).

    \item 
    In Section \ref{sec_related_work}, we provide a detailed comparison with related works on both fixed- and high-dimensional spectral algorithms, including consistency, minimax optimality, and benign overfitting. 
   
        In Proposition \ref{prop_learning_curve_of_sequence}, we observe that the learning curves mentioned above can be accurately recovered from a much simpler sequence model.
        Undoubtedly, studying sequence model is far more tractable and provides a clear, decoupled perspective on the learning dynamics  
    {\color{red} ?implication}
    

    \item In Section \ref{sec_general}, we consider kernels defined on general domains in $\mathbb{R}^{d}$. By adopting the spectral-scaling assumption and hyper-contractivity assumption, we determine the learning curve of large-dimensional KRR in Theorem \ref{theorem_learn_curve_general}. 
    Our results indicate that the new phenomena observed in large-dimensional kernel regression still exist even without the specific properties associated with inner-product kernels defined on $\mathbb{S}^{d-1}$, and further suggest that existing large-dimensional KRR results for spherical inner-product kernels may extend to more general kernels.

    
\end{itemize}
\fi 





\iffalse
    \dhnote{too detailed}
    In Section \ref{sec_related_work}, we provide a detailed comparison of related works. Specifically, we first recall results of learning curves of fixed-dimensional spectral algorithms \citep{li2024generalization, li2023asymptotic}. 
    Then, we recollect existing results in large-dimensional spectral algorithms, including 
    consistency of KRR / KGF \citep{ghorbani2021linearized, Donhauser_how_2021, mei2022generalization, xiao2022precise, misiakiewicz_spectrum_2022, hu2022sharp}), 
    minimax optimality \citep{lu2023optimal, zhang2024optimal, lu2024saturation, lu2024pinsker}, 
    and benign overfitting of kernel interpolation \citep{Liang_Just_2019, liang2020multiple, barzilai2023generalization, aerni2023strong, zhang2024phase}.
    Finally, motivated by recent results \citep{cheng2022dimension, misiakiewicz2024non} on deriving tractable approximations for excess risk of KRR based on sequence model, we determined the learning curves of equivalent sequence models in Proposition \ref{prop_learning_curve_of_sequence}.
\fi 
    





















\fi


\subsection{Notations}\label{subsec_notations}
For any integer $q \geq 2$, we denote the norm in $L^q:=L^q(\mathcal{X}, \rho_{\mathcal{X}})$ by $\|\cdot\|_{L^q}$. 
% For any integer $q=2,3, \cdots$, let's denote the norm in $L^q:=L^q(\mathcal{X}, \rho_{\mathcal{X}})$ as $\|\cdot\|_{L^q}$. 
For \( x \in [0, \infty] \), denote by \( \lfloor x \rfloor \) the greatest integer less than or equal to \( x \) if \( x < \infty \), with the convention  \( \lfloor x \rfloor = \infty \) if \( x = \infty \). Similarly, let \( \lceil x \rceil \) denote the smallest integer greater than or equal to \( x \) if \( x < \infty \), and \( \lceil x \rceil = \infty \) if \( x = \infty \).
%We use $z \sim \rho$ to denote that $z$ follows the distribution $\rho$.

We use asymptotic notations $O_d(\cdot),~o_d(\cdot),~\Omega_d(\cdot)$ and $\Theta_d(\cdot)$ to describe the behavior of sequences as $d \to \infty$.
For instance, we say two sequences of (deterministic) quantities $\{U(d), V(d)\}_{d=1}^{\infty}$ satisfy $U(d) =  o_{d}(V(d))$ if and only if for any $\varepsilon > 0$, there exists a constant $D_{\varepsilon}$ that depends only on $\varepsilon$ and the absolute positive constants listed in Definition A.1, such that for any $d > D_{\varepsilon}$, we have $U(d)< \varepsilon V(d)$. 
We also use the notation $V(d) \gg U(d)$ and $U(d)\ll V(d)$ if $U(d) = o_{d}(V(d))$.


For random variables, the probabilistic asymptotic notations $O_{d, \mathbb{P}}(\cdot), o_{d, \mathbb{P}}(\cdot), \Omega_{d, \mathbb{P}}(\cdot), \Theta_{d, \mathbb{P}}(\cdot)$ are also used. For instance, we say two sequences of random variables $\{X(d), Y(d)\}_{d=1}^{\infty}$ satisfy $ X(d) =  O_{d, \mathbb{P}}(Y(d)) $ if and only if for any $\delta > 0$, there exist two constants $C_{\delta} $ and $ D_{\delta}$ that depend only on $\delta$ and the absolute positive constants listed in Definition A.1,  such that for any $d > D_{\delta}$, we have $\mathbb{P}\left( |X(d)| > C_{\delta} |Y(d)| \right) < \delta$.

For symmetric matrix \( A\) of dimension \( k \), and any positive integer \( j \leq k \), let \( \lambda_j(A) \) denote the \( j \)-th largest eigenvalue of \( A \). 
In particular, let \( \lambda_{\text{max}}(A) \) and \( \lambda_{\text{min}}(A) \) denote the largest and smallest eigenvalues of \( A \), respectively. 
For another symmetric matrix \( \tilde{A} \) of the same dimension, we write $\tilde{A}\geq A$, if $\tilde{A}-A$ is positive semi-definite. 
We write $\tilde{A}=\Omega_{d, \mathbb{P}}(1) A$ if and only if for any $\delta > 0$, there exist two constants $C_{\delta}>0$ and $ D_{\delta}$ that depend only on $\delta$ and the absolute positive constants listed in Definition A.1,  such that for any $d > D_{\delta}$, we have $\tilde{A} \geq C_{\delta}A$ with probability at least $1-\delta$.
For any non-symmetric matrix $U$, define $\|U\|_{2} := \sqrt{\|UU^{\top}\|_{2}}$.









\section{Problem setting on the sphere}\label{sec_settings}




Suppose that we observe $n$ i.i.d. samples $(x_i, y_i)$ for $i =1,2,\ldots, n$ from the model:
\begin{equation}\label{equation:true_model}
    y=f_{\star}(x)+\epsilon,
\end{equation}
where $x_i$'s are sampled from $\rho_{\calX}$, which is the uniform distribution on $\mathcal{X} =\mathbb{S}^{d-1}\subset \bbR^{d}$,
$y \in \mathcal{Y} \subset \mathbb{R}$,
$f_{\star}$ is the regression function defined on $\mathcal{X}$, and  $\epsilon_1, \ldots, \epsilon_n$ are conditionally independent given $(x_1, \ldots, x_n)$, with conditional mean 0 and variance $\sigma^2$. 
Sections~\ref{sec_settings} and \ref{sec_main_results} focus on this spherical design, whereas Section~\ref{sec_general} later considers KRR on general domains. 
In high dimensions, standardized covariates are often concentrated near a thin spherical shell \citep{talagrand1996new}, so the spherical model is a natural reference setting. 
The spherical setting is also analytically convenient, because inner-product kernels admit an explicit spectral block structure that can be analyzed sharply.



We denote the response vector by $Y=\left(y_1, \ldots, y_n\right)^{\top} \in \mathbb{R}^n$ and the design matrix by $X=\left[x_1, \ldots, x_n\right]^{\top} \in \mathbb{R}^{n \times d}$.
In this work, we consider the large-dimensional setting where the dimension $d$ grows with the sample size  $n$. 
Specifically, the sample size and the dimension satisfy the following assumption:

\begin{assumption}\label{assump_asymptotic}
    There exist positive constants $c_1$, $c_2$, and $\gamma$ such that the sample size satisfies
\begin{align}\label{Asym}
    %n =\alpha d^{\gamma}(1+o_d(1)),
    c_{1} d^{\gamma} \leq n \leq c_{2} d^{\gamma},
\end{align}
and we denote $\ell_{\gamma} := \lfloor \gamma\rfloor$.
\end{assumption}







\subsection{Inner product kernels}\label{subsec_sphere_kernel}


Throughout this paper, we consider the following inner product kernels.
A continuous inner product kernel $\mathscr{K}$ defined on $\mathbb{S}^{d-1}$ is given by 
\begin{equation*}
%\label{eqn:def_inner_product_kernel}
    \mathscr{K}(x, x^\prime) = \Phi(\left\langle x, x^\prime \right\rangle), \forall~ x, x^\prime \in \mathbb{S}^{d-1},
\end{equation*}
where  $\Phi:[-1,1] \to \mathbb{R}$ is a continuous function satisfying the following assumption. 
% %independent of $d$.
% \dhrev{
% To avoid unnecessary notation, let us make the following assumption on the function $\Phi(\cdot)$.
% }{How is this avoiding notation? Can we just say: We make the following assumption on the function $\Phi(\cdot)$}

\begin{assumption}\label{assu:coef_of_inner_prod_kernel} 
$\Phi(t) \in \mathcal{C}^{\infty} \left([-1,1]\right)$ is a fixed function (independent of $d$) and there exists a non-negative sequence of absolute constants $\{a_i\}_{i = -1}^{\infty}$ such that
    \begin{displaymath}
        \Phi(t) = \sum\nolimits_{i=0}^\infty a_i t^i
        \quad \text{ and } \quad 
        \sum\nolimits_{i=0}^\infty a_i \leq a_{-1}, 
    \end{displaymath}
    where the constant $a_{-1}$ denotes an upper bound on the sum of coefficients. 
    Furthermore, we require  
    $a_{i} > 0$ for any $0 \leq i \leq \ell_{\gamma} + 2$, where $\ell_{\gamma}$ is defined in Assumption \ref{assump_asymptotic}.
\end{assumption}


    %$p$ is given in Definition \ref{def:abs_constants} and 



The primary purpose of Assumption \ref{assu:coef_of_inner_prod_kernel} is to ensure the positive definiteness of the kernel and to derive the asymptotic order for the 
%first $\ell_{\gamma} + 3$ 
eigenvalues $\mu_0, \cdots, \mu_{\ell_{\gamma} + 2}$ defined later in (\ref{spherical_decomposition_of_inner_main}). 
First, Theorem 1.b in \cite{gneiting2013strictly} states that an inner product kernel on the sphere is positive definite in all dimensions if and only if all coefficients $a_{i}$ are non-negative. 
Second, the strict positivity of the coefficients $a_k$ for $0 \leq k \leq \ell_{\gamma} + 2$ implies that the eigenvalues  satisfy $\mu_k = \Theta_d( d^{-k} )$ for all $k$ in this range.
%
% having a power series with the first $\ell_{\gamma} + 3$ coefficients $a_k > 0$ allows us to deduce that the eigenvalues satisfy $\mu_k = \Theta_d( d^{-k} )$, where $k \leq \ell_{\gamma} + 2$
%
Furthermore, our main results extend to settings where certain coefficients in $\{a_{j}\}_{j \geq 0}$ are zero. 
For example, one can consider the two-layer NTK defined in \cite{Bietti_on_2019}, where $a_j=0$ for any odd $j\geq 3$.
For more discussion of Assumption \ref{assu:coef_of_inner_prod_kernel}, readers can refer to the discussions below Assumption 4 in \cite{lu2023optimal}, below 
Assumption 1 in \cite{lu2024saturation}, and below Assumption 1 in \cite{zhang2024phase}.


The analyticity of $\Phi$ is a standard assumption in the related literature (e.g., \cite{liang2020multiple, xiao2022precise, hu2022sharp, lu2023optimal, zhang2024optimal, lu2024saturation, zhang2024phase, lu2024pinsker}). Indeed, many commonly used inner product kernels defined on $\mathbb{S}^{d-1}$ satisfy this assumption:
\begin{itemize}
    \item ReLU NTK: As studied in \cite{Bietti_deep_2021}, the associated $\Phi$ is analytic (see page 6 therein). Moreover, their Corollary 3 guarantees that all coefficients are positive.

    \item Gaussian kernel: When defined on the sphere, this kernel is analytic with positive coefficients.
    
    \item Random feature kernel: Consider kernels of the following form \citep{ghorbani2021linearized, mei2022generalization}:
    $$
    \mathscr{K}(x, x^{\prime}) := \mathbb{E}_{w}[\sigma(\langle w, x\rangle)\sigma(\langle w, x^{\prime}\rangle)],
    $$
    where $\sigma(\cdot)$ is an activation function. As shown in Equation (2) of \cite{Bietti_deep_2021}, $\Phi$ admits a power series expansion with $a_j \geq 0$ for all $j$  (see also \cite{daniely2016toward}).
    
\end{itemize}












Define the integral operator as
\begin{equation}\label{def_T}
    T(f)(x)=\int \mathscr{K}(x, x^{\prime}) f(x^{\prime}) ~\mathsf{d} \rho_{\calX}(x^{\prime}). 
\end{equation}
By Theorem 1.b in \cite{gneiting2013strictly}, Assumption \ref{assu:coef_of_inner_prod_kernel} implies that the kernel $\mathscr{K}$ is positive definite. Moreover, both $\mathscr{K}$ and the trace of $T$ are bounded:
\begin{equation*}
%\label{eqn:maximum_of_kernel}
\begin{aligned}
    K_{\max}&:=\sup_{x \in \mathcal X} \mathscr{K}(x, x) 
    = \sup_{x \in \mathcal X} \Phi(\left\langle x, x \right\rangle) 
    \leq \sum\nolimits_{i=0}^{\infty} a_i \leq a_{-1}\\
    \text{Tr}(T)&:=\int \mathscr{K}(x, x)  ~\mathsf{d} \rho_{\calX}(x) \leq a_{-1}. 
\end{aligned}
\end{equation*}
Consequently, standard results (e.g., \cite{steinwart2012mercer}) ensure that $T$ is a positive, self-adjoint, trace-class, and compact operator.  
%
The celebrated Mercer's theorem then assures that the decomposition
\begin{align}\label{eqn:mercer_decomp}
    \mathscr{K}(x, x^{\prime})=\sum_{j=1}^{\infty}\lambda_{j}\phi_{j}(x)\phi_{j}( x^{\prime}),
\end{align}
where the eigenvalues satisfy $\lambda_{1} \geq \lambda_{2} \geq \ldots \geq 0$, and the corresponding eigenfunctions are $\{\phi_{j}(\cdot)\}_{j\geq 1}$. 
%form an orthonormal basis in the $L^2(\calX, \rho_{\calX})$ function space.
%Calculating these eigenvalues is generally not straightforward. For example, \cite{jianfa2022generalization, li2023statistical} determined the eigenvalue decay rates $\lambda_j \asymp j^{-(d+1)/d}$ of neural-tangent kernels (NTK) on $\mathbb{R}^{d}$ with fixed dimension $d$ under certain conditions. Fortunately, 
Furthermore, since $\mathscr{K}$ is an inner product kernel defined on the sphere, the Funk-Hecke formula provides a more explicit spectral decomposition:
\begin{equation}\label{spherical_decomposition_of_inner_main}
\begin{aligned}
\mathscr{K}(x,x^\prime) = \sum_{k=0}^{\infty} \mu_{k} \sum_{j=1}^{N(d, k)} \psi_{k, j}(x) \psi_{k, j}\left(x^\prime\right),
\end{aligned}
\end{equation}
where for each $k\geq 0$, $\mu_{k}$ is the distinct eigenvalue with multiplicity $N(d, k)$ associated with degree $k$, and 
$\{\psi_{k, j}(\cdot): j=1,\ldots, N(d, k)\}$ are spherical harmonic polynomials of degree $k$. 
The dimension of the degree-$k$ subspace is given by $N(d,0)=1$ and $N(d, k) = \frac{2k+d-2}{k} \cdot \frac{(k+d-3)!}{(d-2)!(k-1)!}$ for $k\geq 1$. 
Notably, in contrast to the ordered sequence $\{\lambda_j\}$, the sequence of distinct eigenvalues $\{\mu_{k}\}_{k \geq 0}$ is not necessarily non-increasing. 
For more details on inner product kernels, see, e.g.,  \citep{gallier2009notes, ghorbani2021linearized}.

The following proposition, adapted from \cite{lu2024pinsker}, characterizes the asymptotic behavior of the leading eigenvalues $\lambda_j$ (or equivalently, the distinct eigenvalues $\mu_k$).


\begin{proposition}\label{prop:inner_edr}
    Suppose Assumptions 
\ref{assump_asymptotic} and 
\ref{assu:coef_of_inner_prod_kernel}  hold for an absolute constant $\gamma>0$.
    Then,
\begin{itemize}
    \item For any integer $k  = 0, 1, \ldots, \ell_{\gamma} +2$, we have
    \iffalse
\begin{equation*}
\begin{aligned}
 \mu_{k} = a_{k} k! d^{-k}(1+o_d(1)) \quad \text{ and } \quad
N(d, k) = \frac{d^{k}}{k!}(1+o_d(1)).
\end{aligned}
\end{equation*}
\fi

\begin{equation*}
\begin{aligned}
 \mu_{k} = \Theta_d( d^{-k}) \quad \text{ and } \quad
N(d, k) = \Theta_d( d^k ).
\end{aligned}
\end{equation*}
For ease of notation, denote
$N_{-1}=0$ and $N_{k} = \sum_{k^{\prime}=0}^{k}N(d, k^{\prime}) = \sum_{k^{\prime}=0}^{k}\Theta_d(d^{k^{\prime}}) = \Theta_d(d^k)$. 


\item There exists a constant \( \mathfrak{C} \) depending only on \( \gamma, a_0, \ldots, a_{\ell_{\gamma} + 2} \), such that for any \( d \geq \mathfrak{C} \), we have for any $0 \leq k \leq \ell_{\gamma}+1$: 
    \begin{align*}
    &  \lambda_{N_{k-1}+1} = \lambda_{N_{k-1}+2}= \cdots=\lambda_{N_{k}} = \mu_k,\\
     &  \{ 
         \phi_{N_{k-1}+1}, \phi_{N_{k-1}+2}, \ldots, \phi_{N_{k}}\} =\{\psi_{k,1}, \ldots, \psi_{k,N(d,k)}\}.
    \end{align*}
\end{itemize}
\end{proposition}

Proposition \ref{prop:inner_edr} reveals the spectral block structure of large-dimensional inner product kernels. 
In comparison, most fixed-dimensional kernels do not have such properties. We list some examples:
\begin{itemize}
    \item  For Sobolev kernels defined on a bounded domain with smooth boundary with smoothness $r>d/2$, we have $C_1(d) j^{-2r/d} \leq \lambda_j \leq C_2(d) j^{-2r/d}$ \citep{edmunds1996function, zhang2023optimality}.


    \item For ReLU NTK defined on $\mathcal{S}^{d-1}$, we have $C_1(d) j^{-d/(d-1)} \leq \lambda_j \leq C_2(d) j^{-d/(d-1)}$ \citep{Bietti_on_2019, Bietti_deep_2021}.


    \item  For Gaussian kernels defined on $\mathcal{S}^{d-1}$, we have $C_1(d) e^{-C_3(d) j} \leq \lambda_j \leq C_2(d) e^{-C_4(d) j}$ (\cite{Minh_Mercer_2006}).
    %\cite{rasmussen2003gaussian}
% \dhnote{
% Gaussian kernels on $\mathbb{R}^d$ is NOT an inner product kernel. 
% Is there a result for Gaussian kernel on $S^{d-1}$? This is the one we list earlier. }
\end{itemize}


In the large-dimensional setting considered throughout this paper, we assume $d \geq \mathfrak{C}$ such that the results of Proposition \ref{prop:inner_edr} hold. 





\subsection{The interpolation space}\label{subsec:interpolation_space}

For any $s \geq 0$, the interpolation space $[\mathcal{H}]^s$ associated with the kernel $K$ is defined as
\begin{equation*}
  [\mathcal{H}]^s := 
  \Big\{
  \sum\nolimits_{j=1}^{\infty} f_j \phi_{j}(\cdot): \left(\lambda_j^{-s / 2} f_j\right)_{j} \in \ell^2 
  \Big\} 
  \subseteq L^{2}(\mathcal{X}, \rho_{\mathcal{X}}),
\end{equation*}
where the eigenpairs $(\lambda_j, \phi_j)$ are defined in \eqref{eqn:mercer_decomp}. 
Equipped with the norm
\begin{equation*}\begin{aligned}
\Big\|\sum\nolimits_{j=1}^{\infty} f_j  \phi_j
\Big\|_{[\mathcal{H}]^s} :=
\Big(\sum\nolimits_{j=1}^{\infty} \lambda_j^{-s} f_j^2
\Big)^{1 / 2}, 
\end{aligned}
\end{equation*}
$[\mathcal{H}]^s $ is a separable Hilbert space with orthonormal basis $ \{\lambda_{j}^{s/2} \phi_{j}\}_{j=1}^\infty$. 
Generally speaking, functions in $[\mathcal{H}]^s$ become smoother as $s$ increases (see, e.g., the example of Sobolev spaces in \cite{edmunds1996function, zhang2023optimality}). 
The two most interesting interpolation spaces are $[\calH]^{0} \subseteq L^{2}$ and $[\calH]^{1} = \calH$. 


%Moreover, since $K$ is an inner product kernel defined on $\mathbb{S}^{d-1}$, we have $[\calH]^{0} = L^{2}$ if and only if $a_j >0$ for any $j=0, 1, \ldots$, where the absolute constants $\{a_j\}_{j \geq 0}$ are defined in Assumption \ref{assu:coef_of_inner_prod_kernel} (see, e.g., \cite{steinwart2012mercer, fischer2020_SobolevNorm}).


Standard kernel regression analysis typically assumes the well-specified case where the true regression function $f_{\star}$ falls into the RKHS $\calH$ (e.g., \cite{Caponnetto2006OptimalRF, caponnetto2007optimal, Yao2007OnES, raskutti2014early, Liang_Just_2019, lu2023optimal}). 
However, subsequent research has suggested that the RKHS assumption might be too restrictive, prompting interest in the performance of kernel regression in the misspecified case with $s \in (0, 1)$ (see \citet{fischer2020_SobolevNorm}, \citet{zhang2023optimality,zhang2023optimality_2}, \citet{zhang2024optimal}, and \citet{lu2024pinsker,lu2024saturation}). 
Recently, several studies on large-dimensional kernel regression have considered the extreme case where $s=0$ (e.g., \cite{ghorbani2021linearized, mei2021learning, mei2022generalization, misiakiewicz_learning_2021}).
% To fully capture the performance of large-dimensional kernel regression and provide a unified explanation for previous work, 
To provide a unified analysis of large-dimensional kernel regression, 
we assume that $f^*$ falls in a ball of radius $\sqrt{R}$ in $[\calH]^{s}$:
%, where $s$ and $R$ are positive absolute constants given in Definition \ref{def:abs_constants}.
\begin{assumption}\label{assump_function_calss}
There exist absolute constants $s \geq 0$ and $R>0$, such that we have
    \begin{equation}\label{eqn:function_calss}
    f_{\star} \in 
    %\sqrt{R}[\calB]^{s} := 
    \left\{f \in [\calH]^{s} \mid \|f\|_{[\calH]^{s}} \leq \sqrt{R}\right\}.
\end{equation}
%where $0< R = R(d) < \infty$ can be bounded or go to infinity as $d$ goes to infinity.
\end{assumption}







\iffalse
From Proposition \ref{prop:inner_edr}, for any $d \geq \mathfrak{C}$, we can denote 
\begin{equation*}
%\label{eqn_decompose_reg_func}
    f_{\star}(x) = \sum_{k=0}^{\ell_{\gamma}+1} \sum_{j=1}^{N(d,k)} \theta_{k,j} \psi_{k,j}(x) + \sum_{j =B_{\ell_{\gamma}+1}+1}^{\infty} f_j \phi_{j}(x),
\end{equation*}
where $\sum_{k=0}^{\ell_{\gamma}+1} \sum_{j=1}^{N(d,k)} \mu_k^{-s} \theta_{k,j}^2+ \sum\nolimits_{j=B_{\ell_{\gamma}+1}+1}^{\infty} \lambda_j^{-s} f_j^2 \leq R$.
\dhmargin{Therefore}{Why?}, throughout this paper, we always assume that $d \geq \mathfrak{C}$.
\fi



% Notice that if $d$ is upper bounded, then the asymptotic notations defined in Section \ref{subsec_notations} is meaningless. Therefore, in this paper, we always assume that $d$ is greater than a certain constant, and without loss of generality, we assume that $d \geq \mathfrak{C}$ such that results in Proposition \ref{prop:inner_edr} hold. 
% Then, from Assumption 

From Assumption~\ref{assump_function_calss} and Proposition \ref{prop:inner_edr}, we can denote 
\begin{equation}\label{eqn_decompose_reg_func}
    f_{\star}(x) = \sum_{k=0}^{\ell_{\gamma}+1} \sum_{j=1}^{N(d,k)} \theta_{k,j} \psi_{k,j}(x) + \sum_{j =N_{\ell_{\gamma}+1}+1}^{\infty} f_j \phi_{j}(x),
\end{equation}
where $\sum_{k=0}^{\ell_{\gamma}+1} \sum_{j=1}^{N(d,k)} \mu_k^{-s} \theta_{k,j}^2+ \sum\nolimits_{j=N_{\ell_{\gamma}+1}+1}^{\infty} \lambda_j^{-s} f_j^2 = \|f\|_{[\calH]^{s}}^2 \leq R$.
















\subsection{Spectral algorithms}

Below we recall the definition of analytic filter functions (see \cite{bauer2007_RegularizationAlgorithms, li2024generalization, lu2024saturation} for further details).
KRR, KGF, iterated ridge, and kernel gradient descent are all covered by this filter framework. 

\begin{definition}[Analytic filter functions]
  \label{def:filter}
  Let $\left\{ \reg : (0, K_{\max}] \to \R_{\geq 0} \mid \lambda \in [0, 1] \right\}$ be a family of analytic functions indexed by the regularization parameter $\lambda$.
  We call $\left\{ \reg \mid \lambda \in [0, 1]  \right\}$ an analytic filter family with qualification $\tau \in [1,\infty]$, if there exist positive absolute constants $C_1, C_2, C_3$, and $C_4$, such that:
  \begin{enumerate}
      \item 
      For all $\lambda\in[0,1]$ and $z\in(0,K_{\max}]$, we have $z \reg(z) \in [0, 1]$ and the function $\rem(z):= 1- z\reg(z)$ is non-increasing in $z$.
      
      \item 
        For any $z\in(0,K_{\max}]$ and any $1\leq t\leq \tau$,  there exists a constant $\mathfrak{C}_3$ only depending on $C_3$ and $t$, such that 
\begin{equation}\label{eq:Filter_case1}
    \begin{aligned}
        &~ C_1 (z+\lambda)^{-1} \leq \reg(z) \leq C_2 (z+\lambda)^{-1}\\
        &~ \rem(z) \leq \mathfrak{C}_3 \lambda^{t}z^{-t}.
    \end{aligned}
\end{equation}

        \item 
   
            If $\tau<\infty$, then there exists a positive constant $\mathfrak{C}_{4}$ only depending on $\tau$ and $\lambda_1$, such that we have
    \begin{align}
    \label{eq:Filter_Rem_finite_case2}
        \rem(\lambda_1) \geq \mathfrak{C}_{4} \lambda^{\tau},
    \end{align}
    where $\lambda_1$ is the largest eigenvalue of $\mathscr{K}$ defined in (\ref{eqn:mercer_decomp}).
        
    % {\color{red}        
    %     If $\tau<\infty$, then for any constant $\mathfrak{C}_4 >0$ only depending on the absolute positive constants listed in Definition \ref{def:abs_constants}, there exists a constant $\mathfrak{C}_4^{\prime}>0$ only depending on $C_4$,  $\mathfrak{C}_4$, and $\tau$, such that we have
    % \begin{align}
    % % \label{eq:Filter_Rem_finite_case2}
    %     \rem(z) \geq \mathfrak{C}_4^{\prime} \lambda^{\tau}z^{-\tau},
    %     \quad 
    %     &\forall \lambda \in [0, 1], ~ z >  \mathfrak{C}_4 \lambda.
    % \end{align}
    % }

    \item 
    
    When $s>0$, we also assume that the analytic filter function $\reg$ satisfies the remaining conditions listed in Appendix C of \cite{lu2024saturation} (i.e., equations (30), (31), (C1), and (C2) in \cite{lu2024saturation}). For convenience, these requirements are restated in Appendix J.1.
    
    % \dhnote{What are these requirements? I cannot find them. }
        
  \end{enumerate}
  \end{definition}

%We also request that $1-z\reg(z)$ is robust against small perturbation.

\begin{remark}
    The conditions restated in Appendix J.1 essentially require that $\reg$ admits an extension to a larger domain, while preserving all relevant properties that hold on the original domain. This allows us to apply the ``analytic functional argument'' developed by \cite{li2024generalization}.
\end{remark}











\iffalse
\begin{assumption}\label{assume_pert_filter}
    Suppose $z=\Omega(\lambda)$ and $\delta = o_d(\lambda)$, then we have
    $$
    1-z\reg(z+\delta) = {\blue \left[ 1-z\reg(z)\right](1+o_d(1)) + O_d\left(\frac{\delta}{z}\right)} = \left[ 1-z\reg(z)\right](1+o_d(1)) +
    o_d\left(\frac{\lambda}{z}\right).
    $$
\end{assumption}
\fi




For any $\reg$ in an analytic filter family, we define the associated spectral algorithm as follows. 
Suppose the Taylor series expansion of the analytic function $\reg$ is $\reg(z)= b_{\lambda, 0} + b_{\lambda, 1} z + b_{\lambda, 2} z^2 + \cdots$; 
we extend this to the matrix domain for any square matrix $A \in \mathbb{R}^{n \times n}$ as 
\begin{equation}\label{eqn_def_filter_matrix_function}
    \reg(A) := b_{\lambda, 0} \mathrm{I}_{n} + b_{\lambda, 1} A + b_{\lambda, 2} A^2 + \cdots. 
\end{equation}
The spectral estimator is then given by 
\begin{align}
  \label{eq:SA}
  \hat{f}_{\lambda}(\cdot) = \frac{1}{n}\mathscr{K}(\cdot, X) \reg\left(\frac{1}{n}K\right) Y,
\end{align}
where the matrix $K$ is defined as $(K)_{i, j} := \mathscr{K}(x_i, x_j)$; see, e.g., \cite{bauer2007_RegularizationAlgorithms, li2024generalization}. 

\iffalse
For any $y \in \mathbb{R}$, let $K_{x} : \R \to \calH$ be given by $K_{x}(y) = y \cdot \mathscr{K}(x,\cdot)$, whose adjoint $K_{x}^* : \calH \to \R$ is given by $K_{x}^*(f) = \ang{\mathscr{K}(x,\cdot),f}_{\calH} = f(x)$.
Moreover, we denote by $T_{x} = K_{x} K_{x}^*$ and $T_{X} = \frac{1}{n}\sum_{i=1}^n T_{x_i}$.
We also define the sample basis function
\begin{align}
  \label{eq:GZ}
  \hat{g}_Z = \frac{1}{n} \sum\nolimits_{i=1}^n K_{x_i}(y_i) = \frac{1}{n} \sum\nolimits_{i=1}^n y_i \cdot \mathscr{K}(x_i,\cdot).
\end{align}
Now, the estimator of the spectral algorithm is defined by
\begin{align}
  \label{eq:SA}
  \hat{f}_{\lambda} := \reg(T_{X}) \hat{g}_Z.
\end{align}
\fi

    % \lambda = \Theta_d\left(d^{-u}\right),

    

Throughout, we assume that the regularization parameter $\lambda$ follows a polynomial scaling in the dimension $d$. 
\begin{assumption}\label{assume_lambda}
    There exist constants $c_3,c_4>0$ and a constant exponent $u \in (0, \infty]$, such that we have
    $$
    c_3 d^{-u} \leq \lambda \leq c_4 d^{-u},
    $$
    where the case $u = \infty$ corresponds to $\lambda=0$ (i.e., the interpolation limit). Define  $\ell_{\lambda} = \lfloor u \rfloor$ and $\tilde{\ell}=\min\{\ell_{\gamma}, \ell_{\lambda}\}$.
\end{assumption}

The exponent $u$ and the index $\tilde{\ell}$ will be used to formulate the following two technical assumptions. 

%for determining the learning curves.


Definition \ref{def:filter} only imposes conditions on the scalar function $\reg$ defined on $(0, K_{\max}]$. 
However, these conditions are insufficient for the learning curve analysis that depends on evaluating $\reg$ at random kernel matrices as in \eqref{eq:SA}. 
To avoid unnecessary technical complications in the proof, we introduce the following assumption, which will be verified in Section~\ref{subsec_example_spe} for standard analytic filters.


%whereas controlling the order of the learning curve requires us to bound the spectral decomposition of the matrix function $\reg: \mathbb{R}^{n \times n} \to \mathbb{R}^{n \times n}$. 

\begin{assumption}\label{assume_pert_filter}

When $0 < u < \gamma <1$, we require that $\rem$ defined in Definition \ref{def:filter} satisfies
$$
    \rem(K/n) = \Omega_{d, \mathbb{P}}(1)\rem(K_{0}/n), 
    $$
    where %$ \Psi_{0, 1} = \left(\psi_{0, 1}(x_{1}), \cdots, \psi_{0, 1}(x_{n}) \right)^{\top} \in \mathbb{R}^{n}$ and $K_{0} = \mu_0 \Psi_{0, 1} \Psi_{0, 1}^{\top}$.
    $K_{0} = \mu_0 \mathbf{1} \mathbf{1}^{\top}$.
    
% \dhnote{
% Is $K_{0}$ just a constant matrix?
% }

    
    % are defined in Appendix \ref{append_notation}.
    

    % \dhnote{You didn't define $K_{0}$ but defined $K_{\leq 0}$; should be more explicit. }
    % \dhnote{In the main text, can we put the definition of $K$ here (you have already used it in \eqref{eq:SA}) and use words to describe $K_0$? }
    % {\green
    % [Have defined $K$ and $K_0$ above, also add a remark in Appendix A.]
    % }
\end{assumption}


\iffalse
We also require that $\rem$ satisfies the following assumption:
    \begin{itemize}
        
    \item  If $0 < u < \gamma <1$, then we have
    $$
    \rem(K/n) = \Omega_{d, \mathbb{P}}(1)\rem(K_{0}/n), 
    $$
    where the matrices $K:=\mathscr{K}(X, X)$ and $K_{0}$ are defined in Appendix \ref{append_notation}.
    \end{itemize} 
\fi





\subsubsection{Examples of analytic spectral algorithms}\label{subsec_example_spe}

We now illustrate the concept of analytic filter functions using several classical spectral algorithms.

% Certain analytic filter functions can be used to construct many commonly used spectral algorithms.
% We provide several examples as follows.


\begin{example}[Kernel ridge regression]
  \label{example:KRR}
For kernel ridge regression (KRR), the filter is  
  \begin{align*}
    \reg^{\krr}(z) = \frac{1}{z+\lambda},\quad \rem^{\krr}(z) = \frac{\lambda}{z+\lambda},\quad \tau=1.
  \end{align*}
\end{example}

\begin{example}[Kernel gradient flow]
\label{example:gradient_flow}
  % \label{example:GradientFlow}
For kernel gradient flow with time parameter $\lambda^{-1}$, the filter is
  \begin{align*}
    \reg^{\gf}(z) = \frac{1-e^{-z / \lambda}}{z},\quad \rem^{\gf}(z) = e^{-z / \lambda}, \quad \tau=\infty.
  \end{align*}
\end{example}



\begin{example}[Iterated ridge regression]
  \label{example:IteratedRidge}
  Let $q \geq 1$ be fixed.
  The $q$-fold iterated ridge estimator corresponds to the filter 
  \begin{align*}
    \reg^{\mathrm{IT}, q}(z) = \frac{1}{z}\left[ 1 - \frac{\lambda^q}{(z+\lambda)^q} \right],
    \quad
    \psi^{\mathrm{IT}, q}_\lambda(z) =\frac{\lambda^q}{(z+\lambda)^q},\quad \tau=q.
  \end{align*}
\end{example}





\begin{example}[Kernel gradient descent]
  \label{example:GradientDescent}
  Kernel gradient descent may be viewed as a time-discretized version of kernel gradient flow.
Fix a step size $\eta>0$, and apply gradient descent to the empirical loss for $t$ iterations. 
The resulting estimator is associated with the spectral filter
  \begin{align*}
    \reg^{\mr{GD}}(z) &= \eta \sum_{k=0}^{t-1} (1-\eta z)^{k} = \frac{1-(1-\eta z)^t}{z},\quad \lambda = (\eta t)^{-1}, \\
    \rem^{\mr{GD}}(z) &= (1-\eta z)^t, \quad\tau=\infty.
  \end{align*}
\end{example}



The filter functions in the above examples 
satisfy Assumption~\ref{assume_pert_filter}; see 
Appendix J.3 
for a detailed verification.


\subsection{Nondegenerate signal}


% Finally, we need the following assumption to obtain the corresponding lower bound of the excess risk.


To derive matching lower bounds for the excess risk, we require that the signal has nontrivial energy in the degrees around $\tilde{\ell}$ defined in Assumption~\ref{assume_lambda}.


\begin{assumption}\label{assump_source condition_lower_bound}
There exists an absolute constant $c_{0} > 0$ such that for any $ d \geq \mathfrak{C}$, we have
\begin{equation*}
%\label{ass of fi}
        \mu_{m}^{-s} \sum\limits_{j=1}^{N(d, m)}  \theta_{m,j}^{2} \ge c_{0}, ~~m=\tilde{\ell}, \tilde{\ell}+1;
   \end{equation*}
Moreover, if $s> 2\tau$, we further assume that
$$
\sum\limits_{k=0}^{\tilde{\ell}} \sum\limits_{j =1}^{N(d, k)}  \theta_{k, j}^{2} \ge c_{0}.
$$
\end{assumption}


Assumption~\ref{assump_source condition_lower_bound} ensures that $s$ is the effective smoothness level of $f_{\star}$ for the purpose of lower bounding the excess risk.
Assumption \ref{assump_function_calss} alone allows the possibility that $f_{\star}\in[\mathcal{H}]^{t}$ for some $t>s$, in which case the excess risk may decay faster and a sharp lower bound at smoothness level $s$ cannot be identified. 
Assumption~\ref{assump_source condition_lower_bound} rules out this scenario and allows us to obtain a matching lower bound.
Following the discussion in Section~2.5 in the supplementary material of \cite{zhang2024phase}, we can show that Assumptions \ref{assump_function_calss} and \ref{assump_source condition_lower_bound} together imply that
\begin{equation}\label{implication of assumption 2}
\left\{
\begin{gathered}
f_{\star} \in [\mathcal{H}]^{s} \quad \text{for any } d; \\
\forall t > s,\quad f_{\star} \notin [\mathcal{H}]^{t} \quad \text{when } d \text{ is sufficiently large}.
\end{gathered}
\right.
\end{equation}
Moreover, we can further show that functions violating Assumption \ref{assump_source condition_lower_bound} are easier to learn, in the sense that they admit faster convergence rates. Therefore, Assumption~\ref{assump_source condition_lower_bound} is essential for characterizing the worst-case convergence rate of the excess risk over the class $[\mathcal{H}]^{s}$.
Similar assumptions have been adopted when establishing lower bounds for excess risk. Examples include (8) in \cite{Cui2021GeneralizationER} and Assumption 3 in \cite{li2023asymptotic} for fixed-dimensional settings, and Assumption 5 in \cite{zhang2024optimal}, Assumption 2 in \cite{lu2024saturation}, and Assumption 2 in \cite{zhang2024phase} for large-dimensional settings. 

\begin{remark}
    Several lower bound results in the literature remain valid, in the sense that the same proofs apply, if their respective non-degeneracy conditions are replaced by Assumption \ref{assump_source condition_lower_bound}. 
    These include Lemmas~21--23 of \cite{zhang2024optimal}, Lemma~D.14 of \cite{lu2024saturation}, and Lemmas~5.4.3--5.4.4 of \cite{zhang2024phase}. 
    %
% \dhnote{How about this: Several lower bound results in the literature remain valid, in the sense that the same proofs apply, if their respective non-degeneracy conditions are replaced by Assumption \ref{assump_source condition_lower_bound}. 
% These include Lemmas~21--23 of \cite{zhang2024optimal}, Lemma~D.14 of \cite{lu2024saturation}, and Lemmas~5.4.3--5.4.4 of \cite{zhang2024phase}. 
% }
    % By examining Lemmas 21--23 in \cite{zhang2024optimal}, Lemma D.14 in \cite{lu2024saturation}, and Lemmas 5.4.3 and 5.4.4 in \cite{zhang2024phase}, we observe that even if the corresponding assumptions are replaced by Assumption \ref{assump_source condition_lower_bound}, their conclusions still hold.
\end{remark}








\section{Main results}\label{sec_main_results}

This section is devoted to determining the convergence rate of the excess risk of $\hat{f}_{\lambda}$ under the setup in Section~\ref{sec_settings}. 
We begin with the bias--variance decomposition:
\begin{align}\label{eq bias var decomposition}
    E_{x, \epsilon} \left[ \left(\hat{f}_{\lambda}(x) - f_{\star}(x) \right)^{2} \right] =  \mathrm{var}(\hat{f}_{\lambda}) + \mathrm{bias}^2(\hat{f}_{\lambda}),
\end{align}
where $E_{x, \epsilon}$ denotes expectation with respect to the new sample $x$ and the random noise terms $\epsilon_i$ in the training samples,
\begin{equation}\label{eq var term}
    \mathrm{var}(\hat{f}_{\lambda}) := E_{x, \epsilon}\left[\left(\hat{f}_{\lambda}(x)-E_\epsilon\left(\hat{f}_{\lambda}(x)\right)\right)^2\right],
\end{equation}
and
\begin{equation}\label{eq bias term}
    \mathrm{bias}^2(\hat{f}_{\lambda}) := E_x\left[\left(E_\epsilon\left(\hat{f}_{\lambda}(x)\right)-f_{\star}(x)\right)^2\right].
\end{equation}

By determining the convergence rate of each term in \eqref{eq bias var decomposition}, our main result is the complete learning curve of large-dimensional spectral algorithms.

\begin{theorem}\label{theorem_learn_curve}
    Suppose Assumptions \ref{assump_asymptotic}, \ref{assu:coef_of_inner_prod_kernel}, 
\ref{assump_function_calss},
\ref{assume_lambda}, 
\ref{assume_pert_filter},
and \ref{assump_source condition_lower_bound} hold with $\tau \leq \infty$, $s \geq 0$, and $u \in (0, \infty]$. 
Denote $\tilde{s}=\min\{s, 2\tau\}$. Then the excess risk of the spectral algorithm estimator in (\ref{eq bias var decomposition}) satisfies
\begin{equation*}
\begin{aligned}
    E_{x, \epsilon} \left[ \left(\hat{f}_{\lambda}(x) - f_{\star}(x) \right)^{2} \right]
= \Theta_{d, \mathbb{P}} \left( 
d^{\gamma - \tilde{\ell} - 1 - 2  \max\{\gamma - u, 0\}}
+ d^{\tilde{\ell} - \gamma}
+ d^{-(\tilde{\ell} + 1)  s}
+ \mathbf{1}\{\tau<\infty\} d^{-2  \tau  u + (2  \tau - \tilde{s})  \tilde{\ell}}
\right).
\end{aligned}
\end{equation*}
\end{theorem}

\iffalse
\begin{equation*}
\begin{aligned}
    E_{x, \epsilon} \left[ \left(\hat{f}_{\lambda}(x) - f_{\star}(x) \right)^{2} \right]
= \Theta_{d, \mathbb{P}} \left( \frac{n}{d^{\tilde{\ell}+1}(n\lambda+1)^{2}} + \frac{d^{\tilde{\ell}}}{n} + 
    d^{-(\tilde{\ell}+1)s}
    +\mathbf{1}\{\tau<\infty\}\lambda^{2\tau}d^{(2\tau-\tilde{s})\tilde{\ell} }\right).
\end{aligned}
\end{equation*}
\fi



\iffalse
\dhnote{Why isn't $\lambda$ replaced by $d^{-u}$? $\tilde{l}$ also depends on $u$. }
\dhnote{
Similarly, why not replace $n$ by $d^{\gamma}$
}


{\green You are right that it will be clearer to express $\lambda$ and $n$ by $d^{-u}$ and $d^{\gamma}$. I will change all these notations in the main results and the corresponding proof.}
\fi


Theorem \ref{theorem_learn_curve} fully characterizes the convergence rate of the excess risk in different asymptotic regimes $n = \Theta_d( d^{\gamma} )$, different smoothness $s \geq 0$, different filter qualification $\tau$, and different choices of the regularization parameter $\lambda=\Theta_d(d^{-u})$. Therefore, it provides the complete learning curve for large-dimensional spectral algorithms. 


Compared with interpolation-only or optimally tuned analyses, a key difficulty is the regime $\ell_\lambda < \ell_\gamma$ (where $0<\lambda \ll n^{-1}$). 
In this regime, the proof can no longer be organized solely around the interpolation cutoff $\ell_\gamma$, because the block $\ell_\lambda < k \le \ell_\gamma$ is neither absorbed into the low-degree fitted block nor negligible as a tail term. 
One must isolate and control this intermediate spectral block in both the bias and the variance. 
The remaining ingredients combine the large-dimensional kernel-matrix concentration results of \cite{ghorbani2021linearized} with the analytic functional framework for spectral filters. For the reader's convenience, Appendix~C provides a proof sketch for a representative case.

The full proof is deferred to the appendices. Appendix~D bounds the variance term in \eqref{eq var term}, Appendix~E bounds the bias term in \eqref{eq bias term}, and Appendix~F combines these bounds with Proposition~B.1 to prove Theorem~\ref{theorem_learn_curve}.



% To highlight the key new steps, we provide a proof sketch for a representative case in Appendix C. 
% Apart from 
% large-dimensional kernel matrix concentration results of
% \cite{ghorbani2021linearized}, 
% the main technical novelty of the paper is the treatment of the under-regularized case $\ell_\lambda < \ell_\gamma$. 
% In this case, the proof can no longer be organized solely at the interpolation cutoff. 
% Instead, one must track an intermediate spectral block between $\ell_\lambda+1$ and $\ell_\gamma$ in the analysis. 
% This creates a new difficulty absent from interpolation-only or optimally tuned analyses.



% The proof of Theorem \ref{theorem_learn_curve} is deferred to the appendices and is organized as follows: We will first bound the variance term (\ref{eq var term}) and the bias term (\ref{eq bias term}) in Appendix D and Appendix E, respectively. Then, combining with results in \cite{lu2024saturation} (see, e.g., Proposition B.1), we prove Theorem \ref{theorem_learn_curve} in Appendix F.





% We note that the upper bound relies only on Assumptions \ref{assump_asymptotic}--\ref{assume_pert_filter}, while Assumption \ref{assump_source condition_lower_bound} is needed solely for the matching lower bound. 


% We discuss the interpretation and several implications of Theorem \ref{theorem_learn_curve} in the following subsections. 

Theorem~\ref{theorem_learn_curve} will be interpreted in the following subsections from the perspectives of phase transitions, benign overfitting, saturation, and sequence-model equivalence.


\subsection{Three regimes of the learning curve} 
The rate in Theorem \ref{theorem_learn_curve} comprises two components: the variance error rate $d^{\gamma - \tilde{\ell} - 1 - 2  \max\{\gamma - u, 0\}}
+ d^{\tilde{\ell} - \gamma}$ and the bias error rate $d^{-(\tilde{\ell} + 1)  s}
+ \mathbf{1}\{\tau<\infty\} d^{-2  \tau  u + (2  \tau - \tilde{s})  \tilde{\ell}}$.
In general, as $u$ (the rate of the regularization parameter $\lambda$) increases, the variance error rate gradually decreases while the bias error rate increases. 
By balancing the rates of the bias and variance terms, one can identify the optimal choice of $u$, denoted as $u^{\prime} = u^{\prime}(\tau)$. 
The quantity $u^{\prime}$ is conceptually important, but its explicit expression is not involved in our discussion, so we defer its expression to Proposition F.4 in the appendix.  

% or \cite{lu2024saturation}

Using $u^{\prime}$ and $\gamma$, we divide the learning curve provided by Theorem \ref{theorem_learn_curve} into three regimes as $u$ increases (equivalently, $\lambda$ decreases):
\begin{itemize}
    \item {\bf Over-regularized regime: } When $u < u^{\prime}$, there is too much regularization and the bias term dominates. In this regime, the excess risk rate is non-increasing in $u$ and the learning curve has many plateaus.
    
    \vspace{3pt}
    
    \item {\bf Under-regularized regime: } When $u \in (u^{\prime}, \gamma)$, 
    the rate is non-decreasing in $u$ and the learning curve has many plateaus.

    \vspace{3pt}
    
    \item {\bf Interpolation regime: } When $u \geq \gamma$, the rate remains unchanged.
\end{itemize}

% \dhnote{The above cases may be better dived by the regimes of $u$. }

\iffalse
\begin{remark}
An interesting observation is that the bias error rate can still be dominant in the under-regularized and interpolation regime when $s$ is too small.
%Though the learning curve is non-decreasing in the under-regularized regime, one of the bias error rate $d^{-(\tilde{\ell}+1)s}$ can still be the dominant term when $s$ is too small. 
For example, consider the case $\gamma = 0.9$, $s=0.5$, and $\tau=\infty$. In this case, we have $u^{\prime}=0.45$. For any $u \geq 0.45$, the variance term is $d^{\gamma - \tilde{\ell} - 1 - 2  \max\{\gamma - u, 0\}}
+ d^{\tilde{\ell} - \gamma} = o_d(d^{-(\tilde{\ell}+1)s})$, so it is dominated by the bias term.
\end{remark}
\fi




From the above description, we can conclude that the overall learning curve of the large-dimensional spectral algorithm does not follow a simple U-shape. Specifically,  as illustrated in Figure \ref{fig_theory_benign_overfit} (d):
\begin{itemize}
    \item Several plateaus appear in the over-regularized and under-regularized regimes;

    \vspace{3pt}
    
    \item In the interpolation regime, the learning curve is essentially a straight horizontal line over $u\in [\gamma, \infty)$, meaning that the generalization performance of the spectral algorithm is equivalent to that of kernel interpolation (i.e., $u=\infty$).
\end{itemize}


% Theorem \ref{theorem_learn_curve} also offers a fresh perspective on understanding the benign overfitting phenomenon and the saturation effect discussed in \citep{zhang2024optimal, lu2024saturation, zhang2024phase}. 
% \dhnote{does this belong to the next part?}


% {\green [The following three paragraphs have been modified. Specifically, the second paragraph is copied from the Introduction section.]}


\begin{remark}
Theorem \ref{theorem_learn_curve} reveals an interesting phenomenon.
When $s$ is sufficiently small, the error can remain dominated by the bias even for large $u$ in the under-regularized and interpolation regimes. 
For example, consider the case $\gamma = 0.9$, $s=0.5$, and $\tau=\infty$. In this case, we have $u^{\prime}=0.45$. 
    For any $u \geq 0.45$, the variance term is of order $d^{\gamma - \tilde{\ell} - 1 - 2  \max\{\gamma - u, 0\}}
    + d^{\tilde{\ell} - \gamma} = o_d(d^{-(\tilde{\ell}+1)s})$, so it is negligible relative to the bias term.
    This suggests that, when $s$ is sufficiently small, spectral algorithms can still achieve strong performance, and even benign overfitting, well into the under-regularized and interpolation regime. Section~\ref{sec: benigh overfitting} makes this statement precise by characterizing benign overfitting.
\end{remark}





\subsection{Full characterization of benign overfitting}\label{sec: benigh overfitting}

Theorem \ref{theorem_learn_curve} yields a sharp characterization of benign overfitting across the under-regularized and interpolation regimes. 
To state it, recall the minimax benchmark
\[
d^{p-\gamma}+d^{-(p+1)s},
\qquad
p:=\left\lfloor \frac{\gamma}{s+1}\right\rfloor,
\]
which is the optimal rate of excess risk over the class $[\calH]^s$ (see Proposition B.2).
Benign overfitting occurs when the estimator is under regularized and the risk is minimax rate-optimal. 
We further recall a threshold function from \citet{zhang2024phase}:
\begin{equation}\label{eqn_threshold_s}
\Gamma(\gamma) := \begin{cases}
\infty,  & \gamma \in (0, 0.5], \\
(\gamma-\ell_{\gamma})/\ell_{\gamma}, & \gamma \in (\ell_{\gamma}, \ell_{\gamma}+0.5] ~\text{ and }~ \ell_{\gamma} \geq 1, \\
(\ell_{\gamma}+1-\gamma)/(\ell_{\gamma}+1), & \gamma \in (\ell_{\gamma}+0.5, \ell_{\gamma}+1),\\
0, & \gamma =1, 2, \cdots .
\end{cases}
\end{equation}

The following corollary compares excess risk of the spectral algorithm relative to the minimax rate according to $s\leq \Gamma(\gamma)$  and $s>\Gamma(\gamma)$. 



\iffalse

%%%%%%%%%%%%%%%%%%%%%%%%%%%%%%%%%%%%%%%
{\color{red} [The following three paras should be deleted]} Existing work has primarily focused on the benign overfitting phenomenon in large-dimensional kernel ridge regression within the interpolation regime (see, e.g., \cite{Liang_Just_2019, liang2020multiple, aerni2023strong, barzilai2023generalization, zhang2024phase}). For example, \cite{zhang2024phase} determined the convergence rate of the excess risk of the large-dimensional kernel interpolation algorithm, i.e., the case $\lambda = 0$ (see, e.g., Proposition \ref{prop_intepolation}). 
\dhnote{why ref this prop?}
{\green [Proposition \ref{prop_intepolation} restates Corollary 1 in \cite{zhang2024phase}]}
They further found that benign overfitting exhibits in the interpolation regime (i.e., kernel interpolation is minimax rate-optimal) if and only if $0<s \leq \Gamma(\gamma)$, where
    \begin{equation}
        \begin{aligned}
    \Gamma(\gamma) := \begin{cases} 
    \infty,  & \gamma \in (0, 0.5], \\ 
    (\gamma-\ell_{\gamma})/\ell_{\gamma}, & \gamma \in (\ell_{\gamma}, \ell_{\gamma}+0.5] ~\text{ and }~ \ell_{\gamma} \geq 1, \\ 
    (\ell_{\gamma}+1-\gamma)/(\ell_{\gamma}+1), & \gamma \in (\ell_{\gamma}+0.5, \ell_{\gamma}+1)\\
    0, & \gamma =1, 2, \cdots.
    \end{cases}
    \end{aligned}
    \end{equation}
However, these prior studies do not fully resolve whether benign overfitting extends to the under-regularized regime. 
The motivation for studying benign overfitting stems from the goal of ensuring good generalization even under weak regularization (e.g., when the stopping time in gradient descent is very large). 
In practice, computational limitations often yield solutions in the under-regularized regime rather than exact interpolation. Therefore, if benign overfitting were possible in the under-regularized regime while interpolation fails for $s > \Gamma(\gamma)$, previous conclusions regarding benign overfitting would be incomplete.
\dhnote{The last sentence is very confusing... what does "while interpolation fails for $s > \Gamma(\gamma)$" mean? What are "previous conclusions"? What is considered "complete" in order to say "incomplete".  }

\dhnote{"In contrast" need A vs B; what are A and B here?}
In contrast, Theorem \ref{theorem_learn_curve} leads to the following corollary.
%%%%%%%%%%%%%%%%%%%%%%%%%%%%%%%%%%%%%%%%%
\fi 


\begin{theorem}\label{coroll_benign_overfit}
    Suppose all assumptions in Theorem \ref{theorem_learn_curve} hold and define \(p := \lfloor \gamma/(s+1) \rfloor\). 
    Recall that $u^{\prime} = u^{\prime}(\tau)$ is the optimal choice of $u$ that balances the rates of the bias and variance terms (see Proposition F.4 for its expression). 
    Then we have the following statements:
    
    \begin{itemize}
        \item[(i)] Suppose $0 \leq s \leq \Gamma(\gamma)$. Then for any $u \in [u^{\prime}, \infty]$, the excess risk of the spectral algorithm estimator in (\ref{eq bias var decomposition}) satisfies
\begin{equation*}
\begin{aligned}
    E_{x, \epsilon} \left[ \left(\hat{f}_{\lambda}(x) - f_{\star}(x) \right)^{2} \right]
= \Theta_{d, \mathbb{P}} \left( 
d^{p-\gamma} + 
    d^{-(p+1)s}
\right).
\end{aligned}
\end{equation*}

\item[(ii)] Suppose $s > \Gamma(\gamma)$. Then there exists $\tilde{u} \in (u^{\prime}, \gamma)$ such that for $\tilde{\lambda}=d^{-\tilde{u}}$, we have
\begin{equation*}
\begin{aligned}
    E_{x, \epsilon} \left[ \left(\hat{f}_{\tilde{\lambda}}(x) - f_{\star}(x) \right)^{2} \right]
 \gg  
d^{p-\gamma} + 
    d^{-(p+1)s} \quad \text{ in probability.}
\end{aligned}
\end{equation*}

    \end{itemize}
\end{theorem}


\iffalse %dh 2026-03-26: reconstruct the context 
Theorem~\ref{coroll_benign_overfit} has two implications.
Firstly, 
recall that the minimax rate is \(d^{p-\gamma} + 
    d^{-(p+1)s}\) (see, e.g., Proposition B.2). Therefore, when the regression function $f_{\star}$ is less smooth w.r.t. $\mathcal{H}$, benign overfitting in large-dimensional spectral algorithms can occur not only in the interpolation regime but also in the under-regularized regime. 
    This observation aligns with empirical findings on neural networks (e.g., Figure 2 in \cite{zhang2021understanding}, and \cite{zhang2016understanding, belkin2018understand, belkin2019does, belkin2021fit, nakkiran2021deep}): even without strong smoothness assumptions on the regression function, neural networks generalize well when the training time $t = \lambda^{-1}$ meets or exceeds the optimal stopping time. Our results thus offer a new theoretical perspective on benign overfitting in neural networks.
Secondly, the condition for benign overfitting remains consistent across both the under-regularized and interpolation regimes: it occurs if and only if $0 < s \leq \Gamma(\gamma)$.
\fi

When $0 < s \le \Gamma(\gamma)$, the theorem shows that once the regularization level enters the range $u \ge u^{\prime}$, the excess risk remains on the minimax scale for every further choice of $u$, including both the under-regularized regime $u \in (u^{\prime},\gamma)$ and the interpolation limit $u=\infty$. 
Therefore, in the less smooth regime, benign overfitting for large-dimensional spectral algorithms occurs not only at interpolation, but throughout the under-regularized regime as well. 
This observation aligns with empirical findings on neural networks, where even without strong smoothness assumptions on the regression function, neural networks generalize well when the training time $t = \lambda^{-1}$ meets or exceeds the optimal stopping time; see \cite{zhang2021understanding, zhang2016understanding, belkin2018understand, belkin2019does, belkin2021fit, nakkiran2021deep}.
In this sense, our result offers a kernel-limit perspective relevant to benign-overfitting phenomena discussed in the neural-network literature.

% In this sense, our result thus offers a new theoretical perspective on benign overfitting in neural networks.


Moreover, the threshold $\Gamma(\gamma)$ separates benign and non-benign behavior in both the under-regularized and interpolation regimes.
Existing work has primarily focused on the benign overfitting phenomenon in large-dimensional kernel ridge regression within the interpolation regime. 
See, for example, \cite{Liang_Just_2019, liang2020multiple, aerni2023strong, barzilai2023generalization, zhang2024phase}.
In particular, \cite{zhang2024phase} derived the convergence rate of the excess risk for kernel interpolation when $\lambda=0$ (restated in Proposition B.3) and established that benign overfitting occurs if and only if $0 < s \leq \Gamma(\gamma)$. 
Theorem~\ref{coroll_benign_overfit} shows that this same condition is exactly what allows the benign behavior to persist beyond interpolation and throughout the whole region $u \geq u^{\prime}$. When $s > \Gamma(\gamma)$, this persistence fails, since there exists some $\tilde u \in (u^{\prime},\gamma)$ for which the excess risk is strictly above the minimax rate.





\iffalse % dh 2026-03-26: this part is absorbed in the interpretation and in the introduction. 
Existing work has primarily focused on the benign overfitting phenomenon in large-dimensional kernel ridge regression within the interpolation regime. 
See, for example, \cite{Liang_Just_2019, liang2020multiple, aerni2023strong, barzilai2023generalization, zhang2024phase}.
In particular, \cite{zhang2024phase} derived the convergence rate of the excess risk for kernel interpolation (restated in Proposition \ref{prop_intepolation}) and established that benign overfitting occurs (i.e., kernel interpolation is minimax rate-optimal) if and only if $0 < s \leq \Gamma(\gamma)$, where
\begin{equation}\label{eqn_threshold_s}
        \begin{aligned}
    \Gamma(\gamma) := \begin{cases} 
    \infty,  & \gamma \in (0, 0.5], \\ 
    (\gamma-\ell_{\gamma})/\ell_{\gamma}, & \gamma \in (\ell_{\gamma}, \ell_{\gamma}+0.5] ~\text{ and }~ \ell_{\gamma} \geq 1, \\ 
    (\ell_{\gamma}+1-\gamma)/(\ell_{\gamma}+1), & \gamma \in (\ell_{\gamma}+0.5, \ell_{\gamma}+1)\\
    0, & \gamma =1, 2, \cdots.
    \end{cases}
    \end{aligned}
    \end{equation}

However, these results do not fully address whether benign overfitting can persist into the under-regularized regime. 
Importantly, studying benign overfitting of spectral algorithms in the under-regularized regime is closely aligned with a common practical setting, namely early stopping, as considered in seminal works on benign overfitting. 
For example, \cite{zhang2016understanding} trained the neural networks with gradient descent, and stopped training at a finite time $t=1/\lambda \geq T$, which corresponds to a small but nonzero $\lambda$.
This suggests that the benign overfitting phenomenon they observed likely occurs in the under-regularized regime.



In contrast to results for only kernel interpolation (i.e., $\lambda=0$) in \cite{zhang2024phase}, Theorem \ref{theorem_learn_curve} allows us to characterize learning curves across both the under-regularized and the interpolation regimes.
We summarize its implication for benign overfitting in the following corollary. 
\fi 


\iffalse
\begin{itemize}
    \item Recall that the minimax rate is \(d^{p-\gamma} + 
    d^{-(p+1)s}\) (see, e.g., Proposition B.2). Therefore, when the regression function $f_{\star}$ is less smooth w.r.t. $\mathcal{H}$, benign overfitting in large-dimensional spectral algorithms can occur not only in the interpolation regime but also in the under-regularized regime.\\ 
    This observation aligns with empirical findings on neural networks (e.g., Figure 2 in \cite{zhang2021understanding}, and \cite{zhang2016understanding, belkin2018understand, belkin2019does, belkin2021fit, nakkiran2021deep}): even without strong smoothness assumptions on the regression function, neural networks generalize well when the training time $t = \lambda^{-1}$ meets or exceeds the optimal stopping time. Our results thus offer a new theoretical perspective on benign overfitting in neural networks.

    \item The condition for benign under-regularized remains consistent across both the under-regularized and interpolation regimes: it occurs if and only if $0 < s \leq \Gamma(\gamma)$.
        
\end{itemize}
\fi



The following example provides further insight into why benign overfitting only occurs for small $s$.
Consider the Sobolev RKHS $W^{m,2}(\mathcal{X})$ with $m > d/2$ and the extreme case where $u=\infty$ (i.e., kernel interpolation). The interpolation space of $W^{m,2}(\mathcal{X})$ satisfies $[W^{m,2}(\mathcal{X})]^s \cong W^{m s, 2}(\mathcal{X})$, meaning that larger values of $s$ correspond to higher-order differentiability and smoothness of functions in the interpolation space.
For small $s$, the regression function contains substantial high-frequency components. In this case, kernel interpolation that passes through all data points primarily captures these high-frequency features of the underlying signal, leading to an excess risk that matches the minimax optimal rate. 
In contrast, for large $s$, the regression function is inherently smooth. However, noise in the training data forces the kernel interpolation estimator to oscillate sharply to fit every data point exactly. This results in spurious high-frequency fluctuations between training points, causing significant deviation from the smooth regression function.






Figure \ref{fig_theory_benign_overfit} illustrates the behavior of learning curves when $s \leq \Gamma(\gamma)$ and $s > \Gamma(\gamma)$. As shown in Figure \ref{fig_theory_benign_overfit} (a) and (c), when $s \leq \Gamma(\gamma)$, the learning curve remains flat throughout both the under-regularized regime (indicated by hatched patterns) and the interpolation regime. In contrast, when $s > \Gamma(\gamma)$, the learning curve increases in the under-regularized regime.





\begin{figure}[btp]
% \vskip 0.05in
\centering
\subfigure[]{\includegraphics[width=0.48\columnwidth]{figures/zeta_plot_for_s=0.5_gamma=1.5.pdf}}
\subfigure[]{\includegraphics[width=0.48\columnwidth]{figures/zeta_plot_for_s=1.5_gamma=1.5.pdf}}

\vspace{-10pt}

\subfigure[]{\includegraphics[width=0.48\columnwidth]{figures/zeta_plot_for_s=0.13_gamma=2.6.pdf}}
\subfigure[]{\includegraphics[width=0.48\columnwidth]{figures/zeta_plot_for_s=0.2_gamma=2.6.pdf}}

\vspace{-10pt}

\caption{
A graphical representation of the learning curves of large dimensional spectral algorithms with $\tau=1$ (blue) and $\tau=4$ (orange) obtained in Theorem \ref{theorem_learn_curve};
and the minimax rate (red) obtained in Proposition B.2.
Specifically, for any given $(s, \gamma, \tau, u)$, we set $\zeta = \max\{\gamma - \tilde{\ell} - 1 - 2  \max\{\gamma - u, 0\}, \tilde{\ell} - \gamma, -(\tilde{\ell} + 1)  s, \ln(\mathbf{1}\{\tau<\infty\}) + -2  \tau  u + (2  \tau - \tilde{s})  \tilde{\ell}\}$ for the spectral algorithm, and set $\zeta = \max\{p-\gamma, -(p+1)s\}$ for minimax rate.
In subfigure (a) and (c), $s \leq \Gamma(\gamma)$; while in subfigure (b) and (d), $s > \Gamma(\gamma)$. 
Under-regularized regimes are marked by hatched patterns with corresponding colors.
}
\label{fig_theory_benign_overfit}
% \vskip 0.05in
\end{figure}




\iffalse
 In contrast, Theorem~\ref{coroll_benign_overfit} demonstrates that when the regression function $f_{\star}$ is less smooth w.r.t. $\mathcal{H}$, benign overfitting in large-dimensional spectral algorithms can occur not only in the interpolation regime but also in the under-regularized regime. 
We note that this observation closely resembles the benign overfitting phenomenon observed in numerical experiments on neural networks (see, e.g., Figure 2 in \cite{zhang2021understanding}, and \cite{zhang2016understanding, belkin2018understand, belkin2019does, belkin2021fit, nakkiran2021deep}): 
Even without requiring the regression function to be relatively smooth, neural networks generalizes equally well when the training time $t=\lambda^{-1}$ equals or exceeds the optimal stopping time. 
Therefore, our results provide a new perspective on investigating benign overfitting in neural networks.
\fi









\subsection{Explaining the saturation effect}

From Theorem \ref{theorem_learn_curve} and the minimax benchmark (Proposition B.2) we obtain the \emph{saturation effect} in large dimensions: spectral algorithms with qualification \(\tau < s\) fail to attain the rate of the information-theoretic lower bound on the excess risk, whereas those with qualification \(\tau \geq s\) do. Specifically, when \(\tau < s\), no matter how carefully one tunes $u$, the convergence rate of the excess risk is much slower than the minimax rate $d^{p-\gamma}+d^{-(p+1)s}$, where $p:=\lfloor \gamma / (s+1) \rfloor$; while when \(\tau \geq s\), the convergence rate of the excess risk attains the above minimax rate with $u = u^{\prime}$ defined in Proposition F.4.

While the same phenomenon was previously reported in \citet{zhang2024optimal, lu2024saturation}, those authors did not explain why this dichotomy occurs. 
Our Theorem \ref{theorem_learn_curve} fills this gap by revealing the geometric mechanism behind the saturation effect: 
\begin{itemize}
    \item []
    \textit{A smaller qualification parameter $\tau$ makes the learning curve decrease more slowly in the over-regularized regime, which shifts the optimal choice $u^{\prime}(\tau)$ to the right and raises the minimum excess risk.}
\end{itemize}



\iffalse

Previous work \citep{zhang2024optimal, lu2024saturation} characterized the excess risk of optimally tuned spectral algorithms and demonstrated the potential \emph{saturation effect}: 
spectral algorithms with qualification \(\tau < s\) fail to attain the rate of the information-theoretic lower bound on the excess risk, whereas those with qualification \(\tau \geq s\) do.
However, these results do not explain why this dichotomy occurs.



By characterizing the full learning curve as a function of $u$,  
Theorem \ref{theorem_learn_curve} reveals 
\fi

\begin{figure}[btp]
% \vskip 0.05in
\centering
\subfigure[]{\includegraphics[width=0.48\columnwidth]{figures/zeta_plot_for_s=2.5_gamma=6.0.pdf}}
\subfigure[]{\includegraphics[width=0.48\columnwidth]{figures/zeta_plot_for_s=3.0_gamma=3.0.pdf}}


\caption{
Another graphical representation of the learning curve of large dimensional spectral algorithms with $\tau=1$ (blue),
$\tau=2$ (green),
and $\tau=4$ (orange) obtained in Theorem \ref{theorem_learn_curve};
and the minimax rate (red) obtained in Proposition B.2. 
Dashed vertical lines indicate the optimal rate of regularization parameter $\lambda$ with qualification $\tau$, that is, the global minimum point $u^{\prime}(\tau)$ of the learning curve.
}
\label{fig_theory_saturation}
% \vskip 0.05in
\end{figure}

We use Figure~\ref{fig_theory_saturation} to illustrate the above geometric mechanism. 
Consider two learning curves with $s \in \{2.5, 3\}$ and qualification parameters $(\tau_1, \tau_2)=(1, 4)$, and let $u^{\prime}(\tau_i)$, $i=1,2$, denote the corresponding global minimum points. 
Because $\tau_1$ is smaller, its descent is slower, so $u^{\prime}(\tau_1) > u^{\prime}(\tau_2)$.
Now observe from the figure that the $\tau_1$-curve meets the $\tau_2$-curve exactly at $u^{\prime}(\tau_1)$, the point where the $\tau_1$-curve attains its own minimum.
At the point $u^{\prime}(\tau_1)$, the $\tau_2$-curve is already in its increasing regime where the variance term already dominates. 
Since the variance term is shared among all $\tau$, the minimum of the $\tau_1$-curve cannot pass below the $\tau_2$-curve. 
Hence the minimum excess risk achieved by $\tau_1$ is strictly larger than that achieved by $\tau_2$. 









\iffalse

{\color{red}
Previous work \citep{zhang2024optimal, lu2024saturation} characterized the excess risk of optimally tuned spectral algorithms and demonstrated the potential \emph{saturation effect}: 
spectral algorithms with qualification \(\tau < s\) fail to attain the rate of the information-theoretic lower bound on the excess risk, whereas those with qualification \(\tau \geq s\) do. 
Theorem \ref{theorem_learn_curve} complements these results by characterizing the entire learning curve as a function of $u$. 
This fuller description reveals the geometric mechanism behind the saturation effect: 
\begin{quote}
    \textit{a smaller qualification parameter $\tau$ makes the learning curve decrease more slowly in the over-regularized regime, which shifts the optimal choice $u^{\prime}(\tau)$ to the right and raises the minimum excess risk.}
\end{quote}


% \begin{figure}[btp]
% % \vskip 0.05in
% \centering
% \subfigure[]{\includegraphics[width=0.48\columnwidth]{figures/zeta_plot_for_s=2.5_gamma=6.0.pdf}}
% \subfigure[]{\includegraphics[width=0.48\columnwidth]{figures/zeta_plot_for_s=3.0_gamma=3.0.pdf}}


% \caption{
% Another graphical representation of the learning curve of large dimensional spectral algorithms with $\tau=1$ (\textcolor{myblue}{blue}),
% $\tau=2$ (\textcolor{mygreen}{green}),
% and $\tau=4$ (\textcolor{myorange}{orange}) obtained in Theorem \ref{theorem_learn_curve};
% and the minimax rate ({\color{red} red}) obtained in Proposition B.2. 
% Dashed vertical lines indicate the optimal rate of regularization parameter $\lambda$ with qualification $\tau$, that is, the global minimum point $u^{\prime}(\tau)$ of the learning curve.
% }
% \label{fig_theory_saturation_copied}
% % \vskip 0.05in
% \end{figure}


To see this geometric mechanism, consider two learning curves with qualification parameters $\tau_2 > \tau_1 \geq 1$, and let $u^{\prime}(\tau_i)$, $i=1,2$, denote the corresponding global minimum points. 
As shown in Figure~\ref{fig_theory_saturation}, both learning curves decrease in the over-regularized regime and increase in the under-regularized and interpolation regimes. 
The smaller qualification $\tau_1$ produces a flatter descent in the over-regularized regime, and therefore yields a higher curve.
As a consequence, the minimizer $u^{\prime}(\tau_1)$ occurs at a larger value of $u$ than $u^{\prime}(\tau_2)$. 
At the point $u^{\prime}(\tau_1)$, the $\tau_2$-curve is already in its increasing regime where the variance term already dominates. 
Since the variance term is shared among all $\tau$, the minimum of the $\tau_1$-curve cannot go across the $\tau_2$-curve. 
Hence the minimum excess risk achieved by $\tau_1$ is strictly larger than that achieved by $\tau_2$. 

% More precisely, when $\tau_2\geq s$, Proposition F.4 yields $p\leq u^{\prime}(\tau_2)<p+1$ for $p=$

%todo: shorten the following part or just remove
When $\tau_2 \geq s > \tau_1$, previous work implies that the optimally tuned spectral algorithm with qualification $\tau_2$ attains the minimax rate, whereas the optimally tuned algorithm with qualification $\tau_1$ does not.
\dhnote{I don't think we can say much on $s>\tau_1$, can we?}
In summary, Theorem \ref{theorem_learn_curve} explains this dichotomy through the geometry of the full learning curves: smaller qualification leads to a slower descent, a rightward shift of the optimal tuning parameter, and a higher minimum excess risk. 
This helps explain why results derived for KRR, which corresponds to $\tau=1$, need not carry over directly to spectral algorithms with larger qualification, such as KGF.
}

\fi 




\begin{comment} % [2026-04-23 remove]
\begin{remark}
The saturation effect is especially relevant for understanding lazy-training limits of neural networks. 
Indeed, the convergence of wide neural networks in the lazy regime to KGF has motivated substantial interest in large-dimensional spectral algorithms. 
However, most studies have focused on KRR (with $\tau=1$) rather than KGF (with $\tau=\infty$). 
Since KGF has higher qualification than KRR, the saturation effect shows that conclusions based on large-dimensional KRR may directly extend to large-dimensional neural networks, even in the lazy regime.
% The results of \citet{zhang2024optimal, lu2024saturation} revealed that previous results on large-dimensional KRR may not directly extend to large-dimensional neural networks, even in the lazy regime.
% Our results provide an addition geometric explanation the generalization advantage of neural networks over KRR.


%

%
% In recent years, the observation that neural networks in the ``lazy regime'' converge to the kernel gradient flow (KGF) has significantly stimulated research on large-dimensional spectral algorithms. 
% However, limited by existing tools, most studies have focused on kernel ridge regression (KRR) rather than KGF. 
% The results of \citet{zhang2024optimal, lu2024saturation} revealed that previous results on large-dimensional KRR (with \(\tau = 1\)) may not directly extend to large-dimensional neural networks, even in the lazy regime.
\end{remark}
\end{comment}

\iffalse % dh 2026-03-26: reduce the intro and put the literature in remark
\hrule 

In recent years, the observation that neural networks in the ``lazy regime'' converge to the kernel gradient flow (KGF) has significantly stimulated research on large-dimensional spectral algorithms. 
However, limited by existing tools, most studies have focused on kernel ridge regression (KRR) rather than KGF. 
Recent work by \cite{zhang2024optimal, lu2024saturation} derived the convergence rate for optimally tuned spectral algorithms and demonstrated a corresponding saturation effect: spectral algorithms with qualification \(\tau < s\) fail to attain the rate of the information-theoretic lower bound on the excess risk, whereas those with qualification \(\tau \geq s\) do. 
Consequently, they concluded that previous results on large-dimensional KRR (with \(\tau = 1\)) may not directly extend to large-dimensional neural networks, even in the lazy regime.



\begin{figure}[btp]
% \vskip 0.05in
\centering
\subfigure[]{\includegraphics[width=0.48\columnwidth]{figures/zeta_plot_for_s=2.5_gamma=6.0.pdf}}
\subfigure[]{\includegraphics[width=0.48\columnwidth]{figures/zeta_plot_for_s=3.0_gamma=3.0.pdf}}


\caption{
Another graphical representation of the learning curve of large dimensional spectral algorithms with $\tau=1$ (\textcolor{myblue}{blue}),
$\tau=2$ (\textcolor{mygreen}{green}),
and $\tau=4$ (\textcolor{myorange}{orange}) obtained in Theorem \ref{theorem_learn_curve};
and the minimax rate ({\color{red} red}) obtained in Proposition B.2. 
Dashed vertical lines indicate the optimal rate of regularization parameter $\lambda$ with qualification $\tau$, that is, the global minimum point $u^{\prime}(\tau)$ of the learning curve.
}
\label{fig_theory_saturation}
% \vskip 0.05in
\end{figure}



% {\green [ave deleted the itemize]} 
Theorem \ref{theorem_learn_curve} clarifies how the smaller $\tau_1$ yields a higher learning curve and leads to the saturation effect.
Let $\tau_2 > \tau_1 \geq 1$ and let $u^{\prime}(\tau_i)$, $i=1, 2$ denote the corresponding global minimum points of their learning curves.
First, as illustrated in Figure~\ref{fig_theory_saturation}, the learning curve decreases monotonically in the over-regularized regime and increases monotonically in the under-regularized and interpolation regimes. A smaller $\tau_1$ leads to a flatter descent in the over-regularized regime, which results in a higher learning curve.
Second, note that $u^{\prime}(\tau_1)$ lies in the region where the $\tau_2$-curve is already increasing. Therefore, if $\tau_2 \geq s > \tau_1$, the minimum excess risk for $\tau_1$ is strictly higher than that for $\tau_2$, which implies the saturation effect.
\fi 


\iffalse
\begin{itemize}
    \item As shown in Figure~\ref{fig_theory_saturation}, the learning curve decreases monotonically in the over-regularized regime and increases monotonically in the under-regularized and interpolation regimes. 
    A smaller $\tau_1$ leads to a flatter descent in the over-regularized regime, which results in a higher learning curve.

    \item Notice that $u^{\prime}(\tau_1)$ lies in the region where the $\tau_2$-curve is already increasing. Therefore, if $\tau_2 \geq s > \tau_1$, the minimum excess risk for $\tau_1$ is strictly higher than that for $\tau_2$, which implies the saturation effect.
    
\end{itemize}
\fi


\iffalse
Theorem \ref{theorem_learn_curve} further clarifies the impact of qualification \(\tau\) on the learning curve and explains the emergence of saturation effect. 
Consider $\tau_2 > \tau_1 \geq 1$, and let $u^{\prime}(\tau_i)$, $i=1, 2$ denote the corresponding global minimum points of their learning curves.
As illustrated in Figure~\ref{fig_theory_saturation}, in the over-regularized regime ($u<u^{\prime}(\tau_i)$), the learning curve decreases monotonically; while in the under-regularized / interpolation regimes ($u>u^{\prime}(\tau_i)$), the curve increases monotonically with values that are independent of $\tau_i$ (Actually, Proposition \ref{prop_e_3} implies that when $u>u^{\prime}(\tau)$, the only term in the convergence rate that depends on \(\tau\) is dominated by other three terms).
Since the smaller $\tau_1$ causes a slower decline in the over-regularized regime, we have $u^{\prime}(\tau_1)>u^{\prime}(\tau_2)$, indicating that $u^{\prime}(\tau_1)$ is in the under-regularized / interpolation regimes of the $\tau_2$'s learning curve.
Notably, when $\tau_2\geq s > \tau_1$, $u^{\prime}(\tau_1)$ falls in a regime where $\tau_2$'s learning curve has already begun to rise, ensuring that the global minimum of $\tau_1$'s learning curve is higher than the one of $\tau_2$'s.
\fi






\subsection{Equivalence with the sequence model}

Another consequence of Theorem \ref{theorem_learn_curve} is that, in the large-dimensional regime, the learning curve for kernel regression agrees with that of an associated sequence model up to constant factors. This gives a precise sense in which the sequence model serves as an effective proxy for kernel regression under sufficient regularization.

Consider the sequence model
\begin{equation}\label{eqn_sequence_model}
    z_j = f_j + \xi_j,\quad j=1, 2, \cdots,
\end{equation}
where $(z_j)_{j=1}^{\infty}$ is the observation sequence, $f_j$ is the coefficient of $f_{\star}$ with respect to $\phi_j$, that is, we have $f_j = \langle f_{\star}, \phi_j \rangle_{L^2}$, $j=1, 2, \cdots$. Moreover, $\xi_j$, $j=1, 2, \cdots$ are (not necessarily independent) random variables with mean zero and variance $\sigma^2 / n$.
Suppose $\lambda_j$'s are known and consider the parametrization $f_j = \lambda_j^{1/2} \beta_j$, $j=1, 2, \cdots$. 
As shown in Proposition I.3, 
the least squares estimator under ridge regularization is given by $\hat{f}_j^{\lambda} = \lambda_j \reg^{\krr}(\lambda_j) z_j$, $j=1, 2, \cdots$. 
Similarly, the gradient flow estimator is $\hat{f}_j^{\lambda} = \lambda_j \reg^{\gf}(\lambda_j) z_j$. 
More generally, for any filter function $\reg$ in Definition \ref{def:filter},
we define the estimator of $(f_j)_{j=1}^{\infty}$ as:
\begin{equation}\label{eqn_est_sequence_model}
    \hat{f}_j^{\lambda} := \lambda_j \reg(\lambda_j) z_j,\quad j=1, 2, \cdots
\end{equation}
whose excess risk is
\begin{equation}\label{eqn_excess_risk_sequence}
    \begin{aligned}
        E_{\xi}\left[\sum_{j=1}^{\infty}(\hat{f}_j^{\lambda} - f_j)^2\right]
        =
        \frac{\sigma^2}{n}\sum_{j =1}^\infty \left[ \lambda_j \reg(\lambda_j) \right]^2
        + \sum\limits_{j=1}^{\infty} \left( \rem(\lambda_j) f_{j}\right)^{2}.
    \end{aligned}
\end{equation}
Note that the same parametrization and excess risk $E_{\xi}\left[\sum_{j=1}^{\infty}(\hat{f}_j^{\lambda} - f_j)^2\right]$ are also considered in recent work that focuses on the concentration of the excess risks of KRR and ridge regression in sequence models (see, e.g., \cite{cheng2022dimension, misiakiewicz2024non}). 


The next proposition provides the learning curve of the estimator in \eqref{eqn_est_sequence_model} and clarifies its relationship with the learning curve for kernel regression in Theorem \ref{theorem_learn_curve}.


\begin{proposition}\label{prop_learning_curve_of_sequence}
    Adopt all the notations and assumptions in Theorem \ref{theorem_learn_curve}. 
    %Further suppose that $\lambda = \Omega_d(n^{-1})$.
    Then the excess risk of $(\hat{f}_j^{\lambda})_{j=1}^{\infty}$ satisfies
\begin{equation}\label{eqn_prop_learning_curve_of_sequence_0}
\begin{aligned}
     E_{\xi}\left[\sum_{j=1}^{\infty}(\hat{f}_j^{\lambda} - f_j)^2\right]
= 
\Theta_{d} \left( 
d^{2u - \gamma - \ell_{\lambda} - 1}
+ d^{\ell_{\lambda} - \gamma}
+ d^{-(\ell_{\lambda} + 1)  s}
+ \mathbf{1}\{\tau<\infty\} d^{-2  \tau  u + (2  \tau - \tilde{s})  \ell_{\lambda}}
\right).
\end{aligned}
\end{equation}
Consequently, Theorem \ref{theorem_learn_curve} implies
\begin{equation}\label{eqn_prop_learning_curve_of_sequence}
    E_{x, \epsilon} \left[ \left(\hat{f}_{\lambda}(x) - f_{\star}(x) \right)^{2} \right] = \Theta_{d, \mathbb{P}} \left(
E_{\xi}\left[\sum_{j=1}^{\infty}(\hat{f}_j^{\max\{\lambda, n^{-1}\}} - f_j)^2\right]
\right).
\end{equation}
\end{proposition}


Proposition \ref{prop_learning_curve_of_sequence} shows that for $\lambda=\Omega_d(n^{-1})$, the learning curve of a spectral algorithm in kernel regression has the same asymptotic order as that of the corresponding estimator in the associated sequence model.
In this regime, the sequence model captures the statistical behavior of large-dimensional spectral algorithms up to constant factors, which reinforces the idea that the sequence model serves as a theoretically tractable proxy for kernel regression. See Appendix B.3.1 for more discussions. 
The appearance of $\max\{\lambda,n^{-1}\}$ also clarifies that this equivalence is a sufficiently-regularized phenomenon rather than an exact identity for arbitrarily small $\lambda$.



\iffalse
\subsection{Experiments}



In this experiment, we numerically compute the values of three functions, denoted by $f_1$, $f_2$, and $f_3$, at a set of sample points $\{x_i\}_{i=1}^{n}$. The functions are constructed such that they belong to the interpolation spaces $[\mathcal{H}]^{0.5}$, $[\mathcal{H}]^{1}$, and $[\mathcal{H}]^{1.5}$, respectively, where $\mathcal{H}$ is a reproducing kernel Hilbert space (RKHS) associated with a spherical kernel. For computational convenience, we employ the radial basis function (RBF) kernel defined as
\[
k(x,y)=\exp\left(-\frac{\|x-y\|^2}{2\sigma^2}\right).
\]
Throughout the experiment, the following steps are performed.

\subsubsection{Kernel Matrix Construction}
Given $n$ sample points $\{x_i\}_{i=1}^{n}$ in $\mathbb{R}^d$, we construct the kernel matrix $K \in \mathbb{R}^{n \times n}$ where
\[
K_{ij}=k(x_i,x_j).
\]
This matrix is symmetric and positive semidefinite by Mercer’s theorem \cite{Mercer1909}.

\subsubsection{Nyström Approximation}
To approximate the eigen-decomposition of the kernel integral operator, we employ the Nyström method. The kernel matrix $K$ is decomposed as
\[
K = U \Sigma U^\top,
\]
where $\Sigma = \operatorname{diag}(\sigma_1,\dots,\sigma_n)$ approximates the eigenvalues $\{\lambda_j\}$ and the columns of $U$ approximate the evaluations $\{\phi_j(x_i)\}$ of the corresponding eigenfunctions at the sample points \cite{Williams2001}.

\subsubsection{Function Construction in Interpolation Spaces}
The interpolation spaces $[\mathcal{H}]^s$ are defined as
\[
[\mathcal{H}]^s := \Big\{ \sum_{j=1}^{\infty} f_j \phi_{j}(\cdot): \left(\lambda_j^{-s / 2} f_j\right)_{j} \in \ell_2 \Big\}.
\]
In order to ensure that the constructed function $f(\cdot)$ lies in $[\mathcal{H}]^s$, we choose coefficients of the form
\[
f_j = \lambda_j^{s/2}\, b_j,
\]
where $\{b_j\}$ is a square-summable sequence. In our numerical implementation, we truncate the infinite series to the first $M$ terms (with $M \le n$). For the three cases, we have:
\begin{itemize}
    \item For $s=0.5$: $f_j^{(0.5)} = \lambda_j^{0.25}\, b_j^{(0.5)}$.
    \item For $s=1$: $f_j^{(1)} = \lambda_j^{0.5}\, b_j^{(1)}$.
    \item For $s=1.5$: $f_j^{(1.5)} = \lambda_j^{0.75}\, b_j^{(1.5)}$.
\end{itemize}

\subsubsection{Computation of Function Values}
Using the approximate eigenfunctions obtained from the Nyström method, the value of the function at a sample point $x_i$ is computed as
\[
f_s(x_i) \approx \sum_{j=1}^{M} \lambda_j^{s/2}\, b_j^{(s)}\, U_{ij},
\]
or in matrix notation,
\[
f_s \approx U \, \operatorname{diag}\left(\lambda_1^{s/2}, \dots, \lambda_M^{s/2}\right) \, b^{(s)},
\]
where $s \in \{0.5,1,1.5\}$ and $b^{(s)}$ is a vector of coefficients chosen such that $\|b^{(s)}\|_2 < \infty$. This approach follows standard practices in kernel approximation and interpolation in RKHS \cite{Rasmussen2006,Scholkopf2002,Cucker2007}.

\subsubsection{Summary}
The experimental procedure involves:
\begin{enumerate}
    \item Constructing the kernel matrix from the sample points.
    \item Approximating the spectral decomposition using the Nyström method.
    \item Constructing the functions in the desired interpolation spaces by appropriately weighting the eigenvalues.
    \item Computing the function values at the sample points.
\end{enumerate}
This method provides a flexible numerical framework for controlling the smoothness of the functions via the parameter $s$.











\fi








% [too long as section title:] Learning curve and benign overfitting of large-dimensional KRR on general domains

% \section{Large-dimensional KRR on general domains}
\section{An extension for KRR on general domains}
\label{sec_general}




Section~\ref{sec_main_results} establishes the full learning curves of large-dimensional analytic spectral algorithms for inner-product kernels on $\mathbb{S}^{d-1}$, where the spherical inner-product structure yields explicit low-degree spectral block properties.
This raises the question of whether these low-degree spectral block properties are the key ingredient behind the large-dimensional learning curve beyond the spherical setting.


As the first attempt, we study this question for KRR, which is the main object in the existing large-dimensional kernel regression literature. 
We consider kernels on general domains in $\mathbb{R}^d$ under structural assumptions that match the low-degree spectral scaling derived in the spherical case, together with a hyper-contractivity condition on the associated eigenfunctions.
Under these assumptions, Theorem~\ref{theorem_learn_curve_general} establishes the learning curve of large-dimensional KRR on general domains. 
The result has the same form as in the spherical KRR case, and therefore the characterization of benign overfitting in Theorem~\ref{coroll_benign_overfit} continues to hold within this class of kernels.












\iffalse
\hrule 
In this section, we will extend our main results in Theorem \ref{theorem_learn_curve} and determine the learning curve of large-dimensional kernel ridge regression with kernels defined on general domains.


Recall that, most existing results for the learning curve of large-dimensional spectral algorithms focus on the kernel ridge regression method with inner-product kernels defined on the hyper-sphere $\mathbb{S}^{d-1}$. See Appendix \ref{sec_related_work} for details. 

\fi 


\iffalse [DH: Relocated to sec 2]
The widely-adopted assumption of spherical distribution is reasonable in many practical scenarios. 
In high-dimensional spaces, normalized data vectors often concentrate around a thin spherical shell due to the concentration of measure phenomenon \citep{talagrand1996new}. 
As a result, when the inputs are standardized to have unit norm, their empirical distribution tends to approximate the uniform distribution on the sphere. 
This makes the spherical model a convenient and accurate abstraction for analyzing the asymptotic behavior of spectral algorithms.
\fi 

\iffalse [Suppressed]
Meanwhile, we believe that it remains important to investigate more general kernels and more general input distributions, since it not only broadens the applicability of our theoretical results but also helps to identify whether the new phenomena of large-dimensional kernel regression are merely bought by the nice properties associated with inner-product kernels defined on $\mathbb{S}^{d-1}$.
\fi

\subsection{Notation and assumptions for general domains}


We introduce the settings for this section, which parallel those in Section \ref{sec_settings}. 
The main components are as follows: 
Let $\mathcal{X} \subset \bbR^{d}$ be a measurable space, and let $\rho_{\mathcal{X}}$ be a probability measure on $\mathcal{X}$.
Let $\tilde{\mathscr{K}}$ be a continuous and positive-definite kernel on $\mathcal{X} \times \mathcal{X}$ with respect to $\rho_{\mathcal{X}}$, and it admits a Mercer decomposition
\begin{equation}\label{eqn_mercer_general}
    \tilde{\mathscr{K}}(x,x^\prime) = \sum_{k=0}^{\infty} \sum_{j=1}^{\tilde{N}(d, k)} \tilde{\lambda}_{k, j}  \tilde{\psi}_{k, j}(x) \tilde{\psi}_{k, j}\left(x^\prime\right),
\end{equation}
where the eigenvalues $\{\tilde{\lambda}_{0, 1}, \cdots, \tilde{\lambda}_{0, \tilde{N}(d, 0)}, \tilde{\lambda}_{1, 1}, \cdots, \tilde{\lambda}_{1, \tilde{N}(d, 1)}, \cdots\}$ form a non-increasing sequence,
% \dhnote{Note that this bounded sum condition is not enough! We need $\sup_{d}[...]<\infty$ (constant upper bound)}
and the corresponding eigenfunctions of $\tilde{\lambda}_{k, j}$ are $\tilde{\psi}_{k, j}$ with $k=0, 1, \cdots$ and $j=1, \cdots, \tilde{N}(d, k)$. 
Furthermore, we denote $\tilde{\calH}$ as the RKHS associated with the kernel $\tilde{\mathscr{K}}$, and $[\tilde{\calH}]^{s}$ as the corresponding interpolation space. Finally, we denote the regression function as
    $$
        f_{\star}(x) = \sum_{k=0}^{\infty} \sum_{j=1}^{\tilde{N}(d,k)} \theta_{k,j} \tilde{\psi}_{k,j}(x).
    $$

    
\iffalse
\begin{itemize}
    \item Let $\mathcal{X} \subset \bbR^{d}$ be a measurable space, and let $\rho_{\mathcal{X}}$ be a probability measure on $\mathcal{X}$.

    \item Let $\tilde{\mathscr{K}}$ be a continuous and positive-definite kernel on $\mathcal{X} \times \mathcal{X}$ with respect to $\rho_{\mathcal{X}}$, and it admits a Mercer decomposition
\begin{equation*}
    \tilde{\mathscr{K}}(x,x^\prime) = \sum_{k=0}^{\infty} \sum_{j=1}^{\tilde{N}(d, k)} \tilde{\lambda}_{k, j}  \tilde{\psi}_{k, j}(x) \tilde{\psi}_{k, j}\left(x^\prime\right),
\end{equation*}
where the eigenvalues $\{\tilde{\lambda}_{0, 1}, \cdots, \tilde{\lambda}_{0, \tilde{N}(d, 0)}, \tilde{\lambda}_{1, 1}, \cdots, \tilde{\lambda}_{1, \tilde{N}(d, 1)}, \cdots\}$ form a non-increasing sequence with bounded sum 
$\sum_{k=0}^{\infty} \sum_{j=1}^{\tilde{N}(d, k)} \tilde{\lambda}_{k, j}\leq \tilde{\mathfrak{C}}$, {\green where $\tilde{\mathfrak{C}}<\infty$ is an absolute constant,}
\dhnote{Note that this bounded sum condition is not enough! We need $\sup_{d}[...]<\infty$ (constant upper bound)}
and the corresponding eigenfunctions of $\tilde{\lambda}_{k, j}$ is $\tilde{\psi}_{k, j}$ with $k=0, 1, \cdots$ and $j=1, \cdots, \tilde{N}(d, k)$.

    \item We denote $\tilde{\calH}$ as the RKHS associated with $\tilde{\mathscr{K}}$, and denote the corresponding interpolation space as $[\tilde{\calH}]^{s}$.

    \item We denote the regression function as
    $$
        f_{\star}(x) = \sum_{k=0}^{\infty} \sum_{j=1}^{\tilde{N}(d,k)} \theta_{k,j} \tilde{\psi}_{k,j}(x).
    $$
    
    
\end{itemize}
\fi

Next, we replace certain assumptions in Section \ref{sec_settings} with the following ones.


\begin{assumption}\label{assump_asymptotic_general}
    We assume that there exist positive absolute constants $c_1$, $c_2$, and a non-integer $\gamma>0$, such that the sample size satisfies
% \begin{align*}
    %n =\alpha d^{\gamma}(1+o_d(1)),
    $c_{1} d^{\gamma} \leq n \leq c_{2} d^{\gamma}$,
% \end{align*}
and we denote $\ell_{\gamma} := \lfloor \gamma\rfloor < \gamma$.
\end{assumption}


In the spherical setting, Proposition \ref{prop:inner_edr} identifies the low-degree spectral-block scaling of inner-product kernels on $\mathbb{S}^{d-1}$.
For KRR on general domains, the next assumption abstracts the corresponding spectral scaling needed in our learning-curve analysis.


\begin{assumption}\label{assump_eigenval_general}
    We assume that the sum of the eigenvalues is bounded, that is, there exists an absolute constant $\tilde{\mathfrak{C}}$ such that $\sum_{k=0}^{\infty} \sum_{j=1}^{\tilde{N}(d, k)} \tilde{\lambda}_{k, j}\leq \tilde{\mathfrak{C}}$. Moreover, for any $k  = 0, 1, \cdots, 2\ell_{\gamma} +2$, we assume that
    $$\tilde{N}(d, k) = \Theta_d(d^{k}) \quad \text{ and } \quad 
        \tilde{\lambda}_{k, j} = \Theta_d(d^{-k}), \quad \text{ for } 1\leq j \leq \tilde{N}(d, k).
    $$
\end{assumption}


Assumption \ref{assump_eigenval_general} holds for several types of common kernels in large dimensions. For example:
\begin{itemize}
    \item Under Assumption \ref{assu:coef_of_inner_prod_kernel}, inner product kernels defined on $\mathbb{S}^{d-1}$ with uniform distribution (see, e.g., Proposition \ref{prop:inner_edr});

    \item Under Assumption \ref{assu:coef_of_inner_prod_kernel}, random feature kernels defined on the hypercube $\{-1, 1\}^{d}$ with uniform distribution (see, e.g., Appendix D in \cite{mei2022generalization});

    \item Gaussian kernels defined on $\mathbb{R}^d$ with standard Gaussian distribution (see, e.g., \cite{rasmussen2003gaussian});

    \item Product Laguerre kernels defined on $[0, \infty)^{d}$ with the product of independent Gamma distributions, see Appendix I.4.1 for a detailed definition.
    %(we will provide a discussion in Appendix \ref{appendix_Laguerre_polynomials}).
    
\end{itemize}

We provide a detailed
verification in Appendix I.4.
Meanwhile, to the best of our knowledge, we do not find any other form of eigen-decay patterns within the class of kernels currently understood in large dimensions. 





The third assumption is known as the hyper-contractivity assumption.


\begin{assumption}\label{assump_hyper_contractivity}
    For any integer $q \geq 1$, there exists a constant $C=C(q, \ell_{\gamma})$ such that
$$
\sup_{k \leq 2\ell_{\gamma}+1} \sup_{j \leq \tilde{N}(d, k)} \|\tilde{\psi}_{k, j}\|_{L^{2q}} \leq C(q, \ell_{\gamma}).
$$
\end{assumption}

Assumption \ref{assump_hyper_contractivity} is originally adopted in \cite{mei2022generalization, misiakiewicz2024non} to replace the addition formula property holding for spherical harmonics (see, e.g., Proposition 1.18 in \cite{gallier2009notes}). 
Assumption \ref{assump_hyper_contractivity} holds for several types of common kernels in large dimensions, such as (i) inner-product kernels defined on $\mathbb{S}^{d-1}$ with uniform distribution and (ii) random feature kernels defined on the hypercube $\{-1, 1\}^{d}$ with uniform distribution (see, e.g., Appendix E in \cite{mei2022generalization}). 
These two examples also satisfy Assumption 8, and hence provide concrete settings covered by Theorem~\ref{theorem_learn_curve_general}.



\begin{remark}
    \cite{misiakiewicz2024non} also considers a hyper-contractivity assumption \textit{on the unknown true regression function $f_{\star}$}, which posits that $\left\|f_{\star}\right\|_{L^q} \leq (C q)^{(\ell_{\gamma}+1) / 2}\left\|f_{\star}\right\|_{L^2}$ for all integers $q \geq 2$. In Appendix I.5, we use a counterexample to demonstrate that the hyper-contractivity of eigen-functions does not imply the hyper-contractivity of the unknown true regression function.
    Therefore, the hyper-contractivity assumption on $f_{\star}$ is not weaker than Assumption \ref{assump_hyper_contractivity}.
\end{remark}









The next two assumptions are for the regression functions. 
Similar to Assumptions \ref{assump_function_calss} and \ref{assump_source condition_lower_bound}, they imply that $f_{\star}$ is exactly in $[\tilde{\calH}]^{s}$.

\begin{assumption}\label{assump_function_calss_general}
There exist absolute constants $s \geq 0$ and $R>0$, such that we have
    \begin{equation*}
    f_{\star} \in 
    \left\{f \in [\tilde{\calH}]^{s} \mid \|f\|_{[\tilde{\calH}]^{s}} \leq \sqrt{R}\right\}.
\end{equation*}
\end{assumption}

\begin{assumption}\label{assump_source condition_lower_bound_general}
Suppose there exists an absolute constant $c_{0} > 0$ such that for any $ d \geq \mathfrak{C}$, we have
\begin{equation*}
        \sum\limits_{j=1}^{\tilde{N}(d, m)} \tilde{\lambda}_{m, j}^{-s} \theta_{m,j}^{2} \ge c_{0}, ~~m=\tilde{\ell}, \tilde{\ell}+1;
   \end{equation*}
Moreover, if $s> 2$, we further assume that
$\sum\limits_{k=0}^{\tilde{\ell}} \sum\limits_{j =1}^{\tilde{N}(d, k)}  \theta_{k, j}^{2} \ge c_{0}$.
\end{assumption}






















\subsection{Learning curve of large-dimensional KRR in general domains}

The following theorem determines the learning curve of large-dimensional KRR with general kernels defined on general domains.

\begin{theorem}\label{theorem_learn_curve_general}
    Suppose Assumptions \ref{assume_lambda}, 
    \ref{assump_asymptotic_general}, 
    \ref{assump_eigenval_general}, 
    \ref{assump_hyper_contractivity},
\ref{assump_function_calss_general},
and \ref{assump_source condition_lower_bound_general} hold with $s \geq 0$ and $u \in (0, \infty]$. 
Denote $\tilde{s}=\min\{s, 2\}$. Then the excess risk of the KRR estimator in (\ref{eq bias var decomposition}) satisfies
\begin{equation*}
\begin{aligned}
    E_{x, \epsilon} \left[ \left(\hat{f}_{\lambda}(x) - f_{\star}(x) \right)^{2} \right]
= \Theta_{d, \mathbb{P}} \left( 
d^{\gamma - \tilde{\ell} - 1 - 2  \max\{\gamma - u, 0\}}
+ d^{\tilde{\ell} - \gamma}
+ d^{-(\tilde{\ell} + 1)  s}
+ d^{-2  u + (2 - \tilde{s})  \tilde{\ell}}
\right).
\end{aligned}
\end{equation*}
Consequently, results in Theorem~\ref{coroll_benign_overfit} still hold under the above assumptions and conditions.
\end{theorem}

% \dhrev{
% Theorem \ref{theorem_learn_curve_general} indicates that the new phenomena observed by prior works in large-dimensional kernel ridge regression (as discussed in Appendix B) are not limited to spherical inner-product kernels.
% Instead, we can obtain learning curves of large-dimensional KRR with kernels in general domains, as long as the spectral-scaling and hyper-contractivity assumptions hold.
% Therefore, Theorem \ref{theorem_learn_curve_general} justifies the simplified spherical kernel models adopted in prior works and in Theorem \ref{theorem_learn_curve}, and further suggests that other large-dimensional KRR results with spherical inner-product kernels may also extend to more general kernels.
% }{
% Theorem \ref{theorem_learn_curve_general} shows that the rate characterization obtained in Theorem~\ref{theorem_learn_curve} is not tied to the exact sphere kernel setting. 
% Rather, it extends to KRR on general domains for kernels satisfying Assumptions \ref{assump_eigenval_general} and \ref{assump_hyper_contractivity}.
% In this sense, the spherical analysis captures structural features that are sufficient for the same large-dimensional behavior to persist beyond the sphere.
% This supports the use of spherical kernel models as a tractable basic setting for studying large-dimensional KRR. 
% }

% \dhnote{
% The point is to avoid backing up something we don't prove ("indicates that the new phenomena observed by prior works"). 
% }

Theorem \ref{theorem_learn_curve_general} shows that the rate characterization obtained in Theorem~\ref{theorem_learn_curve} is not tied to the exact sphere kernel setting. 
Rather, it extends to KRR on general domains for kernels satisfying Assumptions \ref{assump_eigenval_general} and \ref{assump_hyper_contractivity}.
In this sense, the spherical analysis captures structural features that are sufficient for the same large-dimensional behavior to persist beyond the sphere.
This supports the use of spherical kernel models as a tractable basic setting for studying large-dimensional KRR. 



\begin{remark}\label{remark_compare_with_mei_mis}
Two existing results \citep{mei2022generalization, misiakiewicz2024non} also provide tight bounds on excess risk of large-dimensional KRR with general kernels defined on general domains. 
Specifically, \cite{mei2022generalization} showed that 
    \begin{equation*}
    \left|E_{x, \epsilon} \left[ \left(\hat{f}_{\lambda}(x) - f_{\star}(x) \right)^{2} \right]-
    \text{E}_{\text{approx}}
    \right| =o_{d, \mathbb{P}}\left(1\right),
\end{equation*}
where $\text{E}_{\text{approx}}$ is a certain deterministic quantity. However, as we have discussed in Appendix B.2, such results are uninformative when $s>0$ since the excess risk is $o_{d, \mathbb{P}}(1)$.
Later, \cite{misiakiewicz2024non} developed a tighter approximation (see, e.g., Theorem 1 and Theorem 2 of \cite{misiakiewicz2024non})
    \begin{equation*}
    E_{x, \epsilon} \left[ \left(\hat{f}_{\lambda}(x) - f_{\star}(x) \right)^{2} \right] = \text{E}_{\text{approx}} \cdot \left(1 + o_{d, \mathbb{P}}\left(1\right)\right).
\end{equation*}
    However, their assumptions are either difficult to verify (Assumption 1 of \cite{misiakiewicz2024non}), or are too restrictive when $s>0$, as we have discussed in Appendix B.3.
On the contrary, for kernel ridge regression methods, Theorem \ref{theorem_learn_curve_general} not only extends Theorem \ref{theorem_learn_curve} from kernels on the sphere to kernels on general domains. It also extends the results of \citet{mei2022generalization} and \citet{misiakiewicz2024non} by determining the convergence rate of the excess risk when $s>0$.
\end{remark}





\section{Numerical experiments}\label{sec_experiments}
We use several experiments to illustrate Theorem \ref{theorem_learn_curve}.
Consider the following two inner-product kernels:
\begin{itemize}
    \item[(i)] Neural Tangent Kernel (NTK) of a two-layer ReLU neural network:
    $$
        \mathscr{K}^{\mathrm{ntk}}(x, x^\prime) := \Phi(\langle x, x^{\prime} \rangle), ~~x, x^{\prime} \in \mathbb{S}^{d-1},
    $$
    where $\Phi(t)=\left[\sin{(\arccos t)}+2(\pi-\arccos t)t\right]/ (2 \pi)$.
    
    \item[(ii)] RBF kernel with a fixed bandwidth:
    $$
        \mathscr{K}^{\mathrm{rbf}}(x,x^{\prime}) = \exp{\left(-\frac{\|x-x^{\prime}\|_{2}^{2}}{2}\right)}, ~~x, x^{\prime} \in \mathbb{S}^{d-1}.
    $$
\end{itemize}
The RBF kernel satisfies Assumption \ref{assu:coef_of_inner_prod_kernel}. For the NTK, the coefficients of $\Phi(\cdot)$, $\{a_{j}\}_{j=0}^{\infty}$, satisfy $ a_{j} > 0, j \in \{0, 1\} \cup \{2,4,6,\ldots\}$ and $ a_{j} = 0, j \in \{3,5,7,\ldots\}$ (see, e.g., \cite{lu2023optimal}). 
As noted after Assumption \ref{assu:coef_of_inner_prod_kernel}, our results can be extended to inner product kernels with certain zero coefficients $a_j$. Specifically, for any $\gamma>0$, as long as $ a_{j} > 0$ for $j = \ell_{\gamma}, \ell_{\gamma}+1$, the proof and convergence rate remain the same. Therefore, for $\gamma<2$ in our experiments, the theoretical convergence rates for the NTK can be directly obtained from Theorem \ref{theorem_learn_curve}. 

We consider the following data generation process:
$y_{i} = f_{\star}(x_{i}) + \epsilon_{i}, ~~ i = 1, \cdots, n,$ 
where each $x_{i}$ is independently sampled from the uniform distribution on $\mathbb{S}^{d-1}$, and $\epsilon_{i} \overset{\text{i.i.d.}}{\sim} \mathcal{N}(0,1)$.
%
%
For different training sample sizes $n$,  we set the corresponding dimension as $d=\lfloor n^{1/\gamma}\rfloor$. For each kernel and dimension $d$, we consider the following regression function $f_{\star}$:
\begin{equation}\label{experiment true function}
    f_{\star}(x) = 
    \left\{\begin{matrix}
\sum_{i=1}^{3} \mathscr{K}(\xi_{i},x), & \quad \text{ if } s=1; \vspace{5pt}\\
\sum_{i=1}^{3} \sum_{k=0}^{2} d^{k(1-s)/2} P_{k, d}(\langle\xi_{i}, x\rangle), & \quad \text{ if } s \neq 1;
\end{matrix}\right.
\end{equation}
where $P_{k, d}(\cdot)$ is the $k$-th normalized Gegenbauer polynomial with parameter $(d-1)/2$ (see, e.g., Appendix L.2).
%
%
The points $\xi_{1}, \xi_{2}, \xi_{3}$ are independently sampled from the uniform distribution on $\mathbb{S}^{d-1}$, subject to the constraint 
    \begin{equation}\label{eqn_req_inner_of_xi_in_exp}
        \sum_{i \neq i^{\prime}}^{3} \left| \langle \xi_{i}, \xi_{i^{\prime}} \rangle \right| \leq 2.9,
    \end{equation}
    which excludes nearly collinear configurations. 
    Once sampled, the triplet remains fixed for each experimental run (i.e., for each combination of $n$, $d$, and $s$), but is resampled across different runs. 
One can verify that $f_{\star}$ satisfies Assumption \ref{assump_function_calss} and Assumption \ref{assump_source condition_lower_bound} (see, e.g., Appendix L.2).

    

\begin{remark}
    The condition (\ref{eqn_req_inner_of_xi_in_exp}) is very mild: the LHS is of order $O_{d, \mathbb{P}}(d^{-1})$. To see this, note that for any $1 \leq i \neq i^{\prime} \leq 3$, the inner products $\langle \xi_{i}, \xi_{i^{\prime}} \rangle$ have mean zero and variance $(d+1)^{-1}$, thus for any $\varepsilon>0$, from Chebyshev's inequality we have $\mathbb{P}(\sum_{i \neq i^{\prime}}^{3} \left| \langle \xi_{i}, \xi_{i^{\prime}} \rangle \right| \geq \varepsilon) \leq 6 \max_{1 \leq i \neq i^{\prime} \leq 3} \mathbb{P}( \left| \langle \xi_{i}, \xi_{i^{\prime}} \rangle \right| \geq \varepsilon) \leq 6/[\varepsilon^2(d+1)]$. 
\end{remark}

% {\color{red}
% I don't quite get the meaning. Do you mean:\\
% v1: 
% $\xi_{1}, \xi_{2}, \xi_{3}$ are fixed within each repetition, but may vary across repetitions.
% \\
% or 
% v2: $\xi_{1}, \xi_{2}, \xi_{3} \in \mathbb{S}^{d-1}$ are fixed points, shared across all repetitions.
% }


For each $n$, we repeat the experiments 50 times and present the average generalization error of (a) kernel gradient flow (KGF); (b) kernel ridge regression (KRR) on $1,000$ testing samples.
% \dhnote{Should we unify the name KGF? We didn't use GF elsewhere.}
We fit a least squares line on the log scale to the excess risk as a function of the dimension, which yields an estimate of the convergence rate (i.e., the slope reported in the figure legends).

% We used logarithmic least squares to fit the excess risk with respect to the sample size, resulting in an estimate of the convergence rate (i.e., the slope in the legends of the figures). 

\begin{remark}
    We select KGF and KRR for experiments due to their dual significance. 
    On one hand, they are central and widely studied instances of spectral algorithms for kernel regression.
    On the other hand, the class of analytic spectral algorithms is characterized by a qualification parameter $\tau \in [1, \infty]$, with KRR and KGF serving as canonical representatives of the two ends, $\tau=1$ and $\tau =\infty$. 
    Hence, their agreement with the predicted slopes shows that the large-dimensional asymptotic picture in Theorem~\ref{theorem_learn_curve} is already visible at finite sample sizes for the two canonical endpoint filters, and provides empirical support for the broader theory of analytic spectral algorithms.
\end{remark}








\noindent {\bf Type 1 experiments:} We consider four parameter settings of  $(\gamma, s, u)$ as follows:
\begin{itemize}
    \item $(\gamma, s, u) =(1.5, 1.5, 0.5)$, $n$ from $2600$ to $5000$, with intervals $200$, $d = \lfloor n^{2/3} \rfloor$.

    \item $(\gamma, s, u) =(0.8, 1.0, 2.0)$, $n$ from $1000$ to $3000$, with intervals $100$, $d = \lfloor n^{5/4} \rfloor$.

    \item $(\gamma, s, u) =(1.0, 1.0, 2.0)$, $n$ from $1000$ to $3000$, with intervals $100$, $d = n$.

    \item $(\gamma, s, u) =(1.2, 2.0, 1.0)$, $n$ from $1600$ to $4000$, with intervals $150$, $d = \lfloor n^{5/6} \rfloor$.
    
\end{itemize}
We set $\lambda = C_{\lambda} d^{-u}$ with $C_{\lambda} \in \{0.001, 0.01, 0.1, 0.2, 0.4, 0.5, 0.7, 0.8, 1, 10\}$.
For each method and each setting, we report the results obtained with the best performing choice of $C_{\lambda}$ that consistently has the smallest generalization error, and separately report the results for all tested values of $C_{\lambda}$. 
% For each pattern of $(\gamma, s, u)$, we first report the results with the best choice of $C_{\lambda}$ that has the smallest average generalization error, and then report results with all choices of $C_{\lambda}$.
As shown in 
Figure \ref{experiment_1_settings_1} and 
Figure \ref{experiment_1_settings_2}, as well as Figures 7 and 8 in Appendix L.1,
the empirical convergence rates agree with the theoretical rates predicted by Theorem \ref{theorem_learn_curve}.










\begin{figure}
\centering
\subfigure{\includegraphics[width=0.45\columnwidth]{figures_experiments_1/plot_setting_1_best_C_NTK.pdf}}
\subfigure{\includegraphics[width=0.45\columnwidth]{figures_experiments_1/plot_setting_1_best_C_RBF.pdf}}
\caption{
Type 1 experiments with parameters $(\gamma, s, u) = (1.5, 1.5, 0.5)$ and regularization $\lambda = C_{\lambda} d^{-u}$.
The left panel uses the NTK kernel, and the right panel uses the RBF kernel.
In each panel, we plot $\ln(\text{Excess risk})$ versus $\ln(d)$ for KRR ($\tau=1$) and KGF ($\tau=\infty$), with each method using its optimal $C_{\lambda}$ value (selected separately).
Dashed lines show the least-squares fits. 
The legend reports the fitted slopes (with their 68\% confidence intervals), and the theoretical predictions from Theorem \ref{theorem_learn_curve}.
Results for additional $C_{\lambda}$ values are provided in Figure 9 in Appendix L.1.
}
\label{experiment_1_settings_1}
\end{figure}
\begin{figure} 
\centering
\subfigure{\includegraphics[width=0.45\columnwidth]{figures_experiments_1/plot_setting_2_best_C_NTK.pdf}}
\subfigure{\includegraphics[width=0.45\columnwidth]{figures_experiments_1/plot_setting_2_best_C_RBF.pdf}}
\caption{Type 1 experiments with $(\gamma, s, u) = (0.8, 1.0, 2.0)$. The setup and analysis are the same as in Figure \ref{experiment_1_settings_1}. Results for additional $C_{\lambda}$ values are shown in Figure 10 in Appendix L.1.}
\label{experiment_1_settings_2}
\end{figure}



\noindent {\bf Type 2 experiments:} 
In the second type of experiments, we fix $\gamma = 1.5$ and consider $s \in \{1, 2\}$. 
We vary $n$ from $3000$ to $5000$ (with gaps $200$) and set $d = \lfloor n^{2/3} \rfloor$. 
% \dhnote{
% I am confused; in the plots, there is no $n$ or $d$. Are they fixed?
% }
% {\green 
For any given $(\gamma, s, u)$, similar to Type 1 experiments (see, e.g., Figure \ref{experiment_1_settings_1}), we fit a least squares line on the log scale to the excess risk as a function of the dimension, which yields an estimate of the convergence rate.
We then report the estimated convergence rates with $u \in [0.1, 2.5]$, and compare them with the theoretical convergence rates given in Theorem \ref{theorem_learn_curve}. 
As shown in Figure \ref{experiment_2_s_1} and Figure \ref{experiment_2_s_2}, the estimated convergence rates align well with the theoretical convergence rates.
% }




\begin{figure}[htb]
\centering
\subfigure{\includegraphics[width=0.45\columnwidth]{figures_experiments_2/plot_exp2_slopes_s_1.0_NTK.pdf}}
\subfigure{\includegraphics[width=0.45\columnwidth]{figures_experiments_2/plot_exp2_slopes_s_1.0_RBF.pdf}}

\caption{
Comparison of the experimental and theoretical convergence rates for Type 2 experiments with $(\gamma, s) = (1.5, 1)$. 
The left panel uses the NTK kernel, and the right panel uses the RBF kernel.
For each value of $u$, the experimental slope (circles with 68\% confidence intervals shown as error bars) is obtained by fitting a least squares line to $\ln(\text{Excess risk})$ against $\ln(d)$;
the theoretical slope (dashed lines with triangles) is predicted by Theorem \ref{theorem_learn_curve} (using $\tau=1$ for KRR and $\tau=\infty$ for KGF).
%The inset at the lower right reports the mean absolute deviation between the experimental and theoretical slopes over $u$. 
}
\label{experiment_2_s_1}
\end{figure}
\begin{figure}
\centering
\subfigure{\includegraphics[width=0.45\columnwidth]{figures_experiments_2/plot_exp2_slopes_s_2.0_NTK.pdf}}
\subfigure{\includegraphics[width=0.45\columnwidth]{figures_experiments_2/plot_exp2_slopes_s_2.0_RBF.pdf}}

\caption{
Type 2 experiments with $(\gamma, s) = (1.5, 2)$. The setup and analysis are the same as in Figure \ref{experiment_2_s_1}. 
}
\label{experiment_2_s_2}
\end{figure}


% \input{sections/conclusion}


\clearpage


\section*{Acknowledgments}
D.H. is the corresponding author. 
D.H. is supported by Singapore MOE AcRF Tier-1 Grant A-8004149-00-00.
Q.L. is supported by the National Natural Science Foundation of China (9237012, 11971257) and Tsinghua University Dushi Program (20251080059). 
Y.X. is supported by the SUSTech-NUS Joint Research Program and the National Natural Science Foundation of China (12571312). 
AI-assisted tools were used only for language polishing. 





%Bibliography
\bibliographystyle{chicago}
\bibliography{reference.bib} 

\clearpage

\appendix



\begin{appendix}

The appendix is organized as follows: 
Appendix \ref{append_notation} introduces all notations used throughout this Appendix; 
Appendix \ref{sec_related_work} provides a detailed comparison between our results and the existing literature;
Appendix \ref{append_sketch_proof_main_results} provides a sketch of proof of Theorem 3.1;
Appendix \ref{append_variance}, \ref{append_bias}, and \ref{append_learning_curve} are the proof of Theorem 3.1;
Appendix \ref{append_kernel_app_sphere} restates and develops several approximation results for kernels defined on the sphere;
Appendix \ref{append_kernel_app_general} generalizes results in Appendix \ref{append_kernel_app_sphere} to kernels on general domains and prove Theorem 4.2;
Appendix \ref{appendix_proof_other_results} proves other results in the main text;
Appendix \ref{appendix_filter_property} discusses properties of the filter function defined in Definition 2.2; 
Appendix \ref{appendix_auxiliary} lists some auxiliary lemmas;
and Appendix \ref{appendix_experiments} provides further experiments and supporting proofs for Section 5.










\section{Notations}\label{append_notation}


In this section, we first specify all the absolute constants on which our asymptotic notations (e.g., $O_{d}()$ and $O_{d, \mathbb{P}}()$) depend. 
That is, any implicit constants appearing in our results depend only on the absolute constants listed below. 
Next, we introduce the notations used throughout this Appendix, most of which are adopted from \cite{ghorbani2021linearized, zhang2024phase}.


%All other constants this paper defines depend on the following absolute constants.


\begin{definition}\label{def:abs_constants}

    We list all the absolute constants used in this paper:
    
    \begin{itemize}

    \item \(\sigma^2\): Upper bound on variance of the noise in (1).

    \vspace{3pt}
    
    \item \(c_1, c_2, \gamma > 0\): Positive constants in the asymptotic framework (2). 

    \vspace{3pt}

    \item \(a_{-1}; a_0, a_1, \ldots, a_{\lfloor \gamma \rfloor + 2} > 0\) : The sum and the first $(\lfloor \gamma \rfloor+ 3)$ coefficients of the Taylor expansion of \(\Phi(\cdot)\) as specified in Assumption 2.

    \vspace{3pt}

    \item \(s \geq 0\), $R>0$: Coefficients in the source condition  (6).

    \vspace{3pt}

    \item $c_{0} >0$: A constant given in Assumption 6.

    \vspace{3pt}
    
    \item $C_i > 0$, $i=1, \cdots, 4$: The constants related to the filter function defined in Definition 2.2.

    \vspace{3pt}

    \item $\tau^{\prime}$: A finite constant defined in (\ref{eqn_def_of_tau_prime}).

    \vspace{3pt}

    \item $\Delta u :=(u - \ell_{\lambda})\mathbf{1}\{u<\infty\} \in [0, 1)$: $u$ is the decay rate of the regularization parameter $\lambda$ in Assumption 4, and $\ell_{\lambda} = \lfloor u \rfloor$.

    \vspace{3pt}

    \item \(c_3, c_4 > 0\): Positive constants in Assumption 4. 

    \vspace{3pt}

    \item $\tilde{\ell}=\min\{\ell_{\gamma}, \ell_{\lambda}\}$ defined in Assumption 4, where $\ell_{\gamma} = \lfloor \gamma \rfloor$.
    
    \end{itemize}
\end{definition}






%The following notations are mainly borrowed from \cite{ghorbani2021linearized, zhang2024phase}.


% We always assume that $d \geq \mathfrak{C}$, where $\mathfrak{C}$ is a constant given in Proposition 2.1.




\paragraph*{Notations on eigenfunctions and spherical harmonics}
Denote $\psi(x) = (\phi_j(x))^{\top}_{j \ge 1} \in \mathbb{R}^{\infty}$, where $\phi_j(\cdot)$ are the eigenfunctions given in (4). For any $k =0, 1, \cdots$ and $j \leq N(d, k)$, denote $N_{k} = \sum_{k^{\prime}=0}^{k}N(d, k^{\prime})$, $ \Psi_{k, j} = \psi_{k,j}(X) = \left(\psi_{k,j}(x_{1}), \cdots, \psi_{k,j}(x_{n}) \right)^{\top} \in \mathbb{R}^{n}$, and
\begin{displaymath}
    \Psi_{k} = \left(\Psi_{k,1}, \cdots, \Psi_{k,N(d,k)} \right) \in \mathbb{R}^{n \times N(d,k)}.
\end{displaymath}
Denote $\Psi =  \left(\psi(x_{1}), \cdots \psi(x_{n}) \right)^{\top} \in \mathbb{R}^{n \times \infty}$. For $k =0, 1, \cdots$, denote 
$$
\Psi_{\leq k} = \left(\Psi_{0}, \cdots, \Psi_{k} \right) \in \mathbb{R}^{n \times N_{k}},
$$
then Proposition 2.1 and the Funk-Hecke formula in (5) together imply that for any $k \leq \ell_{\gamma}+1$, we have $\Psi = \left( \Psi_{\leq k}, \Psi_{>k} \right)$, where $\Psi_{>k} = \left(\Psi_{j_1}, \Psi_{j_2}, \cdots \right)$ with $\{j_1, j_2, \cdots\} \subset \{k+1, k+2, \cdots\}$.\\
When $\ell_{\lambda}<\ell_{\gamma}$, we further denote $\Psi_{\ell_{\lambda}+1, \ell_{\gamma}} = \left(\Psi_{\ell_{\lambda}+1}, \cdots, \Psi_{\ell_{\gamma}} \right)$.


\paragraph*{Notations on eigenvalues}
For any $k = 0, 1, \cdots$, define the diagonal matrix
\begin{displaymath}
    \Sigma_{k} = \mu_{k} \mathrm{I}_{N(d,k)} \in \mathbb{R}^{N(d,k) \times N(d,k)}.
\end{displaymath}

% \dhnote{Excluding $\Sigma_0$??}

Denote $\Sigma = \text{Diag}\{ \Sigma_0, \Sigma_1, \Sigma_2, \cdots\}$. For $k = 0, 1, \cdots$, denote 
$$
\Sigma_{\leq k} = \text{Diag} \left\{\Sigma_{0}, \cdots, \Sigma_{k} \right\} \in \mathbb{R}^{N_{k} \times N_{k}}.
$$
Then Proposition 2.1 and the Funk-Hecke formula in (5) imply that for any $k \leq \ell_{\gamma}+1$, we have $\Sigma = \text{Diag}  \left\{ \Sigma_{\leq k}, \Sigma_{>k} \right\}$, where 
$\Sigma_{>k} = \text{Diag} \left\{\Sigma_{j_1}, \Sigma_{j_2}, \cdots \right\}$ with $\{j_1, j_2, \cdots\} \subset \{k+1, k+2, \cdots\}$.\\
When $\ell_{\lambda}<\ell_{\gamma}$, we further denote $\Sigma_{\ell_{\lambda}+1, \ell_{\gamma}} = \text{Diag} \left\{\Sigma_{\ell_{\lambda}+1}, \cdots, \Sigma_{\ell_{\gamma}} \right\}$.




\paragraph*{Notations on the regression function}
Recall that in (7) we define
\begin{equation*}
    f_{\star}(x) =\theta^{\top}  \psi(x)=\sum_{k=0}^{\ell_{\gamma}+1} \sum_{j=1}^{N(d,k)} \theta_{k,j} \psi_{k,j}(x) + \sum_{j =N_{\ell_{\gamma}+1}+1}^{\infty}f_j \phi_{j}(x),
\end{equation*}
where $\theta = \left( \theta_{0,1}, \theta_{1,1}, \cdots, \theta_{\ell_{\gamma}+1, N(d, \ell_{\gamma}+1)}, f_{N_{\ell_{\gamma}+1}+1}, f_{N_{\ell_{\gamma}+1}+2}, \cdots\right)^{\top}$.\\
For any $k =0, 1, \cdots, \ell_{\gamma}+1$, denote 
\begin{displaymath}
    \theta_{k} = \left( \theta_{k,1}, \cdots, \theta_{k,N(d,k)} \right)^{\top} \in \mathbb{R}^{N(d,k)},
\end{displaymath}
then we can denote $ \theta = \left( \theta^{ \top}_{\leq \tilde{\ell}}, \theta^{\top}_{>\tilde{\ell}} \right)^{\top}$, where
\begin{displaymath}
    \theta_{\leq \tilde{\ell}} = \left( \theta_{0}^{ \top}, \cdots, \theta_{\tilde{\ell}}^{\top} \right)^{\top} \in \mathbb{R}^{N_{\tilde{\ell}}}.
\end{displaymath}



\paragraph*{Notations on the kernel matrix}
Denote the kernel matrix, which is also called the gram matrix, by $K := \mathscr{K}(X, X) \in \mathbb{R}^{n \times n}$, where for $i,j \leq n$, its $(i, j)$-th entry is $(K)_{ij} = \mathscr{K}(x_{i},x_{j})$. From Lemma \ref{lemma lambda max min}, we have $\lambda_{\text{min}}\left( K\right) = \Omega_{d, \mathbb{P}}(1)$, hence the event $E_0 := \{K \text{ is an invertible matrix }\}$ occurs with probability $1-o_{d, \mathbb{P}}(1)$. In the remainder of this paper, our proofs are conditioned on the event $E_0$.\\
For $k \in \{\ell_{\gamma}, \tilde{\ell}\}$, define
$$
K_{\leq k} := \Psi_{\leq k} \Sigma_{\leq k} \Psi_{\leq k}^{\top}, 
\qquad 
K_{>k} := \Psi_{>k} \Sigma_{>k} \Psi_{>k}^{\top}. 
$$
Specifically, since $\psi_{0, 1}(x) \equiv 1$, we have $K_{\leq 0} = \mu_0 \Psi_{0, 1} \Psi_{0, 1}^{\top} = \mu_0 \mathbf{1} \mathbf{1}^{\top} = K_0$, where $K_{0}$ is defined in Assumption 5.

Define $M_{>k} := K_{>k} + n\lambda \mathrm{I}_n$, and
$$
M := K + n\lambda \mathrm{I}_n.
$$
When $\ell_{\lambda}<\ell_{\gamma}$, we further denote $ K_{\ell_{\lambda}+1, \ell_{\gamma}} := \Psi_{\ell_{\lambda}+1, \ell_{\gamma}} \Sigma_{\ell_{\lambda}+1, \ell_{\gamma}} \Psi_{\ell_{\lambda}+1, \ell_{\gamma}}^{\top}$.




\paragraph*{Discussion about the filter function} For notational simplicity, we denote $\varphi_{\lambda, K}:=n^{-1}\reg(K/n) \in \mathbb{R}^{n \times n}$. 
For any orthogonal matrix $B\in \mathbb{R}^{n\times n}$ and any diagonal matrix $S\in \mathbb{R}^{n\times n}$, from (10) we have
\begin{align}
    \label{eq:Filter_Rem_finite_case3}
    \reg(BSB^{\top})=B\reg(S)B^{\top}.
    \end{align}
Therefore, $\varphi_{\lambda, K}$ is a symmetric matrix, and from Definition 2.2 we have
\begin{equation}\label{eqn:spe_properties}
    \begin{aligned}
         &C_1 M^{-1} =   \frac{C_1}{n}\left(\frac{K}{n}+\lambda \mathrm{I}_n \right)^{-1} \leq \varphi_{\lambda, K} \leq  C_2 M^{-1}.
         \\
        % &\left[\mathrm{I}_{n} -  K \varphi_{\lambda, K}  \right] \leq C_3 \lambda^{\tau^{\prime}}\left(\frac{K}{n}\right)^{-\tau^{\prime}}.
    \end{aligned}
\end{equation}








\section{Discussion on related works}\label{sec_related_work}

In this section, we provide a detailed comparison between our results and the existing literature to better illustrate the novelty of our results and proofs.

\subsection{Learning curves of fixed-dimensional spectral algorithms}\label{subsec_learning_curve_fix_d}

The learning curves of fixed-dimensional spectral algorithms became popular and have been well-studied within recent few years.
Several early works attempted to describe the learning curve of KRR using heuristic arguments \cite{Bordelon_Spectrum_2020, Cui2021GeneralizationER, mallinar2022benign}. However, their results relied on the unrealistic Gaussian design assumption, that is, they assumed that the eigenfunctions $\{\phi_j(x)\}_{j=0}^{\infty}$ in (4) are normal random variables.


% {\green [Have deleted the itemize]} 
To provide a rigorous proof of the learning curve of KRR, \cite{li2023asymptotic} rewrote the variance (14) and the bias (15) into integral operator forms. Specifically, recall the definition of $T$ in (3) and define its empirical version $T_{X}$ as follows: For any $y \in \mathbb{R}$, let $K_{x} : \R \to \calH$ be given by $K_{x}(y) = y \cdot \mathscr{K}(x,\cdot)$, whose adjoint $K_{x}^* : \calH \to \R$ is given by $K_{x}^*(f) = \ang{\mathscr{K}(x,\cdot),f}_{\calH} = f(x)$.
Moreover, we denote by $T_{x} = K_{x} K_{x}^*$ and $T_{X} = \frac{1}{n}\sum_{i=1}^n T_{x_i}$.
\cite{li2023asymptotic} then proved the following approximation on variance \& bias:

\noindent {\bf Variance term: }
\begin{equation}\label{eqn_approx_variance_fixed_d_learning_curve}
    \begin{aligned}
    \mathrm{var}(\hat{f}_{\lambda})
    =&~ \frac{\sigma^2}{n} \int_{\mathcal{X}} \left\|\reg^{\krr} \left(T_{X}\right) \mathscr{K}(x,\cdot)\right\|_{L^2,n}^2 \mathrm{d} \rho_{\calX}(x)\\
    \stackrel{(\mathbf{A})}{\approx} &~
    \frac{\sigma^2}{n}\int_{\mathcal{X}} \left\|{\reg^{\krr} \left(T\right) \mathscr{K}(x,\cdot)}\right\|_{L^2, n}^{2} \mathrm{d}\rho_{\calX}(x)
    \stackrel{\mathbf{B}}{\approx}
    \frac{\sigma^2}{n}\int_{\mathcal{X}} \left\|{\reg^{\krr} \left(T\right) \mathscr{K}(x,\cdot)}\right\|_{L^2}^{2} \mathrm{d}\rho_{\calX}(x);
    \end{aligned}
\end{equation}
where $(\mathbf{A})$ is obtained via the key observation that 
\begin{equation}\label{eqn_key_observation_in_approx_A_fixed_d_learning_curve}
    \reg^{\krr}(T_{X})-\reg^{\krr}(T) = \reg^{\krr}(T_{X})(T-T_{X})\reg^{\krr}(T),
\end{equation}
and concentration of $T-T_{X}$, and $\mathbf{B}$ is obtained by a standard concentration of the $L^2$-norm (see, e.g., Proposition C.9 in \cite{li2022saturation}).


\noindent  {\bf Bias term: } 
Denote $g(\cdot) := \int_{\mathcal{X}} \mathscr{K}(\cdot,x) f_{\star}(x) ~\mathsf{d} \rho_{\calX}(x)$ and ${g}_{X}(\cdot) :=n^{-1} \sum_{i=1}^n K_{x_i} f_{\star}(x_i)$. Furthermore, denote $f_{\lambda} = \reg^{\krr}(T) g$ and notice that we have $E_\epsilon(\hat{f}_{\lambda}) = \reg^{\krr}(T_{X}) {g}_{X}$.
\iffalse
that
$$
    \left\|\tilde{f}_\lambda - f_{\lambda}\right\|_{L_{2}} \leq \left\| T^{\frac{1}{2}} \reg^{\krr}(T)^{\frac{1}{2}}\right\| \cdot \left\| \reg^{\krr}(T)^{-\frac{1}{2}} \reg^{\krr}(T_{X}) \reg^{\krr}(T)^{-\frac{1}{2}} \right\| \cdot \left\| \reg^{\krr}(T)^{\frac{1}{2}} \left( \tilde{g} - \reg^{\krr}(T_{X})^{-1} f_{\lambda} \right) \right\|_{\mathcal{H}},
$$
where the last term can be decomposed as 
$$
\reg^{\krr}(T)^{\frac{1}{2}} \left[ \tilde{g} - \reg^{\krr}(T_{X})^{-1} f_{\lambda} \right]  = \reg^{\krr}(T)^{\frac{1}{2}}\left[\left(\tilde{g}-T_{X} f_\lambda\right)-\left(g-T f_\lambda\right)\right]
$$
\fi
They established the following concentration:
\begin{equation}\label{eqn_approx_bias_fixed_d_learning_curve}
    \|E_\epsilon(\hat{f}_{\lambda}) - {f}_{\lambda}\|_{L^2} := \|\reg^{\krr}(T_{X}) {g}_{X} - \reg^{\krr}(T) g\|_{L^2} = o_{\mathbb{P}}\left(\| {f}_{\lambda} - f_{\star}\|_{L^2}\right),
\end{equation}
and hence triangle inequality implies $\mathrm{bias}^2(\hat{f}_{\lambda}) = \|E_\epsilon(\hat{f}_{\lambda}) - f_{\star}\|_{L^2} \approx \| {f}_{\lambda} - f_{\star}\|_{L^2}$.
    



Later, \cite{li2024generalization} extended the work of \cite{li2023asymptotic} to analytic spectral algorithms. They replaced the key observation (\ref{eqn_key_observation_in_approx_A_fixed_d_learning_curve}) with the following analytic functional argument:
$$
\reg(T_X)-\reg(T) = \frac{1}{2\pi i}\oint_{D_{\lambda}} \varphi_{z}^{\krr}(T) (T_{X}-T) \varphi_{z}^{\krr}(T_{X}) \reg(z) \mathsf{d} z,
$$
where the definition of $D_{\lambda}$ is given in Assumption 3 in \cite{li2024generalization}. Similarly to other approximations in \cite{li2023asymptotic}, they derived the learning curve for fixed-dimensional spectral algorithms.


% {\color{red}"both KRR and spectral algorithms" sounds strange since KRR is a spectral algorithm. Maybe "both KRR and more general spectral algorithms"?}

We note that the above approximations for KRR and other analytic spectral algorithms have certain limitations: they require the regularization parameter $\lambda$ to be bounded below. For fixed dimensions and eigenvalues satisfying $\lambda_j \asymp j^{-\beta}$, \citep{li2024generalization, li2023asymptotic} observed that the approximation holds if and only if $\lambda = \Omega(n^{-\beta})$. Consequently, for fixed-dimensional KRR and spectral algorithms, 
tight bounds on the learning curve can be established when $\lambda = \Omega(n^{-\beta})$, while only a lower bound $(\ln(n))^{-4}$ can be established when $\lambda = o(n^{-\beta})$.


Unfortunately, the above limitations prevent us from employing a similar approach to derive the complete learning curves for large-dimensional KRR or spectral algorithms. In the large-dimensional setting, Theorem 2 in \cite{zhang2024optimal} and Appendix 4 in \cite{lu2024saturation} showed that the approximations hold if $u \leq u^{\prime}$ (ignoring logarithmic factors in $d$ for $\lambda$), which only covers the over-regularized regime of the learning curves.




\iffalse
\begin{table}[h!]
\centering
\renewcommand{\arraystretch}{1.6}
\begin{tabular}{|c|c|c|c|}
\hline
u 
& $(0, u^{\prime}(\tau)]$ 
& $[u^{\prime}(\tau), u^{\prime}(1)]$ 
& $[u^{\prime}(1), \infty)$ \\
\hline

$s=0$ 
& \multicolumn{3}{c|}{Theorem D.1 : $\Theta_{d, \mathrm{P}}(1)$} \\
\hline

$s>0, \gamma \geq 1$ 
& \makecell{Theorem 4.1 and 4.2 in \cite{lu2024saturation}:\\
$Bias_{res}=o(Var_{rate} + Bias_{rate})$}
& \makecell{{\color{red} Upper bound: Proposition D.2:} \\
$Bias_{res}=O(Var_{rate} + Bias_{rate})$\\
Lower bound: Theorem D.1}
& \makecell{Theorem D.1, Proposition G.13,\\ 
and Proposition G.14} \\
\hline

$s>0, \gamma < 1$
& Theorem D.1 
& \multicolumn{2}{c|}{Theorem D.1, Proposition G.13, and Proposition G.14} \\
\hline
\end{tabular}
\caption{Bounds for the bias term. From \cite{lu2024saturation} we have $Bias \leq Bias_{main} + Bias_{res}$, $Bias \geq Bias_{main} - Bias_{res}$, and $Bias_{main}=\Theta(Bias_{rate})$}
\end{table}
\fi







\iffalse
and at least one of the following conditions hold (see, e.g., Theorem 2 in \cite{zhang2024optimal} or Appendix 4 in \cite{lu2024saturation}):
\begin{equation}\label{eqn_condition_saturation}
    \text{(i) } \tau =\infty, \quad \text{(ii) }s>1/(2\tau), \quad \text{(iii) } \gamma > ((2\tau+1)s) / (2\tau(1+s)).
\end{equation}
\fi






\subsection{Recent advances in large-dimensional spectral algorithms}\label{subsec_related_work_large_dimension_spe}

\paragraph*{Consistency of KRR / KGF}


Several recent studies \citep{ghorbani2021linearized, Donhauser_how_2021, mei2022generalization, xiao2022precise, misiakiewicz_spectrum_2022, hu2022sharp} focus on the consistency of large-dimensional KRR and KGF with $s=0$. Roughly speaking, they showed the following approximations for KRR and KGF:
\begin{equation}\label{eqn_restate_thm_4_in_linear}
    \left|E_{x, \epsilon} \left[ \left(\hat{f}_{\lambda}(x) - f_{\star}(x) \right)^{2} \right]-
    \text{E}_{\text{approx}}
    \right| =o_{d, \mathbb{P}}\left(1\right),
\end{equation}
where $\text{E}_{\text{approx}}$ is a certain deterministic quantity (see, e.g., Theorem 4 in \cite{ghorbani2021linearized} or Theorem 2 in \cite{xiao2022precise}).



However, when \( s > 0 \), the RHS in (\ref{eqn_restate_thm_4_in_linear}) makes these results not precise enough to provide an exact convergence rate of the excess risk, and the approximation of the excess risk $E_{x, \epsilon} [ (\hat{f}_{\lambda}(x) - f_{\star}(x) )^{2} ]$, especially its convergence rate, when $s>0$ could be much harder since the excess risk is $o_{d, \mathbb{P}}(1)$.
% {\color{red}"the precise estimates when $s>0$ could be much harder" what does "estimte" mean here?}






\paragraph*{Convergence rate of optimally tuned spectral algorithms and corresponding minimax bound}

It is challenging to develop results on the excess risk of spectral algorithms when $s>0$. We summarize several studies as follows.


\cite{lu2023optimal} determined the minimax rate of KGF for $s=1$, that is, $f_{\star}$ falls exactly in the RKHS. However, the proof in \cite{lu2023optimal} relies heavily on empirical process theory, requiring bounds on the empirical loss and its difference from the expected one (the excess risk). Consequently, this approach does not generalize to cases where $s\neq 1$.


Later, \cite{zhang2024optimal, lu2024saturation, zhang2025learning} determined the convergence rate of optimally tuned spectral algorithms by modifying the techniques introduced in Appendix \ref{subsec_learning_curve_fix_d}:


\begin{proposition}[Restatement of 
Theorem 4.1 and Theorem 4.2 in \cite{lu2024saturation}]\label{prop_nips_left_learning_curve}
  % {\color{red}
  % Let $u^{\prime} = u^{\prime}(\tau)$ be the rate of the optimal regularization parameter $\lambda$ with qualification $\tau \leq \infty$ and $s>0$ be the coefficient of source condition in}
%
Let $u^{\prime} = u^{\prime}(\tau)$ be the rate of the optimal regularization parameter $\lambda$ with qualification $\tau \leq \infty$ and $s>0$ be the coefficient of source condition in
\cite{lu2024saturation} (ignoring the logarithm term).
  Suppose Assumptions 1,2,3,4,5,
and 6 hold with $\tau \leq \infty$, $s>0$, and $u \in (0, u^{\prime})$. 
Denote $\tilde{s}=\min\{s, 2\tau\}$.
Further, suppose at least one of the following conditions hold:
  $$
  \text{(i) } \tau =\infty, \quad \text{(ii) }s>1/(2\tau), \quad \text{(iii) } \gamma > ((2\tau+1)s) / (2\tau(1+s));
  $$
  Then the excess risk of the spectral algorithm estimator in (13) satisfies
\begin{align*}
E_{x, \epsilon} \left[ \left(\hat{f}_{\lambda}(x) - f_{\star}(x) \right)^{2} \right]
%\Theta_{d, \mathbb{P}} \left( \frac{n}{d^{\ell_{\lambda}+1}(n\lambda+1)^{2}} + \frac{d^{\ell_{\lambda}}}{n} + d^{-(\ell_{\lambda} + 1)  s}+ \mathbf{1}\{\tau<\infty\}\lambda^{ 2  \tau} d^{(2  \tau - \tilde{s})  \ell_{\lambda}} \right).
= \Theta_{d, \mathbb{P}} \left( 
d^{2u - \gamma - \ell_{\lambda} - 1}
+ d^{\ell_{\lambda} - \gamma}
+ d^{-(\ell_{\lambda} + 1)  s}
+ \mathbf{1}\{\tau<\infty\} d^{-2  \tau  u + (2  \tau - \tilde{s})  \ell_{\lambda}}
\right).
\end{align*}
\end{proposition}


On one hand, Proposition \ref{prop_nips_left_learning_curve} provides the learning curves for spectral algorithms in the over-regularized regime for certain values of $(\tau, s, \gamma)$. Moreover, since $u^{\prime} < \gamma$ (refer to Proposition \ref{prop_e_1}), for any $u \in (0, u^{\prime})$, from Assumption 4 we have $\tilde{\ell}=\ell_{\lambda}$. Therefore, Theorem 3.1 fully recovers the results in Proposition \ref{prop_nips_left_learning_curve}.




On the other hand, we note that Proposition \ref{prop_nips_left_learning_curve} does not apply to the under-regularized or interpolation regimes, and it does not cover the case where $\tau<\infty$ and both $s$ and $\gamma$ are small. We believe that these excluded cases are of equal importance:
\begin{itemize}
    \item As shown in Corollary 3.3 and its subsequent discussion, the benign overfitting phenomenon observed in large-dimensional kernel interpolation and neural networks suggests that the learning curves for small $s$ merit further study. 

    \item The above studies \cite{ghorbani2021linearized, Donhauser_how_2021, mei2022generalization, xiao2022precise, misiakiewicz_spectrum_2022, hu2022sharp} that considered the extreme case where $s=0$ motivate us to investigate the learning curves for all $s$, to provide a unified explanation of these results.

    \item Many large-dimensional datasets in machine learning, especially image datasets, often have dimensions that far exceed the sample size, further encouraging our focus on small $\gamma$. 
\end{itemize}
In summary, it is necessary to refine the conclusions of Proposition \ref{prop_nips_left_learning_curve} to derive the learning curves for large-dimensional spectral algorithms.
However, as discussed in Appendix \ref{subsec_learning_curve_fix_d}, the proof of Proposition \ref{prop_nips_left_learning_curve} relies on the validity of approximations such as (\ref{eqn_approx_variance_fixed_d_learning_curve}) and (\ref{eqn_approx_bias_fixed_d_learning_curve}), which prevents us from obtaining a complete learning curve.






Finally, \cite{lu2024pinsker} determined the exact asymptotic minimax bound for large-dimensional kernel regression problems, known as the Pinsker bound:



\begin{proposition}[Restatement of Theorem 3.1 in \cite{lu2024pinsker}]\label{prop_pinsker_constant}
    Let $\mathcal{P}$ consist of all the distributions $\rho_{f_{\star}}$ on $\mathcal{X} \times \mathcal{Y}$ given by (1) such that Assumption 
1,2, and 3 hold for some $\gamma, s>0$. 
Furthermore, assume that there exists an absolute constant $\alpha>0$, such that we have $n = \alpha d^{\gamma} (1+o_{d}(1))$. Define \(p := \lfloor \gamma/(s+1) \rfloor\).
Then, we have
    \begin{equation*}
        \inf_{\hat{f}} \sup_{
        \rho_{f_{\star}} \in \mathcal{P}} \mathbb{E}_{(X, Y) \overset{\mathcal{D}}{\sim} \rho_{f_{\star}}^{\otimes n}}\left[
    \|\hat{f}-f_{\star}\|_{L^2}^2
    \right]
    =
    \calC^{\star} \cdot \left( 
    d^{p-\gamma} + 
    d^{-(p+1)s}
    \right)(1+o_{d}(1)),
    \end{equation*}
    where $\hat{f}$ is any estimator of $f_{\star}$, measurable with respect to the observed data set $(X, Y)$, and the definition of the Pinsker constant $\calC^{\star}$ can be found in Theorem 3.1 in \cite{lu2024pinsker}.
\end{proposition}






\paragraph*{Benign overfitting phenomenon of KRR in the interpolation regime}


A substantial body of recent work has analyzed large-dimensional KRR in the interpolation regime (i.e., when $\lambda=O_d(n^{-1})$), with particular focus on the case $\lambda=0$ (i.e., kernel interpolation). 
% {\color{red} I don't quite get the word "especially " and what is the difference between "large-dimensional KRR in the interpolation regime" and "kernel interpolation (i.e., KRR with $\lambda=0$) in large dimensions"? }
For example, \cite{Liang_Just_2019} showed the consistency of kernel interpolation with $s=\gamma=1$; subsequent studies \cite{liang2020multiple, barzilai2023generalization} provides upper bounds on the excess risk of kernel interpolation with $\gamma >0$. However, it is usually difficult to derive the matching lower bound, which is not provided in the above studies (also refer to Section 4 in \cite{zhang2024phase} for a detailed discussion).


A tight characterization on the excess risk  (with matching upper and lower bounds) of KRR was provided in \cite{aerni2023strong}. They considered a specific type of convolutional kernel on the discrete hypercube and a special form of the true regression function $f_{\star}(x) = x_{1}x_{2}\cdots x_{L^{*}}$, where $L^{*}$ is an absolute constant that formulated in Theorem 1 in \cite{aerni2023strong}. 






To our knowledge, the results most similar to Theorem 3.1 are provided by \cite{zhang2024phase}. They determined the exact convergence rates of kernel interpolation stated as follows. 

\begin{proposition}[Restatement of Corollary 1 in \cite{zhang2024phase}]\label{prop_intepolation}
Let $\hat{f}_{\mathrm{inter}}(\cdot):= \hat{f}_{0}(\cdot) = \mathscr{K}(\cdot, X) \left(\mathscr{K}(X, X)\right)^{-1} Y$ be the kernel interpolation estimator.
  Suppose Assumptions 1,2,3, and 6 hold with $s \geq 0$. 
Denote $\tilde{s}=\min\{s, 2\tau\}$.
  Then the excess risk of the kernel interpolation estimator satisfies
\begin{align*}
E_{x, \epsilon} \left[ \left(\hat{f}_{0}(x) - f_{\star}(x) \right)^{2} \right]
=
\Theta_{d, \mathbb{P}} \left( 
d^{\gamma - \ell_{\gamma} - 1}
+ d^{\ell_{\gamma} - \gamma}
+ d^{-(\ell_{\gamma} + 1)  s}
\right).
\end{align*}
\end{proposition}


% {\green [Have deleted the itemize]}
It is worth mentioning that the proofs in \cite{zhang2024phase} are difficult to extend to obtain (i) the learning curves for KRR in the over-regularized and under-regularized regimes, or (ii) the learning curves for more general spectral algorithms, even when restricting attention to the interpolation regime. 
To illustrate this, let's restate their proof steps.

They first partition the matrix version of the equality given in (5) into two parts (see Appendix \ref{append_notation} for further details):
$$
K = \Psi_{\leq \ell_{\gamma}} \Sigma_{\leq \ell_{\gamma}} \Psi_{\leq \ell_{\gamma}}^{\top} + \Psi_{>\ell_{\gamma}} \Sigma_{>\ell_{\gamma}} \Psi_{>\ell_{\gamma}}^{\top},
$$
where $\ell_{\gamma} = \lfloor \gamma \rfloor$, and then showed that
\begin{equation}\label{eqn_approx_phase_1}
    \left\|\Psi_{>\ell_{\gamma}} \Sigma_{>\ell_{\gamma}}^j \Psi_{>\ell_{\gamma}}^{\top} - \kappa_j\mathrm{I}_n \right\|_{\mathrm{op}} = o_{d, \mathbb{P}}(\kappa_j), \quad \kappa_j = \sum_{k=\ell_{\gamma}+1}^\infty \mu_k^j N(d, k),\quad j=1,2;
\end{equation}
and that
\begin{equation}\label{eqn_approx_phase_2}
\left\| n^{-1}{\Psi_{\leq \ell_{\gamma}}^{\top} \Psi_{\leq \ell_{\gamma}} } - \mathrm{I}_{N_{\ell_{\gamma}}} \right\|_{\mathrm{op}} = o_{d, \mathbb{P}}(1).
\end{equation}
In the interpolation regime, we have $n\lambda =o_d(1)$. Therefore, they established the following approximations for $M:= K+n\lambda \mathrm{I}_n$:
\begin{equation}\label{eqn_approx_phase_3}
\lambda_{\text{max}}\left( M^{-1}\right) = O_{d, \mathbb{P}}((n\lambda+1)^{-1}) = O_{d, \mathbb{P}}(1) \quad \text{ and } \quad
\mathrm{tr}\left( M^{-1} / n\right)=\Omega_{d, \mathbb{P}}((n\lambda+1)^{-1}) = \Omega_{d, \mathbb{P}}(1).
\end{equation}
For the variance term, they decompose it into two parts:
    \begin{equation}
\label{eq var tr}
    \begin{aligned}
        \mathrm{var}\left( \hat{f}_{\lambda} \right) 
    =  
    \sigma^{2} \cdot \mathrm{tr} \left( M^{-1} \Psi_{>\ell_{\gamma}} \Sigma_{>\ell_{\gamma}}^2 \Psi_{>\ell_{\gamma}}^{\top} M^{-1}\right)  ~+~ \sigma^{2} \cdot \mathrm{tr} \left( M^{-1} \Psi_{\leq \ell_{\gamma}} \Sigma_{\leq \ell_{\gamma}}^{2} \Psi_{\leq \ell_{\gamma}}^{\top} M^{-1}\right),
    \end{aligned}
\end{equation}
so that they can use (\ref{eqn_approx_phase_1}), (\ref{eqn_approx_phase_2}), and (\ref{eqn_approx_phase_3}) to obtain the tight bound of the variance.

For the bias term, they decompose it into five parts, separating the contributions from indices $\leq \ell_{\gamma}$ and $> \ell_{\gamma}$. By employing the Sherman-Morrison-Woodbury (SMW) formula (see, e.g., Proposition \ref{prop_smw_formula}), they convert $M^{-1}$ into $M_{>\ell_{\gamma}}^{-1} := (\Psi_{>\ell_{\gamma}} \Sigma_{>\ell_{\gamma}} \Psi_{>\ell_{\gamma}}^{\top} + n\lambda \mathrm{I}_n)^{-1}$, which by (\ref{eqn_approx_phase_1}) is approximately the identity matrix and thus easier to handle. For example, for the first term $B_{1, 1}$, they obtain
    \begin{equation}\label{eqn_appro_bias_1_1_phase}
     \begin{aligned}
        \left\|B_{1, 1}\right\|_{2}^{2} 
        :=&~ \left\| \left( \Sigma_{\le \ell_{\gamma}} - \Sigma_{\le \ell_{\gamma}} \Psi^{\top}_{\le \ell_{\gamma}} M^{-1} \Psi_{\le \ell_{\gamma}} \Sigma_{\le \ell_{\gamma}}\right)  \Sigma_{\le \ell_{\gamma}}^{-1}  \theta_{\le \ell_{\gamma}} \right\|_{2}^{2}\\
        \overset{\text{ SMW formula } }{=} &~
        \left\| \left( \Sigma_{\le \ell_{\gamma}}^{-1} + \Psi_{\le \ell_{\gamma}}^{\top} M_{>\ell_{\gamma}}^{-1} \Psi_{\le \ell_{\gamma}}  \right)^{-1}  \Sigma_{\le \ell_{\gamma}}^{-1}  \theta_{\le \ell_{\gamma}} \right\|_{2}^{2}.
    \end{aligned}
    \end{equation}
    Similarly, using (\ref{eqn_approx_phase_1}), (\ref{eqn_approx_phase_2}), and (\ref{eqn_approx_phase_3}), they obtained a tight bound on the bias.



\iffalse
We restate their proof steps as follows:
\begin{itemize}
    \item They first partition the matrix version of the equality given in (5) into two parts:
$$
n^{-1} K = \Psi_{\leq \ell_{\gamma}} \Sigma_{\leq \ell_{\gamma}} \Psi_{\leq \ell_{\gamma}}^{\top} + \Psi_{>\ell_{\gamma}} \Sigma_{>\ell_{\gamma}} \Psi_{>\ell_{\gamma}}^{\top},
$$
(see Appendix \ref{append_notation} for further details), and  then showed that
\begin{equation}\label{eqn_approx_phase_1}
    \left\|\Psi_{>\ell_{\gamma}} \Sigma_{>\ell_{\gamma}}^j \Psi_{>\ell_{\gamma}}^{\top} - \kappa_j\mathrm{I}_n \right\|_{\mathrm{op}} = o_{d, \mathbb{P}}(\kappa_j), \quad \kappa_j = \sum_{k=\ell_{\gamma}+1}^\infty \mu_k^j N(d, k),\quad j=1,2;
\end{equation}
and that
\begin{equation}\label{eqn_approx_phase_2}
\left\| n^{-1}{\Psi_{\leq \ell_{\gamma}}^{\top} \Psi_{\leq \ell_{\gamma}} } - \mathrm{I}_{N_{\ell_{\gamma}}} \right\|_{\mathrm{op}} = o_{d, \mathbb{P}}(1).
\end{equation}

    \item In the interpolation regime, we have $n\lambda =o_d(1)$. Therefore, they established the following approximations for $M:= (K+n\lambda \mathrm{I}_n)$:
\begin{equation}\label{eqn_approx_phase_3}
\lambda_{\text{max}}\left( M^{-1}\right) = O_{d, \mathbb{P}}((n\lambda+1)^{-1}) = O_{d, \mathbb{P}}(1) \quad \text{ and } \quad
\mathrm{tr}\left( M^{-1}\right)=\Omega_{d, \mathbb{P}}((n\lambda+1)^{-1}) = \Omega_{d, \mathbb{P}}(1).
\end{equation}

\item \cite{zhang2024phase} then prove the following approximation:
\begin{itemize}
    \item {\bf Variance term: } They decompose the variance into two terms:
    \begin{equation}
\label{eq var tr}
    \begin{aligned}
        \mathrm{var}\left( \hat{f}_{\lambda} \right) 
    =  
    \sigma^{2} \cdot \mathrm{tr} \left( M^{-1} \Psi_{>\ell_{\gamma}} \Sigma_{>\ell_{\gamma}}^2 \Psi_{>\ell_{\gamma}}^{\top} M^{-1}\right)  ~+~ \sigma^{2} \cdot \mathrm{tr} \left( M^{-1} \Psi_{\leq \ell_{\gamma}} \Sigma_{\leq \ell_{\gamma}}^{2} \Psi_{\leq \ell_{\gamma}}^{\top} M^{-1}\right),
    \end{aligned}
\end{equation}
and from (\ref{eqn_approx_phase_1}), (\ref{eqn_approx_phase_2}), and (\ref{eqn_approx_phase_3}), they obtained the tight bound of the variance.


    \item {\bf Bias term: } They decompose the bias into five parts, separating the contributions from indices $\leq \ell_{\gamma}$ and $> \ell_{\gamma}$. By employing the Sherman–Morrison–Woodbury (SMW) formula (see, e.g., Proposition \ref{prop_smw_formula}), they convert $M^{-1}$ into $M_{>\ell_{\gamma}}^{-1} := (\Psi_{>\ell_{\gamma}} \Sigma_{>\ell_{\gamma}} \Psi_{>\ell_{\gamma}}^{\top} + n\lambda \mathrm{I}_n)^{-1}$, which by (\ref{eqn_approx_phase_1}) is approximately the identity matrix and thus easier to handle. For example, for the first term $B_{1, 1}$, they obtain
    \begin{equation}\label{eqn_appro_bias_1_1_phase}
     \begin{aligned}
        \left\|B_{1, 1}\right\|_{2}^{2} 
        =&~ \left\| \left( \Sigma_{\le \ell_{\gamma}} - \Sigma_{\le \ell_{\gamma}} \Psi^{\top}_{\le \ell_{\gamma}} M^{-1} \Psi_{\le \ell_{\gamma}} \Sigma_{\le \ell_{\gamma}}\right)  \Sigma_{\le \ell_{\gamma}}^{-1}  \theta_{\le \ell_{\gamma}} \right\|_{2}^{2}\\
        \overset{\text{ SMW formula } }{=} &~
        \left\| \left( \Sigma_{\le \ell_{\gamma}}^{-1} + \Psi_{\le \ell_{\gamma}}^{\top} M_{>\ell_{\gamma}}^{-1} \Psi_{\le \ell_{\gamma}}  \right)^{-1}  \Sigma_{\le \ell_{\gamma}}^{-1}  \theta_{\le \ell_{\gamma}} \right\|_{2}^{2}.
    \end{aligned}
    \end{equation}
    Similarly, using (\ref{eqn_approx_phase_1}), (\ref{eqn_approx_phase_2}), and (\ref{eqn_approx_phase_3}), they obtained a tight bound on the bias.

\end{itemize}

\end{itemize}
\fi




% {\green [Have deleted the itemize]} 
However, the following difficulties exist in extending the above results to the two cases we mentioned above:
(i) In the over-regularized and under-regularized regimes, we have $\ell_{\lambda} \leq \ell_{\gamma}$. Consequently, the indices in the decompositions of the variance (\ref{eq var tr}) and the bias (\ref{eqn_appro_bias_1_1_phase}) need to be replaced by $\ell_\lambda$. Under these circumstances, the preliminary approximations (\ref{eqn_approx_phase_1}) and (\ref{eqn_approx_phase_3}) no longer hold.
(ii) When KRR is replaced by other spectral algorithms, the matrix $n^{-1} \reg^{\krr}(K/n) := M^{-1}$ is replaced by $n^{-1} \reg(K/n)$. However, the SMW formula does not necessarily hold for $n^{-1} \reg(K/n)$. 

\iffalse
\begin{itemize}
    \item[(i)] Regarding the KRR algorithm: In the over-regularized and under-regularized regimes, we have $\ell_{\lambda} \leq \ell_{\gamma}$. Consequently, the indices in the decompositions of the variance (\ref{eq var tr}) and the bias (\ref{eqn_appro_bias_1_1_phase}) need to be replaced by $\ell_\lambda$. Under these circumstances, the preliminary approximations (\ref{eqn_approx_phase_1}) and (\ref{eqn_approx_phase_3}) no longer hold.
    %, so the proofs in \cite{aerni2023strong, zhang2024phase} cannot be directly extended to derive the learning curves for KRR in the over-regularized / under-regularized regimes.

    \item[(ii)] When KRR is replaced by other spectral algorithms, the matrix $n^{-1} \reg^{\krr}(K/n) := M^{-1}$ is replaced by $n^{-1} \reg(K/n)$. However, the SMW formula does not necessarily hold for $n^{-1} \reg(K/n)$. 
    %Therefore, we cannot directly extend the proofs in \cite{aerni2023strong, zhang2024phase} to obtain the learning curves for spectral algorithms, even in the interpolation regime. 


\end{itemize}
\fi

 % {\color{red} Variance: V22, Bias: Assumption 5}







\iffalse
\paragraph*{Concentration of the excess risks of KRR and ridge regression (RR) in sequence models}

Two recent studies (\cite{cheng2022dimension, misiakiewicz2024non}) have established similarities between 
(i) the excess risk of kernel ridge regression (KRR) in the nonparametric regression model (1) and (ii) the excess risk of ridge regression (RR) on the Gaussian sequence model. 
Specifically, they consider the following two settings:
\begin{itemize}
    \item[(i)]  First, they consider the nonparametric regression model (1)
    with a regression function $f_{\star} = \sum_j \theta_j \phi_j \in L^2$.
    Specifying a kernel $K$ with eigenvalues $\lambda_j$'s, they then consider the KRR estimator with regularization parameter $\lambda$. 
    
    \item[(ii)] They also consider the Gaussian sequence model $z_j=\theta_j+ \xi_j, j=1,2, \cdots$, with 
    a specific variance of the Gaussian noise $\xi_j$.
    Let $\lambda_{\text{RR}} = \lambda_{\text{RR}}(\lambda)$ be given as in (7) of \cite{cheng2022dimension},  $\hat{\theta}_{\lambda_{\text{RR}}}$ be the RR estimator at regularization level $\lambda_{\text{RR}}$ defined in (19) of \cite{cheng2022dimension}, and $R_n(\lambda_{\text{RR}})$ be the corresponding excess risk.
    
\end{itemize}

\cite{cheng2022dimension, misiakiewicz2024non} then built the following approximations on the excess risks of KRR and RR:

\begin{proposition}[Restate Theorem 2 in \cite{misiakiewicz2024non}]\label{prop_restate_misia}
    Suppose Assumption \ref{assump_asymptotic} and \ref{assu:coef_of_inner_prod_kernel} hold for some $\gamma>0$. 
    Further suppose there exists an absolute constant $C$, such that we have $\left\|\mathrm{P}_{>\ell} f_{\star}\right\|_{L^2} \geq\left\|f_{\star}\right\|_{L^2} / C$, and for any integer $q \geq 2$,  we have $\left\|f_{\star}\right\|_{L^q} \leq (C q)^{q(\ell+1) / 2}\left\|f_{\star}\right\|_{L^2}$. Then under certain conditions, we have
    $$
    E_{x} \left[ \left(\hat{f}_{\lambda}(x) - f_{\star}(x) \right)^{2} \mid \epsilon_1, \cdots, \epsilon_n \right] = R_n(\lambda_{\text{RR}}) \left(1 + o_{d, \mathbb{P}}\left(1\right)\right).
    $$
\end{proposition}


{\color{red}
It is obvious that the sequence model is much simpler than kernel regression models. Therefore, one might wonder whether Proposition \ref{prop_restate_misia} could help us derive the learning curves for large-dimensional KRR. However, we face the following obstacles when we try to do so:
\begin{itemize}
    \item \cite{cheng2022dimension, misiakiewicz2024non} did not discuss how to compute the convergence rate of $R_n(\lambda_{\text{RR}})$. In fact, they only provide upper bounds for $R_n(\lambda_{\text{RR}})$ under specific conditions (see, e.g., Proposition 2.2 in \cite{cheng2022dimension}). Worse still, they do not offer an explicit expression between $\lambda_{\text{RR}}$ and $\lambda$, which makes it difficult to determine the convergence rate of $R_n(\lambda_{\text{RR}})$.

    \item The assumptions in \cite{cheng2022dimension, misiakiewicz2024non} are too restrictive. For instance, \cite{cheng2022dimension} requires that the eigenfunctions have zero mean, i.e., \(\mathbb{E}\phi_j = 0\) for all \(j = 0, 1, \cdots\). However, most kernel functions, including spherical inner product kernels, have a constant eigenfunction and thus do not satisfy this assumption. Similarly, if \(s>0\), functions in \([\calH]^{s}\) with nonzero \(L^2\) norms do not meet the conditions of Proposition \ref{prop_restate_misia} since \(\|\mathrm{P}_{>\ell} f_{\star}\|_{L^2} \to 0\); moreover, most functions with zero \(L^2\) norm have nonzero \(L^{q}\) norms for certain \(q\), thereby violating the assumption \(\|f_{\star}\|_{L^q} \leq (Cq)^{q(\ell+1)/2}\|f_{\star}\|_{L^2}\).
    
\end{itemize}
}


{\color{red}
17 heavily rely on the Convex concentration on eigenfunctions, and 61 replace it with the following stronger assumption in inner product kernel regression cases
$$
\left\|\mathrm{P}_{>\ell} f_{\star}\right\|_{L^2} \geq\left\|f_{\star}\right\|_{L^2} / C \text{ and } \left\|f_{\star}\right\|_{L^q} \leq (C q)^{(\ell+1) / 2}\left\|f_{\star}\right\|_{L^2}
$$
However, if \(s>0\), functions in \([\calH]^{s}\) with nonzero \(L^2\) norms do not meet the conditions of Proposition \ref{prop_restate_misia} since \(\|\mathrm{P}_{>\ell} f_{\star}\|_{L^2} \to 0\); moreover, most functions with zero \(L^2\) norm have nonzero \(L^{q}\) norms for certain \(q\), thereby violating the assumption \(\|f_{\star}\|_{L^q} \leq (Cq)^{q(\ell+1)/2}\|f_{\star}\|_{L^2}\).
When $s=0$, the risk do not converge to $0$, hence they actually attain $|R-R_{seq}|=o_d(1)$.



}
\fi












\subsection{Comparison for the learning curves of kernel regression model and sequence model}\label{sec_sequence_model}


Sequence models have been widely considered as simplifications for the analysis of kernel regression models. For example, the celebrated work by \cite{pinsker1980optimal} and \cite{nussbaum1985spline} and subsequent ones show that the exact asymptotic of the minimax risks, i.e., the Pinsker bounds, for Gaussian sequence models and fixed-dimensional kernel regression models over \( [\mathcal{H}]^{s} \) for \( s > 1 \) are identical; \cite{brown1996asymptotic, carter2006continous, reiss2008asymptotic} further showed the classical Le Cam equivalence between above kernel regression models and the Gaussian sequence models;
\cite{li2024improving, li2025diagonal} focus on the sequence model and kernel regression model separately, showing that the order of eigen-functions impacts regression outcomes, and that over-parameterization enhances the generalization capability.
Recently, \cite{cheng2022dimension, misiakiewicz2024non} showed the concentration of the excess risks of KRR in kernel regression models and ridge regression in sequence models.




Proposition 3.4 in Section 3 shows that, under sufficient regularization, the excess risk of large-dimensional kernel regression has the same order as that of the associated sequence model. This complements recent comparison results such as \cite{cheng2022dimension, misiakiewicz2024non}, which study sharper approximations for KRR under stronger assumptions on the regression function and on the sequence-model noise structure. In contrast, our comparison is derived as a consequence of the full learning-curve theorem and applies uniformly across all $s\ge 0$ for the analytic spectral algorithms covered by Definition 2.2.


More specifically, the approximation result of \cite{misiakiewicz2024non} is sharper for KRR when its additional assumptions hold, but those assumptions become restrictive in the large-dimensional regimes covered here. From this perspective, Theorem~3.1 and Proposition~3.4 trade pointwise sharpness for broader regime coverage and a unified treatment of spectral algorithms beyond KRR.


\begin{comment}
    
In the next subsection, we will first determine the learning curve of equivalent sequence models with true parameters defined by the coefficients of the regression function $f_{\star}$ in (\ref{eqn_decompose_reg_func}). Then, we will restate and compare results in \cite{cheng2022dimension, misiakiewicz2024non}.



\subsubsection{Learning curves of equivalent sequence models}

\iffalse
In this subsection, we will consider a specific type of sequence model with true parameters defined by the coefficients of the regression function $f_{\star}$ in (\ref{eqn_decompose_reg_func}). 
Then, we will use the filter function defined in Definition 2.2 to construct estimators of true parameters. 
Finally, we will show that the learning curve for kernel regression model given in Theorem 3.1 equals the learning curve of the above sequence model equals with certain regularization.
\fi

Let us consider the sequence model
\begin{equation}\label{eqn_sequence_model}
    z_j = f_j + \xi_j,\quad j=1, 2, \cdots,
\end{equation}
where $(z_j)_{j=1}^{\infty}$ is the observation, $f_j$ is the coefficient of $f_{\star}$ with respect to $\phi_j$, that is, we have $f_j = \langle f_{\star}, \phi_j \rangle_{L^2}$, $j=1, 2, \cdots$. Moreover, $\xi_j$, $j=1, 2, \cdots$ are (not necessarily independent) random variables with mean zero and variance $\sigma^2 / n$.

Now, suppose $\lambda_j$'s are known and consider the parametrization $f_j = \lambda_j^{1/2} \beta_j$, $j=1, 2, \cdots$. Notice that, applying the ridge regression with regularization parameter $\lambda$ or gradient flow with early-stopping time $t=\lambda^{-1}$ will derive estimators $\hat{f}_j^{\lambda} = \lambda_j \reg(\lambda_j) z_j$, $j=1, 2, \cdots$, where $\reg = \reg^{\krr}$ or $\reg^{\gf}$ are defined in Section \ref{subsec_example_spe} (see, e.g., Proposition \ref{prop_sequence_estimators}). Therefore, we define the estimator of $(f_j)_{j=1}^{\infty}$ as:
\begin{equation}\label{eqn_est_sequence_model}
    \hat{f}_j^{\lambda} := \lambda_j \reg(\lambda_j) z_j,\quad j=1, 2, \cdots,
\end{equation}
where $\reg$ is the filter function defined in Definition 2.2.
Finally, the (expected) excess risk of $(\hat{f}_j^{\lambda})_{j=1}^{\infty}$ is defined as:
\begin{equation}\label{eqn_excess_risk_sequence}
    \begin{aligned}
        E_{\xi}\left[\sum_{j=1}^{\infty}(\hat{f}_j^{\lambda} - f_j)^2\right] := &~ \sum_{j=1}^{\infty}\mathrm{var}(\hat{f}_j^{\lambda}) + \sum_{j=1}^{\infty}\left(E_{\xi}\left[\hat{f}_j^{\lambda}\right] - f_j\right)^2\\
        = &~
        \frac{\sigma^2}{n}\sum_{j =1}^\infty \left[ \lambda_j \reg(\lambda_j) \right]^2 + \sum\limits_{j=1}^{\infty} \left( \rem(\lambda_j) f_{j}\right)^{2},
    \end{aligned}
\end{equation}
where $E_{\xi}$ means taking expectation on the noises $\xi_j$, $j=1, 2, \cdots$.


\begin{remark}
    The same parametrization that $f_j = \lambda_j^{1/2} \beta_j$, $j=1, 2, \cdots$ and the same excess risk $E_{\xi}\left[\sum_{j=1}^{\infty}(\hat{f}_j^{\lambda} - f_j)^2\right]$ are also considered in recent work that focus on the concentration of the excess risks of KRR and ridge regression in sequence models (see, e.g., \cite{cheng2022dimension, misiakiewicz2024non}). 
    %We will provide a detailed discussion in the next subsection.
\end{remark}



The following proposition is a corollary of Lemma D.14 in \cite{lu2024saturation}. It provides the learning curve of the estimator of sequence model and clarifies its relationship with the learning curve of kernel regression model in Theorem 3.1.


% \dhnote{
% This result relies on the approximation of $\lambda_j$ for large-dimensional inner product kernel. The text here seems ignore the specific form of $\lambda_j$ and reads like the result hold for all $\lambda_j$. 
% }

\begin{proposition}\label{prop_learning_curve_of_sequence}
    Adopt all the notations and assumptions in Theorem 3.1. 
    %Further suppose that $\lambda = \Omega_d(n^{-1})$.
    Then the excess risk of $(\hat{f}_j^{\lambda})_{j=1}^{\infty}$ satisfies
\begin{equation}\label{eqn_prop_learning_curve_of_sequence_0}
\begin{aligned}
     E_{\xi}\left[\sum_{j=1}^{\infty}(\hat{f}_j^{\lambda} - f_j)^2\right]
= 
% \left\{\begin{matrix}
% \Theta_{d} \left( 
% d^{2u - \gamma - \ell_{\lambda} - 1}
% + d^{\ell_{\lambda} - \gamma}
% + d^{-(\ell_{\lambda} + 1)  s}
% + \mathbf{1}\{\tau<\infty\} d^{-2  \tau  u + (2  \tau - \tilde{s})  \ell_{\lambda}}
% \right) & \lambda > 0, \\
% +\infty & \lambda = 0. \\
% \end{matrix}\right.
\Theta_{d} \left( 
d^{2u - \gamma - \ell_{\lambda} - 1}
+ d^{\ell_{\lambda} - \gamma}
+ d^{-(\ell_{\lambda} + 1)  s}
+ \mathbf{1}\{\tau<\infty\} d^{-2  \tau  u + (2  \tau - \tilde{s})  \ell_{\lambda}}
\right).
\end{aligned}
\end{equation}
Consequently, Theorem 3.1 implies
\begin{equation}\label{eqn_prop_learning_curve_of_sequence}
    E_{x, \epsilon} \left[ \left(\hat{f}_{\lambda}(x) - f_{\star}(x) \right)^{2} \right] = \Theta_{d, \mathbb{P}} \left(
E_{\xi}\left[\sum_{j=1}^{\infty}(\hat{f}_j^{\max\{\lambda, n^{-1}\}} - f_j)^2\right]
\right).
\end{equation}
\end{proposition}


For any regularization parameter $\lambda = \Omega_d(n^{-1})$, (21) determine the same convergence rate of the excess risks of spectral algorithms applied to the kernel regression model and those applied to the sequence model. 
This means that, under sufficient regularization, the statistical behavior of the two estimators is identical up to constant factors, reinforcing the idea that the sequence model serves as a theoretically tractable proxy for kernel regression in this regime.
\end{comment}



\iffalse
This proposition offers two key insights:
\begin{itemize}
    \item On one hand, for any regularization parameter $\lambda = \Omega_d(n^{-1})$, (21) establishes the asymptotic equivalence between the learning curves of spectral algorithms applied to the kernel regression model and those applied to the sequence model. This means that, under sufficient regularization, the statistical behavior of the two estimators is identical up to constant factors, reinforcing the idea that the sequence model serves as a theoretically tractable proxy for kernel regression in this regime.

    \item On the other hand, when $\lambda = o_d(n^{-1})$, the excess risk of the sequence model becomes significantly larger than that of the kernel regression model. Especially, the excess risk of sequence model diverges in the extreme case where $\lambda=0$.
    
\end{itemize}



\begin{remark}
    The above discrepancy when $\lambda = o_d(n^{-1})$ can be explained as follows. Both models aim to recover the sequence $\{f_j :=\langle f_{\star}, \phi_j \rangle_{L^2}\}_{j=1}^{\infty}$. However, given finite samples $((x_1, y_1), \cdots, (x_n, y_n))$, we do not have direct access to $f_j$, and hence cannot generate the idealized data $z_j=f_j+\xi_j$. Instead, we can only observe a noisy empirical estimate:
    $$
        \hat{z}_j := \frac{1}{n} \sum_{i=1}^{n} y_i \phi_j(x_i) + \xi_j, \quad j \geq 1,
    $$
    and the resulting estimator takes the form $\lambda_j \reg(\lambda_j) \hat{z}_j$. Notice that, the expectation $E_{\epsilon}[n^{-1} \sum_{i=1}^{n} y_i \phi_j(x_i)] = n^{-1} \sum_{i=1}^{n} f_{\star}(x_i) \phi_j(x_i)$ is not necessarily an unbiased estimate of $f_j$. When $\lambda$ is too small, the estimation error $\hat{f}_j^{\lambda} - \lambda_j \reg(\lambda_j) \hat{z}_j$ becomes non-negligible compared to $\hat{f}_j^{\lambda}$ itself. Therefore, the excess risk of sequence model becomes much larger than the excess risk of kernel regression model.
\end{remark}
\fi










\subsubsection{Literature review: Concentration of the excess risks of KRR and ridge regression in sequence models}\label{subsec_related_work_sequence_model}

Recent studies \citep{cheng2022dimension, misiakiewicz2024non} have established approximations between 
(i) the excess risk of kernel ridge regression (KRR) in the nonparametric regression model (1) and (ii) the excess risk of ridge regression (RR) on the Gaussian sequence model. 

% {\green [Have deleted the itemize]} 
Specifically, they first consider the nonparametric regression model (1)
    with a regression function $f_{\star} = \sum_j f_j \phi_j \in L^2$.
    Specifying a kernel $K$ with eigenvalues $\lambda_j$'s, they then consider the KRR estimator in (11) with regularization parameter $\lambda$. 
    They also consider the Gaussian sequence model 
    $$
    \tilde{z}_j=f_j+ \tilde{\xi}_j, j=1,2, \cdots, 
    $$
    where $\tilde{\xi}_j \sim_{\text{i.i.d.}} \calN (0, [\sigma^2 + E_{\tilde{\xi}_j}[\sum_{j=1}^{\infty}(\tilde{f}_j^{\tilde{\lambda}} - f_j)^2]] / n)$.
    Here, $\tilde{\lambda}$ is the unique non-negative solution of
\begin{equation}\label{eqn_lambda_star}
    n \left(1-\frac{\lambda}{\tilde{\lambda}}\right) = \sum_{j=1}^{\infty} \frac{\lambda_j}{\lambda_j + \tilde{\lambda}},
\end{equation}
and $\tilde{f}_j^{\tilde{\lambda}} := \lambda_j \varphi_{\tilde{\lambda}}^{\krr}(\lambda_j) \tilde{z}_j$ is the RR estimator with regularization parameter $\tilde{\lambda}$ for $j=1, 2, \cdots$. 

\begin{remark}
    Notice that the above sequence model is defined implicitly, since the variance of the noise depends on the excess risk, which itself depends on the noise.
\end{remark}



\cite{misiakiewicz2024non} then built the following approximations on the excess risks of KRR and RR:

\begin{proposition}[Restate Theorem 2 in \cite{misiakiewicz2024non}]\label{prop_restate_misia}
    Suppose Assumption 1 and 2 hold for some $\gamma>0$. 
    Further suppose there exists an absolute constant $C$, such that we have $\left\|\mathrm{P}_{>\ell_{\gamma}} f_{\star}\right\|_{L^2} \geq \left\|f_{\star}\right\|_{L^2} / C$, and for all integers $q \geq 2$,  we have $\left\|f_{\star}\right\|_{L^q} \leq (C q)^{(\ell_{\gamma}+1) / 2}\left\|f_{\star}\right\|_{L^2}$. Then under certain conditions, we have
    \begin{equation}\label{eqn_prop_restate_misia}
        E_{x, \epsilon} \left[ \left(\hat{f}_{\lambda}(x) - f_{\star}(x) \right)^{2} \right] = E_{\tilde{\xi}}\left[\sum_{j=1}^{\infty}(\tilde{f}_j^{\tilde{\lambda}} - f_j)^2\right] \cdot \left(1 + o_{d, \mathbb{P}}\left(1\right)\right),
    \end{equation}
where $\mathrm{P}_{>\ell_{\gamma}}$ means the projection onto spherical harmonics with degree $>\ell_{\gamma}$, and $E_{\tilde{\xi}}$ means taking expectation on the noises $\tilde{\xi}_j$, $j=1, 2, \cdots$.
\end{proposition}


% {\green [Have deleted the itemize]} 
With more detailed assumptions on the regression function, Eq. (\ref{eqn_prop_restate_misia}) provides a sharper approximation on the excess risks than Eq. (21) when the algorithm is specified to KRR. To clarify our novelty, we present a direct comparison:

When $s>0$, Eq. (\ref{eqn_prop_restate_misia}) relies on assumptions that become increasingly restrictive as $d$ grows. Specifically, let $f_{\star}(\cdot) = \sum_{j=0}^{\infty} f_j \phi_j(\cdot)$ be any regression function satisfying Assumption 3 and with nonzero \(L^2\) norm. Then Proposition 2.1 implies
    $$
        \left\|\mathrm{P}_{>\ell_{\gamma}} f_{\star}\right\|_{L^2}^2 = \sum_{j=N_{\ell_{\gamma}}+1}^{\infty} f_j^2 \leq \mu_{\ell_{\gamma}+1}^s \sum_{j=N_{\ell_{\gamma}}+1}^{\infty} \lambda_j^{-s} f_j^2 \overset{\text{Assumption 3}}{\leq} O_d \left(d^{-s(\ell_{\gamma}+1)}\right),
    $$
    thus requiring that $\left\|f_{\star}\right\|_{L^q}^2$ is of order at most $d^{-s(\ell_{\gamma}+1)}$ for any $q=2, 3, \cdots$. Therefore, our result applies to a significantly broader class of large-dimensional learning scenarios where the assumptions behind Eq. (\ref{eqn_prop_restate_misia}) may not hold.

For any regression function in $[\calH]^{s}$ with $s \geq 0$ that satisfies all conditions in Proposition \ref{prop_restate_misia}, Eq. (\ref{eqn_prop_restate_misia}) is consistent with Eq. (21). To illustrate this, first notice that we have $\tilde{\lambda} = \Theta_d(\lambda)$ and $\mathrm{var}(\tilde{\xi}_j) = \Theta_d(\mathrm{var}(\xi_j))$: Solving (\ref{eqn_lambda_star}) we have
    \begin{equation}\label{eqn_mis_implies_ours_1}
        \tilde{\lambda} = 
        \left\{\begin{matrix}
\left[1+\Theta_d(d^{u-\gamma})\right]\lambda = \Theta_d(\lambda), &  \lambda \gg n^{-1}\\
\Theta_d(n^{-1}), & \lambda = O_d(n^{-1}); \\
\end{matrix}\right.
    \end{equation}
and
    \begin{equation}\label{eqn_mis_implies_ours_2}
    \sigma^2 + E_{\tilde{\xi}}\left[\sum_{j=1}^{\infty}(\tilde{f}_j^{\tilde{\lambda}} - f_j)^2\right] = \Theta_d(\sigma^2);
    \end{equation}

Therefore, if (\ref{eqn_prop_restate_misia}) holds, then it implies (21) since:
    \begin{equation*}
        \begin{aligned}
           &~ E_{x, \epsilon} \left[ \left(\hat{f}_{\lambda}(x) - f_{\star}(x) \right)^{2} \right] \overset{(\ref{eqn_prop_restate_misia})}{=}
   \Theta_{d, \mathbb{P}} \left( E_{\tilde{\xi}}\left[\sum_{j=1}^{\infty}(\tilde{f}_j^{\tilde{\lambda}} - f_j)^2\right] \right)\\
   \overset{(19) \text{ and } (\ref{eqn_mis_implies_ours_2})}{=} &~
   \Theta_{d, \mathbb{P}} \left(E_{\xi}\left[\sum_{j=1}^{\infty}(\hat{f}_j^{\tilde{\lambda}} - f_j)^2\right]\right)
   \overset{(20) \text{ and } (\ref{eqn_mis_implies_ours_1})}{=} 
   \Theta_{d, \mathbb{P}} \left(
E_{\xi}\left[\sum_{j=1}^{\infty}(\hat{f}_j^{\max\{\lambda, n^{-1}\}} - f_j)^2\right]
\right).
        \end{aligned}
    \end{equation*}





In summary, Theorem 3.1 strictly generalizes Proposition \ref{prop_restate_misia},
covering the cases for any $s \geq 0$ and for certain analytical spectral algorithms including KRR.




\begin{remark}
    The assumptions in Proposition \ref{prop_restate_misia} are pivotal to its proof technique and are intrinsically difficult to relax. Notice that the two assumptions,
$$
\left\|\mathrm{P}_{>\ell_{\gamma}} f_{\star}\right\|_{L^2} \geq\left\|f_{\star}\right\|_{L^2} / C 
\quad \text{ and } \quad 
\left\|f_{\star}\right\|_{L^q} \leq (C q)^{(\ell_{\gamma}+1) / 2}\left\|f_{\star}\right\|_{L^2} \text{ for all integers } q \geq 2,
$$
are sufficient conditions for the convex concentration on eigenfunctions in \cite{cheng2022dimension}, which allows them to apply the Hanson-Wright inequality in Lemma 2.1 of \cite{cheng2022dimension}.
In contrast, our proof employs a different analytical framework that does not rely on those restrictive premises.
\end{remark}









\section{Proof sketch for Theorem 3.1 in a representative regime}\label{append_sketch_proof_main_results}



To highlight the main ideas without carrying the full case analysis, we focus on the case where
\begin{equation}\label{eqn_sketch_0}
    \begin{aligned}
    &~ u \in [u^{\prime}(1), \gamma) \text{ (that is, the under-regularized regime), }\\
    &~ \ell_{\lambda}+1 < \ell_{\gamma} < \gamma, \quad s>0, \quad \text{ and } \tau < \infty
\end{aligned}
\end{equation}

This case already contains the main additional difficulty in proving the full learning curve.
Indeed, once $\ell_{\lambda}<\ell_{\gamma}$, the proof must be organized at the cutoff $\ell_{\lambda}$ rather than only at $\ell_{\gamma}$.  
This creates the intermediate block between degrees $\ell_{\lambda}+1$ and $\ell_{\gamma}$, whose contribution must be tracked explicitly, especially in the bias analysis. 
This is the key new obstacle that does not appear in the same form in interpolation-level arguments or in proofs restricted to the optimally tuned regime.



The remaining cases are handled by the same general strategy, together with appropriate case-dependent modifications of several known results from the literature.
For example, for the case where $u \in [u^{\prime}(\tau), u^{\prime}(1))$, the bias bound requires a variant of the argument in Appendix D.4.2 of \cite{lu2024saturation}; see Proposition \ref{prop_verify_bias_special}. 
The complete proof treats all cases systematically.


% \hrule 
% The remaining cases can be resolved by adapting the argument below, together with routine generalizations of known results from the literature.
% For example, for the case where $u \in [u^{\prime}(\tau), u^{\prime}(1))$, we need a modification of Appendix D.4.2 in \cite{lu2024saturation} to bound the bias, see Proposition \ref{prop_verify_bias_special}.

% The complete proof in the paper treats all cases systematically.


The proof has three ingredients. 
First, we combine kernel approximation results with a decomposition of $K$ at the cutoff index  $\ell_\lambda$. 
Second, we use this decomposition to bound the variance and the bias separately. 
Third, we compare the resulting exponents to identify the dominant term in each range of $u$. 




\paragraph*{Key kernel approximation results}
We first collect some kernel approximation results, and readers can refer to Appendix \ref{append_kernel_app_sphere} for details.

From \cite{ghorbani2021linearized}, we have
\begin{equation}\label{eqn_sketch_1}
    \left\|\Psi_{>\ell_{\gamma}} \Sigma_{>\ell_{\gamma}}^j \Psi_{>\ell_{\gamma}}^{\top} - \kappa_j\mathrm{I}_n \right\|_{\mathrm{op}} = o_{d, \mathbb{P}}(\kappa_j), \quad \kappa_j = \sum_{k=\ell_{\gamma}+1}^\infty \mu_k^j N(d, k),\quad j=1,2;
\end{equation}
and that
\begin{equation}\label{eqn_sketch_2}
\left\| n^{-1}{\Psi_{\leq \ell_{\gamma}}^{\top} \Psi_{\leq \ell_{\gamma}} } - \mathrm{I}_{N_{\ell_{\gamma}}} \right\|_{\mathrm{op}} = o_{d, \mathbb{P}}(1);
\end{equation}
see Lemma \ref{lemma psi psi top} and Lemma \ref{lemma psi top psi}. 
Recall that we have the decomposition 
$$
K = \underbrace{\Psi_{\leq \ell_{\lambda}} \Sigma_{\leq \ell_{\lambda}} \Psi_{\leq \ell_{\lambda}}^{\top}}_{K_{\leq \ell_{\lambda}}} + 
\underbrace{\Psi_{\ell_{\lambda}+ 1, \ell_{\gamma}} \Sigma_{\ell_{\lambda}+ 1, \ell_{\gamma}} \Psi_{\ell_{\lambda}+ 1, \ell_{\gamma}}}_{K_{\ell_{\lambda}+ 1, \ell_{\gamma}}}
+ \underbrace{\Psi_{>\ell_{\gamma}} \Sigma_{>\ell_{\gamma}} \Psi_{>\ell_{\gamma}}^{\top}}_{K_{>\ell_{\gamma}}}.
$$
Combining the block decomposition of $K$ with Weyl's inequality, we obtain several key properties of the kernel matrices $M := K + n\lambda \mathrm{I}_n$, $K_{>\ell_{\lambda}}$, and $M_{>\ell_{\lambda}}:=K_{>\ell_{\lambda}} + n\lambda \mathrm{I}_n$, which will be used throughout the variance and bias analysis.
Specifically, we have 
\begin{equation}\label{eqn_sketch_3}
\left\{
\begin{aligned}
    &~ \lambda_{\text{min}}\left( M\right) = \Omega_{d, \mathbb{P}}(n\lambda), \quad \mathrm{tr}(M / n)=O_{d, \mathbb{P}}(n\lambda);\\
    &~ \text{For any } N_{\ell_{\lambda}}+1, \cdots, N_{\ell_{\lambda}+1}, \text{ we have } \lambda_i(M) = O_{d, \mathbb{P}}(n\lambda);\\
    &~ \lambda_{\text{max}}(M_{>\ell_{\lambda}}) \asymp \lambda_{\text{min}}(M_{>\ell_{\lambda}}) = \Theta_{d, \mathbb{P}}(n\lambda), \quad \lambda_{\text{max}}(K_{>\ell_{\lambda}}) =
 O_{d, \mathbb{P}}\left({d^{\gamma-\ell_{\lambda}-1}}\right);\\
 &~ \text{Denote } A = n^{-1}\left(\Sigma_{\leq \ell_{\lambda}}^{-1} + \Psi_{\leq \ell_{\lambda}}^{\top} M_{>\ell_{\lambda}}^{-1} \Psi_{\leq \ell_{\lambda}}\right), \text{ then } \lambda_{\text{max}} \left(  A^{-1} \right) \asymp \lambda_{\text{min}} \left(  A^{-1} \right) = \Theta_{d, \mathbb{P}}(n\lambda);
\end{aligned}
\right.
\end{equation}
see Lemma \ref{lemma lambda max min}. 
For the bias, we also need the following results
\begin{equation}\label{eqn_sketch_4}
   \left\{ 
   \begin{aligned}
       & \left\| \Psi_{>\ell_{\lambda}} \theta_{>\ell_{\lambda}} \right\|_{2}^{2} = ~ O_{d, \mathbb{P}}\left( n  \| \theta_{>\ell_{\lambda}} \|_{2}^{2}\right); \\
       & \| \theta_{>\ell_{\lambda}} \|_{2}^{2}  = ~ \Theta_d \left(  d^{-(\ell_{\lambda}+1)s} \right); \\
      &  \left\| \Sigma_{\leq \ell_{\lambda}}^{-\tau}  \theta_{\leq \ell_{\lambda}} \right\|_{2}^{2}  = ~ \Theta_{d} \left(  d^{(2\tau-\min\{s, 2\tau\})\ell_{\lambda} } \right);
    \end{aligned}
    \right.
\end{equation}
see Lemma \ref{lemma ge l l2}. Finally, based on the Sherman-Morrison-Woodbury formula, for $k \in \{\ell_{\gamma}, \ell_{\lambda}\}$, we have 
\begin{equation}\label{eqn_sketch_5}
\begin{aligned}
    M^{-1} \Psi_{\leq k} \Sigma_{\leq k}^{1/2} 
        =  
        M_{>k}^{-1} \Psi_{\leq k} \Sigma_{\leq k}^{1/2} \left( \mathrm{I}_{N_{k}} +  \Sigma_{\leq k}^{1/2}  \Psi_{\leq k}^{\top} M_{>k}^{-1} \Psi_{\leq k} \Sigma_{\leq k}^{1/2}\right)^{-1},\\
        \Sigma_{\leq k} - \Sigma_{\leq k} \Psi^{\top}_{\leq k} M^{-1} \Psi_{\leq k}\Sigma_{\leq k}
        = 
        \left( \Sigma_{\leq k}^{-1} + \Psi_{\leq k}^{\top} M_{>k}^{-1} \Psi_{\leq k}  \right)^{-1};
\end{aligned}
\end{equation}
see Lemma \ref{lemma trans in V2}. 



\paragraph*{Bounding the variance}

Since $\mathscr{K}(x,X) = \psi(x)^{\top} \Sigma \Psi^{\top}$, $\varphi_{\lambda, K}:=n^{-1}\reg(K/n)$, and $\epsilon_1, \cdots, \epsilon_n$ are independent with variance $\sigma^2$, 
the variance term (14) can be decomposed according to the three spectral blocks
\[
>\ell_\gamma,\qquad \ell_\lambda+1,\ldots,\ell_\gamma,\qquad \le \ell_\lambda.
\]
More precisely,
\begin{equation*}
%\label{eq var tr}
    \begin{aligned}
        \mathrm{var}\left( \hat{f}_{\lambda} \right) &= E_{x, \epsilon}\left[\left(\mathscr{K}(x,X) \varphi_{\lambda, K} \epsilon \right)^2\right] = E_{x, \epsilon} \left[\left( \psi(x)^{\top} \Sigma^{1/2} \Sigma^{1/2} \Psi^{\top} \varphi_{\lambda, K} \epsilon \right)^2\right]  \\
    =&~  \sigma^{2} \cdot \mathrm{tr} \left( \varphi_{\lambda, K} \Psi \Sigma^{2} \Psi^{\top} \varphi_{\lambda, K} \right)\\
    =&~
    \sigma^{2} \cdot \mathrm{tr} \underbrace{\left(\varphi_{\lambda, K} \Psi_{> \ell_{\gamma}} \Sigma_{> \ell_{\gamma}}^{2} \Psi_{> \ell_{\gamma}}^{\top} \varphi_{\lambda, K}\right)}_{ V_{1}}\\
    &~+
    \sigma^{2} \cdot \mathrm{tr} \underbrace{\left(  \Psi_{\ell_{\lambda}+1, \ell_{\gamma}} \Sigma_{\ell_{\lambda}+1, \ell_{\gamma}}^{2} \Psi_{\ell_{\lambda}+1, \ell_{\gamma}}^{\top} \varphi_{\lambda, K}^2 \right)}_{V_{2, 1}}
    ~+~
    \sigma^{2} \cdot \mathrm{tr} \underbrace{\left(  \Sigma_{\leq \ell_{\lambda}} \Psi_{\leq \ell_{\lambda}}^{\top}  \varphi_{\lambda, K}^2 \Psi_{\leq \ell_{\lambda}} \Sigma_{\leq \ell_{\lambda}} \right)}_{V_{2, 2}}.
    \end{aligned}
\end{equation*}


We can then bound these three trace terms using (\ref{eqn_sketch_1}), (\ref{eqn_sketch_2}), (\ref{eqn_sketch_3}), and (\ref{eqn_sketch_5}); see Appendix \ref{append_variance}. 

To illustrate the new step, we take $\mathrm{tr}(V_{2,1})$ as an example. 
This term results from the intermediate block between $\ell_\lambda$ and $\ell_\gamma$, which appears only when $\ell_\lambda<\ell_\gamma$. Controlling this block is what allows the variance analysis to extend beyond the interpolation-level decomposition.


\iffalse
\color{red}
$\mathrm{tr}(V_{2,1})$ can be bounded from above and from below as follows: 
\begin{equation*}
    \begin{aligned}
        \mathrm{tr} \left( V_{2, 1} \right)
        \overset{(\ref{eqn:spe_properties})}{\leq} &~
        C_2^2 n \lambda_{\text{max}}(M^{-2})
        \cdot \mathrm{tr} \left( \frac{\Psi_{\ell_{\lambda}+1, \ell_{\gamma}} \Sigma_{\ell_{\lambda}+1, \ell_{\gamma}}^{2} \Psi_{\ell_{\lambda}+1, \ell_{\gamma}}^{\top}}{n} \right) \overset{*}{=} O_{d, \mathbb{P}}\left(\frac{1}{d^{\ell_{\lambda}+1}n\lambda^2}\right),
    \end{aligned}
\end{equation*}
and
\begin{equation*}
    \begin{aligned}
        \mathrm{tr} \left( V_{2, 1} \right) 
        \geq &~
        \mu_{\ell_{\lambda}+1}^{2} \mathrm{tr} \left(  \Psi_{\ell_{\lambda}+1}  \Psi_{\ell_{\lambda}+1}^{\top} \varphi_{\lambda, K}^2 \right) \geq \mu_{\ell_{\lambda}+1}^{2} \sum_{i=1}^{n} \lambda_{i}(\Psi_{\ell_{\lambda}+1}  \Psi_{\ell_{\lambda}+1}^{\top}) \lambda_{n-i+1}(\varphi_{\lambda, K}^2)\\
        \overset{*}{\geq} &~ 
        \Omega_{d, \mathbb{P}}(\mu_{\ell_{\lambda}+1}^{2}) \sum_{i=N_{\ell_{\lambda}} + 1}^{N_{\ell_{\lambda}+1}} \lambda_{n-i+1}(\varphi_{\lambda, K}^2) \overset{(\ref{eqn:spe_properties})}{\geq} \Omega_{d, \mathbb{P}}(\mu_{\ell_{\lambda}+1}^{2}) \sum_{i=N_{\ell_{\lambda}} + 1}^{N_{\ell_{\lambda}+1}} \lambda_{i}(M^{-2})\\
        \overset{*}{=} &~ \Omega_{d, \mathbb{P}}\left(\frac{1}{d^{\ell_{\lambda}+1}n\lambda^2}\right),
    \end{aligned}
\end{equation*}
where each step marked by $*$ uses only the estimates (\ref{eqn_sketch_1})-(\ref{eqn_sketch_3}), and (\ref{eqn_sketch_5}). 
\fi 
\color{black}



$\mathrm{tr}(V_{2,1})$ can be bounded from above as
\begin{equation*}
    \begin{aligned}
        \mathrm{tr}(V_{2,1})
        \overset{(\ref{eqn:spe_properties})}{\le}
        C_2^2 n \lambda_{\max}(M^{-2})
        \cdot
        \mathrm{tr}\!\left(
        \frac{\Psi_{\ell_\lambda+1,\ell_\gamma}\Sigma_{\ell_\lambda+1,\ell_\gamma}^2
        \Psi_{\ell_\lambda+1,\ell_\gamma}^\top}{n}
        \right)
        \overset{(*)}{=}
        O_{d,\mathbb P}\!\left(\frac{1}{d^{\ell_\lambda+1}n\lambda^2}\right), 
    \end{aligned}
\end{equation*}
where the $(*)$ step uses the estimates in (\ref{eqn_sketch_1})-(\ref{eqn_sketch_3}), and (\ref{eqn_sketch_5}). 
These same estimates are used repeatedly in the proof of the lower bound as well. 
Combining the upper and lower bounds, we conclude that 
$
\mathrm{tr}(V_{2,1})
=
\Theta_{d,\mathbb P}\!\left(\frac{1}{d^{\ell_\lambda+1}n\lambda^2}\right)$. 






\paragraph*{Bounding the bias} 


To keep the sketch focused, we do not reproduce the full decomposition of $\mathrm{bias}^2(\hat f_\lambda)$ into the five terms $B_1,\ldots,B_5$; see Appendix \ref{append_bias}. 
Most terms can be controlled using the same approximation and estimates in (\ref{eqn_sketch_1})--(\ref{eqn_sketch_5}); see Appendix \ref{append_bias}. 
A genuinely new argument is needed for the term
\[
B_{1}:=\Sigma_{\leq \ell_{\lambda}} \Psi^{\top}_{\leq \ell_{\lambda}} \varphi_{\lambda, K} \Psi_{\leq \ell_{\lambda}} \theta_{\leq \ell_{\lambda}} - \theta_{\leq \ell_{\lambda}},
\]
which is the contribution associated with the low-degree block.
Because of the appearance of $\varphi_{\lambda, K}$, this term involves both the intermediate block $K_{\ell_\lambda+1,\ell_\gamma}$ and the low-degree feature matrix $\Psi_{\leq \ell_{\lambda}}$. 
This interaction is absent in interpolation-only analyses organized at $\ell_\gamma$. 
Controlling this term is one of the main new steps in the proof.


Using the approximation (\ref{eqn_sketch_2}) and an algebraic identity in Proposition \ref{prop_matrix_prop_1_of_filter}, we can obtain
\begin{equation*}
    \begin{aligned}
        \|B_1\|_{2}^{2}   \leq  &~ C \left(
        \underbrace{\left\| \rem(K/n)  
                Z\theta_{\leq \ell_{\lambda}} \right\|_{2}^{2}}_{B_{1,3}}~+
                \underbrace{\left\| K_{\ell_{\lambda}+1, \ell_{\gamma}} K^{-1} Z \theta_{\leq \ell_{\lambda}} \right\|_{2}^{2}}_{B_{1,4}}~+
                \underbrace{\left\| K_{>\ell_{\gamma}} K^{-1} Z \theta_{\leq \ell_{\lambda}} \right\|_{2}^{2}}_{B_{1,5}}\right), 
    \end{aligned}
\end{equation*}
where $C=O_{d, \mathbb{P}}\left(1\right)$; see (\ref{eqn_b1_spe_origin}) and (\ref{eq:filter_implication}) for detailed calculations.
The term $B_{1,4}$ is the most delicate since it involves the interaction matrix $K_{\ell_{\lambda}+1, \ell_{\gamma}}$. 



The key step is to convert $B_{1,4}$ into a cross-block Gram term. 
Using $K^{-1}K_{\le \ell_\lambda}=\mathrm{I}-K^{-1}K_{>\ell_\lambda}$, separating the harmless $K_{>\ell_\gamma}$ contribution, and writing
\[
K_{\ell_\lambda+1,\ell_\gamma}
=
n\tilde Z\Sigma_{\ell_\lambda+1,\ell_\gamma}\tilde Z^\top,
\qquad
\tilde Z:=n^{-1/2}\Psi_{\ell_\lambda+1,\ell_\gamma},
\]
we obtain
\[
B_{1,4}
\le
O_{d,\mathbb P}\!\left(
d^{(2-\min\{s,2\})\ell_\lambda-2\gamma}
\left[
n\left\|   \underbrace{Z^\top \tilde Z\,\Sigma_{\ell_\lambda+1,\ell_\gamma}^2\,\tilde Z^\top Z}_{B_{1,6}}   \right\|_2
+1
\right]
\right).
\]



This form is useful because it isolates the interaction between the low-degree and intermediate-degree feature matrices in a form that is amenable to concentration. 
In particular, $B_{1, 6}$ is identified as a sum of independent random rank‑one matrices, and a standard matrix Bernstein inequality shows that $B_{1, 6}=o_{d,\mathbb{P}}(n^{-1})$.
This is the crucial estimate behind the control of $B_{1,4}$ and the point where the proof goes beyond interpolation-only arguments.



\iffalse
Specifically, let $z_i^\top$ and $\tilde z_i^\top$ denote the $i$th rows of $Z$ and $\tilde Z$, respectively. Then
\[
Z^\top \tilde Z\,\Sigma_{\ell_\lambda+1,\ell_\gamma}
=
\sum_{i=1}^n z_i\bigl(\Sigma_{\ell_\lambda+1,\ell_\gamma}\tilde z_i\bigr)^\top .
\]
Each summand is an independent random rank-one matrix with mean zero, since it depends only on $x_i$ and couples two orthogonal degree blocks. Therefore a matrix Bernstein inequality gives
\[
\left\|Z^\top \tilde Z\,\Sigma_{\ell_\lambda+1,\ell_\gamma}\right\|_2
=
o_{d,\mathbb P}(n^{-1/2}),
\]
and hence
\[
B_{1,6}=o_{d,\mathbb P}(n^{-1}).
\]
This is the crucial estimate behind the control of $B_{1,4}$.
This is the point where the proof goes beyond interpolation-only arguments.
\fi 


\iffalse
\color{red}
\hrule 
% \iffalse [2026-04-18]
The bias term can be broken down as
\begin{equation*}
    \begin{aligned}
        \mathrm{bias}^2(\hat{f}_{\lambda})
        =&~
    \| \underbrace{\Sigma_{\leq \ell_{\lambda}} \Psi^{\top}_{\leq \ell_{\lambda}} \varphi_{\lambda, K} \Psi_{\leq \ell_{\lambda}} \theta_{\leq \ell_{\lambda}} - \theta_{\leq \ell_{\lambda}}  }_{B_{1}}+  \underbrace{\Sigma_{\leq \ell_{\lambda}} \Psi^{\top}_{\leq \ell_{\lambda}} \varphi_{\lambda, K} \Psi_{>\ell_{\lambda}} \theta_{>\ell_{\lambda}}}_{B_{2}} \|_{2}^{2}\\
    &~+
     \|\underbrace{\theta_{>\ell_{\lambda}}}_{B_{3}} - \underbrace{ \Sigma_{>\ell_{\lambda}} \Psi^{\top}_{>\ell_{\lambda}} \varphi_{\lambda, K} \Psi_{\leq \ell_{\lambda}} \theta_{\leq \ell_{\lambda}} }_{B_{4}} - \underbrace{  \Sigma_{>\ell_{\lambda}} \Psi^{\top}_{>\ell_{\lambda}} \varphi_{\lambda, K} \Psi_{>\ell_{\lambda}} \theta_{>\ell_{\lambda}} }_{ B_{5}}\|_{2}^{2}.
    \end{aligned}
\end{equation*}


% Our proof in Appendix \ref{append_bias} involves several non‑trivial observations. We sketch the upper bound for $\| B_{1} \|_{2}^{2}$ as an illustration. 



We then bound these five terms using (\ref{eqn_sketch_1}), (\ref{eqn_sketch_2}), (\ref{eqn_sketch_3}), (\ref{eqn_sketch_4}), and (\ref{eqn_sketch_5}); see Appendix \ref{append_bias}. 
We illustrate the novel step through the upper bound for $\|B_1\|_2^2$. 

Recall the identity
\begin{equation*}
    \mathrm{I}-A\reg(A+B) = A(A+B)^{-1}\rem(A+B) + B(A+B)^{-1},
\end{equation*}
for positive semi-definite matrices $A$ and $B$;
see Proposition \ref{prop_matrix_prop_1_of_filter}.
Using this identity with $A=K_{\leq \ell_{\lambda}} / n$ and $B=K_{> \ell_{\lambda}} / n$ and the shorthand $Z = n^{-1/2}\Psi_{\leq \ell_{\lambda}}$, we rewrite and bound $\| B_{1} \|_{2}^{2}$ as
\begin{equation*}
    \begin{aligned}
        \|B_1\|_{2}^{2} \overset{(\ref{eqn_sketch_2})}{=} &~ 
        \Theta_{d, \mathbb{P}}\left(\frac{1}{n}\right)
                \left\| \left( \mathrm{I}_{n} - \Psi_{\leq \ell_{\lambda}}\Sigma_{\leq \ell_{\lambda}} \Psi^{\top}_{\leq \ell_{\lambda}} \varphi_{\lambda, K} \right) 
                \Psi_{\leq \ell_{\lambda}}\theta_{\leq \ell_{\lambda}} \right\|_{2}^{2}\\
                = &~
                \Theta_{d, \mathbb{P}}\left(1\right) \left\| \left( \mathrm{I}_{n} - A \reg(A+B)\right) Z\theta_{\leq \ell_{\lambda}}\right\|_{2}^{2}
                \\
            (\because \|A(A+B)^{-1}\|_2 \leq 1) \quad   \leq  &~  O_{d, \mathbb{P}}\left(1\right)
        \underbrace{\left\| \rem(K/n)  
                Z\theta_{\leq \ell_{\lambda}} \right\|_{2}^{2}}_{B_{1,3}}\\ 
                &~+
                O_{d, \mathbb{P}}\left(  1\right)
                \underbrace{\left\| K_{\ell_{\lambda}+1, \ell_{\gamma}} K^{-1} Z \theta_{\leq \ell_{\lambda}} \right\|_{2}^{2}}_{B_{1,4}}\\
                &~+
                O_{d, \mathbb{P}}\left(  1\right)
                \underbrace{\left\| K_{>\ell_{\gamma}} K^{-1} Z \theta_{\leq \ell_{\lambda}} \right\|_{2}^{2}}_{B_{1,5}};
    \end{aligned}
\end{equation*}
see (\ref{eqn_b1_spe_origin}) and (\ref{eq:filter_implication}) for detailed calculations.
The term $B_{1,4}$ is the most delicate since it involves the interaction matrix $K_{\ell_{\lambda}+1, \ell_{\gamma}}$. 



The non‑trivial step is to recognize the following algebraic transformation:
$$
B_{1,4} \leq O_{d, \mathbb{P}}\left( d^{(2-\min\{s, 2\})\ell_{\lambda}-2\gamma } \left[ n\underbrace{\left\| Z^{\top} \tilde{Z} \Sigma_{\ell_{\lambda}+1, \ell_{\gamma}}^2 \tilde{Z}^{\top} Z  \right\|_{2}}_{B_{1,6}}+1 \right]\right),
$$
where $\tilde{Z}=n^{-1/2}\Psi_{\ell_{\lambda}+1, \ell_{\gamma}}$. 
With $B_{1, 6}$ identified as a sum of independent random rank‑one matrices, a standard matrix Bernstein inequality shows that $B_{1, 6}=o_{d,\mathbb{P}}(n^{-1})$.
\color{black}

\fi




\paragraph*{Final comparison of the bias and variance rates}

For the regime in \eqref{eqn_sketch_0}, the previous steps show that the variance contributes two candidate orders, while the bias contributes three candidate orders. 
For notational convenience, we denote the corresponding exponents by $v_1(u),v_2(u)$ and $b_1(u),b_2(u),b_3(u)$, respectively; see \eqref{eqn_order_in_bias_and_variance}. 
Thus
\begin{equation*}
    \begin{aligned}
        \mathrm{var}(\hat{f}_{\lambda})
        &= \Theta_{d,\mathbb P}\!\left(d^{v_1(u)}+d^{v_2(u)}\right),\\
        \mathrm{bias}^{2}(\hat{f}_{\lambda})
        &\le O_{d,\mathbb P}\!\left(d^{b_1(u)}+d^{b_2(u)}+d^{b_3(u)}\right).
    \end{aligned}
\end{equation*}
Moreover, the lower bound proved in Theorem~\ref{thm_bias_spe} shows that
\begin{equation*}
    \mathrm{bias}^{2}(\hat{f}_{\lambda})
    \ge \Omega_{d,\mathbb P}\!\left(d^{b_1(u)}\right)
    \qquad\text{whenever } b_1(u)\ge \max\{b_2(u),b_3(u)\}.
\end{equation*}

Therefore, the only remaining step is to compare these exponents. 
Proposition \ref{prop_e_2} and Proposition \ref{prop_e_3} show that, in the present regime,
\begin{equation*}
    \max\{v_1(u),v_2(u),b_1(u)\}\ge \max\{b_2(u),b_3(u)\}.
\end{equation*}
In other words, the secondary bias exponents $b_2(u)$ and $b_3(u)$ never dominate the leading order. 
Hence the excess risk is determined by the two variance exponents together with the principal bias exponent $b_1(u)$, and we obtain
\begin{equation*}
    E_{x,\epsilon}\!\left[\left(\hat f_{\lambda}(x)-f_{\star}(x)\right)^2\right]
    = \Theta_{d,\mathbb P}\!\left(d^{v_1(u)}+d^{v_2(u)}+d^{b_1(u)}\right).
\end{equation*}
This is the final synthesis step that turns the separate bias and variance bounds into the learning-curve formula in the regime \eqref{eqn_sketch_0}; see Appendix \ref{append_learning_curve} for the complete case-by-case comparison.


\iffalse
\color{red}
For the specific regime considered in (\ref{eqn_sketch_0}), we have shown the bounds for the variance and bias:
\begin{equation*}
    \begin{aligned}
        \mathrm{var}(\hat{f}_{\lambda}) = &~ \Theta_{d, \mathbb{P}}\left( d^{v_1(u)} + d^{v_2(u)} \right)\\
        \mathrm{bias}^{2}(\hat{f}_{\lambda}) \leq &~ O_{d, \mathbb{P}}\left(d^{b_1(u)} + d^{b_2(u)}+d^{b_3}\right)\\
        \mathrm{bias}^{2}(\hat{f}_{\lambda}) \geq &~ \Theta_{d, \mathbb{P}} \left( d^{b_1(u)}\right), \text{ if } b_1(u) \geq \max\{b_2(u), b_3(u)\};
    \end{aligned}
\end{equation*}
where the exponents $v_i(u)$ and $b_i(u)$ are defined in (\ref{eqn_order_in_bias_and_variance}).

We then show the following comparison in Proposition \ref{prop_e_2} and Proposition \ref{prop_e_3}:
$$
    \max\{v_1(u), v_2(u), b_1(u)\} \geq \max\{b_2(u), b_3(u)\}.
$$

Consequently, we obtain the tight bound for the excess risk:
$$
E_{x, \epsilon} \left[ \left(\hat{f}_{\lambda}(x) - f_{\star}(x) \right)^{2} \right]
= 
\mathrm{var}(\hat{f}_{\lambda}) + \mathrm{bias}^{2}(\hat{f}_{\lambda})
=
\Theta_{d, \mathbb{P}}(d^{v_1(u)} + d^{v_2(u)} + d^{b_1(u)});
$$
see Appendix \ref{append_learning_curve} for details.
\color{black}
\fi










\section{Variance}\label{append_variance}

In this section, our goal is to prove the following theorem, which bounds the variance term in (13).

\begin{theorem}\label{thm_variance_spe}
    Suppose Assumptions 1, 2, and 4 hold.
    Then the variance term (14) of the spectral algorithm estimator $\hat{f}_{\lambda}$ satisfies
    \begin{displaymath}
        \mathrm{var}(\hat{f}_{\lambda}) = \Theta_{d, \mathbb{P}} \left( 
        d^{\gamma - \tilde{\ell} - 1 - 2  \max\{\gamma - u, 0\}}
+ d^{\tilde{\ell} - \gamma}
        %\frac{n}{d^{\tilde{\ell}+1}(n\lambda+1)^{2}} + \frac{d^{\tilde{\ell}}}{n} 
        \right).
    \end{displaymath} 
\end{theorem}


\begin{proof}
Because $\mathscr{K}(x,X) = \psi(x)^{\top} \Sigma \Psi^{\top}$ and that $\epsilon_1, \cdots, \epsilon_n$ are independent with variance $\sigma^2$, the variance term (14) can be expressed as
\begin{equation*}
%\label{eq var tr}
    \begin{aligned}
        \mathrm{var}\left( \hat{f}_{\lambda} \right) &= E_{x, \epsilon}\left[\left(\mathscr{K}(x,X) \varphi_{\lambda, K} \epsilon \right)^2\right] = E_{x, \epsilon} \left[\left( \psi(x)^{\top} \Sigma \Psi^{\top} \varphi_{\lambda, K} \epsilon \right)^2\right]  \\
    &=  \sigma^{2} \cdot \mathrm{tr} \left( \varphi_{\lambda, K} \Psi \Sigma^{2} \Psi^{\top} \varphi_{\lambda, K} \right)\\
    (\because \text{Lemma~\ref{lemma psi psi top}})  &=~
    \sigma^{2} \cdot \mathrm{tr} \underbrace{\left(\kappa_{2}  \varphi_{\lambda, K} \left(\mathrm{I}_{n} + \Delta_{2}\right) \varphi_{\lambda, K}\right)}_{ V_{1}}  ~+~ \sigma^{2} \cdot \mathrm{tr} \underbrace{\left( \varphi_{\lambda, K} \Psi_{\leq \ell_{\gamma}} \Sigma_{\leq \ell_{\gamma}}^{2} \Psi_{\leq \ell_{\gamma}}^{\top} \varphi_{\lambda, K}\right)}_{ V_{2}},
    \end{aligned}
\end{equation*}
\iffalse
From Lemma \ref{lemma psi psi top}, we can decompose Eq.~\eqref{eq var tr} as
\begin{equation}\label{eq var tr decomposition}
    \mathrm{var}\left( \hat{f}_{\lambda} \right) / \sigma^{2} = \mathrm{tr} \underbrace{\left(\kappa_{2}  \varphi_{\lambda, K} \left(\mathrm{I}_{n} + \Delta_{2}\right) \varphi_{\lambda, K}\right)}_{ V_{1}}  ~+~ \mathrm{tr} \underbrace{\left( \varphi_{\lambda, K} \Psi_{\leq \ell_{\gamma}} \Sigma_{\leq \ell_{\gamma}}^{2} \Psi_{\leq \ell_{\gamma}}^{\top} \varphi_{\lambda, K}\right)}_{ V_{2}},
\end{equation}
\fi
where $ \kappa_{2} = \Theta_{d}\left( d^{-(\ell_{\gamma}+1)} \right)$ and $\|\Delta_{2}\|_{\mathrm{op}} = o_{d, \mathbb{P}}(1)$. We then bound $\mathrm{tr}(V_1)$ and $\mathrm{tr}(V_2)$ separately.




\noindent {\bf Bounding $\mathrm{tr}(V_1)$. }
On the one hand, we have
\begin{equation*}
    \begin{aligned}
        \mathrm{tr}(V_1) \leq &~ n\lambda_{\text{max}}\left( V_{1} \right) = n\kappa_{2} \cdot \lambda_{\text{max}}\left( \varphi_{\lambda, K} \left( \mathrm{I}_{n} + \Delta_{2} \right)  \varphi_{\lambda, K}\right)  \\
        (\because \text{Prop.~\ref{prop_weyl_ine}})\quad \leq  &~n\kappa_{2} \cdot \left[ \lambda_{\text{max}}\left( \varphi_{\lambda, K}^2\right) + \lambda_{\text{max}}\left( \varphi_{\lambda, K} \Delta_{2} \varphi_{\lambda, K}\right)  \right]  \\
        \leq  &~ n\kappa_{2} \cdot \lambda_{\text{max}}\left( \varphi_{\lambda, K}^2\right) \left( 1 + o_{d, \mathbb{P}}(1) \right)\\
        \overset{(\ref{eqn:spe_properties})}{\leq}  &~ C_2^2 n\kappa_{2} \cdot \lambda_{\text{min}}\left( M\right)^{-2} \left( 1 + o_{d, \mathbb{P}}(1) \right)\\
        (\because \text{Lemma~\ref{lemma lambda max min}}) \quad =&~
        O_{d, \mathbb{P}}\left(\frac{n}{d^{\ell_{\gamma}+1}(n\lambda+1)^2}\right);
    \end{aligned}
\end{equation*}
On the other hand, we have
\begin{equation*}
    \begin{aligned}
        \mathrm{tr}(V_1) \geq &~  \kappa_{2} \cdot \mathrm{tr}\left( \varphi_{\lambda, K}^2\right) \left( 1 - o_{d, \mathbb{P}}(1) \right)\\
        \overset{(\ref{eqn:spe_properties})}{\geq}&~
        C_1^2\kappa_{2} \cdot \mathrm{tr}\left( M^{-2}\right) \left( 1 - o_{d, \mathbb{P}}(1) \right)\\
       \overset{(\ast)}{ \geq} &~  C_1^2 \kappa_{2} \cdot n \left[ \mathrm{tr}\left( \frac{M}{n} \right) \right]^{-2} \left( 1 -o_{d, \mathbb{P}}(1) \right)\\
                (\because \text{Lemma~\ref{lemma lambda max min}(i)})\quad =&~
        \Omega_{d, \mathbb{P}}\left(\frac{n}{d^{\ell_{\gamma}+1}(n\lambda+1)^2}\right),
    \end{aligned}
\end{equation*}
where the third inequality $(\ast)$ can be shown using the Jensen's inequality as in Remark~\ref{rem:variance-Jensen-inequ}. 
Therefore, $\mathrm{tr}(V_1) = \Theta_{d, \mathbb{P}}\left(\frac{n}{d^{\ell_{\gamma}+1}(n\lambda+1)^2}\right)$.

\begin{remark}\label{rem:variance-Jensen-inequ}
    Suppose $A = \text{diag}(a_1, \cdots, a_n) \in \mathbb{R}^{n \times n}$ is a positive diagonal matrix. 
    Then, the Jensen's inequality $f\left( \frac{a_1 + \cdots + a_n}{n} \right) \leq \frac{f(x_1) + \cdots + f(x_n)}{n}$ for convex function \( f(x) = x^{-2}, x>0 \) implies
    $$
    \left( \text{tr}\left( \frac{A}{n} \right) \right)^{-2} = \left( \frac{a_1 + \cdots + a_n}{n} \right)^{-2} \leq \frac{a_1^{-2} + \cdots + a_n^{-2}}{n} = \frac{\text{tr}(A^{-2})}{n}.
    $$
Since $M$ is positive definite, we can show $(\ast)$ with eigendecomposition.
\end{remark}


\noindent {\bf Bounding $\mathrm{tr}(V_2)$. }
We bound $\mathrm{tr}(V_2)$ in two cases: (i) $\ell_{\gamma} \leq \ell_{\lambda}$, and (ii) $\ell_{\gamma} > \ell_{\lambda}$.

\noindent (i) On the one hand, when $\ell_{\gamma} \leq \ell_{\lambda}$, we have
\begin{equation}\label{eqn_v_2_large_regular_para}
    \begin{aligned}
        \mathrm{tr} \left( V_{2} \right)
        =&~
        \mathrm{tr} \left( \Sigma_{\leq \ell_{\gamma}} \Psi_{\leq \ell_{\gamma}}^{\top}  \varphi_{\lambda, K}^2 \Psi_{\leq \ell_{\gamma}} \Sigma_{\leq \ell_{\gamma}} \right)
        \overset{(\ref{eqn:spe_properties})}{=} \Theta_d\left(1\right) \cdot
        \mathrm{tr} \left( \Sigma_{\leq \ell_{\gamma}} \Psi_{\leq \ell_{\gamma}}^{\top}  M^{-2} \Psi_{\leq \ell_{\gamma}} \Sigma_{\leq \ell_{\gamma}} \right)\\
        \overset{\text{Lemma } \ref{lemma trans in V2}}{=}&~
        \Theta_d\left(\frac{1}{n}\right) \cdot \mathrm{tr} \left[ \left( \frac{\Sigma_{\leq \ell_{\gamma}}^{-1}}{n} + \frac{\Psi_{\leq \ell_{\gamma}}^{\top} M_{>\ell_{\gamma}}^{-1} \Psi_{\leq \ell_{\gamma}}}{n}   \right)^{-1} \frac{\Psi_{\leq \ell_{\gamma}}^{\top} M_{>\ell_{\gamma}}^{-2} \Psi_{\leq \ell_{\gamma}}}{n} \left( \frac{\Sigma_{\leq \ell_{\gamma}}^{-1}}{n} + \frac{\Psi_{\leq \ell_{\gamma}}^{\top} M_{>\ell_{\gamma}}^{-1} \Psi_{\leq \ell_{\gamma}}}{n}   \right)^{-1}  \right]\\
        \overset{\text{Lemma } \ref{lemma lambda max min}}{=}&~
        \Theta_{d, \mathbb{P}}\left(\frac{(n\lambda+1)^2}{n}\right) \cdot \mathrm{tr} \left[  \frac{\Psi_{\leq \ell_{\gamma}}^{\top} M_{>\ell_{\gamma}}^{-2} \Psi_{\leq \ell_{\gamma}}}{n}   \right]\\
        \overset{\text{Lemma } \ref{lemma lambda max min}}{=}&~
        \Theta_{d, \mathbb{P}}\left(\frac{1}{n}\right) \cdot \mathrm{tr} \left[  \frac{\Psi_{\leq \ell_{\gamma}}^{\top} \Psi_{\leq \ell_{\gamma}}}{n}   \right]
        %\overset{{\green \text{Proposition } 2.1, } \text{Lemma } \ref{lemma psi top psi} \text{ and Lemma } \ref{lemma psi top psi_int}}{=}
        =
        \Theta_{d, \mathbb{P}}\left(\frac{d^{\ell_{\gamma}}}{n}\right),
    \end{aligned}
\end{equation}
% {\green
where the last equality follows from Lemma \ref{lemma psi top psi} and \ref{lemma psi top psi_int}, together with the fact that $N_{\ell_{\gamma}}:= \sum_{k=0}^{\ell_{\gamma}} N(d, k) = \Theta(d^{\ell_{\gamma}})$ (see Proposition 2.1).
% }



% \dhnote{The last equation is $n^{-1}B_\ell\leq n^{-1}\sum_{k=0}^\ell N(d,k)\lesssim n^{-1}\sum_{k=0}^\ell d^k / k! \lesssim n^{-1}d^{\ell}$; am I right?}

% {\green [You are right, I added the Proposition 2.1, which shows that $N(d,k) = \Theta(d^k)$, $k=0, \cdots, \ell$]}

\noindent (ii) On the other hand, when ${\ell_{\gamma}} > \ell_{\lambda}$, similar to (\ref{eqn_v_2_large_regular_para}), we have
\begin{equation*}
    \begin{aligned}
        \mathrm{tr} \left( \Sigma_{\leq \ell_{\lambda}} \Psi_{\leq \ell_{\lambda}}^{\top}  \varphi_{\lambda, K}^2 \Psi_{\leq \ell_{\lambda}} \Sigma_{\leq \ell_{\lambda}} \right)
        =
        \Theta_{d, \mathbb{P}}\left(\frac{d^{\ell_{\lambda}}}{n}\right),
    \end{aligned}
\end{equation*}
hence
\begin{equation*}
    \begin{aligned}
    \mathrm{tr} \left( V_{2} \right)
        =&~
        \mathrm{tr} \left( \varphi_{\lambda, K} \Psi_{\ell_{\lambda}+1, \ell_{\gamma}} \Sigma_{\ell_{\lambda}+1, \ell_{\gamma}}^{2} \Psi_{\ell_{\lambda}+1, \ell_{\gamma}}^{\top} \varphi_{\lambda, K} \right)
        +
        \mathrm{tr} \left( \Sigma_{\leq \ell_{\lambda}} \Psi_{\leq \ell_{\lambda}}^{\top}  \varphi_{\lambda, K}^2 \Psi_{\leq \ell_{\lambda}} \Sigma_{\leq \ell_{\lambda}} \right)\\
        =&~
        \mathrm{tr} \underbrace{\left(  \Psi_{\ell_{\lambda}+1, \ell_{\gamma}} \Sigma_{\ell_{\lambda}+1, \ell_{\gamma}}^{2} \Psi_{\ell_{\lambda}+1, \ell_{\gamma}}^{\top} \varphi_{\lambda, K}^2 \right)}_{V_{2, 1}}
        +
        \Theta_{d, \mathbb{P}}\left(\frac{d^{\ell_{\lambda}}}{n}\right).
    \end{aligned}
\end{equation*}
For the upper bound of $\mathrm{tr}(V_{2, 1})$, we have
\begin{equation*}
    \begin{aligned}
        \mathrm{tr} \left( V_{2, 1} \right) 
        \leq &~
        C_2^2 n\lambda_{\text{max}}(M^{-2})
        \cdot \mathrm{tr} \left( \frac{\Psi_{\ell_{\lambda}+1, \ell_{\gamma}} \Sigma_{\ell_{\lambda}+1, \ell_{\gamma}}^{2} \Psi_{\ell_{\lambda}+1, \ell_{\gamma}}^{\top}}{n} \right)\\
        \overset{\text{ Lemma } \ref{lemma lambda max min}}{=} &~
        O_{d, \mathbb{P}}\left(\frac{n}{(n\lambda+1)^2}\right) \cdot \mathrm{tr} \left( \frac{\Psi_{\ell_{\lambda}+1, \ell_{\gamma}} \Sigma_{\ell_{\lambda}+1, \ell_{\gamma}}^{2} \Psi_{\ell_{\lambda}+1, \ell_{\gamma}}^{\top}}{n} \right)\\
        \overset{\text{Lemma } \ref{lemma psi top psi} \text{ and } \ref{lemma psi top psi_int}}{=}&~
        O_{d, \mathbb{P}}\left(\frac{n}{(n\lambda+1)^2}\right) \cdot \mathrm{tr} \left( \Sigma_{\ell_{\lambda}+1, \ell_{\gamma}}^{2} \right)\\
        =&~
        O_{d, \mathbb{P}}\left(\frac{n}{d^{\ell_{\lambda}+1}(n\lambda+1)^2}\right).
    \end{aligned}
\end{equation*}
Next, we give the lower bound of $\mathrm{tr}(V_{2, 1})$.
From Lemma \ref{lemma psi top psi} and \ref{lemma psi top psi_int}, we can denote the SVD decomposition of $ \Psi_{\ell_{\lambda}+1} $ as $ \Psi_{\ell_{\lambda}+1} = n^{1/2} ~ O H V^{\top} $, where $ O \in \mathbb{R}^{n \times n}$ and $ V \in \mathbb{R}^{N(d, \ell_{\lambda}+1) \times N(d, \ell_{\lambda}+1)}$ are orthogonal matrices; and $HH^{\top}=\text{Diag}\{h_1^2 \geq \cdots \geq h_n^2 \geq 0\}$ is a diagonal matrix with $h_1^2=O_{d, \mathbb{P}}(1)$ and 

\begin{equation}\label{eqn_lower_bound_h_in_variance}
    \begin{aligned}
        h_{N(d, \ell_{\lambda}+1)}^2=\Omega_{d, \mathbb{P}}(1), \quad \text{ if } n > N(d, \ell_{\lambda}+1);\\
        h_{n}^2=\Omega_{d, \mathbb{P}}(1), \quad \text{ if } n \leq N(d, \ell_{\lambda}+1).
    \end{aligned}
\end{equation}

%$h_{\lfloor\min\{n, N(d, \ell_{\lambda}+1)\}/2\rfloor}^2=\Omega_{d, \mathbb{P}}(1)$. 
Denote the SVD decomposition of $\varphi_{\lambda, K}^2$ as $\varphi_{\lambda, K}^2=PDP^{\top}$, where $D=\text{Diag}\{d_1 \geq \cdots \geq d_n > 0\}$ is a diagonal matrix. 

Since $\ell_{\gamma} = \lfloor \gamma \rfloor > \ell_{\lambda}$, $N_{\ell_{\lambda}} = \Theta(d^{\ell_{\lambda}})$, and $N(d, \ell_{\lambda}+1)= \Theta(d^{\ell_{\lambda}+1})$, without loss of generality we assume that
$$
N_{\ell_{\lambda}} + 1 < \min\{n, N(d, \ell_{\lambda}+1)\} \leq N_{\ell_{\lambda}+1}.
$$

Recall that we denote $HH^{\top}=\text{Diag}\{h_1^2 \geq \cdots \geq h_n^2 \geq 0\}$ and $D=\text{Diag}\{d_1 \geq \cdots \geq d_n > 0\}$. Therefore, we have
\begin{equation*}
    \begin{aligned}
    \mathrm{tr} \left( V_{2, 1} \right) 
        \geq &~
        \mathrm{tr} \left(  \Psi_{\ell_{\lambda}+1} \Sigma_{\ell_{\lambda}+1}^{2} \Psi_{\ell_{\lambda}+1}^{\top} \varphi_{\lambda, K}^2 \right)
        =
        \mu_{\ell_{\lambda}+1}^{2} \mathrm{tr} \left(  \Psi_{\ell_{\lambda}+1}  \Psi_{\ell_{\lambda}+1}^{\top} \varphi_{\lambda, K}^2 \right)\\
        =&~
        n \mu_{\ell_{\lambda}+1}^{2} \mathrm{tr} \left(  OHH^{\top}O^{\top} PDP^{\top} \right) = n \mu_{\ell_{\lambda}+1}^{2}  \mathrm{tr} \left(  HH^{\top}O^{\top} PD P^{\top} O\right)\\
        %\overset{\text{Proposition } \ref{prop_ri}}
        {\geq} &~
        n \mu_{\ell_{\lambda}+1}^{2}\sum_{i=1}^{n} h_i^2 d_{n-i+1}
        \overset{(\ref{eqn_lower_bound_h_in_variance})}{=}
        \Omega_{d, \mathbb{P}}\left(\frac{n}{d^{2\ell_{\lambda}+2}}\right) 
        {
        \sum_{i=N_{\ell_{\lambda}} + 1}^{\min\{n, N(d, \ell_{\lambda}+1)\}}
        }
        d_{n-i+1}\\
        =&~
        \Omega_{d, \mathbb{P}}\left(\frac{n}{d^{2\ell_{\lambda}+2}}\right) {
        \sum_{i=N_{\ell_{\lambda}} + 1}^{\min\{n, N(d, \ell_{\lambda}+1)\}}
        } \lambda_{n-i+1}(\varphi_{\lambda, K}^2)\\
        \overset{(\ref{eqn:spe_properties})}{\geq}&~
        \Omega_{d, \mathbb{P}}\left(\frac{n}{d^{2\ell_{\lambda}+2}}\right) {
        \sum_{i=N_{\ell_{\lambda}} + 1}^{\min\{n, N(d, \ell_{\lambda}+1)\}}
        } \left[\lambda_i(M)\right]^{-2}\\
        \overset{\text{Lemma } \ref{lemma lambda max min}}{=}&~
        \Omega_{d, \mathbb{P}}\left(\frac{n}{d^{2\ell_{\lambda}+2}}\right) {
        \sum_{i=N_{\ell_{\lambda}} + 1}^{\min\{n, N(d, \ell_{\lambda}+1)\}}
        } \left[n\lambda + 1\right]^{-2}\\
        =&~
        \Omega_{d, \mathbb{P}}\left(\frac{n}{d^{\ell_{\lambda}+1}(n\lambda+1)^2}\right).
    \end{aligned}
\end{equation*}
% {\green
where the third line comes from von Neumann's trace inequality (see also Proposition \ref{prop_ri}).
% }


% \dhnote{
% The 2nd inequality is not \text{Proposition } \ref{prop_ri} but the von Neumann's trace theorem. 
% }


% {\green 
% [Your first suggestion is right that it is the direct use of the von Neumann's trace theorem.
% (\href{https://en.wikipedia.org/wiki/Trace_inequality#Von_Neumann's_trace_inequality_and_related_results}{Wiki link}).]
% }


% \dhnote{
% Proposition G.7 does not even have any matrices. The correct proposition should state a lower bound on the Trace of $AB$ for two PSD matrices. 
% }


% \dhnote{
% I think the summation over $i$ may have an issue and the count is unclear. 
% }

% {\green
% You are right. Though both $N(d, k)$ and $B_{k}$ are of order $d^{k}$, the summation should begin with $B_{k}+1$, and then all corresponding eigenvalues can be $\mu_{k+1}$.
% }




Combining all the above, we have
\begin{equation*}
    \begin{aligned}
      \mathrm{var}\left( \hat{f}_{\lambda} \right) =&~ \sigma^2 \left(\mathrm{tr}(V_1) + \mathrm{tr} \left( V_{2} \right) \right)
        =
        \Theta_{d, \mathbb{P}}\left(\frac{n}{d^{\tilde{\ell}+1}(n\lambda+1)^2} + \frac{d^{\tilde{\ell}}}{n}\right)\\
        = &~
        \Theta_{d, \mathbb{P}}\left(
        d^{\gamma - \tilde{\ell} - 1 - 2  \max\{\gamma - u, 0\}}
+ d^{\tilde{\ell} - \gamma}
\right),
    \end{aligned}
\end{equation*}
and we get the desired results.
\end{proof}
% $\hfill\square$





\section{Bias}\label{append_bias}


We first define the following notation:
\begin{equation}\label{eqn_def_of_tau_prime}
    \begin{aligned}
  \tau^{\prime} :=
  \begin{cases}
\tau, &  \tau < \infty;\\
s, &  \tau = \infty \text{ and } s \geq 1;\\
1, &  \tau = \infty \text{ and } s \in \{0\} \cup [\frac{1}{2}, 1);\\
(2s)^{-1}, &  \tau = \infty \text{ and } 0 < s < \frac{1}{2}.
  \end{cases}
  \end{aligned}
\end{equation}
This notation will be used in the proof of the bias term in the following way: we will apply (9) in the proof for $t=\tau^{\prime}$. 
We introduce the finite constant $\tau^{\prime}$ to ensure that 
    the constants implicit in the asymptotic notations $O_d, o_d, \Theta_d, \Omega_d$ will remain bounded uniformly in $\tau$ (even when $\tau = \infty$).

    



In this section, our goal is to prove the following theorem, which bounds the bias term in (13).


\begin{theorem}\label{thm_bias_spe}
Suppose Assumptions 1--6 hold.
Then, the bias term (15) of the spectral algorithm estimator $\hat{f}_{\lambda}$ satisfies:
    \begin{equation*}
\begin{aligned}
    \mathrm{bias}^{2}(\hat{f}_{\lambda})  =  &~
    \|B_1+B_2\|_{2}^{2} + \|B_3-B_4-B_5\|_{2}^{2},
\end{aligned} 
\end{equation*}
where
\begin{itemize}

\item[(i)] We have
    \begin{equation*}
    \begin{aligned}
                \|B_1\|_{2}^{2}  =&~ 
                O_{d, \mathbb{P}}\left(d^{-2u + (2-\min\{s, 2\})\tilde{\ell} } + d^{(2-\min\{s, 2\})\tilde{\ell}-2\gamma }\right);
    \end{aligned}
\end{equation*}

    \item[(ii)]
%     \dhnote{
% the operation     $\Omega_{d}(a_n)-O_{d}(b_n)$ is not well-defined; please rewrite in the way of Prop. \ref{prop_verify_bias_special}. 
%     }
%     {\green
    If $0 < u < \gamma <1$, then for any $\tau \leq \infty$, we have
        \begin{equation*}
        \begin{aligned}
                \|B_1\|_{2}^{2}  \leq &~ 
                O_{d, \mathbb{P}}\left(  d^{-2\tau^{\prime} u} \right) + \Delta_1,\\
                \|B_1\|_{2}^{2}  \geq &~ 
                \Omega_{d, \mathbb{P}}\left(  d^{-2\tau u} \right)- \Delta_1, \quad \text{ if } \tau<\infty,
        \end{aligned}
        \end{equation*}
        where $0 \leq \Delta_1 = O_{d, \mathbb{P}}\left( d^{-2\gamma }\right)$;
    % }

\item[(iii)] We have
\begin{equation*}
    \begin{aligned}
                \|B_2\|_{2}^{2} =&~ 
            O_{d, \mathbb{P}}\left(  d^{-(\tilde{\ell}+1)s}\right);
    \end{aligned}
\end{equation*}

\item[(iv)] We have
\begin{equation*}
    \begin{aligned}
                \|B_3-B_5\|_{2}^{2} 
        =
        \Theta_{d, \mathbb{P}} \left( d^{-(\tilde{\ell}+1)s}\right);
    \end{aligned}
\end{equation*}

\item[(v)] We have
\begin{equation*}
    \begin{aligned}
                &~\|B_4\|_{2}^{2} = O_{d, \mathbb{P}} \left(  d^{(2-\min\{s, 2\})\tilde{\ell} - (\gamma  + \ell_{\gamma} + 1)\mathbf{1}\{\ell_{\gamma} \leq \ell_{\lambda}\} - 2(\ell_{\lambda}+1)\mathbf{1}\{\ell_{\gamma} > \ell_{\lambda}\} } \right).
    \end{aligned}
\end{equation*}
\end{itemize}
\end{theorem}

    \begin{remark}
        We remark that the upper bound in Theorem~\ref{thm_bias_spe} is not tight in certain cases, for which we will provide a tight bound in Proposition \ref{prop_verify_bias_special}.
    \end{remark}
    
\begin{proof}
We consider the following two cases: (I) we have $\gamma > \ell_{\gamma}$, or we have $\ell_{\lambda} < \gamma = \ell_{\gamma}$, or we have $\ell_{\lambda} \geq \gamma = \ell_{\gamma}$ and $n \geq 2N(d, \ell_{\gamma})$;
%$\alpha \gamma ! \geq 2$; 
(II) we have $\ell_{\lambda} \geq \gamma = \ell_{\gamma}$ and $n < 2N(d, \ell_{\gamma})$.
%$\alpha \gamma ! < 2$.

\paragraph{Case (I). } Suppose that we have $\gamma > \ell_{\gamma}$, or we have $\ell_{\lambda} < \gamma = \ell_{\gamma}$, or we have $\ell_{\lambda} \geq \gamma = \ell_{\gamma}$ and $n \geq 2N(d, \ell_{\gamma})$. 
We first prove that \begin{equation}\label{eqn_bound_eigen_psi_top_psi}
    \lambda_{\text{max}}(\Psi_{\leq \tilde{\ell}}^{\top}\Psi_{\leq \tilde{\ell}}) = \Theta_{d, \mathbb{P}}(n) \quad \text{ and } \quad \lambda_{\text{min}}(\Psi_{\leq \tilde{\ell}}^{\top}\Psi_{\leq \tilde{\ell}}) = \Theta_{d, \mathbb{P}}(n).
\end{equation}
Eq.~\eqref{eqn_bound_eigen_psi_top_psi} can be seen by considering the following: 
\begin{itemize}
    \item Suppose $\ell_{\lambda} \geq \ell_{\gamma}$. Since $\tilde{\ell}=\ell_{\gamma}$, from Lemma \ref{lemma psi top psi} and Lemma \ref{lemma psi top psi_int}, we have Eq. \eqref{eqn_bound_eigen_psi_top_psi} holds. 

    \item Suppose  $\ell_{\lambda}<\ell_{\gamma}<\gamma$. Recall that $\Psi_{\leq \ell_{\gamma}} = \left[\Psi_{\leq \ell_{\lambda}}, \Psi_{\ell_{\lambda}+1, \ell_{\gamma}}\right]$. 
    Since $\tilde{\ell}=\ell_{\lambda}$, from Lemma \ref{lemma psi top psi}, we have
\begin{align*}
&~ \left\| \frac{1}{n}
\Psi_{\leq \ell_\lambda}^{\top}\,\Psi_{\leq \ell_\lambda} - 
\mathrm{I}_{N_{\ell_\lambda}} \right\|_{\mathrm{op}}\\
    \leq &~ \left\| \frac{1}{n}
\begin{pmatrix}
  \Psi_{\leq \ell_\lambda}^{\top}\,\Psi_{\leq \ell_\lambda}
  & 
  \Psi_{\leq \ell_\lambda}^{\top}\,\Psi_{\ell_\lambda+1,\ell_{\gamma}}
  \\[1em]
  \Psi_{\ell_\lambda+1,\ell_{\gamma}}^{\top}\,\Psi_{\leq \ell_\lambda}
  & 
  \Psi_{\ell_\lambda+1,\ell_{\gamma}}^{\top}\,\Psi_{\ell_\lambda+1,\ell_{\gamma}}
\end{pmatrix} - \begin{pmatrix}
\mathrm{I}_{N_{\ell_\lambda}} & 0 \\[0.5em]
0 & \mathrm{I}_{\,N_{\ell_{\gamma}}-N_{\ell_\lambda}}
\end{pmatrix} \right\|_{\mathrm{op}} = o_{d, \mathbb{P}}(1),
\end{align*}
    which implies Eq. (\ref{eqn_bound_eigen_psi_top_psi}).

    \item Suppose $\ell_{\lambda}<\ell_{\gamma} = \gamma$. Similar to the last case, we can use Lemma \ref{lemma psi top psi_int} to show Eq. (\ref{eqn_bound_eigen_psi_top_psi}).
\end{itemize}






% \dhnote{
% I am not sure about the case $\ell_{\lambda} < \gamma = \ell$ and how $\lambda_{\text{min}}(\Psi_{\leq \tilde{\ell}}^{\top}\Psi_{\leq \tilde{\ell}}) = \Theta_{d, \mathbb{P}}(n)$ is obtained. 

% Lemma G3 and G4 can only apply to $\ell$ and $\ell-1$. If $\ell_\lambda\leq \ell-2$, I don't see how to lower bound $\lambda_{\text{min}}(\Psi_{\leq \tilde{\ell}}^{\top}\Psi_{\leq \tilde{\ell}})$. 
% }


% \dhnote{
% The same issue occurs in \eqref{eqn_finite_tau_bound}. 
% }

% {\green
% [I have modified the proof to explain how to obtain Eq. (\ref{eqn_bound_eigen_psi_top_psi}) in different cases of Case (i).]
% }




The bias term Eq.~(15) can be decomposed as
\begin{equation}\label{eq derivation of bias_spe}
    \begin{aligned}
        \mathrm{bias}^2(\hat{f}_{\lambda}) &= E_{x} \left[\left( \mathscr{K}(x,X) \varphi_{\lambda, K} f_{\star}(X) -f_{\star}(x)\right)^2\right]\\
    =&~ E_{x} \left[\left( \psi(x)^{\top} \Sigma \Psi^{\top} \varphi_{\lambda, K} \Psi \theta - \psi(x)^{\top} \theta\right)^2\right]\\
    =&~ \| \Sigma \Psi^{\top} \varphi_{\lambda, K} \Psi \theta - \theta  \|_{2}^{2}\\
    =&~
    \| \Sigma_{\leq \tilde{\ell}} \Psi^{\top}_{\leq \tilde{\ell}} \varphi_{\lambda, K} \Psi \theta - \theta_{\leq \tilde{\ell}}  \|_{2}^{2} + \| \Sigma_{>\tilde{\ell}} \Psi^{\top}_{>\tilde{\ell}} \varphi_{\lambda, K} \Psi \theta - \theta_{>\tilde{\ell}}  \|_{2}^{2}\\
    =&~
    \| \underbrace{\Sigma_{\leq \tilde{\ell}} \Psi^{\top}_{\leq \tilde{\ell}} \varphi_{\lambda, K} \Psi_{\leq \tilde{\ell}} \theta_{\leq \tilde{\ell}} - \theta_{\leq \tilde{\ell}}  }_{B_{1}}+  \underbrace{\Sigma_{\leq \tilde{\ell}} \Psi^{\top}_{\leq \tilde{\ell}} \varphi_{\lambda, K} \Psi_{>\tilde{\ell}} \theta_{>\tilde{\ell}}}_{B_{2}} \|_{2}^{2}\\
    &~+
     \|\underbrace{\theta_{>\tilde{\ell}}}_{B_{3}} - \underbrace{ \Sigma_{>\tilde{\ell}} \Psi^{\top}_{>\tilde{\ell}} \varphi_{\lambda, K} \Psi_{\leq \tilde{\ell}} \theta_{\leq \tilde{\ell}} }_{B_{4}} - \underbrace{  \Sigma_{>\tilde{\ell}} \Psi^{\top}_{>\tilde{\ell}} \varphi_{\lambda, K} \Psi_{>\tilde{\ell}} \theta_{>\tilde{\ell}} }_{ B_{5}}\|_{2}^{2}.
    \end{aligned}
\end{equation}
Before we bound these terms separately, we need the following calculations.

% \dhnote{The following seems to hold only if $\tilde{\ell}=\ell$ or $\ell-1$; what if $\ell_\lambda\leq \ell-2$? 
% }

% {\green
% [I have modified the proof to explain how to obtain Eq. (\ref{eqn_bound_eigen_psi_top_psi}) in different cases of Case (i).]
% }


Denote $Z = n^{-1/2}\Psi_{\leq \tilde{\ell}}$. 
%
% \color{red}[This part is rewritten.]
% From Lemma \ref{lemma psi top psi} and \ref{lemma psi top psi_int}, we can denote the SVD decomposition of $Z$ as $Z=B (H, 0)^{\top} V^{\top}$, where $B \in \mathbb{R}^{n \times n}$ and $V \in \mathbb{R}^{B_{\tilde{\ell}} \times B_{\tilde{\ell}}}$ are orthogonal matrices, $H = \text{Diag}\{h_1, \cdots, h_{B_{\tilde{\ell}}}\}$ 
% is a diagonal matrix, and $h_1, \cdots, h_{B_{\tilde{\ell}}} = \Theta_{d, \mathbb{P}}(1)$.
%
% \color{blue}
% [New]
Suppose the SVD decomposition of $Z$ is $Z=U (H, 0)^{\top} V^{\top}$, where $U \in \mathbb{R}^{n \times n}$, $V \in \mathbb{R}^{N_{\tilde{\ell}} \times N_{\tilde{\ell}}}$ are orthogonal matrices, and $H = \text{Diag}\{h_1, \cdots, h_{N_{\tilde{\ell}}}\}$ 
is a diagonal matrix. 
Since Eq. \eqref{eqn_bound_eigen_psi_top_psi} holds, we have $h_1, \cdots, h_{N_{\tilde{\ell}}} = \Theta_{d, \mathbb{P}}(1)$.
% \color{black}



Hence, 
\begin{equation}\label{eqn_finite_tau_bound}
\begin{aligned}
    Z \Sigma_{\leq \tilde{\ell}} Z^{\top} =&~
    U \text{Diag}\{ HV^{\top}\Sigma_{\leq \tilde{\ell}} VH, 0_{n-N_{\tilde{\ell}}} \} U^{\top}\\
    =&~
    U_{N_{\tilde{\ell}}} S_0 U_{N_{\tilde{\ell}}}^{\top}
\end{aligned}
\end{equation}
where we have defined the following 
$$
S_0 =~ HV^{\top}\Sigma_{\leq \tilde{\ell}} VH.
$$


%where $\|\Delta_{4}\|_{\mathrm{op}} = o_{d, \mathbb{P}}\left(\lambda_{\min}(\Sigma_{\leq \tilde{\ell}})\right)$.
\iffalse
\\
For any $i \leq B_{\tilde{\ell}}$, from Proposition 2.1 and (\ref{eqn:regular_para_upper}), we have $\lambda_i = \Omega_{d, \mathbb{P}}( \lambda + n^{-1})$. 
Therefore from (8) we have
\begin{equation}\label{eqn_matrix_ine_3}
    \begin{aligned}
       h_i^2\lambda_i\reg(h_i^2\lambda_i + \kappa_1 / n)
        = &~
         \Theta_{d, \mathbb{P}}(1) ((h_i^2\lambda_i + \kappa_1 / n + \lambda)\reg(h_i^2\lambda_i + \kappa_1 / n)) = \Theta_{d, \mathbb{P}}(1).
    \end{aligned}
\end{equation}
\fi

\iffalse
Similarly, when $ \tilde{\ell}=\ell \leq \ell_\lambda$, we can denote the SVD decomposition of $Z\Sigma_{\leq \tilde{\ell}} Z^{\top} +K_{>\tilde{\ell}}/n$ as
\begin{equation}\label{eqn_finite_tau_bound_2}
    Z\Sigma_{\leq \tilde{\ell}} Z^{\top} +K_{>\tilde{\ell}}/n = \tilde{B} \text{Diag}\left\{\tilde{H}, \tilde{G} \right\} \tilde{B}^{\top},
\end{equation}
where $\tilde{B} \in \mathbb{R}^{n \times n}$ is an orthogonal matrix, and $\tilde{H} = \text{Diag}\{\tilde{h}_1, \cdots, \tilde{h}_{B_{\tilde{\ell}}}\}, \tilde{G} = \text{Diag}\{\tilde{g}_1, \cdots, \tilde{h}_{n-B_{\tilde{\ell}}}\}$ are diagonal matrices. Moreover, from Lemma \ref{lemma psi psi top}, \ref{lemma psi top psi}, and \ref{lemma psi top psi_int}, we have $\tilde{h}_j = \Theta_{d, \mathbb{P}}(\lambda_{j}(\Sigma_{\leq \tilde{\ell}}))$ and $\tilde{h}_k = \Theta_{d, \mathbb{P}}(n^{-1})$, with $j \leq B_{\tilde{\ell}}$ and $k \leq n- B_{\tilde{\ell}}$.
\fi






\iffalse
For any $i \leq B_{\tilde{\ell}}$, from Proposition 2.1 and (\ref{eqn:regular_para_upper}), we have $\lambda_i = \Omega_{d, \mathbb{P}}( \lambda + n^{-1})$. 
Therefore,
%for any constants $\kappa_1, \kappa_1 >0$ (we will determine them later) only depending on the absolute positive constants listed in Definition \ref{def:abs_constants}, 
on one hand, from (8)
%, there exists a constant $0<\kappa_1^{\prime}<1$ only depending on $\kappa_1$ and the absolute positive constants listed in Definition \ref{def:abs_constants}, such that 
we have 
\begin{equation}\label{eqn_matrix_ine_1}
    \begin{aligned}
       &~ \left(h_i^2\lambda_i + \frac{\kappa_1}{n}\mathbf{1}\{\lambda_i^{-1} =\Theta_{d, \mathbb{P}}(n)\}\right)\reg\left(h_i^2\lambda_i + \frac{\kappa_1}{n}\right)\\
        \geq &~
         \left(h_i^2\lambda_i + \frac{\kappa_1}{n}\mathbf{1}\{\lambda_i^{-1} =\Theta_{d, \mathbb{P}}(n)\}\right)\left(h_i^2\lambda_i + \frac{\kappa_1}{n}\right)^{-1} \left[1 - C_3 \lambda\left(h_i^2\lambda_i + \frac{\kappa_1}{n}\right)^{-1}\right]\\
        \geq  &~
        \left\{\begin{matrix}
(1-o_{d, \mathbb{P}}(1) ) \left[1 - O_{d, \mathbb{P}}( \lambda\lambda_i^{-1})\right], &\quad \lambda_i^{-1} =o_{d, \mathbb{P}}(n)  \\
1 - O_{d, \mathbb{P}}( \lambda\lambda_i^{-1}), & \quad \lambda_i^{-1} =\Theta_{d, \mathbb{P}}(n)
\end{matrix}\right.
    \end{aligned}
\end{equation}
On the other hand, when $\tau<\infty$, from (9), there exists a constant $\mathfrak{C}_4^{\prime}>0$ only depending on the absolute positive constants listed in Definition \ref{def:abs_constants}, such that we have 
\begin{equation}\label{eqn_matrix_ine_2}
    \begin{aligned}
       h_i^2\lambda_i\reg(h_i^2\lambda_i + \kappa_1 / n)
        \leq &~
         h_i^2\lambda_i(h_i^2\lambda_i + \kappa_1 / n)^{-1} \left[1 - \mathfrak{C}_4^{\prime} \lambda^{\tau}(h_i^2\lambda_i +\kappa_1 / n)^{-\tau}\right]\\
        \leq &~
        1 - \Omega_{d, \mathbb{P}}( \lambda^{\tau}\lambda_i^{-\tau}).
    \end{aligned}
\end{equation}


Finally, for any $i \leq B_{\tilde{\ell}}$, similar as above, 
\fi







\noindent {\bf Bounding $\|B_1\|_{2}^{2}$. }
We consider the following two cases: (B1.i) $\ell_{\gamma} \leq \ell_\lambda$; 
(B1.ii) $\ell_{\gamma} > \ell_\lambda$.

\noindent {\it Proof of case (B1.i). } 



When $\ell_{\gamma} \leq \ell_\lambda$, 
\iffalse
from Lemma \ref{lemma psi psi top}, %there exist constants $\kappa_1, \kappa_1 >0$ only depending on the absolute positive constants listed in Definition \ref{def:abs_constants}, such that we have $\kappa_1 \mathrm{I}_n \leq K_{>\tilde{\ell}} \leq \kappa_1 \mathrm{I}_n$.
we have $K_{>\tilde{\ell}} = \kappa_1 \left(\mathrm{I}_{n} + \Delta_{1}\right)$ with $\kappa_{1} = \Theta_{d}(1)$ and $\|\Delta_{1}\|_{\mathrm{op}} = o_{d, \mathbb{P}}(1)$.
Therefore, 
\fi
from (\ref{eqn_bound_eigen_psi_top_psi}) we have
\iffalse
\begin{equation}\label{eqn_b1_spe_origin}
    \begin{aligned}
                &~\|B_1\|_{2}^{2} = 
                \left\| \left( \mathrm{I}_{N_{\tilde{\ell}}} - \Sigma_{\leq \tilde{\ell}} \Psi^{\top}_{\leq \tilde{\ell}} \varphi_{\lambda, K} \Psi_{\leq \tilde{\ell}}\right)    \theta_{\leq \tilde{\ell}} \right\|_{2}^{2}\\
                =&~ 
                \Theta_{d, \mathbb{P}}\left(\frac{1}{n}\right)
                \left\| \Psi_{\leq \tilde{\ell}}\left( \mathrm{I}_{N_{\tilde{\ell}}} - \Sigma_{\leq \tilde{\ell}} \Psi^{\top}_{\leq \tilde{\ell}} \varphi_{\lambda, K} \Psi_{\leq \tilde{\ell}}\right)    \theta_{\leq \tilde{\ell}} \right\|_{2}^{2}\\
                =&~ 
                \Theta_{d, \mathbb{P}}\left(\frac{1}{n}\right)
                \left\| \left( I_{n} - \Psi_{\leq \tilde{\ell}}\Sigma_{\leq \tilde{\ell}} \Psi^{\top}_{\leq \tilde{\ell}} \varphi_{\lambda, K} \right) 
                \Psi_{\leq \tilde{\ell}}\theta_{\leq \tilde{\ell}} \right\|_{2}^{2}\\
                {=}&~
        \Theta_{d, \mathbb{P}}\left(1\right)
        \left\| \left( \mathrm{I}_n - Z\Sigma_{\leq \tilde{\ell}} Z^{\top} \reg(Z\Sigma_{\leq \tilde{\ell}} Z^{\top} +K_{>\tilde{\ell}}/n) \right) 
                Z\theta_{\leq \tilde{\ell}} \right\|_{2}^{2}\\
                {\leq}&~
        \Theta_{d, \mathbb{P}}\left(1\right)\left\| \left( \mathrm{I}_n - (Z\Sigma_{\leq \tilde{\ell}} Z^{\top} +K_{>\tilde{\ell}}/n) \reg(Z\Sigma_{\leq \tilde{\ell}} Z^{\top} +K_{>\tilde{\ell}}/n) \right) 
                Z\theta_{\leq \tilde{\ell}} \right\|_{2}^{2}\\
                &~+
                \Theta_{d, \mathbb{P}}\left(1\right)\left\| K_{>\tilde{\ell}} 
                \varphi_{\lambda, K}
                Z
                \theta_{\leq \tilde{\ell}} \right\|_{2}^{2}\\
                \overset{(\ref{eqn_finite_tau_bound})}{=}&~
                \Theta_{d, \mathbb{P}}\left(1\right) 
                \left\| \left( \mathrm{I}_n - B S_{0} B^{\top} \reg(B (S_{0}+ \frac{\kappa_1}{n} \mathrm{I}_n) B^{\top}) \right) 
                Z\theta_{\leq \tilde{\ell}} \right\|_{2}^{2}\\
                \overset{(\ref{eq:Filter_Rem_finite_case3})}{=}&~
                \Theta_{d, \mathbb{P}}\left(1\right) 
                \left\| \left( \mathrm{I}_n - B S_{0}  \reg( S_{0}+ \frac{\kappa_1}{n} \mathrm{I}_n ) B^{\top}\right) 
                Z\theta_{\leq \tilde{\ell}} \right\|_{2}^{2}\\
                \overset{(\ref{eqn_matrix_ine_1})}{=}&~
                O_{d, \mathbb{P}}\left(\lambda^{2}\right) 
                \left\| \text{Diag}\{\Sigma_{\leq \tilde{\ell}}^{-1}, 0_{n-B_{\tilde{\ell}}}\}
                B^{\top}B (H+\Delta_3, 0)^{\top} V^{\top}\theta_{\leq \tilde{\ell}} \right\|_{2}^{2}\\
                 =&~
                O_{d, \mathbb{P}}\left(\lambda^{2}\right)
        \left\| \Sigma_{\leq \tilde{\ell}}^{-1}
                \theta_{\leq \tilde{\ell}} \right\|_{2}^{2}\\
                \overset{\text{Lemma } \ref{lemma ge l l2}}{=} &~
                O_{d, \mathbb{P}}\left(  \lambda^{2} d^{(2-\min\{s, 2\})\tilde{\ell} }\right).
    \end{aligned}
\end{equation}
\fi

\begin{equation}\label{eqn_b1_spe_origin}
    \begin{aligned}
                &~\|B_1\|_{2}^{2} = 
                \left\| \left( \mathrm{I}_{N_{\tilde{\ell}}} - \Sigma_{\leq \tilde{\ell}} \Psi^{\top}_{\leq \tilde{\ell}} \varphi_{\lambda, K} \Psi_{\leq \tilde{\ell}}\right)    \theta_{\leq \tilde{\ell}} \right\|_{2}^{2}\\
                =&~ 
                \Theta_{d, \mathbb{P}}\left(\frac{1}{n}\right)
                \left\| \Psi_{\leq \tilde{\ell}}\left( \mathrm{I}_{N_{\tilde{\ell}}} - \Sigma_{\leq \tilde{\ell}} \Psi^{\top}_{\leq \tilde{\ell}} \varphi_{\lambda, K} \Psi_{\leq \tilde{\ell}}\right)    \theta_{\leq \tilde{\ell}} \right\|_{2}^{2}\\
                =&~ 
                \Theta_{d, \mathbb{P}}\left(\frac{1}{n}\right)
                \left\| \left( \mathrm{I}_{n} - \Psi_{\leq \tilde{\ell}}\Sigma_{\leq \tilde{\ell}} \Psi^{\top}_{\leq \tilde{\ell}} \varphi_{\lambda, K} \right) 
                \Psi_{\leq \tilde{\ell}}\theta_{\leq \tilde{\ell}} \right\|_{2}^{2}\\
                {=}&~
        \Theta_{d, \mathbb{P}}\left(1\right)
        \left\| \left( \mathrm{I}_n - Z\Sigma_{\leq \tilde{\ell}} Z^{\top} \reg(Z\Sigma_{\leq \tilde{\ell}} Z^{\top} +K_{>\tilde{\ell}}/n) \right) 
                Z\theta_{\leq \tilde{\ell}} \right\|_{2}^{2}\\
                {\leq}&~
        O_{d, \mathbb{P}}\left(1\right)\left\| \underbrace{\left( \mathrm{I}_n - (Z\Sigma_{\leq \tilde{\ell}} Z^{\top} +K_{>\tilde{\ell}}/n) \reg(Z\Sigma_{\leq \tilde{\ell}} Z^{\top} +K_{>\tilde{\ell}}/n) \right) 
                Z\theta_{\leq \tilde{\ell}} }_{B_{1, 1}}\right\|_{2}^{2}\\
                &~+
                O_{d, \mathbb{P}}\left(1\right)\left\| \underbrace{K_{>\tilde{\ell}} 
                \varphi_{\lambda, K}
                Z
                \theta_{\leq \tilde{\ell}} }_{B_{1, 2}}\right\|_{2}^{2}.
    \end{aligned}
\end{equation}

For $B_{1,1}$:\\
From Lemma \ref{lemma psi psi top}, there exist constants $\mathfrak{C}_5, \mathfrak{C}_6 >0$ such that we have $\mathfrak{C}_5 \mathrm{I}_n \leq K_{>\tilde{\ell}} \leq \mathfrak{C}_6 \mathrm{I}_n$. 
Furthermore, $\tilde{\ell}=\ell_{\gamma}$ since $\ell_{\gamma} \leq \ell_{\lambda}$, and $\lambda_{\min}(\Sigma_{\leq \ell_{\gamma}})=\mu_{\ell_{\gamma}}=\Omega(d^{-\ell_{\gamma}})$ by Proposition~2.1. 
Hence, we have $\lambda_{\max}(K_{>\tilde{\ell}}/n)=O_{d, \mathbb{P}}\left(\lambda_{\min}(\Sigma_{\leq \tilde{\ell}})\right)$. 

% \dhnote{
% So this is $\lambda_{\min}(\Sigma_{\leq \tilde{\ell}})=\Omega_{d,\mathbb{P}}(1/n)$. Why is this true if $\ell_\lambda<\ell-1$?
% }

% {\green
% [We are now proving the case B1.i, which considers the case where $\ell \leq \ell_{\lambda}$]
% }






Therefore, from (\ref{eqn_bound_eigen_psi_top_psi}) we have

\begin{equation}\label{eqn_B11}
    \begin{aligned}
       &~ \|B_{1, 1}\|_2^2 \\
       (\because \rem(z) \text{ is non-increasing  w.r.t } z)\quad \leq&~
       \left\| \left( \mathrm{I}_n - (Z\Sigma_{\leq \tilde{\ell}} Z^{\top}) \reg(Z\Sigma_{\leq \tilde{\ell}} Z^{\top} ) \right) 
                Z\theta_{\leq \tilde{\ell}} \right\|_{2}^{2}\\
                (\because \text{ Eqns. (\ref{eq:Filter_Rem_finite_case3}) and (\ref{eqn_finite_tau_bound})} )\qquad =
        &~\left\| \left[ \mathrm{I}_{N_{\tilde{\ell}}} - S_0 \reg\left(S_0 \right) \right]  H V^{\top}\theta_{\leq \tilde{\ell}} \right\|_{2}^{2}\\
   (\because \Sigma_{\leq \tilde{\ell}} VH H^{-1} V^{\top}  \Sigma_{\leq \tilde{\ell}}^{-1}=\mathrm{I}_{N_{\tilde{\ell}}}) 
   {=}
     &~\left\| \rem\left( S_0 \right) HV^{\top} \left(\Sigma_{\leq \tilde{\ell}} VH H^{-1} V^{\top}  \Sigma_{\leq \tilde{\ell}}^{-1}\right) \theta_{\leq \tilde{\ell}} \right\|_{2}^{2}\\
        \leq &~
        \left\| \rem\left(S_0\right) S_0 \right\|_{2}^{2} \cdot \left\| H^{-1} V^{\top}  \Sigma_{\leq \tilde{\ell}}^{-1} \theta_{\leq \tilde{\ell}} \right\|_{2}^{2}\\
(\because h_1, \cdots, h_{B_{\tilde{\ell}}} = \Theta_{d, \mathbb{P}}(1)) \qquad        \leq &~
        O_{d, \mathbb{P}}\left(1\right)\left\| \rem\left(S_0\right) S_0 \right\|_{2}^{2} \cdot \left\|\Sigma_{\leq \tilde{\ell}}^{-1} \theta_{\leq \tilde{\ell}} \right\|_{2}^{2}\\
        (\because \text{Eq. (8)}) \qquad = 
        &~O_{d, \mathbb{P}}\left(\lambda^{2}\right) \left\| \Sigma_{\leq \tilde{\ell}}^{-1}\theta_{\leq \tilde{\ell}} \right\|_{2}^{2}\\
                (\because \text{Lemma~\ref{lemma ge l l2} (iii)}) \qquad = &~
                O_{d, \mathbb{P}}\left(  \lambda^{2} d^{(2-\min\{s, 2\})\tilde{\ell} }\right);
    \end{aligned}
\end{equation}


% \dhnote{
% I am concern about the last inequality. Is it obvious? It seems to me that it has used  $h_1, \cdots, h_{B_{\tilde{\ell}}} = \Theta_{d, \mathbb{P}}(1)$ and also rely on the condition number of $\rem\left(HV^{\top}\Sigma_{\leq \tilde{\ell}} VH \right) $. 
% }






\iffalse [Backup 2025-09-18, DH]
\textbf{Bounding $B_{1,2}$.}
Eq.~\eqref{eqn_bound_eigen_psi_top_psi} implies that $\lambda_{\min}(Z^{\top}Z)=\Theta_{d,\mathbb{P}}(1)$  and $\lambda_{\max}(Z^{\top}Z)=\Theta_{d,\mathbb{P}}(1)$. Therefore,
\begin{equation}\label{eqn_bound_b_12}
    \begin{aligned}
& \left\| B_{1, 2} \right\|_{2}^{2} \\
(\because \lambda_{\min}(Z^{\top}Z)=\Theta_{d,\mathbb{P}}(1))\qquad  
=&~ O_{d, \mathbb{P}}\left(1\right)\left\| \left(\varphi_{\lambda, K} 
                Z\Sigma_{\leq \tilde{\ell}}
                Z^{\top}\right)Z
                \Sigma_{\leq \tilde{\ell}}^{-1}\theta_{\leq \tilde{\ell}} \right\|_{2}^{2}
                \\ 
        (\because \text{Eq. (\ref{eqn:spe_properties})}) \qquad {=} & O_{d, \mathbb{P}}\left(\frac{1}{n^2}\right)\left\| 
                Z\Sigma_{\leq \tilde{\ell}}^{-1}\theta_{\leq \tilde{\ell}} \right\|_{2}^{2}\\
  (\because \lambda_{\max}(Z^{\top}Z)=\Theta_{d,\mathbb{P}}(1))\qquad               =&~O_{d, \mathbb{P}}\left(\frac{1}{n^2}\right)\left\| 
                \Sigma_{\leq \tilde{\ell}}^{-1}\theta_{\leq \tilde{\ell}} \right\|_{2}^{2}\\
               ( \because\text{Lemma \ref{lemma ge l l2}(iii)}  ) \qquad  {=} &
                O_{d, \mathbb{P}}\left( d^{(2-\min\{s, 2\})\tilde{\ell}-2\gamma }\right).
    \end{aligned}
\end{equation}
\fi 

% \newpage

For $B_{1,2}$:\\
Eq.~\eqref{eqn_bound_eigen_psi_top_psi} implies that $\lambda_{\min}(Z^{\top}Z)=\Theta_{d,\mathbb{P}}(1))$ and $\lambda_{\max}(Z)=\Theta_{d,\mathbb{P}}(1)$. Therefore, 
\iffalse [format]
\begin{equation}\label{eqn_bound_b_12}
    \begin{aligned}
& \left\| B_{1, 2} \right\|_{2}^{2} \\ 
=&~ O_{d, \mathbb{P}}\left(1\right)\left\| \left(\varphi_{\lambda, K} 
                Z\Sigma_{\leq \tilde{\ell}}
                Z^{\top}\right)Z\left(Z^\top Z\right)^{-1}
                \Sigma_{\leq \tilde{\ell}}^{-1}\theta_{\leq \tilde{\ell}} \right\|_{2}^{2}
                \\ 
        (\because \text{Eq. (\ref{eqn:spe_properties})}) \qquad {=} & O_{d, \mathbb{P}}\left(\frac{1}{n^2}\right)\left\| 
                Z\left(Z^\top Z\right)^{-1}\Sigma_{\leq \tilde{\ell}}^{-1}\theta_{\leq \tilde{\ell}} \right\|_{2}^{2}\\
  (\because \lambda_{\max}(Z)=\Theta_{d,\mathbb{P}}(1),~~ \lambda_{\min}(Z^{\top}Z)=\Theta_{d,\mathbb{P}}(1)))\qquad               =&~O_{d, \mathbb{P}}\left(\frac{1}{n^2}\right)\left\| 
                \Sigma_{\leq \tilde{\ell}}^{-1}\theta_{\leq \tilde{\ell}} \right\|_{2}^{2}\\
               ( \because\text{Lemma \ref{lemma ge l l2} (iii)}  ) \qquad  {=} &
                O_{d, \mathbb{P}}\left( d^{(2-\min\{s, 2\})\tilde{\ell}-2\gamma }\right).
    \end{aligned}
\end{equation}
\fi 

Therefore,
\begin{equation}\label{eqn_bound_b_12}
\begin{aligned}
\left\| B_{1, 2} \right\|_{2}^{2}
&= O_{d, \mathbb{P}}\left(1\right)
\left\|
\left(\varphi_{\lambda, K} 
Z\Sigma_{\leq \tilde{\ell}}
Z^{\top}\right)
Z\left(Z^\top Z\right)^{-1}
\Sigma_{\leq \tilde{\ell}}^{-1}\theta_{\leq \tilde{\ell}}
\right\|_{2}^{2} \\
&= O_{d, \mathbb{P}}\left(\frac{1}{n^2}\right)
\left\|
Z\left(Z^\top Z\right)^{-1}
\Sigma_{\leq \tilde{\ell}}^{-1}\theta_{\leq \tilde{\ell}}
\right\|_{2}^{2} \\
&= O_{d, \mathbb{P}}\left(\frac{1}{n^2}\right)
\left\|
\Sigma_{\leq \tilde{\ell}}^{-1}\theta_{\leq \tilde{\ell}}
\right\|_{2}^{2} \\
&= O_{d, \mathbb{P}}\left( d^{(2-\min\{s, 2\})\tilde{\ell}-2\gamma }\right),
\end{aligned}
\end{equation}
where 
\begin{itemize}
    \item the second equality follows from Eq.~\eqref{eqn:spe_properties};
    \item the third equality is because 
$
\lambda_{\max}(Z)=\Theta_{d,\mathbb{P}}(1),
~\lambda_{\min}(Z^{\top}Z)=\Theta_{d,\mathbb{P}}(1)$; 
\item the last equality follows from Lemma~\ref{lemma ge l l2}(iii).
\end{itemize} 





\noindent {\it Proof of case (B1.ii). } 
Suppose $\ell_\lambda < \ell_{\gamma}$. 

\iffalse
from Lemma \ref{lemma lambda max min} (iv) we have $\lambda_{\text{max}}(K_{>\ell_{\lambda}}/n) =
%{\blue O_{d, \mathbb{P}}(d^{-\ell_{\lambda}-1})=} 
o_{d, \mathbb{P}}(\lambda)$. Therefore, on one hand, we have
\begin{equation*}
\begin{aligned}
      \|B_1\|_{2}^{2}  \overset{(\ref{eqn_b1_spe_origin})}{=}&~
      \Theta_{d, \mathbb{P}}\left(1\right)
        \left\| \left( \mathrm{I}_n - Z\Sigma_{\leq \tilde{\ell}} Z^{\top} \reg(Z\Sigma_{\leq \tilde{\ell}} Z^{\top} +K_{>\tilde{\ell}}/n) \right) 
                Z\theta_{\leq \tilde{\ell}} \right\|_{2}^{2}\\
        \overset{\text{Assumption } 5 \text{ and Lemma } \ref{lemma ge l l2}}{=}&~
      %\Theta_{d, \mathbb{P}}\left(1\right)
        %\left\| \left( \mathrm{I}_n - Z\Sigma_{\leq \tilde{\ell}} Z^{\top} \reg(Z\Sigma_{\leq \tilde{\ell}} Z^{\top}) \right) Z\theta_{\leq \tilde{\ell}} \right\|_{2}^{2} + {\blue O_{d, \mathbb{P}}\left(  d^{-2\ell_{\lambda}-2} d^{(2-\min\{s, 2\})\tilde{\ell} }\right)}\\
        %        {=}&~ \Theta_{d, \mathbb{P}}\left(1\right)\left\| \left( \mathrm{I}_n - Z\Sigma_{\leq \tilde{\ell}} Z^{\top} \reg(Z\Sigma_{\leq \tilde{\ell}} Z^{\top}) \right) Z\theta_{\leq \tilde{\ell}} \right\|_{2}^{2} + {\blue O_{d, \mathbb{P}}\left(  d^{-2-\min\{s, 2\}\tilde{\ell} }\right)}\\
         %       {=}&~
      \Theta_{d, \mathbb{P}}\left(1\right)
        \left\| \left( \mathrm{I}_n - Z\Sigma_{\leq \tilde{\ell}} Z^{\top} \reg(Z\Sigma_{\leq \tilde{\ell}} Z^{\top}) \right) 
                Z\theta_{\leq \tilde{\ell}} \right\|_{2}^{2} + o_{d, \mathbb{P}}\left(  \lambda^{2} d^{(2-\min\{s, 2\})\tilde{\ell} }\right).
\end{aligned}
\end{equation*}
As in (27), we can then show that
$$
\left\|B_1\right\|_2^2=O_{d, \mathbb{P}}\left(\lambda^{2 \tau^{\prime}} d^{\left(2 \tau^{\prime}-\min \left\{s, 2 \tau^{\prime}\right\}\right) \tilde{\ell}}\right)+o_{d, \mathbb{P}}\left(  \lambda^{2} d^{(2-\min\{s, 2\})\tilde{\ell} }\right),
$$
and
$$
\left\|B_1\right\|_2^2=\Omega_{d, \mathbb{P}}\left(\lambda^{2 \tau} d^{(2 \tau-\min \{s, 2 \tau\}) \tilde{\ell}}\right)-o_{d, \mathbb{P}}\left(  \lambda^{2} d^{(2-\min\{s, 2\})\tilde{\ell} }\right),
$$
when $\tau<\infty$.

On the other hand, 
\fi

\iffalse
we have
\begin{equation*}
\begin{aligned}
      \|B_1\|_{2}^{2}  \overset{(\ref{eqn_b1_spe_origin})}{=}&~
      \Theta_{d, \mathbb{P}}\left(1\right)
        \left\| \left( \mathrm{I}_n - Z\Sigma_{\leq \tilde{\ell}} Z^{\top} \reg(Z\Sigma_{\leq \tilde{\ell}} Z^{\top} +K_{>\tilde{\ell}}/n) \right) 
                Z\theta_{\leq \tilde{\ell}} \right\|_{2}^{2}\\
        \overset{\text{Assumption } 5}
        {\leq}&~
      O_{d, \mathbb{P}}\left(1\right)
        \left\| Z\Sigma_{\leq \tilde{\ell}} Z^{\top} (K/n)^{-1}\left( \mathrm{I}_n - (K/n) \reg(K/n) \right) 
                Z\theta_{\leq \tilde{\ell}} \right\|_{2}^{2}\\ 
                &~+  
                O_{d, \mathbb{P}}\left(  1\right)
                \left\| (K_{>\tilde{\ell}}/n) (K/n)^{-1} Z\Sigma_{\leq \tilde{\ell}} Z^{\top} 
                Z \Sigma_{\leq \tilde{\ell}}^{-1} \theta_{\leq \tilde{\ell}} \right\|_{2}^{2}\\
                {\leq}&~
                O_{d, \mathbb{P}}\left(1\right)
        \left\| \left( \mathrm{I}_n - (K/n) \reg(K/n) \right) 
                Z\theta_{\leq \tilde{\ell}} \right\|_{2}^{2}\\ 
                &~+
                O_{d, \mathbb{P}}\left(  1\right)
                \left\| K_{\tilde{\ell}+1, \ell_{\gamma}} K^{-1} Z\Sigma_{\leq \tilde{\ell}} Z^{\top}
                Z \Sigma_{\leq \tilde{\ell}}^{-1} \theta_{\leq \tilde{\ell}} \right\|_{2}^{2}\\
                &~+
                O_{d, \mathbb{P}}\left(  1\right)
                \left\| K_{>\ell_{\gamma}} K^{-1} Z\Sigma_{\leq \tilde{\ell}} Z^{\top}
                Z \Sigma_{\leq \tilde{\ell}}^{-1} \theta_{\leq \tilde{\ell}} \right\|_{2}^{2}
                \\
                \overset{\rem(z) \text{ is non-increasing  w.r.t } z}{\leq}&~
      O_{d, \mathbb{P}}\left(1\right)
        \underbrace{\left\|  \left( \mathrm{I}_n - Z\Sigma_{\leq \tilde{\ell}} Z^{\top} \reg(Z\Sigma_{\leq \tilde{\ell}} Z^{\top}) \right) 
                Z\theta_{\leq \tilde{\ell}} \right\|_{2}^{2}}_{B_{1,3}}\\ 
                &~+
                O_{d, \mathbb{P}}\left(  1\right)
                \underbrace{\left\| K_{\tilde{\ell}+1, \ell_{\gamma}} K^{-1} Z\Sigma_{\leq \tilde{\ell}} Z^{\top}
                Z \Sigma_{\leq \tilde{\ell}}^{-1} \theta_{\leq \tilde{\ell}} \right\|_{2}^{2}}_{B_{1,4}}\\
                &~+
                O_{d, \mathbb{P}}\left(  1\right)
                \underbrace{\left\| K_{>\ell_{\gamma}} K^{-1} Z\Sigma_{\leq \tilde{\ell}} Z^{\top}
                Z \Sigma_{\leq \tilde{\ell}}^{-1} \theta_{\leq \tilde{\ell}} \right\|_{2}^{2}}_{B_{1,5}}.
\end{aligned}
\end{equation*}
\fi

Let $A=Z \Sigma_{\leq \tilde{\ell}} Z^{\top}$ and $B=K_{>\tilde{\ell}} / n$. Then $A+B=Z \Sigma_{\leq \tilde{\ell}} Z^{\top}+K_{>\tilde{\ell}} / n=K/n$. 
By
Proposition \ref{prop_matrix_prop_1_of_filter}, we have 
\begin{equation}\label{eq:filter_implication}
\begin{aligned}
   &~ \mathrm{I}_n - Z\Sigma_{\leq \tilde{\ell}} Z^{\top} \reg(Z\Sigma_{\leq \tilde{\ell}} Z^{\top} +K_{>\tilde{\ell}}/n)\\
 =&~ Z\Sigma_{\leq \tilde{\ell}} Z^{\top} (K/n)^{-1}\left( \mathrm{I}_n - (K/n) \reg(K/n) \right)  + (K_{>\tilde{\ell}}/n) (K/n)^{-1}. 
\end{aligned}
 \end{equation}
It follows that
\begin{equation*}
\begin{aligned}
      \|B_1\|_{2}^{2}  =&~
      \Theta_{d, \mathbb{P}}\left(1\right)
        \left\| \left( \mathrm{I}_n - Z\Sigma_{\leq \tilde{\ell}} Z^{\top} \reg(Z\Sigma_{\leq \tilde{\ell}} Z^{\top} +K_{>\tilde{\ell}}/n) \right) 
                Z\theta_{\leq \tilde{\ell}} \right\|_{2}^{2}\\
        \leq &~
      O_{d, \mathbb{P}}\left(1\right)
        \left\| Z\Sigma_{\leq \tilde{\ell}} Z^{\top} (K/n)^{-1}\left( \mathrm{I}_n - (K/n) \reg(K/n) \right) 
                Z\theta_{\leq \tilde{\ell}} \right\|_{2}^{2}\\ 
                &~+  
                O_{d, \mathbb{P}}\left(  1\right)
                \left\| (K_{>\tilde{\ell}}/n) (K/n)^{-1} Z\Sigma_{\leq \tilde{\ell}} Z^{\top} 
                Z (Z^{\top} 
                Z)^{-1} \Sigma_{\leq \tilde{\ell}}^{-1} \theta_{\leq \tilde{\ell}} \right\|_{2}^{2}\\
                 \leq  &~  O_{d, \mathbb{P}}\left(1\right)
        \underbrace{\left\| \left( \mathrm{I}_n - (K/n) \reg(K/n) \right) 
                Z\theta_{\leq \tilde{\ell}} \right\|_{2}^{2}}_{B_{1,3}}\\ 
                &~+
                O_{d, \mathbb{P}}\left(  1\right)
                \underbrace{\left\| K_{\tilde{\ell}+1, \ell_{\gamma}} K^{-1} Z\Sigma_{\leq \tilde{\ell}} Z^{\top}
                Z (Z^{\top} 
                Z)^{-1} \Sigma_{\leq \tilde{\ell}}^{-1} \theta_{\leq \tilde{\ell}} \right\|_{2}^{2}}_{B_{1,4}}\\
                &~+
                O_{d, \mathbb{P}}\left(  1\right)
                \underbrace{\left\| K_{>\ell_{\gamma}} K^{-1} Z\Sigma_{\leq \tilde{\ell}} Z^{\top}
                Z (Z^{\top} 
                Z)^{-1} \Sigma_{\leq \tilde{\ell}}^{-1} \theta_{\leq \tilde{\ell}} \right\|_{2}^{2}}_{B_{1,5}},
\end{aligned}
\end{equation*}
where the first equation follows from the 4th line in Equation~\eqref{eqn_b1_spe_origin}, and the first inequality is because of Equation~\eqref{eq:filter_implication} and $\lambda_{\min}(Z^{\top}Z)=\Theta_{d,\mathbb{P}}(1))$. 
By the same reasoning as in Eq. (\ref{eqn_B11}), we can show that $B_{1,3} = O_{d, \mathbb{P}}\left(  \lambda^{2} d^{(2-\min\{s, 2\})\tilde{\ell} }\right)$. Moreover, we have
\begin{equation}\label{eqn_bias_B1_case2}
    \begin{aligned}
    B_{1,5} =&~ O_{d, \mathbb{P}}\left(  1\right)\left\| K^{-1} Z\Sigma_{\leq \tilde{\ell}} Z^{\top}
                Z (Z^{\top} 
                Z)^{-1} \Sigma_{\leq \tilde{\ell}}^{-1} \theta_{\leq \tilde{\ell}} \right\|_{2}^{2}\\
                 =&~ O_{d, \mathbb{P}}\left(  \frac{1}{n^2}\right)\left\| 
                Z (Z^{\top} 
                Z)^{-1} \Sigma_{\leq \tilde{\ell}}^{-1} \theta_{\leq \tilde{\ell}} \right\|_{2}^{2}   \qquad(\because \text{Eq. (\ref{eqn:spe_properties})})\\
                =&~O_{d, \mathbb{P}}\left(\frac{1}{n^2}\right)\left\| 
                \Sigma_{\leq \tilde{\ell}}^{-1}\theta_{\leq \tilde{\ell}} \right\|_{2}^{2}    \qquad  (\because \lambda_{\max}(Z)=\Theta_{d,\mathbb{P}}(1),~~ \lambda_{\min}(Z^{\top}Z)=\Theta_{d,\mathbb{P}}(1)))\\
                {=} &
                O_{d, \mathbb{P}}\left( d^{(2-\min\{s, 2\})\tilde{\ell}-2\gamma }\right),
\end{aligned}
\end{equation}
where the last equation follows from Lemma \ref{lemma ge l l2} (iii). 

Finally, denote $\tilde{Z}=n^{-1/2}\Psi_{\ell_{\lambda}+1, \ell_{\gamma}}$. 
Recall that $K_{\leq \tilde{\ell}}=Z\Sigma_{\leq \tilde{\ell}} Z^{\top}$. 
Using $K^{-1}K_{\leq \tilde{\ell}}=\mathrm{I}_n-K^{-1}K_{>\tilde{\ell}}$, we have
\begin{equation*}
    \begin{aligned}
    &~ B_{1,4} = \frac{1}{n^2}\left\| K_{\tilde{\ell}+1, \ell_{\gamma}} K^{-1} K_{\leq \tilde{\ell}}
                Z (Z^{\top} 
                Z)^{-1} \Sigma_{\leq \tilde{\ell}}^{-1} \theta_{\leq \tilde{\ell}} \right\|_{2}^{2}\\
    %line 2
    \leq &~ 
    \frac{1}{n^2}\left\| K_{\tilde{\ell}+1, \ell_{\gamma}} 
                Z (Z^{\top} 
                Z)^{-1} \Sigma_{\leq \tilde{\ell}}^{-1} \theta_{\leq \tilde{\ell}} \right\|_{2}^{2}
               \\
             &~  +
                \frac{1}{n^2}\left\| K_{\tilde{\ell}+1, \ell_{\gamma}} K^{-1} K_{>\tilde{\ell}} 
                Z (Z^{\top} 
                Z)^{-1} \Sigma_{\leq \tilde{\ell}}^{-1} \theta_{\leq \tilde{\ell}} \right\|_{2}^{2}\\
    % line 3
%\tilde{\ell}=\ell_{\lambda} \text{ and } 
              (\because  \|K_{\tilde{\ell}+1, \ell_{\gamma}} K^{-1}\|_2 \leq 1)\quad {\leq} &~
                \frac{1}{n^2}\left\| K_{\tilde{\ell}+1, \ell_{\gamma}} 
                Z (Z^{\top} 
                Z)^{-1}\Sigma_{\leq \tilde{\ell}}^{-1} \theta_{\leq \tilde{\ell}} \right\|_{2}^{2}
               \\
               &~  +
                \frac{1}{n^2}\left\| K_{\tilde{\ell}+1, \ell_{\gamma}} K^{-1} K_{\tilde{\ell}+1, \ell_{\gamma}} 
                Z (Z^{\top} 
                Z)^{-1}\Sigma_{\leq \tilde{\ell}}^{-1} \theta_{\leq \tilde{\ell}} \right\|_{2}^{2}\\
                &~+
                \frac{1}{n^2}\left\| K_{>\ell_{\gamma}} 
                Z (Z^{\top} 
                Z)^{-1}\Sigma_{\leq \tilde{\ell}}^{-1} \theta_{\leq \tilde{\ell}} \right\|_{2}^{2}\\
                (\because  \|K_{\tilde{\ell}+1, \ell_{\gamma}} K^{-1}\|_2 \leq 1)\quad {\leq} &~
                \frac{2}{n^2}\left\| K_{\tilde{\ell}+1, \ell_{\gamma}} 
                Z (Z^{\top} 
                Z)^{-1}\Sigma_{\leq \tilde{\ell}}^{-1} \theta_{\leq \tilde{\ell}} \right\|_{2}^{2}
                +
                \frac{1}{n^2}\left\| K_{>\ell_{\gamma}} 
                Z (Z^{\top} 
                Z)^{-1}\Sigma_{\leq \tilde{\ell}}^{-1} \theta_{\leq \tilde{\ell}} \right\|_{2}^{2}\\
            % line 4
                (\because \text{Lemmas } \ref{lemma psi psi top} \text{ and } \ref{lemma ge l l2})\quad {\leq}&~
                \frac{2}{n^2}\left\| K_{\tilde{\ell}+1, \ell_{\gamma}} 
                Z (Z^{\top} 
                Z)^{-1}\Sigma_{\leq \tilde{\ell}}^{-1} \theta_{\leq \tilde{\ell}} \right\|_{2}^{2}
                + O_{d, \mathbb{P}}\left( d^{(2-\min\{s, 2\})\tilde{\ell}-2\gamma }\right)\\
                %%%%% next line
                (\because\text{ Lemmas } \ref{lemma psi top psi} \text{ and } \ref{lemma psi top psi_int}) \quad   {\leq}&~
                O_{d, \mathbb{P}}\left( \frac{1}{n}\right)
                \left\| Z^{\top} \tilde{Z} \Sigma_{\ell_{\lambda}+1, \ell_{\gamma}}^2 \tilde{Z}^{\top} Z  \right\|_{2}
                \left\| (Z^{\top} 
                Z)^{-1}\Sigma_{\leq \tilde{\ell}}^{-1} \theta_{\leq \tilde{\ell}} \right\|_{2}^{2}
                \\   &~
                + O_{d, \mathbb{P}}\left( d^{(2-\min\{s, 2\})\tilde{\ell}-2\gamma }\right)\\
                (\because\text{Lemma } \ref{lemma ge l l2}) \quad  {\leq}&~
                O_{d, \mathbb{P}}\left( d^{(2-\min\{s, 2\})\tilde{\ell}-2\gamma } \left[ n\underbrace{\left\| Z^{\top} \tilde{Z} \Sigma_{\ell_{\lambda}+1, \ell_{\gamma}}^2 \tilde{Z}^{\top} Z  \right\|_{2}}_{B_{1,6}}+1 \right]\right).
\end{aligned}
\end{equation*}
It suffices to show that $B_{1, 6}=o_{d,\mathbb{P}}(n^{-1})$. 

We will use the matrix Bernstein inequality to bound $B_{1, 6}$ as follows. 

% \dhnote{You showed $B_{1, 6} = o_{d, \mathbb{P}}(n^{-1/2})$ below, but this is not enough. You need $B_{1, 6} = O_{d, \mathbb{P}}(n^{-1})$. }

% {\green It is a typo, see below.}

For $i=1,\cdots, n$, denote
\begin{align*}
    h_i =&~ \left[1, \psi_{1,1}(x_i), \cdots, \psi_{\ell_\lambda, N(d, \ell_\lambda)}(x_i) \right]^{\top} \in \mathbb{R}^{N_{\ell_\lambda} \times 1}\\
    \tilde{h}_i =&~ \left[\mu_{\ell_\lambda+1}\psi_{\ell_\lambda+1,1}(x_i), \cdots, \mu_{\ell_{\gamma}}\psi_{\ell_{\gamma}, N(d, \ell_{\gamma})}(x_i) \right]^{\top} \in \mathbb{R}^{(N_{\ell_{\gamma}}-N_{\ell_\lambda}) \times 1}\\
    A_i =&~ n^{-1} h_i \tilde{h}_i^{\top} \in \mathbb{R}^{N_{\ell_\lambda} \times (N_{\ell_{\gamma}}-N_{\ell_\lambda})}.
\end{align*}

    By the addition formula (see, e.g., Proposition 1.18 in \cite{gallier2009notes}, page 56 in \cite{ghorbani2021linearized}, or (39) in \cite{lu2024pinsker}), we have
    \begin{equation}\label{eqn_add_on_tilde_h}
    \begin{aligned}
        \tilde{h}_i^{\top} \tilde{h}_i =  \sum_{k=\ell_{\lambda}+1}^{\ell_{\gamma}} \mu_k^2 \sum_{j=1}^{N(d, k)} \psi_{k, j}^2(x_i)
        = \sum_{k=\ell_{\lambda}+1}^{\ell_{\gamma}} \mu_k^2 N(d, k),
    \end{aligned}
    \end{equation}
    and
    \begin{equation}\label{eqn_add_on_h}
    \begin{aligned}
        \|h_i h_i^{\top}\|_{2} = h_i^{\top} h_i = \sum_{k=1}^{\ell_{\lambda}} \sum_{j=1}^{N(d, k)} \psi_{k, j}^2(x_i)
        = \sum_{k=1}^{\ell_{\lambda}} N(d, k) = N_{\ell_{\lambda}}.
    \end{aligned}
    \end{equation}

Therefore, we have $B_{1, 6} = \|\sum\nolimits_{i=1}^{n}A_i\|_{2}^{2}$, and we can compute $\mathbb{E}A_i =~ 0$, 
\begin{align*}
    \|A_i\|_{2} =&~ \sqrt{\|A_i A_i^{\top}\|_{2}} = \frac{1}{n}\sqrt{\tilde{h}_i^{\top}\tilde{h}_i \|h_i h_i^{\top}\|_{2}}\\
   (\because \eqref{eqn_add_on_tilde_h} \text{ and } \eqref{eqn_add_on_h})\quad =  &~ \frac{1}{n}\sqrt{\sum_{k=\ell_{\lambda}+1}^{\ell_{\gamma}} \mu_k^2 N(d, k) N_{\ell_{\lambda}}}\\
(\because \text{Prop. 2.1})   =  &~ O_d(d^{-1 / 2 - \gamma}),\\
\end{align*}
and similarly
\begin{align*}
\mathbb{E}[A_i A_i^{\top}] \overset{(\ref{eqn_add_on_tilde_h})}{=}&~ 
    n^{-2} \sum_{k=\ell_{\lambda}+1}^{\ell_{\gamma}} \mu_k^2 N(d, k)\mathbb{E} [h_i h_i^{\top}]=
    O_d(d^{-(\ell_\lambda +1)-2\gamma}) \mathrm{I}_{N_{\ell_\lambda}},\\
    %
    %
    \mathbb{E} [A_i^{\top} A_i] \overset{(\ref{eqn_add_on_h})}{=}&~ 
    n^{-2} N_{\ell_{\lambda}} \mathbb{E} [\tilde{h}_i \tilde{h}_i^{\top}] = \frac{N_{\ell_\lambda}}{n^2}\Sigma_{\ell_{\lambda}+1, \ell_{\gamma}}^2=O_d(d^{- (\ell_{\lambda}+2) - 2\gamma}) \mathrm{I}_{N_{\ell_{\gamma}}-N_{\ell_\lambda}}.
\end{align*}
\color{black}
Therefore, from Proposition \ref{prop_non_sym_matrix_bernstein}, we have
\begin{align*}
\mathbb{P}\left(\left\|\sum_{i=1}^n A_i\right\|_2 \geq n^{-1/2} \right) & \leq 2N_{\ell_{\gamma}} \exp\left(-\frac{n^{-1}}{O_d(d^{-\ell_\lambda -1-\gamma}) + O_d(d^{-1/2-3\gamma / 2})}\right) \\
&=\exp \left( O_d(\ell_{\gamma} \log d) - \Theta_d (d^{\min(\ell_\lambda+1, 1/2+\gamma/2} \right) \\
&= o_d(1),
\end{align*}
and hence $B_{1, 6} = \|\sum\nolimits_{i=1}^{n}A_i\|_{2}^{2} = o_{d, \mathbb{P}}(n^{-1})$.


\vspace{10pt}

\noindent {\it Proof of the case $0 \leq u < \gamma <1$. }

If $0 \leq u < \gamma <1$, we have 
$\ell_{\gamma}=\ell_{\lambda}=0$, from (\ref{eqn_finite_tau_bound}), we have
\begin{equation}\label{eqn_lower_bound_H_S0}
    \begin{aligned}
    H =&~ h_1 = \Theta_{d, \mathbb{P}}(1)\\
    \Sigma_{\leq \tilde{\ell}} =&~ \Sigma_{\leq 0} = \mu_0 = \Omega_{d, \mathbb{P}}( 1 )\\
    S_0 =&~ HV^{\top}\Sigma_{\leq \tilde{\ell}} VH = HV^{\top}\Sigma_{\leq 0} VH = \mu_0 h_1^2 = \Theta_{d, \mathbb{P}}( 1 ) = \Omega_{d, \mathbb{P}}(\lambda).
\end{aligned}
\end{equation}



Following the same derivation in Eq. (\ref{eqn_b1_spe_origin}) but with the triangle inequality applied in an opposite direction, when $\tau < \infty$ we have  

\begin{equation*}
    \begin{aligned}
        \|B_{1}\|_{2}^{2} \geq &~ \Omega_{d, \mathbb{P}}\left(1\right)\left\| \rem(K/n)
                Z\theta_{\leq 0} \right\|_{2}^{2} - \left\| B_{1, 2} \right\|_{2}^{2}\\
                = &~ \Omega_{d, \mathbb{P}}\left(1\right)\left\| \rem(K/n)
                Z\theta_{\leq 0} \right\|_{2}^{2} - \Delta_1\\
                \overset{\text{Assumption } 5}{\geq} &~
                \Omega_{d, \mathbb{P}}\left(1\right)\left\| \rem(K_{0}/n)
                Z\theta_{\leq 0} \right\|_{2}^{2} - \Delta_1\\
                { \overset{ (\ref{eqn_finite_tau_bound})}{=}}&~
                \Omega_{d, \mathbb{P}}\left(1\right)\left\| \rem(U \text{Diag}\{S_{0}, 0_{n-1} \} U^{\top})
                U (H, 0)^{\top} V^{\top}\theta_{\leq 0} \right\|_{2}^{2} - \Delta_1\\
                { \overset{(\ref{eq:Filter_Rem_finite_case3})}{=}}
        &~\Omega_{d, \mathbb{P}}\left(1\right)\left[ \rem\left(S_0 \right)  H V^{\top}\theta_{\leq 0} \right]^{2}- \Delta_1\\
        { 
        \overset{(\ref{eqn_lower_bound_H_S0})}{=}
        }&~
        \Omega_{d, \mathbb{P}}\left(1\right)\left[ \rem\left(\mu_0 h_1^2 \right)  h_1 \theta_{\leq 0} \right]^{2}- \Delta_1\\
        % line 5
        {
        \overset{(9) \text{ and } (\ref{eqn_lower_bound_H_S0})}{\geq}
        }
        &~
         \Omega_{d, \mathbb{P}}\left(\lambda^{2\tau}\right) \theta_{\leq 0}^{2}- \Delta_1\\
         = &~ \Omega_{d, \mathbb{P}}\left(\lambda^{2\tau}\right) \left[ \Sigma_{\leq 0}^{-\tau}\theta_{\leq 0} \right]^{2}- \Delta_1\\
                \overset{\text{Lemma } \ref{lemma ge l l2}}{=} &~
                \Omega_{d, \mathbb{P}}\left(  \lambda^{2\tau} \right)- \Delta_1.
    \end{aligned}
\end{equation*}
where the third last and the second last lines use the fact that $\Sigma_{\leq 0} = \mu_0 = \Theta_d(1)$ obtained by Proposition 2.1, and (\ref{eqn_bound_b_12}) implies $0 \leq \Delta_1 = \left\| B_{1, 2} \right\|_{2}^{2} = O_{d, \mathbb{P}}\left( d^{-2\gamma }\right)$.


% {\color{red} where the inequality $\mathbf{(A)}$ follows from (9) and 
% $S_0 = \mu_0 V^2 H^2 = \Omega_{d, \mathbb{P}}( 1 ) = \Omega_{d, \mathbb{P}}( \lambda )$.} \dhnote{What does this equation mean??}

% {\green [In (\ref{eqn_lower_bound_H_S0}) I provided a detailed proof.]}



Similarly to (\ref{eqn_b1_spe_origin}) and above, for any $\tau \leq \infty$, we have
\begin{equation*}
    \begin{aligned}
        \|B_{1}\|_{2}^{2} \leq &~  O_{d, \mathbb{P}}\left(1\right)\left\| \rem(K/n)
                Z\theta_{\leq 0} \right\|_{2}^{2} + \Delta_1\\
        (\because \rem(z) \text{ is non-increasing  w.r.t } z)\quad \leq&~
        O_{d, \mathbb{P}}\left(1\right)\left\| \rem(K_{0}/n)
                Z\theta_{\leq 0} \right\|_{2}^{2} + \Delta_1\\
                \overset{(8)}{\leq} &~
                O_{d, \mathbb{P}}\left(  \lambda^{2\tau^{\prime}} \right) + \Delta_1.
    \end{aligned}
\end{equation*}








% Combining the above three cases:
In summary, we have completed the proof for the following statements about $B_{1}$:
\begin{itemize}
    \item We have
    \begin{equation}\label{eqn_bound_B_1_spe}
    \begin{aligned}
                \|B_1\|_{2}^{2}  =&~ 
                O_{d, \mathbb{P}}\left(d^{-2u + (2-\min\{s, 2\})\tilde{\ell} } + d^{(2-\min\{s, 2\})\tilde{\ell}-2\gamma }\right);
    \end{aligned}
\end{equation}

    \item If $0 < u < \gamma <1$, then for any $\tau \leq \infty$, we have
    \begin{equation*}
    \begin{aligned}
                \|B_1\|_{2}^{2}  \leq &~ 
                O_{d, \mathbb{P}}\left(  d^{-2\tau^{\prime} u} \right) + \Delta_1;\\
                \|B_1\|_{2}^{2}  \geq &~ 
                \Omega_{d, \mathbb{P}}\left(  d^{-2\tau u} \right)- \Delta_1, \quad \text{ if } \tau<\infty,
    \end{aligned}
\end{equation*}
where $0 \leq \Delta_1 = O_{d, \mathbb{P}}\left( d^{-2\gamma }\right)$.

\iffalse
\item When $\ell_\lambda < \ell$ and $\tau<\infty$, we have
\begin{equation*}
    \begin{aligned}
        \left\|B_1\right\|_2^2=&~\Omega_{d, \mathbb{P}}\left(\lambda^{2 \tau} d^{(2 \tau-\min \{s, 2 \tau\}) \tilde{\ell}}\right)-O_{d, \mathbb{P}}\left(  d^{-2-\min\{s, 2\}\tilde{\ell} }\right)\\
        =&~\Omega_{d, \mathbb{P}}\left(\lambda^{2 \tau} d^{(2 \tau-\min \{s, 2 \tau\}) \tilde{\ell}}\right)-o_{d, \mathbb{P}}\left(  \lambda^{2} d^{(2-\min\{s, 2\})\tilde{\ell} }\right).
    \end{aligned}
\end{equation*}
\fi
\end{itemize}

% \newpage 

\noindent {\bf Bounding $\|B_2\|_{2}^{2}$. }
We have
\begin{equation}\label{eqn_bound_B_2_spe}
    \begin{aligned}
                \|B_2\|_{2}^{2} 
=&~  
\theta_{>\tilde{\ell}}^\top \Psi_{>\tilde{\ell}}^\top \varphi_{\lambda, K} ^\top
\Psi_{\leq \tilde{\ell}} 
\Sigma_{\leq \tilde{\ell}} 
{\Sigma_{\leq \tilde{\ell}} \Psi^{\top}_{\leq \tilde{\ell}} \varphi_{\lambda, K} \Psi_{>\tilde{\ell}} \theta_{>\tilde{\ell}}}
                 \\
=&~ \theta_{>\tilde{\ell}}^\top \Psi_{>\tilde{\ell}}^\top \varphi_{\lambda, K} ^\top
\Psi_{\leq \tilde{\ell}} 
\Sigma_{\leq \tilde{\ell}} 
(\Psi_{\leq \tilde{\ell}}^\top \Psi_{\leq \tilde{\ell}} ) (\Psi_{\leq \tilde{\ell}}^\top \Psi_{\leq \tilde{\ell}} )^{-2}(\Psi_{\leq \tilde{\ell}}^\top \Psi_{\leq \tilde{\ell}} ) 
{\Sigma_{\leq \tilde{\ell}} \Psi^{\top}_{\leq \tilde{\ell}} \varphi_{\lambda, K} \Psi_{>\tilde{\ell}} \theta_{>\tilde{\ell}}}                 
                 \\
                \leq &~ \left\|\left(\Psi_{\leq \tilde{\ell}}^\top \Psi_{\leq \tilde{\ell}} \right)^{-2}\right\|_2 
                \left\|(\Psi_{\leq \tilde{\ell}}^\top \Psi_{\leq \tilde{\ell}} ) 
{\Sigma_{\leq \tilde{\ell}} \Psi^{\top}_{\leq \tilde{\ell}} \varphi_{\lambda, K} \Psi_{>\tilde{\ell}} \theta_{>\tilde{\ell}}}   \right\|^2 
%%
%%
%                 \left[\theta_{>\tilde{\ell}}^\top \Psi_{>\tilde{\ell}}^\top \varphi_{\lambda, K} ^\top
% \Psi_{\leq \tilde{\ell}} 
% \Sigma_{\leq \tilde{\ell}} 
% (\Psi_{\leq \tilde{\ell}}^\top \Psi_{\leq \tilde{\ell}} ) (\Psi_{\leq \tilde{\ell}}^\top \Psi_{\leq \tilde{\ell}} ) 
% {\Sigma_{\leq \tilde{\ell}} \Psi^{\top}_{\leq \tilde{\ell}} \varphi_{\lambda, K} \Psi_{>\tilde{\ell}} \theta_{>\tilde{\ell}}}   \right]              
               \\
                \leq &~ \left\|\left(\Psi_{\leq \tilde{\ell}}^\top \Psi_{\leq \tilde{\ell}} \right)^{-2}\right\|_2 \|\Psi_{\leq \tilde{\ell}}\Psi_{\leq \tilde{\ell}}^\top \|_2
                \left\|K_{\leq \tilde{\ell}} \varphi_{\lambda, K} \Psi_{>\tilde{\ell}} \theta_{>\tilde{\ell}} \right\|^2    \\ \overset{(\ref{eqn_bound_eigen_psi_top_psi})}{=}
                &~ 
             O_{d, \mathbb{P}}\left(\frac{1}{n}\right)\left( \Psi_{>\tilde{\ell}} \theta_{>\tilde{\ell}}  \right)^{\top} \left( \varphi_{\lambda, K} K_{\leq \tilde{\ell}}^2   \varphi_{\lambda, K}  \right) \left( \Psi_{>\tilde{\ell}} \theta_{>\tilde{\ell}}  \right)\\
             \overset{(\ref{eqn:spe_properties})}{=}&~
             O_{d, \mathbb{P}} \left( \frac{1}{n} \right) \cdot \|\Psi_{>\tilde{\ell}} \theta_{>\tilde{\ell}} \|_{2}^{2}
            \overset{\text{Lemma } \ref{lemma ge l l2}}{=}
            O_{d, \mathbb{P}}\left(  d^{-(\tilde{\ell}+1)s}\right).
    \end{aligned}
\end{equation}

\noindent {\bf Bounding $\|B_3\|_{2}^{2}$. }
We have
\begin{equation}\label{eqn_bound_B_3_spe}
    \begin{aligned}
                \|B_3\|_{2}^{2} 
        \overset{\text{Lemma } \ref{lemma ge l l2}}{=}&~
        \Theta_{d, \mathbb{P}} \left( d^{-(\tilde{\ell}+1)s}\right).
    \end{aligned}
\end{equation}



\noindent {\bf Bounding $\|B_4\|_{2}^{2}$. }
We have
 % [below is wrong; DH.  2025-09-18]
\begin{equation*}
    \begin{aligned}
   \|B_4\|_{2}^{2} = &~   
  \theta_{\leq \tilde{\ell}}^{\top} \Psi_{\leq \tilde{\ell}}^{\top} \varphi_{\lambda, K}^{\top}\left(\Psi_{>\tilde{\ell}} \Sigma_{>\tilde{\ell}}^2 \Psi_{>\tilde{\ell}}^{\top}\right) \varphi_{\lambda, K} \Psi_{\leq \tilde{\ell}} \theta_{\leq \tilde{\ell}}
  \\
                \overset{(\ref{eqn_bound_eigen_psi_top_psi})}{\leq} &~
        O_{d, \mathbb{P}} \left( n\|\Sigma_{> \tilde{\ell}}\|_{\mathrm{op}} \right)
        \|\Sigma_{>\tilde{\ell}}^{1/2} \Psi^{\top}_{>\tilde{\ell}} \varphi_{\lambda, K} Z \left(\Sigma_{\leq \tilde{\ell}} Z^{\top} Z (Z^{\top} 
                Z)^{-1} \Sigma_{\leq \tilde{\ell}}^{-1}\right) \theta_{\leq \tilde{\ell}}\|_{2}^{2}
        \\
        \leq &~
        O_{d, \mathbb{P}} \left( n\|\Sigma_{> \tilde{\ell}}\|_{\mathrm{op}} \right)  
        \left\| Z\Sigma_{\leq \tilde{\ell}} Z^{\top} \varphi_{\lambda, K} K_{> \tilde{\ell}}  \varphi_{\lambda, K} Z \Sigma_{\leq \tilde{\ell}} Z^{\top}\right\|_{2} 
        \left\| Z (Z^{\top} 
                Z)^{-1} \Sigma_{\leq \tilde{\ell}}^{-1} \theta_{\leq \tilde{\ell}}  \right\|_{2}^{2}\\
       \overset{\text{Lemma } \ref{lemma ge l l2}}{\leq} &~
        O_{d, \mathbb{P}} \left( d^{\gamma-\tilde{\ell}-1} \right)  
        \left\| Z\Sigma_{\leq \tilde{\ell}} Z^{\top} \varphi_{\lambda, K} K_{> \tilde{\ell}}  \varphi_{\lambda, K} Z \Sigma_{\leq \tilde{\ell}} Z^{\top}\right\|_{2} 
        \left\| Z (Z^{\top} 
                Z)^{-1} \Sigma_{\leq \tilde{\ell}}^{-1} \theta_{\leq \tilde{\ell}}  \right\|_{2}^{2}\\
        \overset{(\ref{eqn_bound_eigen_psi_top_psi})}{=}&~
        O_{d, \mathbb{P}} \left( d^{\gamma-\tilde{\ell}-1} \right)  
        \left\| Z\Sigma_{\leq \tilde{\ell}} Z^{\top} \varphi_{\lambda, K} K_{> \tilde{\ell}}  \varphi_{\lambda, K} Z \Sigma_{\leq \tilde{\ell}} Z^{\top}\right\|_{2} 
        \left\| \Sigma_{\leq \tilde{\ell}}^{-1} \theta_{\leq \tilde{\ell}}  \right\|_{2}^{2}\\
        \overset{\text{Lemma } \ref{lemma ge l l2}}{=}&~ 
         O_{d, \mathbb{P}} \left(  d^{(2-\min\{s, 2\})\tilde{\ell} +\gamma - \tilde{\ell}-1 } \right)\left\| Z\Sigma_{\leq \tilde{\ell}} Z^{\top} \varphi_{\lambda, K} K_{> \tilde{\ell}}  \varphi_{\lambda, K} Z \Sigma_{\leq \tilde{\ell}} Z^{\top}\right\|_{2}\\
         {=}&~ 
         O_{d, \mathbb{P}} \left(  d^{(2-\min\{s, 2\})\tilde{\ell} -\gamma - \tilde{\ell}-1 } \right)\left\| K_{\leq \tilde{\ell}} \varphi_{\lambda, K} K_{> \tilde{\ell}}  \varphi_{\lambda, K} K_{\leq \tilde{\ell}}\right\|_{2}.
    \end{aligned}
\end{equation*}


We derive the bound on $\left\| K_{\leq \tilde{\ell}} \varphi_{\lambda, K} K_{> \tilde{\ell}}  \varphi_{\lambda, K} K_{\leq \tilde{\ell}}\right\|_{2}$ in two cases. 

\newpage

\textit{Case 1:} When $\tilde{\ell} = \ell_{\gamma} \leq \ell_{\lambda}$, we have $u\geq \gamma$ and 
    \begin{align*}
        &~\left\| K_{\leq \tilde{\ell}} \varphi_{\lambda, K} K_{> \tilde{\ell}}  \varphi_{\lambda, K} K_{\leq \tilde{\ell}}\right\|_{2}\\
       (\text{Lemma } \ref{lemma psi psi top}) \quad {=} &~O_{d, \mathbb{P}}(1) \left\| K_{\leq \tilde{\ell}} \varphi_{\lambda, K}  \varphi_{\lambda, K} K_{\leq \tilde{\ell}}\right\|_{2} = O_{d, \mathbb{P}}(1) \left\|  \varphi_{\lambda, K} K_{\leq \tilde{\ell}}\right\|_{2}^{2}\\
        (\text{By Eq.~\eqref{eqn:spe_properties}}) \quad {=}
        &~
        O_{d, \mathbb{P}}(1) \left\|  M^{-1} \Psi_{\leq \tilde{\ell}} \Sigma_{\leq \tilde{\ell}}^{1/2} \Sigma_{\leq \tilde{\ell}}^{1/2} \Psi_{\leq \tilde{\ell}}^{\top} \right\|_{2}^{2}
        \\
       (\text{Lemma } \ref{lemma trans in V2})\quad {=}
        &~
        O_{d, \mathbb{P}}(1) \left\|M_{>\tilde{\ell}}^{-1} \Psi_{\leq \tilde{\ell}} \Sigma_{\leq \tilde{\ell}}^{1/2} \left( \mathrm{I}_{N_{\tilde{\ell}}} +  \Sigma_{\leq \tilde{\ell}}^{1/2}  \Psi_{\leq \tilde{\ell}}^{\top} M_{>\tilde{\ell}}^{-1} \Psi_{\leq \tilde{\ell}} \Sigma_{\leq \tilde{\ell}}^{1/2}\right)^{-1} \Sigma_{\leq \tilde{\ell}}^{1/2} \Psi_{\leq \tilde{\ell}}^{\top}\right\|_{2}^{2}\\
        =&~
        O_{d, \mathbb{P}}(d^{-2\gamma}) \left\|M_{>\tilde{\ell}}^{-1} \Psi_{\leq \tilde{\ell}}  \left( \frac{\Sigma_{\leq \tilde{\ell}}^{-1}}{n} +  \Psi_{\leq \tilde{\ell}}^{\top} M_{>\tilde{\ell}}^{-1} \Psi_{\leq \tilde{\ell}} /n \right)^{-1} \Psi_{\leq \tilde{\ell}}^{\top}\right\|_{2}^{2}\\
      (\because \text{Lemma \ref{lemma lambda max min} \& Eq.~\eqref{eqn_bound_eigen_psi_top_psi}}) \quad {=}
        &~
        O_{d, \mathbb{P}}(d^{-\gamma}) \left\|  \left( \frac{\Sigma_{\leq \tilde{\ell}}^{-1}}{n} +  \Psi_{\leq \tilde{\ell}}^{\top} M_{>\tilde{\ell}}^{-1} \Psi_{\leq \tilde{\ell}} /n \right)^{-1} \Psi_{\leq \tilde{\ell}}^{\top}\right\|_{2}^{2}\\
        (\because \text{Lemma \ref{lemma lambda max min} and }u\geq \gamma)\quad  {=}
        &~
         O_{d, \mathbb{P}}(d^{-\gamma}) \left\|  \Psi_{\leq \tilde{\ell}}^{\top}\right\|_{2}^{2} \\
         (\text{Eq.~\eqref{eqn_bound_eigen_psi_top_psi}}) \quad &~ =  O_{d, \mathbb{P}}(1).
    \end{align*}



\textit{Case 2:} When $\ell_{\gamma} > \ell_{\lambda}$, we have
    \begin{align*}
        &~\left\| K_{\leq \tilde{\ell}} \varphi_{\lambda, K} K_{> \tilde{\ell}}  \varphi_{\lambda, K} K_{\leq \tilde{\ell}}\right\|_{2}\\
        \overset{\text{Lemma } \ref{lemma lambda max min}}{=} &~O_{d, \mathbb{P}}\left({d^{\gamma-\ell_{\lambda}-1}}\right) \left\|  \varphi_{\lambda, K} K_{\leq \tilde{\ell}}\right\|_{2}^{2}\\
        \overset{(\ref{eqn:spe_properties})}{=}
        &~
        O_{d, \mathbb{P}}\left({d^{\gamma-\ell_{\lambda}-1}}\right) \left\|  M^{-1} \Psi_{\leq \tilde{\ell}} \Sigma_{\leq \tilde{\ell}}^{1/2} \Sigma_{\leq \tilde{\ell}}^{1/2} \Psi_{\leq \tilde{\ell}}^{\top} \right\|_{2}^{2}
        \\
        \overset{\text{Lemma } \ref{lemma trans in V2}}{=}
        &~
        O_{d, \mathbb{P}}\left({d^{\gamma-\ell_{\lambda}-1}}\right) \left\|M_{>\tilde{\ell}}^{-1} \Psi_{\leq \tilde{\ell}} \Sigma_{\leq \tilde{\ell}}^{1/2} \left( \mathrm{I}_{N_{\tilde{\ell}}} +  \Sigma_{\leq \tilde{\ell}}^{1/2}  \Psi_{\leq \tilde{\ell}}^{\top} M_{>\tilde{\ell}}^{-1} \Psi_{\leq \tilde{\ell}} \Sigma_{\leq \tilde{\ell}}^{1/2}\right)^{-1} \Sigma_{\leq \tilde{\ell}}^{1/2} \Psi_{\leq \tilde{\ell}}^{\top}\right\|_{2}^{2}\\
        =&~
        O_{d, \mathbb{P}}\left({d^{-\gamma-\ell_{\lambda}-1}}\right) \left\|M_{>\tilde{\ell}}^{-1} \Psi_{\leq \tilde{\ell}}  \left( \frac{\Sigma_{\leq \tilde{\ell}}^{-1}}{n} +  \Psi_{\leq \tilde{\ell}}^{\top} M_{>\tilde{\ell}}^{-1} \Psi_{\leq \tilde{\ell}} /n \right)^{-1} \Psi_{\leq \tilde{\ell}}^{\top}\right\|_{2}^{2}\\
        \overset{\text{Lemma \ref{lemma lambda max min} and (\ref{eqn_bound_eigen_psi_top_psi})}}{=}
        &~
        O_{d, \mathbb{P}}\left({d^{-\ell_{\lambda}-1}} (n\lambda)^{-2}\right) \left\|  \left( \frac{\Sigma_{\leq \tilde{\ell}}^{-1}}{n} +  \Psi_{\leq \tilde{\ell}}^{\top} M_{>\tilde{\ell}}^{-1} \Psi_{\leq \tilde{\ell}} /n \right)^{-1} \Psi_{\leq \tilde{\ell}}^{\top}\right\|_{2}^{2}\\
        \overset{\text{Lemma \ref{lemma lambda max min}}}{=}
        &~
         O_{d, \mathbb{P}}\left({d^{-\ell_{\lambda}-1}}\right)) \left\|  \Psi_{\leq \tilde{\ell}}^{\top}\right\|_{2}^{2} \overset{\text{ (\ref{eqn_bound_eigen_psi_top_psi})}}{=}  O_{d, \mathbb{P}}\left({d^{\gamma-\ell_{\lambda}-1}}\right).
    \end{align*}


\bigskip 

Combining both cases, we have
\begin{equation}\label{eqn_bound_B_4_spe}
    \|B_4\|_{2}^{2} = O_{d, \mathbb{P}} \left(  d^{(2-\min\{s, 2\})\tilde{\ell} - (\gamma  + \ell_{\gamma} + 1)\mathbf{1}\{\ell_{\gamma} \leq \ell_{\lambda}\} - 2(\ell_{\lambda}+1)\mathbf{1}\{\ell_{\gamma} > \ell_{\lambda}\} } \right).
\end{equation}





\newpage


\noindent {\bf Bounding $\|B_5\|_{2}^{2}$. }
We have
\begin{equation}\label{eqn_bound_B_5_spe}
    \begin{aligned}
\|B_5\|_{2}^{2} 
        =&~
        \left( \Psi_{>\tilde{\ell}} \theta_{>\tilde{\ell}}  \right)^{\top} 
        \left( \varphi_{\lambda, K}   \Psi_{>\tilde{\ell}} \Sigma_{>\tilde{\ell}}^2 \Psi_{>\tilde{\ell}}^{\top} \varphi_{\lambda, K} \right) 
        \left( \Psi_{>\tilde{\ell}} \theta_{>\tilde{\ell}}  \right)\\
        \leq &~
        \|\Sigma_{> \tilde{\ell}}\|_{\mathrm{op}} \cdot \|\varphi_{\lambda, K} K_{> \tilde{\ell}}  \varphi_{\lambda, K}\|_{2} \cdot \|\Psi_{>\tilde{\ell}} \theta_{>\tilde{\ell}}\|_{2}^{2}\\
        \overset{\text{Lemma } \ref{lemma ge l l2}}{\leq}&~
         O_{d, \mathbb{P}} \left(  d^{-(\tilde{\ell}+1)(s+1) + \gamma} \right) \cdot \|\varphi_{\lambda, K} K  \varphi_{\lambda, K}\|_{2}\\
         =&~
         O_{d, \mathbb{P}} \left(  d^{-(\tilde{\ell}+1)(s+1) + \gamma} \right) \cdot \|\varphi_{\lambda, K}^{1/2} \left(\varphi_{\lambda, K}^{1/2} K \varphi_{\lambda, K}^{1/2}\right)  \varphi_{\lambda, K}^{1/2}\|_{2}\\
         \leq &~
         O_{d, \mathbb{P}} \left(  d^{-(\tilde{\ell}+1)(s+1) + \gamma} \right) \cdot \|\varphi_{\lambda, K}\|_{2}\\
         \overset{(\ref{eqn:spe_properties})}{=} &~
         O_{d, \mathbb{P}} \left(  d^{-(\tilde{\ell}+1)(s+1) + \gamma} \right) \cdot \|M^{-1}\|_{2}\\
         \overset{\text{Lemma } \ref{lemma lambda max min}}{=}&~
         O_{d, \mathbb{P}}\left( \frac{n }{d^{(\tilde{\ell}+1)(s+1)} (n\lambda+1)}\right) = O_{d, \mathbb{P}}\left( \frac{1}{d^{\tilde{\ell}+1}(\lambda +n^{-1})}\cdot d^{-(\tilde{\ell}+1)s} \right)\\
         \overset{(\ref{eqn:regular_para_upper})}{=}&~
         o_{d, \mathbb{P}}\left(   d^{-(\tilde{\ell}+1)s} \right).
    \end{aligned}
\end{equation}
\iffalse
Now we begin to bound $\mathrm{bias}^2(\hat{f}_{\lambda})$. 
From (\ref{eq derivation of bias_spe}),
(\ref{eqn_bound_B_1_spe}), 
(\ref{eqn_bound_B_2_spe}), 
(\ref{eqn_bound_B_3_spe}), 
(\ref{eqn_bound_B_4_spe}), and
(\ref{eqn_bound_B_5_spe}), 
we have
\begin{equation*}
\begin{aligned}
    \mathrm{bias}^{2}(\hat{f}_{\lambda})=  &~
    O_{d, \mathbb{P}}\left(d^{-(\tilde{\ell}+1)s} + \lambda^{2}d^{(2-\min\{s, 2\})\tilde{\ell} }\right)\\
    &~
    + \mathbf{1}\{\gamma > \tilde{\ell}\}O_{d, \mathbb{P}} \left( d^{-\min\{s, 2\}\tilde{\ell} + \tilde{\ell}-\gamma -1 } \right)
    +
        \mathbf{1}\{\gamma = \tilde{\ell}\}
        O_{d, \mathbb{P}}\left(d^{-\min\{s, 2\})\gamma}\right).
\end{aligned} 
\end{equation*}
\fi

\paragraph{\bf Case (II). } When $\ell_{\lambda} \geq \gamma = \ell_{\gamma}$ and $n < 2N(d, \ell_{\gamma})$,
%$\alpha \gamma ! < 2$, for any $d \geq \mathfrak{C}^{\prime}$, a sufficiently large constant, we have $n = \alpha d^{\ell}(1+o_{d}(1))< 2N(d, \ell) = 2d^{\ell}(\ell !)^{-1}(1+o_d(1))$. Therefore, we assume that $d \geq \mathfrak{C}^{\prime}$.
denote ${N}_{\underline{\ell_{\gamma}}}=\sum_{k=0}^{\ell_{\gamma}-1}N(d,k)+n/4 < N_{\ell_{\gamma}}$ and
$$
\Psi_{\leq {\underline{\ell_{\gamma}}}} = \left(\Psi_{1}, \cdots, \Psi_{\ell_{\gamma}-1}, \left(\Psi_{k,1}, \cdots, \Psi_{k,n/4} \right) \right) \in \mathbb{R}^{n \times {N}_{\underline{\ell_{\gamma}}}};
$$
then, Lemma \ref{lemma psi top psi_int} implies that 
\begin{equation*}
    \lambda_{\text{max}}(\Psi_{\leq \underline{\ell_{\gamma}}}^{\top}\Psi_{\leq \underline{\ell_{\gamma}}}) = \Theta_{d, \mathbb{P}}(n)
    \quad \text{ and } \quad
    \lambda_{\text{min}}(\Psi_{\leq \underline{\ell_{\gamma}}}^{\top}\Psi_{\leq \underline{\ell_{\gamma}}}) = \Theta_{d, \mathbb{P}}(n).
\end{equation*}
Similar to Appendix \ref{append_notation}, we can denote $\Psi_{> {\underline{\ell_{\gamma}}}}$, $\Sigma_{\leq {\underline{\ell_{\gamma}}}}$, $\Sigma_{> {\underline{\ell_{\gamma}}}}$, $\theta_{\leq {\underline{\ell_{\gamma}}}}$, $\theta_{> {\underline{\ell_{\gamma}}}}$, $K_{\leq {\underline{\ell_{\gamma}}}}$, and $K_{> {\underline{\ell_{\gamma}}}}$. 

Notice that we have ${N}_{\underline{\ell_{\gamma}}} = \Theta_d(n)=\Theta_d(d^{\ell_{\gamma}}) = \Theta_d(N(d, \ell_{\gamma})) = \Theta_d(N_{\ell_{\gamma}})$, hence from Lemma \ref{lemma psi top psi_int}, similar to the proof of (iv) in Lemma \ref{lemma lambda max min}, we have
\begin{align*}
        \lambda_{\text{max}}(M_{>\underline{\ell_{\gamma}}}) 
        = \Theta_{d, \mathbb{P}}(n\lambda)
        \quad \text{ and } \quad
        \lambda_{\text{min}}(M_{>\underline{\ell_{\gamma}}}) = \Theta_{d, \mathbb{P}}(n\lambda).
    \end{align*}
Similarly, we can show that Lemma \ref{lemma lambda max min}(v), Lemma \ref{lemma ge l l2}, and Lemma \ref{lemma trans in V2} hold with $\tilde{\ell}$ replaced by $\underline{\ell_{\gamma}}$. Therefore, following the above proof for {\bf Case (I)}, we can obtain the same bounds for $\|B_i\|$ with $i=1, \cdots, 5$.
\end{proof}





\subsection{Bounds of Bias term in the under-regularized regimes}

The following proposition is a modification of Appendix D.4.2 in \cite{lu2024saturation}. It provides a tight bound when $u^{\prime}(\tau) \leq u < u^{\prime}(1)$ and $\gamma \geq 1$, where $u^{\prime} = u^{\prime}(\tau)$ is the optimal rate of the regularization parameter $\lambda$ with qualification $\tau \leq \infty$ and coefficient of source condition $s>0$ defined in \cite{lu2024saturation} (ignoring the logarithm term).


\begin{proposition}\label{prop_verify_bias_special}
    Suppose Assumptions 1--6 hold.
Suppose that $\tau \leq \infty$, $s>0$,
\iffalse
one of the following conditions hold:
\begin{itemize}
    \item $\tau=\infty$ and $s > 1$;
    \item $\tau<\infty$ and $s > 0$.
\end{itemize}
Further suppose that 
\fi
$u^{\prime}(\tau) \leq u < u^{\prime}(1)$, and $\gamma \geq 1$. 
Then, the bias term (15) of the spectral algorithm estimator $\hat{f}_{\lambda}$ satisfies:

    \begin{equation*}
\begin{aligned}
    \mathrm{bias}^{2}(\hat{f}_{\lambda})  \leq  &~
    \Theta_d\left( 
    d^{-(\tilde{\ell} + 1)  s}
+ \mathbf{1}\{\tau<\infty\} d^{-2  \tau  u + (2  \tau - \tilde{s})  \tilde{\ell}} 
\right) 
    + 
    \Delta_2,\\
    \mathrm{bias}^{2}(\hat{f}_{\lambda})  \geq  &~
    \Theta_d\left( 
    d^{-(\tilde{\ell} + 1)  s}
+ \mathbf{1}\{\tau<\infty\} d^{-2  \tau  u + (2  \tau - \tilde{s})  \tilde{\ell}}
\right) 
    -
    \Delta_2,
\end{aligned} 
\end{equation*}
where $0 \leq \Delta_2 = o_{d, \mathbb{P}} \left(
d^{\gamma - \tilde{\ell} - 1 - 2  \max\{\gamma - u, 0\}}
+ d^{\tilde{\ell} - \gamma}
+ d^{-(\tilde{\ell} + 1)  s}
+ \mathbf{1}\{\tau<\infty\} d^{-2  \tau  u + (2  \tau - \tilde{s})  \tilde{\ell}}
\right)$.

\end{proposition}
\begin{proof} 
We adopt notations in \cite{lu2024saturation}, and especially:
   \begin{align*}
       &~ \mathcal{N}_{1, \varphi}(\lambda), \mathcal{N}_{2, \varphi}(\lambda), \mathcal{M}_{1, \varphi}(\lambda), 
       \mathcal{M}_{2, \varphi}(\lambda) \text{ in (42) of \cite{lu2024saturation}}\\
       &~ \mathcal{N}_{1}(\lambda) := \mathcal{N}_{1, \varphi^{\krr}}(\lambda) \text{ below (42) of \cite{lu2024saturation}}\\
       &~ \tilde{f}_{\lambda}, f_{\lambda} \text{ in (36) and (38) of \cite{lu2024saturation}}.
   \end{align*}

We also note that the following notations are different in our paper and in \cite{lu2024saturation, zhang2024optimal}:
\begin{itemize}
    \item In \cite{lu2024saturation}, $\ell:=\lfloor u \rfloor$; while in our paper, we denote $\ell_{\lambda}:=\lfloor u \rfloor$.

    \item 
    Consider the case where $0<s<1$. 
    In this case,  $\tilde{s}:=\min(s, 2\tau)=s$, so the notation $\tilde{p} = \lfloor \gamma / (\tilde{s}+1) \rfloor$ used in \cite{zhang2024optimal} and in Proposition \ref{prop_e_1} is the same as the notation $p = \lfloor \gamma / (s+1) \rfloor$ in \cite{lu2024saturation}.
    Notice that one can show that $p= \tilde{p} < \gamma$ (see Proposition \ref{prop_e_1}).

% \dhnote{
% Do you suggest that $p$ here is the $\tilde{p}$ in Proposition \ref{prop_e_1}? If so, state it?
% }

\end{itemize}




Without loss of generality, we assume that $\lambda=d^{-u}$.
We divide the proof into the following three cases: (i) $1 \leq s \leq \tau$, (ii) $\tilde{s}:=\min\{s, 2\tau\} > \tau$, and (iii) $0<s<1$.

\noindent {\bf Case (i): $1 \leq s \leq \tau$. }  
    Notice that:
    \begin{itemize}
        \item From (36) and (41) of \cite{lu2024saturation}, we have that $\mathrm{bias}(\hat{f}_{\lambda})$ equals to the bias term defined in (41) of \cite{lu2024saturation}.

        \item From (60) of \cite{lu2024saturation}, we have that $\mathrm{bias}(\hat{f}_{\lambda}) \leq \left\| f_{\lambda} - f_{\star}\right\|_{L^{2}} 
    +
    \left\| \tilde{f}_{\lambda} - f_{\lambda}\right\|_{L^{2}}$ and $\mathrm{bias}(\hat{f}_{\lambda}) \geq \left\| f_{\lambda} - f_{\star}\right\|_{L^{2}} 
    -
    \left\| \tilde{f}_{\lambda} - f_{\lambda}\right\|_{L^{2}}$. Moreover, the latter two terms can be upper bounded separately:
    \begin{itemize}
        \item From Lemma D.6 and Lemma D.14 in \cite{lu2024saturation}, we have 
        $$
        \left\| f_{\lambda} - f_{\star}\right\|_{L^{2}}^2 = \mathcal{M}_{2, \varphi}(\lambda) = \Theta_d\left( 
    d^{-(\ell_{\lambda}+1)s}
    +\mathbf{1}\{\tau<\infty\}\lambda^{2\tau}d^{(2\tau-\tilde{s})\ell_{\lambda} }\right).
        $$

        \item Suppose the following inequalities hold (which we will prove later)
        \begin{equation}\label{eqn:check_condition_bias_1}
      \begin{aligned}
          \frac{ \mathcal{N}_{1}(\lambda) \mathcal{M}_{1, \varphi}^2(\lambda)}{n^2}
    &= 
    o_{d}\left( \mathcal{M}_{2, \varphi}(\lambda) + \frac{\sigma^2}{n} \calN_{2,\varphi}(\lambda)  \right),\\
    \frac{\mathcal{N}_{1}(\lambda)}{n} \ln(n) (\ln\lambda^{-1})^2 \cdot \sum_{i=1}^\infty \frac{\lambda^2 \lambda_i \reg^2(\lambda_i)}{\lambda + \lambda_i} f_i^2  
    &= 
    o_{d}\left( \mathcal{M}_{2, \varphi}(\lambda) + \frac{\sigma^2}{n} \calN_{2,\varphi}(\lambda)  \right).
      \end{aligned}
  \end{equation}
  Then, on one hand, since $u < u^{\prime}(1) < \gamma$, from Lemma 23 in \cite{zhang2024optimal}, there exists $\epsilon>0$, such that
        $$
        n^{\epsilon - 1}\mathcal{N}_{1}(\lambda) = \Theta_d(n^{\epsilon - 1} \lambda^{-1})= o_d(1);
        $$
    On the other hand, following the proof of Lemma D.7 in \cite{lu2024saturation}, with (63) and (64) of \cite{lu2024saturation} being 
    %replaced by 
    restated as
    (\ref{eqn:check_condition_bias_1}),
 we have 
 \begin{equation*}
        \left\| \tilde{f}_{\lambda} - f_{\lambda}\right\|_{L^{2}}^2 = o_{d, \bbP}\left( \mathcal{M}_{2, \varphi}(\lambda) + \frac{\sigma^2}{n} \calN_{2,\varphi}(\lambda)  \right).
    \end{equation*}
  
    \end{itemize}
        

        

        \item Proposition \ref{prop_e_1} implies that $u < u^{\prime}(1) < \gamma$, hence $\tilde{\ell}=\ell_{\lambda}$.
    \end{itemize}

Therefore, from the above arguments, we only need to verify that (\ref{eqn:check_condition_bias_1}) hold.


\begin{remark}
    We remark that the above argument is the same as that in \cite{lu2024saturation}, where they verify an analogous version of (\ref{eqn:check_condition_bias_1}) 
    via their Lemmas D.21-D.23.
\end{remark}




  

Since $\tau \geq s$, we have $(\lambda d^{\ell_{\lambda}})^{\tau} \leq (\lambda d^{\ell_{\lambda}})^{s}$, hence from Lemma D.14 in \cite{lu2024saturation} we have
\begin{align*}
        \frac{ \mathcal{N}_{1}(\lambda) \mathcal{M}_{1, \varphi}^2(\lambda)}{n}
        &=
        O_d\left(
        \lambda^{2(s-1)} d^{-\gamma+\ell_{\lambda}s} \mathbf{1}\{\tau < \infty\}
        +
        \lambda^{-1} d^{-\gamma -(\ell_{\lambda}+1)(s-1)}
        \right)\\
        \mathcal{N}_{1}(\lambda) \ln(n) (\ln\lambda^{-1})^2 \cdot \sum_{i=1}^\infty \frac{(\lambda)^2 \lambda_i \varphi_{\lambda}^2(\lambda_i)}{\lambda + \lambda_i} f_i^2
        &=
        O_d\left((\ln(d))^3 \right)
    \cdot
    O_d\left(
        \lambda d^{\max\{\ell_{\lambda}(2-s), 0\}} + \lambda^{-1}d^{-s(\ell_{\lambda}+1)}
    \right),
    \end{align*}
Denote $\mathbf{I}=\lambda^{2(s-1)} d^{-\gamma+\ell_{\lambda}s}$,
$\mathbf{II}=\lambda^{-1} d^{-\gamma -(\ell_{\lambda}+1)(s-1)}$,
$\mathbf{III}=\lambda d^{\max\{\ell_{\lambda}(2-s), 0\}}(\ln(d))^3$,\\
and $\mathbf{IV}=\lambda^{-1}d^{-s(\ell_{\lambda}+1)}(\ln(d))^3$. Denote $A_1 = d^{2u-\ell_{\lambda}-1}$, $A_2 = d^{\ell_{\lambda}}$, $A_3 = d^{\gamma-s(\ell_{\lambda}+1)}$. 
Then, (\ref{eqn:check_condition_bias_1}) is implied by the following relationships:
\begin{equation*}
    \begin{aligned}
        n^{-1} \cdot \left(\mathbf{I} + \mathbf{II}\right) =&~ o_{d}\left( \mathcal{M}_{2, \varphi}(\lambda) + \frac{\sigma^2}{n} \calN_{2,\varphi}(\lambda)  \right), \\
        n^{-1} 
    \cdot \left(\mathbf{III} + \mathbf{IV}\right) =&~ o_{d}\left( \mathcal{M}_{2, \varphi}(\lambda) + \frac{\sigma^2}{n} \calN_{2,\varphi}(\lambda)  \right). 
    \end{aligned}
\end{equation*}



\begin{itemize}
    \item Recall that Proposition \ref{prop_e_1} implies that $u < u^{\prime}(1) < \gamma$. Hence, from Lemma D.14 in \cite{lu2024saturation}, we have that 
        \begin{align*}
          &~~  n \cdot \left(\mathcal{M}_{2, \varphi}(\lambda) +  \frac{\sigma^2}{n}\mathcal{N}_{2, \varphi}(\lambda)\right)\\
          = &~
    n \cdot \Theta_{d} \left( \frac{n}{d^{\ell_{\lambda}+1}(n\lambda+1)^{2}} + \frac{d^{\ell_{\lambda}}}{n} + 
    d^{-(\ell_{\lambda}+1)s}
    +\mathbf{1}\{\tau<\infty\}\lambda^{2\tau}d^{(2\tau-\tilde{s})\ell_{\lambda} }\right)\\
    = &~ \Omega_d (A_1 + A_2 + A_3).
        \end{align*}
      
        


    \iffalse
    \item We have
    $$
    n\left( \mathcal{M}_{2, \varphi}(\lambda) + \frac{\sigma^2}{n} \calN_{2,\varphi}(\lambda)  \right) = \Omega_d (A_1 + A_2 + A_3).
    $$
    \fi

    \item We have
    $$
        \mathbf{I}=d^{-2u(s-1) -\gamma +\ell_{\lambda}s} = d^{\ell_{\lambda}} d^{-u(s-1) -\gamma}d^{(\ell_{\lambda}-u)(s-1)} = o_d( d^{\ell_{\lambda}})=o_d(A_2).
    $$

    \item Proposition \ref{prop_e_1} implies that $u < u^{\prime}(1) < \gamma$, hence
    $$
    \mathbf{II} = d^{\ell_{\lambda}}d^{u - \gamma}d^{-s(\ell_{\lambda}+1) + 1} \overset{s \geq 1}{=} o_d( d^{\ell_{\lambda}})=o_d(A_2).
    $$

    \item When $s\geq 2$, we have 
    $$
    \mathbf{III} = d^{-u} (\ln(d))^3 \overset{u>0}{=} o_d( d^{\ell_{\lambda}})=o_d(A_2);
    $$
    When $s<2$, we have
    $$
    \mathbf{III} = d^{-u + \ell_{\lambda}(2-s)} (\ln(d))^3 = d^{\ell_{\lambda}} d^{-u - \ell_{\lambda}(s-1)} (\ln(d))^3 = o_d( d^{\ell_{\lambda}})=o_d(A_2).
    $$
    
    \item Proposition \ref{prop_e_1} implies that $u < u^{\prime}(1) < \gamma$, hence
    $$
    \mathbf{IV}=\lambda^{-1}d^{-s(\ell_{\lambda}+1)}(\ln(d))^3 = o_d( d^{\gamma-s(\ell_{\lambda}+1)} )=o_d(A_3).
    $$
    
\end{itemize}
Combining all these, we obtained (\ref{eqn:check_condition_bias_1}).



\noindent {\bf Case (ii): $\tilde{s}:=\min\{s, 2\tau\} > \tau$. } Similar to {\bf Case (i)}, we only need to show that (\ref{eqn:check_condition_bias_1}) holds.
From Lemma D.14 in \cite{lu2024saturation} we have
\begin{align*}
      &  \frac{ \mathcal{N}_{1}(\lambda) \mathcal{M}_{1, \varphi}^2(\lambda)}{n}
        =
        O_d\left(
        \lambda^{2(\tau-1)} d^{-\gamma+\ell_{\lambda}(2\tau-\tilde{s})}
        +
        \lambda^{-1} d^{-\gamma -(\ell_{\lambda}+1)(s-1)}\mathbf{1}\{s \leq 2\tau\}
        \right),  \\
      &  \mathcal{N}_{1}(\lambda) \ln(n) (\ln\lambda^{-1})^2 \cdot \sum_{i=1}^\infty \frac{(\lambda)^2 \lambda_i \varphi_{\lambda}^2(\lambda_i)}{\lambda + \lambda_i} f_i^2
         =
        O_d\left((\ln(d))^3 \right)
    \cdot
    O_d\left(
        \lambda d^{\max\{\ell_{\lambda}(2-s), 0\}} + \lambda^{-1}d^{-s(\ell_{\lambda}+1)}
    \right).
    \end{align*}
Denote $\mathbf{I}=\lambda^{2(\tau-1)} d^{-\gamma+\ell_{\lambda}(2\tau-\tilde{s})}$,
$\mathbf{II}=\lambda^{-1} d^{-\gamma -(\ell_{\lambda}+1)(s-1)}$,
$\mathbf{III}=\lambda d^{\max\{\ell_{\lambda}(2-s), 0\}}(\ln(d))^3$,\\
and $\mathbf{IV}=\lambda^{-1}d^{-s(\ell_{\lambda}+1)}(\ln(d))^3$.
Denote $A_1 = d^{2u-\ell_{\lambda}-1}$, $A_2 = d^{\ell_{\lambda}}$, $A_3 = d^{\gamma-s(\ell_{\lambda}+1)}$.
Then, (\ref{eqn:check_condition_bias_1}) is implied by the following relationships:
\begin{equation*}
    \begin{aligned}
        n^{-1} \cdot \left(\mathbf{I} + \mathbf{II}\right) =&~ o_{d}\left( \mathcal{M}_{2, \varphi}(\lambda) + \frac{\sigma^2}{n} \calN_{2,\varphi}(\lambda)  \right)\\
        n^{-1} 
    \cdot \left(\mathbf{III} + \mathbf{IV}\right) =&~ o_{d}\left( \mathcal{M}_{2, \varphi}(\lambda) + \frac{\sigma^2}{n} \calN_{2,\varphi}(\lambda)  \right)
    \end{aligned}
\end{equation*}


\begin{itemize}
    \item Recall that Proposition \ref{prop_e_1} implies that $u < u^{\prime}(1) < \gamma$. Hence, from Lemma D.14 in \cite{lu2024saturation}, we have that 
        \begin{align*}
           &~ n \cdot \left(\mathcal{M}_{2, \varphi}(\lambda) +  \frac{\sigma^2}{n}\mathcal{N}_{2, \varphi}(\lambda)\right)\\
           = &~
    n \cdot \Theta_{d} \left( \frac{n}{d^{\ell_{\lambda}+1}(n\lambda+1)^{2}} + \frac{d^{\ell_{\lambda}}}{n} + 
    d^{-(\ell_{\lambda}+1)s}
    +\mathbf{1}\{\tau<\infty\}\lambda^{2\tau}d^{(2\tau-\tilde{s})\ell_{\lambda} }\right)\\
    = &~ \Omega_d (A_1 + A_2 + A_3).
        \end{align*}
    
    \item When $s >2\tau$, we have
    $$
        \mathbf{I}=\lambda^{2(\tau-1)} d^{-\gamma} {=} o_d( d^{\ell_{\lambda}})=o_d(A_2);
    $$
    When $\tau < s \leq 2\tau$, we have $(\lambda d^{\ell_{\lambda}})^{2\tau} \leq (\lambda d^{\ell_{\lambda}})^{s}$, hence
    
    \begin{equation*}
        \begin{aligned}
            \mathbf{I} =&~ (\lambda d^{\ell_{\lambda}})^{2\tau} d^{2u-\gamma -s\ell_{\lambda}} = O_d(d^{2u-us-\gamma}) = O_d(d^{2u-\ell_{\lambda}-1})O_d(d^{\ell_{\lambda}-us}d^{1-\gamma})\\ (\because \gamma \geq 1,~ s> 1) \quad = &~ O_d(d^{2u-\ell_{\lambda}-1}) o_d(d^{\ell_{\lambda}-s})=o_d(A_1).
        \end{aligned}
    \end{equation*}



    
    \item Same as the proof in {\bf Case (i)}, we have
    $$
    \mathbf{II}, \mathbf{III}, \mathbf{IV} = o_d(d^{\ell_{\lambda}} + d^{\gamma-s(\ell_{\lambda}+1)} ) = o_d(A_2 + A_3).
    $$
\end{itemize}
Combining all these, we obtained (\ref{eqn:check_condition_bias_1}).



\noindent {\bf Case (iii): $0< s <1$. } 
The proof is based on the following  argument, which is similar to the one in {\bf Case (i)}:  
% Notice that:
% {\green
%     [Remark: I have updated the following arguments and deleted the duplicate parts]
% }
    \begin{itemize}
        \item []
        
        Suppose there exists a constant $\epsilon>0$ only depending on $s$ and $\gamma$, such that the following inequalities hold (which we will prove later)
\begin{equation}\label{eqn:check_condition_bias_3}
      \begin{aligned}
   \frac{\mathcal{N}_{1}(\lambda)}{n} \ln(n) (\ln\lambda^{-1})^2 \cdot \sum_{i=1}^\infty \frac{\lambda^2 \lambda_i \reg^2(\lambda_i)}{\lambda + \lambda_i} f_i^2  
    &= o_{d} 
    \left( \mathcal{M}_{2, \varphi}(\lambda) + \frac{\sigma^2}{n} \calN_{2,\varphi}(\lambda)  \right);\\
    n^{-2} \mathcal{N}_{1}(\lambda) 
    \left(
    \left\| f_{\lambda} \right\|_{L^{\infty}}^2  
    +
    n^{1-s+\epsilon}
    \right)
    &= o_{d}\left( \mathcal{M}_{2, \varphi}(\lambda) + \frac{\sigma^2}{n} \calN_{2,\varphi}(\lambda)  \right).
      \end{aligned}
  \end{equation}
  
  Then, following the proof of Lemma D.8 in \cite{lu2024saturation}, where (75) and (76) therein are restated as (\ref{eqn:check_condition_bias_3}), we obtain the desired bounds for the bias term $\mathrm{bias}^{2}(\hat{f}_{\lambda})$.
  

\iffalse   
        \item From (36) and (41) of \cite{lu2024saturation}, we have that $\mathrm{bias}(\hat{f}_{\lambda})$ equals to the bias term defined in (41) of \cite{lu2024saturation}.

        \item From (60) of \cite{lu2024saturation}, we have that $\mathrm{bias}(\hat{f}_{\lambda}) \leq \left\| f_{\lambda} - f_{\star}\right\|_{L^{2}} 
    +
    \left\| \tilde{f}_{\lambda} - f_{\lambda}\right\|_{L^{2}}$. Moreover, 
    \begin{itemize}
        \item From Lemma D.6 in \cite{lu2024saturation}, we have 
        $$
        \left\| f_{\lambda} - f_{\star}\right\|_{L^{2}} = \mathcal{M}_{2, \varphi}(\lambda)^{\frac{1}{2}}.
        $$

        \item Since $u < u^{\prime}(1) < \gamma$, from Lemma 23 in \cite{zhang2024optimal}, there exists $\epsilon>0$, such that
        $$
        n^{\epsilon - 1}\mathcal{N}_{1}(\lambda) = \Theta_d(n^{\epsilon - 1} \lambda^{-1})= o_d(1).
        $$

        \item Following the proof of Lemma D.8 in \cite{lu2024saturation}, with (75) and (76) of \cite{lu2024saturation} being replaced by the following condition: 
        there exists a constant $\epsilon>0$ only depending on $s$ and $\gamma$, such that we have

 Then we have 
 \begin{equation*}
        \left\| \tilde{f}_{\lambda} - f_{\lambda}\right\|_{L^{2}}^2 = O_{d, \bbP}\left( \mathcal{M}_{2, \varphi}(\lambda) + \frac{\sigma^2}{n} \calN_{2,\varphi}(\lambda)  \right).
    \end{equation*}

  
    \end{itemize}
        

        \item From Lemma D.14 in \cite{lu2024saturation}, we have that 
        $$
      \mathcal{M}_{2, \varphi}(\lambda) +  \frac{\sigma^2}{n}\mathcal{N}_{2, \varphi}(\lambda) = 
    \Theta_{d} \left( \frac{n}{d^{\ell_{\lambda}+1}(n\lambda+1)^{2}} + \frac{d^{\ell_{\lambda}}}{n} + 
    d^{-(\ell_{\lambda}+1)s}
    +\mathbf{1}\{\tau<\infty\}\lambda^{2\tau}d^{(2\tau-\tilde{s})\ell_{\lambda} }\right).
        $$

        \item Proposition \ref{prop_e_1} implies that $u < u^{\prime}(1) < \gamma$, hence $\tilde{\ell}=\ell_{\lambda}$.

        \fi
    \end{itemize}


Therefore, we only need to establish (\ref{eqn:check_condition_bias_3}). 


Proposition \ref{prop_e_1} implies that $u < u^{\prime}(1) < \gamma$, hence from Lemma D.14 in \cite{lu2024saturation} and Lemma 23 in \cite{zhang2024optimal}, we have
    \begin{align*}
        \mathcal{N}_{1}(\lambda) \ln(n) (\ln\lambda^{-1})^2 \cdot \sum_{i=1}^\infty \frac{(\lambda)^2 \lambda_i \varphi_{\lambda}^2(\lambda_i)}{\lambda + \lambda_i} f_i^2
        &=
        O_{d}\left((\ln(d))^3 \right)
    \cdot
    O_{d}\left(
        \lambda d^{\ell_{\lambda}(2-s)} + \lambda^{-1}d^{-s(\ell_{\lambda}+1)}
    \right)\\
    n^{-1} \mathcal{N}_{1}(\lambda) 
    n^{1-s}
        &=
        O_{d}\left(
        \lambda^{-1} 
         d^{-\gamma s}
        \right).
    \end{align*}

Moreover, notice that (8) implies $\reg(z) \leq C_2 (z+\lambda)^{-1}$. Hence, similar to the proof of Lemma 27 in \cite{zhang2024optimal} (for kernel ridge regression estimator), we can show that
\begin{align*}
    n^{-1} \mathcal{N}_{1}(\lambda)\left\| f_{\lambda} \right\|_{L^{\infty}}^2 =&~ O_{d}\left(
   d^{-\gamma} \lambda^{-1}  \left\| f_{\lambda} \right\|_{L^{\infty}}^2    \right)\\
   =&~ O_{d}\left(
   d^{-\gamma} \lambda^{-1}  \left[ d^{\frac{(1-s)\ell_{\lambda}}{2}} + \lambda^{-1} d^{-\frac{(1+s)(\ell_{\lambda}+1)}{2}} \right]    \right).
\end{align*}




    

Denote $\mathbf{I}=\lambda d^{\ell_{\lambda}(2-s)} (\ln(d))^3$,
$\mathbf{II}=\lambda^{-1}d^{-s(\ell_{\lambda}+1)}(\ln(d))^3$,  
$\mathbf{III}=\lambda^{-1} 
         d^{-\gamma s}$, 
$\mathbf{IV}= 
         d^{u-\gamma + \frac{(1-s)\ell_{\lambda}}{2}}$,
and $\mathbf{V}= 
         d^{2u-\gamma -\frac{(1+s)(\ell_{\lambda}+1)}{2}}$.     
Denote $A_1 = d^{2u-\ell_{\lambda}-1}$, $A_2 = d^{\ell_{\lambda}}$, $A_3 = d^{\gamma-s(\ell_{\lambda}+1)}$.

Note that (\ref{eqn:check_condition_bias_3}) is implied by the following inequalities:
\begin{equation*}
    \begin{aligned}
        n^{-1} \cdot \left(\mathbf{I} + \mathbf{II}\right) =&~ o_{d}\left( \mathcal{M}_{2, \varphi}(\lambda) + \frac{\sigma^2}{n} \calN_{2,\varphi}(\lambda)  \right)\\
        n^{-1}
    \cdot \left(\mathbf{III} + \mathbf{IV} + \mathbf{V}\right) =&~ o_{d}\left( \mathcal{M}_{2, \varphi}(\lambda) + \frac{\sigma^2}{n} \calN_{2,\varphi}(\lambda)  \right)
    \end{aligned}
\end{equation*}




\begin{itemize}
    \item Recall that Proposition \ref{prop_e_1} implies that $u < u^{\prime}(1) < \gamma$. Hence, from Lemma D.14 in \cite{lu2024saturation}, we have that 
        \begin{align*}
           &~ n \cdot \left(\mathcal{M}_{2, \varphi}(\lambda) +  \frac{\sigma^2}{n}\mathcal{N}_{2, \varphi}(\lambda)\right)\\
           = &~
    n \cdot \Theta_{d} \left( \frac{n}{d^{\ell_{\lambda}+1}(n\lambda+1)^{2}} + \frac{d^{\ell_{\lambda}}}{n} + 
    d^{-(\ell_{\lambda}+1)s}
    +\mathbf{1}\{\tau<\infty\}\lambda^{2\tau}d^{(2\tau-\tilde{s})\ell_{\lambda} }\right)\\
    = &~ \Omega_d (A_1 + A_2 + A_3).
        \end{align*}
    
    \item We have
    $$
    \mathbf{I} = d^{-u + \ell_{\lambda}(2-s)} (\ln(d))^3 = d^{\ell_{\lambda}} d^{-u + \ell_{\lambda}(1-s)} (\ln(d))^3 \overset{s>0}{=} o_d( d^{\ell_{\lambda}})=o_d(A_2).
    $$
    \item Proposition \ref{prop_e_1} implies that $u < u^{\prime}(1) < \gamma$, hence
    $$
    \mathbf{II}=\lambda^{-1}d^{-s(\ell_{\lambda}+1)}(\ln(d))^3 = o_d( d^{\gamma-s(\ell_{\lambda}+1)} )=o_d(A_3).
    $$

    \item For $\mathbf{III}$, we analyze three possible cases:
    \begin{itemize}
        \item Suppose $u-\ell_{\lambda} < \gamma s$. Then
        $$
            \mathbf{III} = O_d(d^{u-\gamma s}) = O_d(d^{\ell_{\lambda}}d^{u-\ell_{\lambda}-\gamma s}) = o_d(d^{\ell_{\lambda}})=o_d(A_2).
        $$

        \item Suppose $u-\ell_{\lambda} \geq \gamma s$ and $s>1/2$. Then
        $$\begin{aligned}
        \mathbf{III} &= O_d(d^{2u-\ell_{\lambda}-1}d^{u-\gamma s -2u+\ell_{\lambda}+1})\\
        & = O_d(d^{2u-\ell_{\lambda}-1}d^{-2\gamma s +1} ) \\
        & \overset{\gamma \geq 1, s>1/2}{=}  o_d(d^{2u-\ell_{\lambda}-1})=o_d(A_1).
            \end{aligned}
        $$

        \item Suppose $u-\ell_{\lambda} \geq \gamma s$ and $s \leq 1/2$. We first prove by contradiction that $\ell_{\lambda}<\ell_{\gamma}$: Suppose that $\ell_{\lambda} = \ell_{\gamma}$, then from {\it Case III} in the proof of Proposition \ref{prop_e_1} with $\tau=1$  and $\tilde{p}=p = \lfloor \gamma / (s+1) \rfloor$, we have
        %and $p:=\lfloor \gamma/(s+1)\rfloor = \lfloor u^{\prime}(1)\rfloor$, we have
        $$
        p:=\lfloor \gamma/(s+1)\rfloor \leq \ell_{\gamma} = \ell_{\lambda} \overset{u<u^{\prime}(1)} \leq p,
        %\leq u^{\prime}(1) < p+1 \leq \ell_{\gamma} +1,
        $$
        implying that $\ell_{\lambda}=p$. Hence, from {\it Case III} in the proof of Proposition \ref{prop_e_1}  with $\tilde{p}=p$, we further have
        \begin{align*}
           &~ \gamma s \leq u-\ell_{\lambda} = u- p\\ 
           = &~
           \left\{\begin{matrix}
\frac{\gamma - p - sp}{2} \leq \frac{s}{2} < \gamma s, & p \geq 1 \text{ and } p(s+1) \leq \gamma < ps+p+s;\\
\frac{s}{2}< \gamma s, & p \geq 0 \text{ and } ps+p+s \leq \gamma < ps+p+s+1,
\end{matrix}\right.
        \end{align*}
leading to a contradiction. (Notice that $p(s+1) \leq \gamma < ps+p+s$ with $p=0$ is not allowed since $\gamma \geq 1 >s$). Therefore, we have $\ell_{\lambda}<\ell_{\gamma}$, and hence
        $$
            \mathbf{III} = O_d(d^{\gamma-s(\ell_{\lambda}+1)}d^{u-\gamma s -\gamma + s(\ell_{\lambda}+1)}) = O_d(d^{\gamma-s(\ell_{\lambda}+1)}d^{u-\gamma}) = o_d(d^{\gamma-s(\ell_{\lambda}+1)}) = o_d(A_3).
        $$
        
    \end{itemize}

    \item We have 
    $$
        \mathbf{IV} = d^{\ell_{\lambda}} d^{u-\gamma + \frac{(1-s)\ell_{\lambda}}{2} - \ell_{\lambda}} \overset{u<\gamma, s<1}{=} o_d(d^{\ell_{\lambda}}) = o_d(A_2).
    $$


    \item Finally, we have
    \begin{align*}
        \mathbf{V} =&~  d^{2u-\ell_{\lambda}-1}  d^{2u-\gamma -\frac{(1+s)(\ell_{\lambda}+1)}{2} - 2u + \ell_{\lambda} + 1} = d^{2u-\ell_{\lambda}-1} d^{-\gamma + \frac{(1-s)(\ell_{\lambda}+1)}{2}}\\
        \overset{0<s<1}{=} &~ o_d(d^{2u-\ell_{\lambda}-1}d^{-\frac{2\gamma - \ell_{\lambda} - 1)}{2}}) = o_d(d^{2u-\ell_{\lambda}-1}) = o_d(A_1).
    \end{align*}
        
    

    
\end{itemize}
Combining all the above steps, we obtained (\ref{eqn:check_condition_bias_3}).    
\end{proof}





\section{Learning curve}\label{append_learning_curve}

In this section, we will use Theorem \ref{thm_variance_spe}, Theorem \ref{thm_bias_spe}, Proposition \ref{prop_nips_left_learning_curve}, and Proposition \ref{prop_verify_bias_special} to determine the convergence rate of the excess risk in (13). 


\subsection{Notations and restatement of results}

We first introduce some notations.


Recall that $u$ is an absolute constant given in Assumption 4, and we have $\lambda = \Theta_d\left(d^{-u}\right)$. Denote
\begin{equation}\label{eqn_order_in_bias_and_variance}
    \begin{aligned}
    &~v_1 = v_1(u) = \gamma - \tilde{\ell} - 1 - 2  \max\{\gamma - u, 0\},\\
    &~v_2 = v_2(u) =  \tilde{\ell} - \gamma,\\
    &~b_1 = b_1(u) = -(\tilde{\ell} + 1)  s,\\
    &~b_2 = b_2(u, t) = -2  t  u + (2  t - \min\{s, 2t\})  \tilde{\ell}, \quad 1 \leq t<\infty,\\
    &~v_1^{\prime} = v_1^{\prime}(u) = \gamma - \ell_{\lambda} - 1 - 2(\gamma - u),\\
    &~v_2^{\prime} = v_2^{\prime}(u) =  \ell_{\lambda} - \gamma,\\
    &~b_1^{\prime} = b_1^{\prime}(u) = -(\ell_{\lambda} + 1)  s,\\
    &~b_2^{\prime} = b_2^{\prime}(u, t) = -2  t  u + (2  t - \min\{s, 2t\})  \ell_{\lambda}, \quad 1 \leq t<\infty,\\
    &~b_3 = b_3(u) =  (2 - \min\{s, 2\})  \tilde{\ell} -2  \gamma,\\
    &~b_4 = b_4(u) = (2 - \min\{s, 2\})  \tilde{\ell} -\gamma - \tilde{\ell} - 1.
    \end{aligned}
\end{equation}
Furthermore, when $s>0$,
%denote $u^{\prime, \star} = u^{\prime, \star} (\tau)= \arg\min\{u \geq 0: \max\{v_1^{\prime}(u), v_2^{\prime}(u)\} = \max\{b_1^{\prime}(u), b_2^{\prime}(u, \tau)\}\}$. 
let $u^{\prime} = u^{\prime}(\tau)$ be the optimal rate of the regularization parameter $\lambda$ with qualification $\tau \leq \infty$ defined in \cite{lu2024saturation} (we will also describe $u^{\prime}$ in the proof of Proposition \ref{prop_e_1}). 

\begin{itemize}
    \item Let's first restate results in Theorem \ref{thm_variance_spe} and Theorem \ref{thm_bias_spe} with notations defined in (\ref{eqn_order_in_bias_and_variance}):
\begin{equation}\label{eqn_learning_curve_pre_results_1}
    \begin{aligned}
        \mathrm{var}(\hat{f}_{\lambda}) = &~ \Theta_{d, \mathbb{P}}\left( d^{v_1} + d^{v_2} \right)\\
        \mathrm{bias}^{2}(\hat{f}_{\lambda})  =  &~
    \|B_1+B_2\|_{2}^{2} + \|B_3-B_4-B_5\|_{2}^{2}\\
    \|B_1\|_{2}^{2}  =&~ 
                O_{d, \mathbb{P}}\left(d^{b_2(u, 1)}+d^{b_3}\right)\\
                \|B_2\|_{2}^{2} =&~ 
            O_{d, \mathbb{P}}\left(  d^{b_1}\right)\\
            \|B_3-B_5\|_{2}^{2} 
        =&~
        \Theta_{d, \mathbb{P}} \left( d^{b_1}\right)\\
        (\because \ell_{\gamma} + 1 > \gamma \text{ and } \ell_{\lambda}+1 > u) \quad \|B_4\|_{2}^{2} =&~ o_{d, \mathbb{P}} \left(d^{b_2(u, 1)}+d^{b_3}\right);
    \end{aligned}
\end{equation}
Moreover, if $0 < u < \gamma <1$, we have $\tilde{\ell}:=\min\{\ell_{\gamma}, \ell_{\lambda}\}=0$. 
Therefore, for any $\tau \leq \infty$, we have

\begin{equation}\label{eqn_learning_curve_pre_results_2}
    \begin{aligned}
    \|B_1\|_{2}^{2}  \geq&~ 
                O_{d, \mathbb{P}}\left( d^{b_2(u, \tau^{\prime})}  \right) + \Delta_1\\
                \|B_1\|_{2}^{2}  =&~ 
                \Omega_{d, \mathbb{P}}\left( d^{b_2(u, \tau)} \right)- \Delta_1, \quad \text{ if } \tau<\infty\\
    \|B_4\|_{2}^{2} =&~ O_{d, \mathbb{P}} \left(d^{-\gamma-1}\right) = O_{d, \mathbb{P}} \left(d^{b_4}\right),           
    \end{aligned}
\end{equation}
where $\tau^{\prime}$ is defined as in (\ref{eqn_def_of_tau_prime}), and $0 \leq \Delta_1 = O_{d, \mathbb{P}}\left( d^{b_3}\right)$.


\item 
We next restate the results in Proposition \ref{prop_nips_left_learning_curve} with notations defined in (\ref{eqn_order_in_bias_and_variance}). Suppose all assumptions and conditions in Proposition \ref{prop_nips_left_learning_curve} hold; this implies that $\tau \leq \infty$, $s>0$, $u \in (0, u^{\prime}(\tau))$, and at least one of the following conditions hold: $\text{(i) } \tau =\infty, \quad \text{(ii) }s>1/(2\tau), \quad \text{(iii) } \gamma > ((2\tau+1)s) / (2\tau(1+s))$).
Then for any $u < u^{\prime}(\tau)$, we have
\begin{equation}\label{eqn_learning_curve_pre_results_3}
    \begin{aligned}
&~E_{x, \epsilon} \left[ \left(\hat{f}_{\lambda}(x) - f_{\star}(x) \right)^{2} \right]\\
=&~
\Theta_{d, \mathbb{P}} \left( d^{v_1^{\prime}} + d^{v_2^{\prime}} + d^{b_1^{\prime}} + \mathbf{1}\{\tau<\infty\}d^{b_2^{\prime}(u, \tau)}\right)\\
(\because \text{ Proposition \ref{prop_e_1}}) \quad =&~
\Theta_{d, \mathbb{P}} \left( d^{v_1} + d^{v_2} + d^{b_1} + \mathbf{1}\{\tau<\infty\}d^{b_2(u, \tau)}\right);
    \end{aligned}
\end{equation}




\item Finally, let's restate results in Proposition \ref{prop_verify_bias_special} with notations defined in (\ref{eqn_order_in_bias_and_variance}). Suppose all assumptions and conditions in Proposition \ref{prop_verify_bias_special} hold; this implies that $\tau \leq \infty$, $s>0$, $u^{\prime}(\tau) \leq u < u^{\prime}(1)$, and $\gamma \geq 1$).
Then for any $u^{\prime}(\tau) \leq u < u^{\prime}(1)$, we have

\begin{equation}\label{eqn_learning_curve_pre_results_4}
    \begin{aligned}
        \mathrm{bias}^{2}(\hat{f}_{\lambda}) \leq &~ \Theta_{d} \left( d^{b_1} + \mathbf{1}\{\tau<\infty\}d^{b_2(u, \tau)}\right) + \Delta_2\\
        \mathrm{bias}^{2}(\hat{f}_{\lambda}) \geq &~ \Theta_{d} \left( d^{b_1} + \mathbf{1}\{\tau<\infty\}d^{b_2(u, \tau)}\right) - \Delta_2;
    \end{aligned}
\end{equation}
where $0 \leq \Delta_2 = o_{d, \mathbb{P}} \left( d^{v_1} + d^{v_2} + d^{b_1} + \mathbf{1}\{\tau<\infty\}d^{b_2(u, \tau)}\right)$


\end{itemize}



















In the next two subsections, we will consider the case with $s=0$ in Appendix \ref{appendix_f2} and with $s>0$ in Appendix \ref{appendix_f3}. 

% the cases where
% \begin{itemize}
%     \item[(E.2)] $s=0$.

%     \item[(E.3)] $s>0$.

%     % \item[(E.4)] $\tau=\infty$ and $s>0$.  
% \end{itemize}


\subsection{$s=0$}\label{appendix_f2}

When $s=0$, Theorem \ref{thm_variance_spe} and Theorem \ref{thm_bias_spe} together imply the following results:

\begin{corollary}\label{coroll_learn_curve_case_1}
    Suppose Assumptions 1,2,3,4,5, and 6 hold with $\tau \leq \infty$ and $s=0$. Then, the excess risk of the spectral algorithm estimator in (13) satisfies
\begin{equation*}
    E_{x, \epsilon} \left[ \left(\hat{f}_{\lambda}(x) - f_{\star}(x) \right)^{2} \right] = \Theta_{d, \mathbb{P}} \left( 1 \right).
\end{equation*}
\end{corollary}
\begin{proof}
    Notice that $b_1=0$ and these four quantities are non-positive:  $v_1, v_2,  b_2(u, 1), b_3$. Hence (\ref{eqn_learning_curve_pre_results_1}) implies that 
    $$
    E_{x, \epsilon} \left[ \left(\hat{f}_{\lambda}(x) - f_{\star}(x) \right)^{2} \right] = O_{d, \mathbb{P}} \left( d^{v_1} + d^{v_2} \right) + O_{d, \mathbb{P}} \left( d^{b_1} + d^{b_2(u, 1)} + d^{b_3} \right) = O_{d, \mathbb{P}} \left( 1 \right),
    $$
    and that
    $$
    E_{x, \epsilon} \left[ \left(\hat{f}_{\lambda}(x) - f_{\star}(x) \right)^{2} \right] \geq 
    \|B_3-B_4-B_5\|_{2}^{2} \geq \|B_3-B_5\|_{2}^{2} - \|B_4\|_{2}^{2}
    = \Omega_{d, \mathbb{P}} \left( 1 \right),
    $$
    and we obtain the desired results.
\end{proof}




% \input{appendix/learn_1}



\subsection{$s>0$}\label{appendix_f3}

We will prove the following results.

\begin{corollary}\label{coroll_learn_curve_case_2}
Suppose Assumptions 1,2,3,4,5, and 6 hold with $s>0$. 
Denote $\tilde{s}=\min\{s, 2\tau\}$.
Assume in addition that either $\tau=\infty$, or $s>1/(2\tau)$, or $\gamma > ((2\tau+1)s) / (2\tau(1+s))$.
Then the excess risk of the spectral algorithm estimator in (13) satisfies
\begin{equation*}
\begin{aligned}
    &~E_{x, \epsilon} \left[ \left(\hat{f}_{\lambda}(x) - f_{\star}(x) \right)^{2} \right]\\
=&~ \Theta_{d, \mathbb{P}} \left( \frac{n}{d^{\tilde{\ell}+1}(n\lambda+1)^{2}} + \frac{d^{\tilde{\ell}}}{n} + 
    d^{-(\tilde{\ell}+1)s}
    +\mathbf{1}\{\tau<\infty\}\lambda^{2\tau}d^{(2\tau-\tilde{s})\tilde{\ell} }\right).
\end{aligned}
\end{equation*}
\end{corollary}


\begin{corollary}\label{coroll_learn_curve_case_2_complex}
    Suppose Assumptions 1,2,3,4,5, and 6 hold with $s>0$. 
Denote $\tilde{s}=\min\{s, 2\tau\}$.
Assume in addition that both $\tau<\infty$, and $s \leq 1/(2\tau)$, and $\gamma \leq ((2\tau+1)s) / (2\tau(1+s))$.
Then, the excess risk of the spectral algorithm estimator in (13) satisfies
\begin{equation*}
\begin{aligned}
    &~E_{x, \epsilon} \left[ \left(\hat{f}_{\lambda}(x) - f_{\star}(x) \right)^{2} \right]\\
=&~ \Theta_{d, \mathbb{P}} \left( \frac{n}{d^{\tilde{\ell}+1}(n\lambda+1)^{2}} + \frac{d^{\tilde{\ell}}}{n} + 
    d^{-(\tilde{\ell}+1)s}
    +\lambda^{2\tau}d^{(2\tau-\tilde{s})\tilde{\ell} }\right).
\end{aligned}
\end{equation*}
\end{corollary}



\begin{proof}[Proof of Corollary \ref{coroll_learn_curve_case_2}]
Let's separate the proof into four parts:
\begin{itemize}
    \item [(i)] $u < u^{\prime}(\tau)$; 
    \item [(ii)] $u^{\prime}(\tau) \geq u^{\prime}(1)$ and $u \geq u^{\prime}(\tau)$;
    \item [(iii)] $u^{\prime}(\tau) < u^{\prime}(1)$ and $u \geq u^{\prime}(1)$; 
    \item [(iv)] $u^{\prime}(\tau) < u^{\prime}(1)$ and $u \in [u^{\prime}(\tau), u^{\prime}(1))$.
\end{itemize}





\noindent {\bf Case (i). } Suppose $u < u^{\prime}(\tau)$. (\ref{eqn_learning_curve_pre_results_3}) implies
\begin{align*}
&~E_{x, \epsilon} \left[ \left(\hat{f}_{\lambda}(x) - f_{\star}(x) \right)^{2} \right]
=
\Theta_{d, \mathbb{P}} \left( d^{v_1} + d^{v_2} + d^{b_1} + \mathbf{1}\{\tau<\infty\}d^{b_2(u, \tau)}\right).
\end{align*}

\noindent {\bf Case (ii). } Suppose $u^{\prime}(\tau) \geq u^{\prime}(1)$. For any $u \geq u^{\prime}(\tau) \geq u^{\prime}(1)$, from Proposition \ref{prop_e_2} and Proposition \ref{prop_e_3}, we have
\begin{equation*}
%\label{eqn_tau_finite_b234_bound}
    \max\{v_1(u), v_2(u), b_1(u)\} \geq \max\{\mathbf{1}\{\tau<\infty\} b_2(u, \tau), b_2(u, 1), b_3(u)\},
\end{equation*}
Therefore from (\ref{eqn_learning_curve_pre_results_1}) we have
\begin{itemize}
    \item Upper bound:
    \begin{align*}
E_{x, \epsilon} \left[ \left(\hat{f}_{\lambda}(x) - f_{\star}(x) \right)^{2} \right]
=
O_{d, \mathbb{P}} \left( d^{v_1} + d^{v_2} + d^{b_1}\right)
=
O_{d, \mathbb{P}} \left( d^{v_1} + d^{v_2} + d^{b_1} + \mathbf{1}\{\tau<\infty\}d^{b_2(u, \tau)}\right);
\end{align*}

\item Lower bound:
\begin{align*}
&~ E_{x, \epsilon} \left[ \left(\hat{f}_{\lambda}(x) - f_{\star}(x) \right)^{2} \right]
\geq  \Omega_{d, \mathbb{P}} \left( d^{v_1} + d^{v_2} \right) + \|B_3-B_4-B_5\|_{2}^{2}\\
= &~
\Omega_{d, \mathbb{P}} \left( d^{v_1} + d^{v_2} + d^{b_1}\right)
=
\Omega_{d, \mathbb{P}} \left( d^{v_1} + d^{v_2} + d^{b_1} + \mathbf{1}\{\tau<\infty\}d^{b_2(u, \tau)}\right).
\end{align*}


\end{itemize}



\noindent {\bf Case (iii). } Suppose $u^{\prime}(\tau) < u^{\prime}(1)$. For any $u \geq u^{\prime}(1) \geq u^{\prime}(\tau)$, from Proposition \ref{prop_e_2} and Proposition \ref{prop_e_3}, we have
\begin{equation*}
%\label{eqn_tau_finite_b234_bound}
    \max\{v_1(u), v_2(u), b_1(u)\} \geq \max\{\mathbf{1}\{\tau<\infty\}b_2(u, \tau), b_2(u, 1), b_3(u)\},
\end{equation*}
Therefore, as in {\bf Case (ii)}, we can show that
\begin{align*}
E_{x, \epsilon} \left[ \left(\hat{f}_{\lambda}(x) - f_{\star}(x) \right)^{2} \right]
=
\Theta_{d, \mathbb{P}} \left( d^{v_1} + d^{v_2} + d^{b_1}\right)
=
\Theta_{d, \mathbb{P}} \left( d^{v_1} + d^{v_2} + d^{b_1} + \mathbf{1}\{\tau<\infty\}d^{b_2(u, \tau)}\right).
\end{align*}


\noindent {\bf Case (iv). } Suppose $u^{\prime}(\tau) < u^{\prime}(1)$ and $u \in [u^{\prime}(\tau), u^{\prime}(1))$. From Proposition \ref{prop_e_2} and Proposition \ref{prop_e_3}, we have
\begin{equation}
\label{eqn_tau_finite_b234_bound}
    \max\{v_1(u), v_2(u), b_1(u)\} \geq \max\{\mathbf{1}\{\tau<\infty\}b_2(u, \tau), b_3(u), b_4(u)\}.
\end{equation}
Therefore, 
\begin{itemize}
    \item[(a)] When $\gamma \geq 1$, from the variance bound in (\ref{eqn_learning_curve_pre_results_1}) and the bias bound in (\ref{eqn_learning_curve_pre_results_4}), we have
    \begin{align*}
E_{x, \epsilon} \left[ \left(\hat{f}_{\lambda}(x) - f_{\star}(x) \right)^{2} \right]
=
\Theta_{d, \mathbb{P}} \left( d^{v_1} + d^{v_2} + d^{b_1} + \mathbf{1}\{\tau<\infty\}d^{b_2(u, \tau)}\right);
\end{align*}

    \item[(b)] When $\gamma<1$ and $\tau<\infty$, Proposition \ref{prop_e_1} with $t=1$ further implies $0<u<u^{\prime}(1) <\gamma<1$. Hence from (\ref{eqn_learning_curve_pre_results_1}) and (\ref{eqn_learning_curve_pre_results_2}) we have
    \begin{itemize}
    \item Upper bound:
    \begin{align*}
E_{x, \epsilon} \left[ \left(\hat{f}_{\lambda}(x) - f_{\star}(x) \right)^{2} \right]
\leq &~ O_{d, \mathbb{P}} \left( d^{v_1} + d^{v_2}\right) + \|B_1\|_{2}^{2} + \|B_2\|_{2}^{2} + \|B_3 - B_5\|_{2}^{2} + \|B_4\|_{2}^{2}\\
= &~
O_{d, \mathbb{P}} \left( d^{v_1} + d^{v_2}\right) + O_{d, \mathbb{P}}\left( d^{b_2(u, \tau)}  + d^{b_3} + d^{b_1} + d^{b_4} \right)\\
(\because \text{ (\ref{eqn_tau_finite_b234_bound})}) \quad = &~ O_{d, \mathbb{P}} \left( d^{v_1} + d^{v_2} + d^{b_1} + d^{b_2(u, \tau)}\right);
\end{align*}

\item Lower bound:
    \begin{itemize}
        \item If $b_1>b_4$, then
\begin{align*}
&~E_{x, \epsilon} \left[ \left(\hat{f}_{\lambda}(x) - f_{\star}(x) \right)^{2} \right]
\geq  \Omega_{d, \mathbb{P}} \left( d^{v_1} + d^{v_2} \right) + \|B_3-B_4-B_5\|_{2}^{2}\\
=&~
\Omega_{d, \mathbb{P}} \left( d^{v_1} + d^{v_2} + d^{b_1}\right)=\Omega_{d, \mathbb{P}} \left( d^{v_1} + d^{v_2} + d^{b_1} + d^{b_2(u, \tau)}\right);
\end{align*}

\item If $b_1 \leq b_4$, since Proposition \ref{prop_e_2} implies that $b_4<v_2$, hence (\ref{eqn_tau_finite_b234_bound}) implies
\begin{equation*}
%\label{eqn_tau_finite_b234_bound}
    \max\{v_1(u), v_2(u)\} \geq \max\{b_1(u), b_2(u, \tau)\}.
\end{equation*}
Therefore, 
\begin{align*}
E_{x, \epsilon} \left[ \left(\hat{f}_{\lambda}(x) - f_{\star}(x) \right)^{2} \right]
\geq
\Omega_{d, \mathbb{P}} \left( d^{v_1} + d^{v_2}\right)
=
\Omega_{d, \mathbb{P}} \left( d^{v_1} + d^{v_2} + d^{b_1} + d^{b_2(u, \tau)}\right),
\end{align*}

    \end{itemize}

\end{itemize}

    \item[(c)] Finally, suppose $\gamma<1$ and $\tau = \infty$. Recall that from (\ref{eqn_def_of_tau_prime}), we have
    $$
        \tau^{\prime}:= 
        \left\{\begin{matrix}
s, &  s \geq 1;\\
1, &  s \in [\frac{1}{2}, 1);\\
(2s)^{-1}, &  0 < s < \frac{1}{2}.
\end{matrix}\right.
        $$


 Therefore, from Proposition \ref{prop_e_1} we have
    \begin{align*}
    u^{\prime}(\infty) =
    \left\{\begin{matrix}
    u^{\prime}(s), & \quad s \geq 1;\\
\gamma/2 = u^{\prime}(1) & \quad s \in [\frac{1}{2}, 1) \text{ and } \gamma < s;\\
\min\{\gamma (1+s)-s, \gamma / 2 \} \geq \frac{s}{2} = u^{\prime}(1) & \quad s \in [\frac{1}{2}, 1) \text{ and } s \leq \gamma < 1;\\
\gamma s = \frac{\gamma}{2(2s)^{-1}} = u^{\prime}((2s)^{-1}) & \quad 0 < s < \frac{1}{2} \text{ and } \gamma < s;\\
\min\{\gamma (1+s)-s, \gamma / 2 \} \geq s^2 = \frac{s}{2(2s)^{-1}} = u^{\prime}((2s)^{-1}) & \quad 0 < s < \frac{1}{2} \text{ and } s \leq \gamma < 1;
\end{matrix}\right.
\end{align*}
implying that $u^{\prime}(\infty) \geq u^{\prime}(\tau^{\prime})$. 
%{\green [Remark: you are right that the second line is $\gamma /2$]}
Furthermore, the assumption of {\bf Case (iv)} implies $u \geq u^{\prime}(\infty)\geq u^{\prime}(\tau^{\prime})$. 
Therefore, from Proposition \ref{prop_e_2} and Proposition \ref{prop_e_3}, we have
\begin{equation*}
%\label{eqn_tau_finite_b234_bound}
    \max\{v_1(u), v_2(u), b_1(u)\} \geq \max\{b_2(u, \tau^{\prime}), b_3(u), b_4(u)\}.
\end{equation*}
From the assumption of {\bf Case (iv)}, we have
$u<u^{\prime}(1)$. From Proposition \ref{prop_e_1} we have $0<u<u^{\prime}(1) <\gamma<1$.

    Hence from (\ref{eqn_learning_curve_pre_results_1}) and (\ref{eqn_learning_curve_pre_results_2}) we have
    \begin{itemize}
    \item Upper bound:
    \begin{align*}
E_{x, \epsilon} \left[ \left(\hat{f}_{\lambda}(x) - f_{\star}(x) \right)^{2} \right]
\leq &~ O_{d, \mathbb{P}} \left( d^{v_1} + d^{v_2}\right) + \|B_1\|_{2}^{2} + \|B_2\|_{2}^{2} + \|B_3 - B_5\|_{2}^{2} + \|B_4\|_{2}^{2}\\
= &~
O_{d, \mathbb{P}} \left( d^{v_1} + d^{v_2}\right) + O_{d, \mathbb{P}}\left( d^{b_2(u, \tau^{\prime})}  + d^{b_3} + d^{b_1} + d^{b_4} \right)\\
= &~ O_{d, \mathbb{P}} \left( d^{v_1} + d^{v_2} + d^{b_1}\right);
\end{align*}

\item Lower bound:
    \begin{itemize}
        \item If $b_1>b_4$, then
\begin{align*}
&~E_{x, \epsilon} \left[ \left(\hat{f}_{\lambda}(x) - f_{\star}(x) \right)^{2} \right]
\geq  \Omega_{d, \mathbb{P}} \left( d^{v_1} + d^{v_2} \right) + \|B_3-B_4-B_5\|_{2}^{2}\\
=&~
\Omega_{d, \mathbb{P}} \left( d^{v_1} + d^{v_2} + d^{b_1}\right);
\end{align*}

\item If $b_1 \leq b_4$, since Proposition \ref{prop_e_2} implies that $b_1 \leq b_4<v_2$, hence we have
\begin{align*}
E_{x, \epsilon} \left[ \left(\hat{f}_{\lambda}(x) - f_{\star}(x) \right)^{2} \right]
\geq
\Omega_{d, \mathbb{P}} \left( d^{v_1} + d^{v_2}\right)
=
\Omega_{d, \mathbb{P}} \left( d^{v_1} + d^{v_2} + d^{b_1}\right).
\end{align*}

    \end{itemize}

\end{itemize}
    
    
\end{itemize}
Combining all above cases, we obtain the desired results.
\end{proof}











\vspace{10pt}


\begin{proof}[Proof of Corollary \ref{coroll_learn_curve_case_2_complex}]
Let's separate the proof into two parts:
\begin{itemize}
    \item [(i)] $u < \gamma$; 
    \item [(ii)] $u \geq \gamma$.
\end{itemize}





\noindent {\bf Case (i). } Suppose $u < \gamma$. Since $s \leq 1/(2\tau)$ and
$$
\gamma \leq \frac{2\tau + 1}{2\tau} \frac{s}{s+1} \leq \frac{2\tau + 1}{2\tau} \frac{\frac{1}{2\tau}}{\frac{1}{2\tau}+1} < 1,
$$
%from Proposition \ref{prop_e_1} we have 
hence we have $0 < u < \gamma <1$. Therefore, 
%from Theorem \ref{thm_variance_spe}, Theorem \ref{thm_bias_spe}, and Proposition \ref{prop_e_2}, 
from (\ref{eqn_learning_curve_pre_results_1}) and (\ref{eqn_learning_curve_pre_results_2})
we have 
\begin{itemize}
    \item Upper bound:
    \begin{align*}
E_{x, \epsilon} \left[ \left(\hat{f}_{\lambda}(x) - f_{\star}(x) \right)^{2} \right]
\leq &~ O_{d, \mathbb{P}} \left( d^{v_1} + d^{v_2}\right) + \|B_1\|_{2}^{2} + \|B_2\|_{2}^{2} + \|B_3 - B_5\|_{2}^{2} + \|B_4\|_{2}^{2}\\
= &~
O_{d, \mathbb{P}} \left( d^{v_1} + d^{v_2}\right) + O_{d, \mathbb{P}}\left( d^{b_2(u, \tau)}  + d^{b_3} + d^{b_1} + d^{b_4} \right)\\
(\because \text{ Proposition \ref{prop_e_2}}) \quad = &~ O_{d, \mathbb{P}} \left( d^{v_1} + d^{v_2} + d^{b_1} + d^{b_2(u, \tau)}\right);
\end{align*}

\item Lower bound:
    \begin{itemize}
        \item If $b_1>b_4$ and $b_1 \geq b_2(u, \tau)$, then
\begin{align*}
&~ E_{x, \epsilon} \left[ \left(\hat{f}_{\lambda}(x) - f_{\star}(x) \right)^{2} \right]
\geq  \Omega_{d, \mathbb{P}} \left( d^{v_1} + d^{v_2} \right) + \|B_3-B_4-B_5\|_{2}^{2}\\
\overset{b_1 > b_4}{=} &~ \Omega_{d, \mathbb{P}} \left( d^{v_1} + d^{v_2} + d^{b_1}\right)
=
\Omega_{d, \mathbb{P}} \left( d^{v_1} + d^{v_2} + d^{b_1} + d^{b_2(u, \tau)}\right);
\end{align*}

\item If $b_1 \leq b_4$ and $b_1 \geq b_2(u, \tau)$, since Proposition \ref{prop_e_2} implies that $b_4<v_2$, we have $\max\{b_1, b_2(u, \tau)\}<v_2$, hence
\begin{align*}
E_{x, \epsilon} \left[ \left(\hat{f}_{\lambda}(x) - f_{\star}(x) \right)^{2} \right]
\geq
\Omega_{d, \mathbb{P}} \left( d^{v_1} + d^{v_2}\right)
=
\Omega_{d, \mathbb{P}} \left( d^{v_1} + d^{v_2} + d^{b_1} + d^{b_2(u, \tau)}\right),
\end{align*}


\item If $b_2(u, \tau)>b_3$ and $b_1 < b_2(u, \tau)$, then
\begin{align*}
&~ E_{x, \epsilon} \left[ \left(\hat{f}_{\lambda}(x) - f_{\star}(x) \right)^{2} \right]
\geq   \Omega_{d, \mathbb{P}} \left( d^{v_1} + d^{v_2} \right) + \|B_1+B_2\|_{2}^{2}\\
\overset{b_2(u, \tau)>\max\{b_1, b_3\}}{=} &~
\Omega_{d, \mathbb{P}} \left( d^{v_1} + d^{v_2} + d^{b_2}(u, \tau)\right)
=
\Omega_{d, \mathbb{P}} \left( d^{v_1} + d^{v_2} + d^{b_1} + d^{b_2(u, \tau)}\right),
\end{align*}

\item If $b_2(u, \tau) \leq b_3$ and $b_1 < b_2(u, \tau)$, since Proposition \ref{prop_e_2} implies that $b_3 \leq v_2$, we have $\max\{b_1, b_2(u, \tau)\} \leq v_2$, hence
\begin{align*}
E_{x, \epsilon} \left[ \left(\hat{f}_{\lambda}(x) - f_{\star}(x) \right)^{2} \right]
\geq  \Omega_{d, \mathbb{P}} \left( d^{v_1} + d^{v_2} \right) 
=
\Omega_{d, \mathbb{P}} \left( d^{v_1} + d^{v_2} + d^{b_1} + d^{b_2(u, \tau)}\right).
\end{align*}
    \end{itemize}

\end{itemize}

\noindent {\bf Case (ii). } Suppose $u \geq \gamma$. From Proposition \ref{prop_e_1}  with $t=1$ and $t=\tau$, we have $\gamma > u^{\prime}(1)$ and $\gamma > u^{\prime}(\tau)$, hence from Proposition \ref{prop_e_2} and Proposition \ref{prop_e_3}, we have
\begin{equation*}
%\label{eqn_tau_finite_b234_bound}
    \max\{v_1(u), v_2(u), b_1(u)\} \geq \max\{b_2(u, 1), b_2(u, \tau), b_3(u)\},
\end{equation*}
Therefore, from (\ref{eqn_learning_curve_pre_results_1}) we have
\begin{itemize}
    \item Upper bound:
    \begin{align*}
E_{x, \epsilon} \left[ \left(\hat{f}_{\lambda}(x) - f_{\star}(x) \right)^{2} \right]
=
O_{d, \mathbb{P}} \left( d^{v_1} + d^{v_2} + d^{b_1}\right)
=
O_{d, \mathbb{P}} \left( d^{v_1} + d^{v_2} + d^{b_1} + d^{b_2(u, \tau)}\right);
\end{align*}

\item Lower bound:
\begin{align*}
&~ E_{x, \epsilon} \left[ \left(\hat{f}_{\lambda}(x) - f_{\star}(x) \right)^{2} \right]
\geq  \Omega_{d, \mathbb{P}} \left( d^{v_1} + d^{v_2} \right) + \|B_3-B_4-B_5\|_{2}^{2}\\
=&~
\Omega_{d, \mathbb{P}} \left( d^{v_1} + d^{v_2} + d^{b_1}\right)
=
\Omega_{d, \mathbb{P}} \left( d^{v_1} + d^{v_2} + d^{b_1} + d^{b_2(u, \tau)}\right).
\end{align*}

\end{itemize}
Putting the two bounds together yields the desired result. 
\end{proof}


% \input{appendix/learn_3}







\subsection{Quantities calculation of learning curve results}



% We then introduce the following results. Readers can find the definitions of certain notations, e.g., $b_i = b_i(u)$ and $v_i = v_i(u)$ for $i=1,2$, in Appendix \ref{append_learning_curve}.

\begin{proposition}\label{prop_e_1}
    Let $u^{\prime} = u^{\prime}(t)$ be the optimal rate of the regularization parameter $\lambda$ with qualification $t \in [1, \infty]$ and coefficient of source condition $s>0$ defined in \cite{lu2024saturation} (ignoring the logarithm term). Then, we have $u^{\prime} < \gamma$ and $\tilde{p} \leq u^{\prime} < \tilde{p}+1$, where $\tilde{s} := \min\{s, 2t\}$ and $\tilde{p}:=\lfloor \gamma/(\tilde{s}+1)\rfloor$.\\
    As a direct corollary, for any $u \leq u^{\prime}$, we have $\ell_{\lambda} = \lfloor u \rfloor \leq \ell_{\gamma}$, hence
$$
v_1^{\prime} = v_1, ~ v_2^{\prime} = v_2, ~
b_1^{\prime} = b_1, ~ b_2^{\prime}(u, t) = b_2(u, t).
$$
\end{proposition}
\begin{proof}
According to Corollaries D.15--17 in \cite{lu2024saturation}, we divide the proof into four cases.

\noindent {\bf Case I: $1 \leq s \leq t$}\\
\noindent Denote $p:=\lfloor \gamma/(s+1)\rfloor$.\\
   (1.1) Suppose $p \geq 1$ and $p(s+1) \leq \gamma < ps+p+s$. We have
   $$
   u^{\prime} = p+\frac{1}{2} < p+ps \leq \gamma \quad \text{ and } \quad p \leq u^{\prime} < p+1.
   $$
   (1.2) Suppose $p \geq 1$ and $ps+p+s \leq \gamma < ps+p+s+1$. We have
   $$
   u^{\prime} = \frac{\gamma - (p+1)(s-1)}{2} < p+1 < \gamma \quad \text{ and } \quad p \leq u^{\prime} < p+1.
   $$
   (1.3) Suppose $\gamma < s$. We have
   $$
   u^{\prime} = \min\{\gamma, 1\}/2 < \gamma \quad \text{ and } \quad 0 \leq u^{\prime} < 1.
   $$
   (1.4) Suppose $s \leq \gamma < s+1$. We have
   $$
   u^{\prime} = \frac{\gamma - (s-1)}{2} < \gamma \quad \text{ and } \quad 0 \leq u^{\prime} < 1.
   $$

   
\noindent {\bf Case II: $\tilde{s} > t$}\\
\noindent Denote $\tilde{p}:=\lfloor \gamma/(\tilde{s}+1)\rfloor$.\\
   (2.1) Suppose $\gamma \geq 1$ and $0 \leq \gamma - \tilde{p}(\tilde{s} +1) < t$. We have
   $$
   u^{\prime} = \tilde{p} + \frac{\gamma - \tilde{p}(\tilde{s} +1)}{2t} < \tilde{p}+\frac{1}{2} < \frac{\gamma + 1}{2} \leq \gamma \quad \text{ and } \quad \tilde{p} \leq u^{\prime} < \tilde{p}+1.
   $$  
   (2.2) Suppose $\gamma \geq 1$ and $t \leq \gamma - \tilde{p}(\tilde{s} +1) < \tilde{s} + \tilde{s} /t - 1$. We have
   $$
   u^{\prime} = \tilde{p} + \frac{\gamma - \tilde{p}(\tilde{s} +1)+1}{2(t+1)} \leq 
   \left\{\begin{matrix}
\tilde{p}+\frac{\tilde{s} }{2t} < t +\tilde{p} + \tilde{p}\tilde{s}  \leq \gamma, & \tilde{p} \geq 1  \\
\frac{\gamma}{t + 1} < \gamma, & \tilde{p}=0,
\end{matrix}\right.
   $$  
   and $\tilde{p} \leq u^{\prime} < \tilde{p}+1$.\\
   (2.3) Suppose $\gamma \geq 1$ and $\tilde{s} + \tilde{s} /t - 1 \leq \gamma - \tilde{p}(\tilde{s} +1) < \tilde{s} +1$. We have
   $$
   u^{\prime} = \tilde{p} + \frac{\gamma - \tilde{p}(\tilde{s} +1)+1-\tilde{s} }{2} < \tilde{p}+1<\gamma \quad \text{ and } \quad \tilde{p} \leq u^{\prime} < \tilde{p}+1.
   $$ 
   (2.4) Suppose $\gamma < 1$. We have
   $$
   u^{\prime} = \frac{\gamma}{2} < \gamma \quad \text{ and } \quad 0 \leq u^{\prime} < 1.
   $$ 


\noindent {\bf Case III: $t < \infty$ and $s<1$}\\
\noindent Denote $p:=\lfloor \gamma/(s+1)\rfloor$.\\
(3.1) Suppose $p(s+1) \leq \gamma < ps+p+s$. We have
   $$
   u^{\prime} = \frac{\gamma +2t p - sp-p}{2t} 
   \leq 
   \left\{\begin{matrix}
p+\frac{s}{2t} < p(s+1) \leq \gamma, & p \geq 1  \\
\frac{\gamma}{2t} < \gamma, & p=0,
\end{matrix}\right.
   $$ 
   and $p \leq u^{\prime} < p+1$.\\
   (3.2) Suppose $ps+p+s \leq \gamma < ps+p+s+1$. We have
   $$
   u^{\prime} = p + \frac{s}{2t} < ps+p+s \leq \gamma \quad \text{ and } \quad p \leq u^{\prime} < p+1.
   $$ 



\noindent {\bf Case IV: $t = \infty$ and $s<1$}\\
\noindent Denote $p:=\lfloor \gamma/(s+1)\rfloor$.\\
   (4.1) Suppose $p \geq 1$ and $p(s+1) \leq \gamma < ps+p+s$. We have
   $$
   u^{\prime} = p + \frac{s}{2} < p(s+1) \leq \gamma \quad \text{ and } \quad p \leq u^{\prime} < p+1.
   $$  
   (4.2) Suppose $p \geq 1$ and $ps+p+s \leq \gamma < ps+p+s+1$. We have
   $$
   u^{\prime} = \frac{\gamma + p(1-s)}{2} < \gamma \quad \text{ and } \quad p \leq u^{\prime} < p+1.
   $$  
   (4.3) Suppose $\gamma < s$. We have
   $$
   u^{\prime} = \gamma \cdot \min\left\{s, \frac{1}{2}\right\} < \gamma \quad \text{ and } \quad 0 \leq u^{\prime} < 1.
   $$  
   (4.4) Suppose $s \leq \gamma < s+1$. We have
   $$
   u^{\prime} = \min\{\gamma (1+s)-s, \gamma / 2 \} < \gamma \quad \text{ and } \quad 0 \leq u^{\prime} < 1.
   $$ 
\end{proof}















\begin{proposition}\label{prop_e_2}
For any $u \in (0, \infty]$, we have $b_3(u) \leq v_2(u)$ and $b_4(u) < v_2(u)$. %Moreover, if $s<2$, then we further have $b_5(u) < b_1(u)$.
\end{proposition}
\begin{proof}
    We have
    \begin{align*}
        b_3 \leq &~ 2(\tilde{\ell} - \gamma) \overset{\tilde{\ell} \leq \ell_{\gamma} \leq \gamma}{\leq} \tilde{\ell} - \gamma = v_2,\\
        b_4 \leq &~ \tilde{\ell} - \gamma - 1 < v_2.
    \end{align*}
\iffalse
moreover, when $s<2$, we have
    $$
    b_5 < -(\tilde{\ell} + 1)  s = b_1.
    $$
    \fi
\end{proof}





\begin{proposition}\label{prop_e_3}
    Let $u^{\prime} = u^{\prime}(t)$ be the optimal rate of the regularization parameter $\lambda$ with qualification $t < \infty$ and coefficient of source condition $s>0$ defined in \cite{lu2024saturation} (ignoring the logarithm term). Then, for any $u \geq u^{\prime}$, we have
    $$
    \max\{v_1(u), v_2(u), b_1(u)\} \geq b_2(u, t).
    $$
\end{proposition}
\begin{proof}
For notational simplicity, we denote  $\tilde{s}:=\min\{s, 2t\}$  and $b_2(u)=b_2(u, t)$ for any $u \in [0, \infty]$.
 We divide the proof into three steps:
 \begin{itemize}
     \item[(i)] We first show that for any $u^{\prime} \leq u < \lceil u^{\prime} \rceil$, we have
     $$
    \max\{v_1(u), v_2(u), b_1(u)\} \geq b_2(u);
    $$
     
     \item[(ii)] We then show that we have
        $$
        \max\{v_1(\lceil u^{\prime} \rceil), v_2(\lceil u^{\prime} \rceil), b_1(\lceil u^{\prime} \rceil)\} \geq b_2(\lceil u^{\prime} \rceil);
        $$
    \item[(iii)] Finally, we show that, for any integer $q \geq 0$ and any $u\geq q$, if $\max\{v_1(q), v_2(q), b_1(q)\} \geq b_2(q)$, then we have
    $$
        \max\{v_1(u), v_2(u), b_1(u)\} \geq b_2(u);
    $$
 \end{itemize}
Combining (i), (ii), and (iii) with $q=\lceil u^{\prime} \rceil$, we obtain the desired result.


\noindent {\bf Proof of (i). } 
For any integer $q \geq 0$ and any $q \leq q_1 \leq q_2 < q+1$, we have
\begin{equation}\label{eqn_prop_e3_1}
    \begin{aligned}
        v_1(q_2) &\geq v_1(q_1)\\
        v_1(q+1) &= v_1(q) + \mathbf{1}\{q < \ell_{\gamma}\} + 2(\gamma -\ell_{\gamma})\mathbf{1}\{q = \ell_{\gamma}\}\\
        v_2(q_2) &= v_2(q_1)\\
        v_2(q+1) &= v_2(q) + \mathbf{1}\{q < \ell_{\gamma}\}\\
        b_1(q_2) &= b_1(q_1)\\
        b_1(q+1) &= b_1(q) - s\mathbf{1}\{q < \ell_{\gamma}\}\\
        b_2(q_2) &\leq b_2(q_1)\\
        b_2(q+1) &= b_2(q) - \tilde{s}\mathbf{1}\{q < \ell_{\gamma}\} - 2t \mathbf{1}\{q \geq \ell_{\gamma}\}.
    \end{aligned}
\end{equation}
Therefore, for any $u^{\prime} \leq u < \lceil u^{\prime} \rceil$, we have
     \begin{align*}
         \max\{v_1(u), v_2(u), b_1(u)\} 
         \geq&~
         \max\{v_1(u^{\prime}), v_2(u^{\prime}), b_1(u^{\prime})\}\\
         \overset{\text{Corollary D.15-17 in \cite{lu2024saturation}} }{\geq}&~
         b_2(u^{\prime}) \geq b_2(u).
     \end{align*}






 

\noindent {\bf Proof of (ii). } 
If $u^{\prime}$ is an integer, then from (i) we have $\max\{v_1(\lceil u^{\prime} \rceil), v_2(\lceil u^{\prime} \rceil), b_1(\lceil u^{\prime} \rceil)\} \geq b_2(\lceil u^{\prime} \rceil)$. 

Suppose $u^{\prime}$ is not an integer. 
From Proposition \ref{prop_e_1}, we have $\tilde{p} \leq u^{\prime} < \tilde{p}+1$, where $\tilde{p}:=\lfloor \gamma/(\tilde{s}+1)\rfloor$.
Denote $\tilde{\ell}_{u^{\prime}} := \min\{\ell_{\gamma}, \tilde{p}+1\}$, then we have $\lceil u^{\prime} \rceil = \tilde{p}+1$ and
\begin{equation*}
    \begin{aligned}
    %v_1(\lceil u^{\prime} \rceil) =&~ \gamma - \tilde{\ell}_{u^{\prime}} - 1 - 2  \max\{\gamma - \tilde{p}-1, 0\}\\
    v_2(\lceil u^{\prime} \rceil) =&~  \tilde{\ell}_{u^{\prime}} - \gamma\\
    %b_1(\lceil u^{\prime} \rceil) =&~ -(\tilde{\ell}_{u^{\prime}} + 1)  s\\
    b_2(\lceil u^{\prime} \rceil) =&~ -2  t  (\tilde{p}+1) + (2  t - \tilde{s})  \tilde{\ell}_{u^{\prime}}.
    \end{aligned}
\end{equation*}

\begin{itemize}
    \item If $\gamma < \tilde{p}+1$, then by the definition of $\tilde{p} \leq \ell_{\gamma} = \lfloor \gamma \rfloor$, we have $\tilde{p}=\ell_{\gamma}$ and $\tilde{\ell}_{u^{\prime}} = \ell_{\gamma}$. Hence
$$
b_2(\lceil u^{\prime} \rceil) = -2  t - \tilde{s}\ell_{\gamma} \overset{t \geq 1}{<} -1 < \ell_{\gamma} - \gamma = v_2(\lceil u^{\prime} \rceil);
$$

    \item If $\gamma \geq \tilde{p}+1$, then $\tilde{\ell}_{u^{\prime}} = \tilde{p}+1$.
    Since $\gamma<(\tilde{p}+1)(\tilde{s}+1)$, we have
$$
b_2(\lceil u^{\prime} \rceil) = -\tilde{s}\tilde{p}-\tilde{s} < \tilde{p}+1-\gamma = v_2(\lceil u^{\prime} \rceil).
$$
\end{itemize}
Combining the above two cases, we have
$$
        \max\{v_1(\lceil u^{\prime} \rceil), v_2(\lceil u^{\prime} \rceil), b_1(\lceil u^{\prime} \rceil)\} \geq b_2(\lceil u^{\prime} \rceil). 
        $$





\noindent {\bf Proof of (iii). } 
Suppose that for some integer $q \geq 0$, we have $\max\{v_1(q), v_2(q), b_1(q)\} \geq b_2(q)$. 
By (i), for any $q^{\prime \prime}\in[q, q+1)$ we have
\begin{align}\label{eqn_step_1}
  \max\{v_1(q^{\prime \prime}), v_2(q^{\prime \prime}), b_1(q^{\prime \prime})\} &\geq  b_2(q^{\prime \prime}).
  %  \max\{v_1(q^{\prime \prime}), v_2(q^{\prime \prime}), b_1(q^{\prime \prime})\} &\geq \max\{v_1(q), v_2(q), b_1(q)\} \geq b_2(q) \geq b_2(q^{\prime \prime}), \quad q \leq q^{\prime \prime} < q+1.
\end{align}
Thus, the inequality holds throughout the interval $[q, q+1)$.



Next, consider the integer point $q+1$. We analyze three possible cases:
\begin{itemize}
    \item If $s \leq 2 t$, then $\tilde{s}=s$ and (\ref{eqn_prop_e3_1}) implies 
    \begin{align*}
        &~ \max\{v_1(q+1), v_2(q+1), b_1(q+1)\} \geq \max\{v_1(q), v_2(q), b_1(q)\} - s\mathbf{1}\{q < \ell_{\gamma}\}\\
        \geq &~ b_2(q) - s\mathbf{1}\{q < \ell_{\gamma}\} \geq b_2(q+1).
    \end{align*}
    
    \item If $s>2t$ and $q < \ell_{\gamma}$, then from (\ref{eqn_order_in_bias_and_variance}) we have $b_1(q)=-s(q+1) < -2t q = b_2(q)$, so the inequality $\max\{v_1(q), v_2(q), b_1(q)\} \geq b_2(q)$ must come from the first two components:
    $$
        \max\{v_1(q), v_2(q)\} \geq b_2(q).
    $$
    Therefore, (\ref{eqn_prop_e3_1}) implies
    \begin{align*}
        &~ \max\{v_1(q+1), v_2(q+1), b_1(q+1)\} \geq \max\{v_1(q+1), v_2(q+1)\} \geq \max\{v_1(q), v_2(q)\}\\
        \geq &~ b_2(q) \geq b_2(q+1).
    \end{align*}

    \item If $s>2t$ and $q \geq \ell_{\gamma}$, then (\ref{eqn_prop_e3_1}) implies 
    \begin{align*}
        &~ \max\{v_1(q+1), v_2(q+1), b_1(q+1)\} \geq \max\{v_1(q), v_2(q), b_1(q)\} \\
        \geq &~ b_2(q) \geq b_2(q+1).
    \end{align*}
\end{itemize}
Combining (\ref{eqn_step_1}) and all three cases above, we conclude that if 
$$
\max\{v_1(u), v_2(u), b_1(u)\} \geq b_2(u),
$$
holds at integer $u=q$, it also holds for all $u \in [q, q+1]$.

Finally, by induction over integers $q+1, q+2, \cdots$, we have 
$$
\max\{v_1(u), v_2(u), b_1(u)\} \geq b_2(u),\quad u \geq q,
$$
which completes the proof of (iii).
\end{proof}



















\section{Kernel approximation results for kernels on sphere}\label{append_kernel_app_sphere}


In this section, we will restate and develop several approximation results for kernels defined on the sphere (see, e.g., Section 2.1).


Throughout this section, we assume that all conditions and assumptions in Theorem 3.1 hold.





The following lemma is borrowed from Proposition 3 in \cite{ghorbani2021linearized} as well as Lemma S5 in \cite{zhang2024phase}.



\begin{lemma}\label{lemma psi psi top}
   We have
   \iffalse
    \begin{equation*}
        \sup\limits_{k > \ell_{\gamma} + 1} \left\| \frac{\Psi_{k} \Psi_{k}^{\top} }{N(d, k)} - \mathrm{I}_{n} \right\|_{\mathrm{op}} = o_{d, \mathbb{P}}(1),
    \end{equation*}
    which further implies
    \fi
    \begin{equation}\label{eq k ge l in lemma}
    \begin{aligned}
        \Psi_{>\ell_{\gamma}} \Sigma_{>\ell_{\gamma}} \Psi_{>\ell_{\gamma}}^{\top} &= \kappa_1 \left(\mathrm{I}_{n} + \Delta_{1}\right), \quad \|\Delta_{1}\|_{\mathrm{op}} = o_{d, \mathbb{P}}(1)\\
        \Psi_{>\ell_{\gamma}} \Sigma_{>\ell_{\gamma}}^{2} \Psi_{>\ell_{\gamma}}^{\top} &= \kappa_2 \left(\mathrm{I}_{n} + \Delta_{2}\right), \quad \|\Delta_{2}\|_{\mathrm{op}} = o_{d, \mathbb{P}}(1);
    \end{aligned}
    \end{equation}
    with
    \begin{align*}
    \kappa_{1} = \Theta_{d}(1), \quad \kappa_{2} = \Theta_{d}\left( d^{-(\ell_{\gamma}+1)} \right).
\end{align*}
\end{lemma}

The following two lemmas are borrowed from \cite{ghorbani2021linearized} and \cite{xiao2022precise}. Readers can also find the first lemma in Lemma S6 of \cite{zhang2024phase}.


 



\begin{lemma}[Lemma 11 in \cite{ghorbani2021linearized}; Lemma S6 of \cite{zhang2024phase}]
\label{lemma psi top psi}
When $\gamma > \ell_{\gamma}$, we have
    \begin{displaymath}
         \left\| \frac{\Psi_{\leq \ell_{\gamma}}^{\top} \Psi_{\leq \ell_{\gamma}} }{n} - \mathrm{I}_{N_{\ell_{\gamma}}} \right\|_{\mathrm{op}} = o_{d, \mathbb{P}}(1).
    \end{displaymath}
\end{lemma}


\begin{lemma}[Theorem 1 in  \cite{xiao2022precise}]\label{lemma psi top psi_int}
Suppose $\gamma = \ell_{\gamma}$. Denote $N_{\ell_{\gamma}-1} = \sum_{k^{\prime}=0}^{\ell_{\gamma}-1} N(d,k^{\prime})$, then
\begin{equation}\label{eqn_lemma psi top psi_int_1}
    \left\| \frac{\Psi_{\leq \ell_{\gamma}-1}^{\top} \Psi_{\leq \ell_{\gamma}-1} }{n} - \mathrm{I}_{N_{\ell_{\gamma}-1}} \right\|_{\mathrm{op}} = o_{d, \mathbb{P}}(1).
\end{equation}
    And we have
    
    $$
    \lambda_{\text{max}}\left(\frac{1}{n} \Psi_{\ell_{\gamma}}^{\top} \Psi_{\ell_{\gamma}}\right) = \Theta_{d, \mathbb{P}}(1)
    \quad \text{ and } \quad 
    \lambda_{\text{min}\{n, N(d, \ell_{\gamma})\}}\left(\frac{1}{n} \Psi_{\ell_{\gamma}}^{\top} \Psi_{\ell_{\gamma}}\right) = \Theta_{d, \mathbb{P}}(1),
    $$
    hence 
    $$
    \frac{1}{n} \mathrm{tr}\left(\frac{1}{n} \Psi_{\ell_{\gamma}}^{\top} \Psi_{\ell_{\gamma}}\right) = \Theta_{d, \mathbb{P}}(1).
    $$
    
\end{lemma}


\begin{remark}
    For completeness, we briefly explain how to obtain Lemma \ref{lemma psi top psi} and Lemma \ref{lemma psi top psi_int} from Lemma 11 in \cite{ghorbani2021linearized} and Theorem 1 in \cite{xiao2022precise}. 
    A comparison of notations is given in Table~\ref{tab_comparison_of_notations}. 
    Recall  
    % \dhnote{cite Prop. 2.1 later}
    \begin{itemize}
        \item When $\gamma > \ell_{\gamma}$, Proposition 2.1 implies that $n/(N_{\ell_{\gamma}} \log(N_{\ell_{\gamma}}))   = \Theta_d(d^{\gamma - \ell_{\gamma}}\log^{-1}(d)) \to \infty$ as $d \to \infty$. 
        Denote $\Delta = \frac{\Psi_{\leq \ell_{\gamma}}^{\top} \Psi_{\leq \ell_{\gamma}} }{n} - \mathrm{I}_{N_{\ell_{\gamma}}}$, then Lemma 11 in \cite{ghorbani2021linearized} implies $\mathbb{E}[\|\Delta\|_{\mathrm{op}}] = o_{d}(1)$, and Markov's inequality implies $\|\Delta\|_{\mathrm{op}} = o_{d, \mathbb{P}}(1)$.

        \vspace{3pt}

        \item When $\gamma = \ell_{\gamma}$, following the proof of Lemma 11 in \cite{ghorbani2021linearized} with $N_{\ell_{\gamma}}$ being replaced by $N_{\ell_{\gamma}-1}$, we can similarly show (\ref{eqn_lemma psi top psi_int_1}). The remaining part of Lemma \ref{lemma psi top psi_int} is a direct result of Theorem 1 in \cite{xiao2022precise}.
        
    \end{itemize}

  %   \dhnote{
  % Why $n/(N_{\ell_{\gamma}} \log(N_{\ell_{\gamma}})) \to \infty$ as $d \to \infty$?

  % In the second case, why $n/(N_{\ell_{\gamma}}\to $ a constant?
  %   }
    
\end{remark}



\begin{table}[t]
\centering
\caption{Comparison of notations between our paper and \cite{ghorbani2021linearized, xiao2022precise}}
\label{tab_comparison_of_notations}
\begin{tabular}{l l l l}
\hline
\textbf{Concept / Term} & \textbf{Our Paper} & \cite{ghorbani2021linearized} & \cite{xiao2022precise} \\
\hline
Spherical harmonics & $\psi_{k, j}$ & $Y_{k, j}$ & $Y_{k, j}$ \\
\vspace{3pt}
Sample size & $n$ & $n$ & $m$ \\
\vspace{3pt}
Multiplicity of eigenvalue $\mu_k$ & $N(d, k)$ & $B(d, k)$ & $N(d, k)$ \\
\vspace{3pt}
$\lfloor \gamma \rfloor$ & $\ell_{\gamma}$ & $\ell$ & $r$ \\
\vspace{3pt}
$\sum_{k=0}^{\ell_{\gamma}} N(d, k)$ & $N_{\ell_{\gamma}}$ & B &  \\
\hline
\end{tabular}
\end{table}








% Combining Lemma \ref{lemma psi psi top}, Lemma \ref{lemma psi top psi}, and Lemma \ref{lemma psi top psi_int}







The following proposition is borrowed from Lemma 3 in \cite{xiao2022precise}.

\begin{proposition}\label{prop_lemma_3_in_xiao}
    Suppose $\gamma = \ell_{\gamma}$. 
    Denote $\overline{D}:= (n\lambda + 1) \cdot n^{-1}\Sigma_{\leq \ell_{\gamma}}^{-1} + n^{-1}\Psi_{\leq \ell_{\gamma}}^{\top} \Psi_{\leq \ell_{\gamma}}$,
    then we have 
    $$
    \lambda_{\text{max}} \left(  \overline{D} \right)  = \Theta_{d, \mathbb{P}}(1)
    \quad \text{ and } \quad
    \lambda_{\text{min}} \left(  \overline{D} \right) = \Theta_{d, \mathbb{P}}(1).
    $$
\end{proposition}




Recall that $\tilde{\ell} =\min\{\ell_{\gamma}, \ell_{\lambda}\}$. Hence, from Assumption 4 we have
\begin{equation}\label{eqn:regular_para_upper}
    % 1 = o_d\left(\lambda d^{\ell_{\lambda}+1}\right) \quad \text{ and } \quad 1 = \Omega_d\left(\lambda d^{\ell_{\lambda}}\right).
   \lambda d^{\ell_{\lambda}} = O_{d}(1),\text{ and }  \lambda d^{\ell_{\lambda}+1}\rightarrow \infty. 
\end{equation}
In the proof of the next lemma, we will frequently use (\ref{eqn:regular_para_upper}).


\begin{lemma}\label{lemma lambda max min}
We have the following bounds:
\begin{itemize}
    \item[(i)] $\lambda_{\text{min}}\left( M\right) = \Omega_{d, \mathbb{P}}(n\lambda+1)$, 
$\mathrm{tr}(M / n)=O_{d, \mathbb{P}}(n\lambda+1)$.

    \item[(ii)] $M_{>\ell_{\gamma}} = K_{>\ell_{\gamma}} + n\lambda \mathrm{I}_n = \Theta_{d, \mathbb{P}}(n\lambda + 1) \left(\mathrm{I}_{n} + \Delta_{1}\right)$.

    \item[(iii)] 
    When $\ell_{\lambda}+1 \leq \ell_{\gamma}$, for any $i=N_{\ell_{\lambda}}+1, \cdots, \lfloor\min\{n, N_{\ell_{\lambda}+1}\}\rfloor$,
    % \dhnote{The range of $i$ seems to be a typo; there should be $N_{l_\lambda}$. }
    we have
        $$
        \lambda_i(M) = O_{d, \mathbb{P}}(n\lambda + 1).
        $$
    

    \item[(iv)] When $\ell_{\lambda}<\ell_{\gamma}$, we have %$\lambda_{\text{max}}(K_{>\ell_{\lambda}}) = O_{d, \mathbb{P}}(nd^{-\ell_{\lambda}-1})$ and
    \begin{align*}
        \lambda_{\text{max}}(M_{>\ell_{\lambda}}) &= \Theta_{d, \mathbb{P}}(n\lambda),
        \quad
        \lambda_{\text{min}}(M_{>\ell_{\lambda}}) = \Theta_{d, \mathbb{P}}(n\lambda),
    \end{align*}
    and $\lambda_{\text{max}}(K_{>\ell_{\lambda}}) =
 O_{d, \mathbb{P}}\left({d^{\gamma-\ell_{\lambda}-1}}\right) =  
    o_{d, \mathbb{P}}(n\lambda)$.

    \item[(v)] Denote $A = n^{-1}\Sigma_{\leq \tilde{\ell}}^{-1} + n^{-1}\Psi_{\leq \tilde{\ell}}^{\top} M_{>\tilde{\ell}}^{-1} \Psi_{\leq \tilde{\ell}}$, then we have
    \begin{displaymath}
        \lambda_{\text{max}} \left(  A^{-1} \right)
        = \Theta_{d, \mathbb{P}}(n\lambda + 1)
        \quad \text{ and } \quad
        \lambda_{\text{min}} \left(  A^{-1} \right) = \Theta_{d, \mathbb{P}}(n\lambda + 1).
    \end{displaymath}

\end{itemize}
\end{lemma}
\begin{proof}
\noindent {\bf {(i)}} From Lemma \ref{lemma psi psi top}, we have 
\begin{equation*}
        \begin{aligned}
        \lambda_{\text{min}}\left( M\right) \overset{\text{Proposition } \ref{prop_weyl_ine}}{\geq}&~ \lambda_{\text{min}} \left( \Psi_{\leq \ell_{\gamma}} \Sigma_{\leq \ell_{\gamma}} \Psi_{\leq \ell_{\gamma}}^{\top} \right) + \lambda_{\text{min}} \left( \kappa_{1} \mathrm{I}_{n} \right) + \lambda_{\text{min}} \left( \kappa_{1} \Delta_{1} \right) + n\lambda\\
        \geq &~ \kappa_{1}(1-o_{d, \mathbb{P}}(1)) + n\lambda
        = \Omega_{d, \mathbb{P}}(n\lambda+1).
    \end{aligned}
\end{equation*}
Further, from Lemma \ref{lemma psi psi top}, we have
\begin{equation*}
        \begin{aligned}
\mathrm{tr}\left(\frac{M}{n}\right) 
=&~ \mathrm{tr}\left( \frac{K}{n} +\lambda \mathrm{I}_{n} \right)\\
=&~ \mathrm{tr}\left( \frac{1}{n} \Psi_{\leq \ell_{\gamma}} \Sigma_{\leq \ell_{\gamma}} \Psi_{\leq \ell_{\gamma}}^{\top} \right) + \mathrm{tr}\left( \frac{\kappa_{1}}{n} \mathrm{I}_{n}  \right) + \mathrm{tr}\left( \frac{\kappa_{1}}{n} \mathbf{\Delta}_{1}  \right) + n\lambda   \\
=&~ \mathrm{tr}\left( \frac{1}{n} \Psi_{\leq \ell_{\gamma}}^{\top} \Psi_{\leq \ell_{\gamma}} \Sigma_{\leq \ell_{\gamma}} \right) + \kappa_{1} \left( 1 + o_{d, \mathbb{P}}(1)\right) + n\lambda   \\
\overset{\text{Lemma } \ref{lemma psi top psi} \text{ and Lemma } \ref{lemma psi top psi_int}}{=}&~ O_{d, \mathbb{P}}\left(\mathrm{tr}\left(  \Sigma_{\leq \ell_{\gamma}}\right)\right) 
+ \kappa_{1} \left( 1 + o_{d, \mathbb{P}}(1)\right) + n\lambda   \\
=&~ O_{d, \mathbb{P}}(n\lambda+1),  
    \end{aligned}
\end{equation*}

\noindent {\bf {(ii)}} It is a direct corollary of Lemma \ref{lemma psi psi top}.



\noindent {\bf {(iii)}}  
Denote the SVD decomposition of $ \Psi_{\leq \ell_{\gamma}} $ as $ \Psi_{\leq \ell_{\gamma}} = n^{1/2} ~ O H V^{\top} $, where $ O \in \mathbb{R}^{n \times n}$ and $ V \in \mathbb{R}^{N_{\ell_{\gamma}} \times N_{\ell_{\gamma}}}$ are orthogonal matrices, $H = \left[ H_{\star}; 0 \right]^{\top} \in \mathbb{R}^{n \times N_{\ell_{\gamma}}} $, and $H_{\star} \in \mathbb{R}^{N_{\ell_{\gamma}} \times N_{\ell_{\gamma}}} $ is a diagonal matrix. Then, Lemma \ref{lemma psi top psi} and Lemma \ref{lemma psi top psi_int} imply that both $\lambda_{\text{max}}\left(H_{\star}\right)$ and $  \lambda_{\text{min}}\left(H_{\star}\right)$ are of order $\Theta_{d, \mathbb{P}}(1)$.\\
Recall the assumption that  $\ell_{\lambda}+1 \leq \ell_{\gamma}$. Hence, for any $i=N_{\ell_{\lambda}}+1, \cdots, \lfloor\min\{n, N_{\ell_{\lambda}+1}\}\rfloor$, from Lemma \ref{lemma psi psi top}, we have 
\begin{equation*}
        \begin{aligned}
        \lambda_{i}\left( M\right) 
        = &~ \lambda_{i}\left( \Psi_{\leq \ell_{\gamma}}\Sigma_{\leq \ell_{\gamma}} \Psi_{\leq \ell_{\gamma}}^{\top} + M_{>\ell_{\gamma}}\right)\\
        \overset{\text{Proposition~ \ref{prop_weyl_ine}}}{\leq}&~
        \lambda_{i} \left(  \Psi_{\leq \ell_{\gamma}}\Sigma_{\leq \ell_{\gamma}} \Psi_{\leq \ell_{\gamma}}^{\top}  \right) + \lambda_{1}\left(  M_{>\ell_{\gamma}} \right)\\
        \overset{\text{Part (ii)}}{\leq} &~
        \lambda_{i} \left(  \Psi_{\leq \ell_{\gamma}}\Sigma_{\leq \ell_{\gamma}} \Psi_{\leq \ell_{\gamma}}^{\top}  \right) + O_{d, \mathbb{P}}\left(n\lambda+1\right)\\
        =&~ \lambda_{i} \left(  H_{\star} V^{\top} \Sigma_{\leq \ell_{\gamma}} V H_{\star}^{\top}  \right) + O_{d, \mathbb{P}}\left(n\lambda+1\right)\\
        \leq &~ 
        \lambda_{\text{max}}\left(H_{\star}\right)^2 \lambda_{i} \left(\Sigma_{\leq \ell_{\gamma}} \right) + O_{d, \mathbb{P}}\left(n\lambda+1\right)\\
        =&~ O_{d, \mathbb{P}}(n\mu_{\ell_{\lambda}+1} + n\lambda+1) \overset{(\ref{eqn:regular_para_upper})}{=} O_{d, \mathbb{P}}(n\lambda+1).
    \end{aligned}
\end{equation*}


% \dhnote{The first inequality: You have used the Courant-Fischer minimax theorem. I think this should be mentioned. Perhaps we can combine the minimax theorem or its corollary into Proposition~\ref{prop_weyl_ine}. 
% }


% \dhnote{
% The 2nd line: 
% Just to confirm that you used $i > \ell_\lambda$?
% }

\noindent {\bf {(iv)}} When $\ell_{\lambda}<\ell_{\gamma}$, we have $
1 = o_d(n\lambda)$ and $nd^{-\ell_{\lambda}-1}=\Omega_d(1)$.
As in Case (i), we can show that $\lambda_{\text{min}}(M_{>\ell_{\lambda}}) = \Omega_{d, \mathbb{P}}(n\lambda+1) = \Omega_{d, \mathbb{P}}(n\lambda)$. On the other hand, we have
\begin{equation}\label{eqn_minmax_iv_1}
    \begin{aligned}
        \lambda_{\text{max}}(M_{>\ell_{\lambda}})
        \overset{\text{Proposition } \ref{prop_weyl_ine}}{\leq} &~
        \lambda_{\text{max}}(K_{\ell_{\lambda}+1, \ell_{\gamma}}) +
        \lambda_{\text{max}}(K_{>\ell_{\gamma}}) + n\lambda\\
        \overset{\text{Lemma } \ref{lemma psi psi top}}{=}&~
        \lambda_{\text{max}}(K_{\ell_{\lambda}+1, \ell_{\gamma}}) +
        O_{d, \mathbb{P}}(n\lambda)\\
    \end{aligned}
\end{equation}
 Denote the SVD decomposition of $ \Psi_{\ell_{\lambda}+1, \ell_{\gamma}} $ as $ \Psi_{\ell_{\lambda}+1, \ell_{\gamma}} = n^{1/2} ~ O H V^{\top} $, where $ O \in \mathbb{R}^{n \times n}$ and $ V \in \mathbb{R}^{(N_{\ell_{\gamma}}-N_{\ell_{\lambda}}) \times (N_{\ell_{\gamma}}-N_{\ell_{\lambda}})}$ are orthogonal matrices; $H = \left[ H_{\star}; 0 \right]^{\top} \in \mathbb{R}^{n \times (N_{\ell_{\gamma}}-N_{\ell_{\lambda}})} $ and $H_{\star} \in \mathbb{R}^{(N_{\ell_{\gamma}}-N_{\ell_{\lambda}}) \times (N_{\ell_{\gamma}}-N_{\ell_{\lambda}})} $ is a diagonal matrix. 
 Then, Lemma \ref{lemma psi top psi} and Lemma \ref{lemma psi top psi_int} imply that both $\lambda_{\text{max}}\left(H_{\star}\right)$ and $  \lambda_{\text{min}}\left(H_{\star}\right)$ are of order $\Theta_{d, \mathbb{P}}(1)$. Therefore, 
\begin{equation}\label{eqn_minmax_iv_2}
    \begin{aligned}
        \lambda_{\text{max}}(K_{\ell_{\lambda}+1, \ell_{\gamma}})
        =&~
        \lambda_{\text{max}}\left(\Psi_{\ell_{\lambda}+1, \ell_{\gamma}} \Sigma_{\ell_{\lambda}+1, \ell_{\gamma}} \Psi_{\ell_{\lambda}+1, \ell_{\gamma}}^{\top}\right)\\
         = &~
         O_{d, \mathbb{P}}\left(\frac{n}{d^{\ell_{\lambda}+1}}\right) \overset{(\ref{eqn:regular_para_upper})}{=} o_{d, \mathbb{P}}(n\lambda).
    \end{aligned}
\end{equation}
Combining all these, we have that both $\lambda_{\text{max}}(M_{>\ell_{\lambda}})$ and $ \lambda_{\text{min}}(M_{>\ell_{\lambda}})$ are of order $\Theta_{d, \mathbb{P}}(n\lambda)$.
Finally, as in (\ref{eqn_minmax_iv_1}) and (\ref{eqn_minmax_iv_2}), we can show that $\lambda_{\text{max}}(K_{>\ell_{\lambda}}) = O_{d, \mathbb{P}}(nd^{-\ell_{\lambda}-1}+ 1)=O_{d, \mathbb{P}}(nd^{-\ell_{\lambda}-1})=o_{d, \mathbb{P}}(n\lambda)$, finishing the proof of (iv).




\noindent {\bf {(v)}} We have $A = n^{-1}\Sigma_{\leq \tilde{\ell}}^{-1} + n^{-1}\Psi_{\leq \tilde{\ell}}^{\top} M_{>\tilde{\ell}}^{-1} \Psi_{\leq \tilde{\ell}}$. On the one hand, from Lemma \ref{lemma psi top psi}, Lemma \ref{lemma psi top psi_int}, (ii), and (iv), we have
\begin{equation*}
    \lambda_{\text{max}} \left(  A \right) = O_{d, \mathbb{P}}\left(\frac{d^{\tilde{\ell}}}{n} + \frac{1}{n(\lambda + n^{-1})}\right) \overset{(\ref{eqn:regular_para_upper})}{=} O_{d, \mathbb{P}}\left(\frac{1}{n\lambda + 1}\right);
\end{equation*}
On the other hand, to obtain the lower bound of $\lambda_{\text{min}} \left(  A \right)$, we consider the following three cases: 
(a) $\gamma >\ell_{\gamma}$;
(b) $\gamma = \ell_{\gamma}$ and $\ell_{\lambda} < \ell_{\gamma}$;
(b) $\gamma = \ell_{\gamma}$ and $\ell_{\lambda} \geq \ell_{\gamma}$.



\noindent {\bf Case (a). }
When $\gamma >\ell_{\gamma}$, from Lemma \ref{lemma psi top psi}, (ii), and (iv), we have
$$
\lambda_{\text{min}} \left(  A \right) \overset{\text{Proposition } \ref{prop_weyl_ine}}{\geq} \lambda_{\text{min}} \left(n^{-1}\Psi_{\leq \tilde{\ell}}^{\top} M_{>\tilde{\ell}}^{-1} \Psi_{\leq \tilde{\ell}}\right)
= \Omega_{d, \mathbb{P}}\left(\frac{1}{n\lambda + 1}\right).
$$

\noindent {\bf Case (b). }
When $\gamma = \ell_{\gamma}$ and $\tilde{\ell}= \ell_{\lambda} < \ell_{\gamma}$, 
from Lemma \ref{lemma psi top psi_int} and (iv), we have
$$
\lambda_{\text{min}} \left(  A \right) \geq \lambda_{\text{min}} \left(n^{-1}\Psi_{\leq \ell_{\lambda}}^{\top} M_{>\ell_{\lambda}}^{-1} \Psi_{\leq \ell_{\lambda}}\right)
= \Omega_{d, \mathbb{P}}\left(\frac{1}{n\lambda + 1}\right).
$$



\noindent {\bf Case (c). } When $\gamma = \ell_{\gamma}$ and $\ell_{\lambda} \geq \ell_{\gamma}$, we have
\begin{equation*}
    \begin{aligned}
\lambda_{\text{min}} \left(  A \right) \overset{\text{(ii)}}{\geq} 
&~\lambda_{\text{min}} \left(n^{-1}\Sigma_{\leq \ell_{\gamma}}^{-1} + (n\lambda + 1)^{-1}n^{-1}\Psi_{\leq \ell_{\gamma}}^{\top} \Psi_{\leq \ell_{\gamma}}\right) (1- o_{d, \mathbb{P}}(1))\\
\overset{\text{Proposition } \ref{prop_lemma_3_in_xiao}}{=} &~
\Omega_{d, \mathbb{P}}\left(\frac{1}{n\lambda + 1}\right),
    \end{aligned}
\end{equation*}
and we get the desired results.
\iffalse
that 
$$
\lambda_{\text{max}} \left(  A^{-1} \right)  \asymp \lambda_{\text{min}} \left(  A^{-1} \right) = \Theta_{d, \mathbb{P}}(n\lambda + 1).
$$
\fi
\end{proof}





With some modifications, the following lemma restates Lemma 8, 9, and 10 in \cite{zhang2024phase}.

\begin{lemma}\label{lemma ge l l2}
We have
\begin{itemize}
    \item[(i)] $\left\| \Psi_{>\tilde{\ell}} \theta_{>\tilde{\ell}} \right\|_{2}^{2} = O_{d, \mathbb{P}}\left( n  \| \theta_{>\tilde{\ell}} \|_{2}^{2}\right)$.

    \item[(ii)] $\| \theta_{>\tilde{\ell}} \|_{2}^{2}  = \Theta_d \left(  d^{-(\tilde{\ell}+1)s} \right)$.

    \item[(iii)] $\left\| \Sigma_{\leq \tilde{\ell}}^{-\tau^{\prime}}  \theta_{\leq \tilde{\ell}} \right\|_{2}^{2}  = \Theta_{d} \left(  d^{(2\tau^{\prime}-\min\{s, 2\tau^{\prime}\})\tilde{\ell} } \right)$, $\tau^{\prime} \geq 1$.
\end{itemize}
\end{lemma}
\begin{proof}
(i) and (ii) are direct results of Lemma S8 and S10 in \cite{zhang2024phase}. Now let's prove (iii). We have
\begin{equation}\label{eq proof le l-1}
        \left\| \Sigma_{\leq \tilde{\ell}}^{-\tau^{\prime}}  \theta_{\leq \tilde{\ell}} \right\|_{2}^{2} = 
        \sum\limits_{m = 0}^{\tilde{\ell}} \mu_{m}^{-2\tau^{\prime}} \sum\limits_{j = 1 }^{N(d, m)} (\theta_{m,j})^{2} 
        = \sum\limits_{m = 0}^{\tilde{\ell}} \mu_{m}^{s-2\tau^{\prime}} \mu_{m}^{-s}\sum\limits_{j = 1 }^{N(d, m)} (\theta_{m,j})^{2}.
    \end{equation}
On the one hand, when $0 \leq s \leq 2\tau^{\prime}$, (\ref{eq proof le l-1}) implies
\begin{equation*}
\begin{aligned}
    \left\| \Sigma_{\leq \tilde{\ell}}^{-\tau^{\prime}}  \theta_{\leq \tilde{\ell}} \right\|_{2}^{2} & \overset{\text{Assumption } 3}{=} 
    O_{d} \left(  \mu_{\tilde{\ell}}^{s-2\tau^{\prime}} \right) =
        O_{d} \left(  d^{(2\tau^{\prime}-s)\tilde{\ell} } \right)\\
        \left\| \Sigma_{\leq \tilde{\ell}}^{-\tau^{\prime}}  \theta_{\leq \tilde{\ell}} \right\|_{2}^{2} &\geq \mu_{\tilde{\ell}}^{s-2\tau^{\prime}} \mu_{\tilde{\ell}}^{s}\sum\limits_{j = 1 }^{N(d, \tilde{\ell})} (\theta_{\tilde{\ell},j})^{2} \overset{\text{Assumption } 6}{=}
        \Omega_{d} \left( d^{(2\tau^{\prime}-s)\tilde{\ell} } \right).
\end{aligned}
    \end{equation*}
On the other hand, when $s > 2\tau^{\prime}$,
since $\mu_m = O_d(1)$, $m=0, 1, \cdots, \tilde{\ell}$, (\ref{eq proof le l-1}) implies
\begin{equation*}
\begin{aligned}
    \left\| \Sigma_{\leq \tilde{\ell}}^{-\tau^{\prime}}  \theta_{\leq \tilde{\ell}} \right\|_{2}^{2} &\overset{\text{Assumption } 3}{=} 
    O_{d} \left( 1 \right)\\
        \left\| \Sigma_{\leq \tilde{\ell}}^{-\tau^{\prime}}  \theta_{\leq \tilde{\ell}} \right\|_{2}^{2} &= 
        \Omega_d(1) \cdot \sum\limits_{m = 0}^{\tilde{\ell}}  \sum\limits_{j = 1 }^{N(d, m)} (\theta_{m,j})^{2}
        \overset{\text{Assumption } 6}{=}
        \Omega_d(1),
\end{aligned}
    \end{equation*}
and we obtain the desired results.
\end{proof}




The following two identities are based on the Sherman-Morrison-Woodbury formula (see, e.g., Proposition \ref{prop_smw_formula}).


% {\green
% [
% In the proof for kernel defined on general domains, we need the second line of equality below, so I add it. 
% ]
% }


\begin{lemma}\label{lemma trans in V2}
For $k \in \{\ell_{\gamma}, \ell_{\lambda}\}$, we have
% \begin{equation*}
%     \begin{aligned}
%         M^{-1} \Psi_{\leq k} \Sigma_{\leq k}^{1/2} 
%         = &~ 
%         M_{>k}^{-1} \Psi_{\leq k} \Sigma_{\leq k}^{1/2} \left( \mathrm{I}_{N_{k}} +  \Sigma_{\leq k}^{1/2}  \Psi_{\leq k}^{\top} M_{>k}^{-1} \Psi_{\leq k} \Sigma_{\leq k}^{1/2}\right)^{-1}, \\ 
%         \Sigma_{\leq k} - \Sigma_{\leq k} \Psi^{\top}_{\leq k} M^{-1} \Psi_{\leq k}\Sigma_{\leq k}
%         = &~
%         \left( \Sigma_{\leq k}^{-1} + \Psi_{\leq k}^{\top} M_{>k}^{-1} \Psi_{\leq k}  \right)^{-1}.
%     \end{aligned}
% \end{equation*}

\begin{equation*}
    \begin{aligned}
        M^{-1} \Psi_{\leq k} \Sigma_{\leq k}^{1/2} 
        = &~ 
        M_{>k}^{-1} \Psi_{\leq k} \Sigma_{\leq k}^{1/2} \left( \mathrm{I}_{N_{k}} +  \Sigma_{\leq k}^{1/2}  \Psi_{\leq k}^{\top} M_{>k}^{-1} \Psi_{\leq k} \Sigma_{\leq k}^{1/2}\right)^{-1}, 
    \end{aligned}
\end{equation*}
and 
\begin{equation*}
    \begin{aligned}
    \Sigma_{\leq k} - \Sigma_{\leq k} \Psi^{\top}_{\leq k} M^{-1} \Psi_{\leq k}\Sigma_{\leq k}
        = &~
        \left( \Sigma_{\leq k}^{-1} + \Psi_{\leq k}^{\top} M_{>k}^{-1} \Psi_{\leq k}  \right)^{-1}.
    \end{aligned}
\end{equation*}
        
\end{lemma}
\begin{proof}
Recall that we have $M = \Psi_{\leq k} \Sigma_{\leq k} \Psi_{\leq k}^-{\top} + M_{>k}$. 
\begin{itemize}
    \item Applying the Sherman-Morrison-Woodbury formula with $Z= \Psi_{\leq k} \Sigma_{\leq k}^{1/2}$ and $A=M_{>k}$, we have
    \begin{equation*}
    \begin{aligned}
        M^{-1} \Psi_{\leq k} \Sigma_{\leq k}^{1/2} 
        = &~  (ZZ^{\top} + A)^{-1} Z
        = A^{-1}Z - A^{-1} Z(\mathrm{I}_{N_{k}} + Z^{\top} A^{-1} Z)^{-1} Z^{\top} A^{-1}Z\\
        =&~
        A^{-1}Z \left(\mathrm{I}_{N_{k}} - (\mathrm{I}_{N_{k}} + Z^{\top} A^{-1} Z)^{-1} \left(\mathrm{I}_{N_{k}} + Z^{\top} A^{-1} Z - \mathrm{I}_{N_{k}}\right)\right)
        \\
        =&~
        A^{-1} Z \left( \mathrm{I}_{N_{k}} +  Z^{\top} A^{-1} Z\right)^{-1}
        \\
        =&~
        M_{>k}^{-1} \Psi_{\leq k} \Sigma_{\leq k}^{1/2} \left( \mathrm{I}_{N_{k}} +  \Sigma_{\leq k}^{1/2}  \Psi_{\leq k}^{\top} M_{>k}^{-1} \Psi_{\leq k} \Sigma_{\leq k}^{1/2}\right)^{-1}.
    \end{aligned}
\end{equation*}

\item Applying the Sherman-Morrison-Woodbury formula with $Z= \Psi_{\leq k} M_{> k}^{-1/2}$ and $A=\Sigma_{\leq k}^{-1}$, we have
\begin{equation*}
    \begin{aligned}
        \left( \Sigma_{\leq k}^{-1} + \Psi_{\leq k}^{\top} M_{>k}^{-1} \Psi_{\leq k}  \right)^{-1}
        = &~
        (ZZ^{\top} + A)^{-1}\\
        = &~
        A^{-1} - A^{-1} Z(\mathrm{I}_{B_{k}} + Z^{\top} A^{-1} Z)^{-1} Z^{\top} A^{-1}\\
        = &~
        \Sigma_{\leq k} - \Sigma_{\leq k} \Psi^{\top}_{\leq k} M^{-1} \Psi_{\leq k}\Sigma_{\leq k}.
    \end{aligned}
\end{equation*}
\end{itemize}
\end{proof}





\section{Proof of results 
%with kernels on general domains: Spectral gap and hyper-contrativity
in Section~4
}\label{append_kernel_app_general}

In this section, we will generalize the results for spherical inner-product kernels defined on the sphere $\mathbb{S}^{d}$ to kernels on general domains and prove Theorem 4.2.

Throughout this section, we assume that all conditions and assumptions in Theorem 4.2 hold.

We use the same notations as in Appendix \ref{append_notation}, except that each symbol is augmented with a tilde, e.g., $\Psi \to \tilde{\Psi}$ and $\Sigma \to \tilde{\Sigma}$.



\subsection{Kernel approximation results for general kernels }

The following lemma is a corollary of Theorem 6 in \cite{mei2022generalization}, which generalizes the results in Lemma \ref{lemma psi psi top} and Lemma \ref{lemma psi top psi}.

\begin{lemma}\label{lemma psi psi top_general}
   We have
    \begin{equation}\label{eq k ge l in lemma_general_1}
    \begin{aligned}
        \tilde{\Psi}_{>\ell_{\gamma}} \tilde{\Sigma}_{>\ell_{\gamma}} \tilde{\Psi}_{>\ell_{\gamma}}^{\top} &= \kappa_1 \left(\mathrm{I}_{n} + \Delta_{1}\right), \quad \|\Delta_{1}\|_{\mathrm{op}} = o_{d, \mathbb{P}}(1),
    \end{aligned}
    \end{equation}
    and
    \begin{equation}\label{eq k ge l in lemma_general_2}
    \begin{aligned}
        \tilde{\Psi}_{>\ell_{\gamma}} \tilde{\Sigma}_{>\ell_{\gamma}}^{2} \tilde{\Psi}_{>\ell_{\gamma}}^{\top} &= \kappa_2 \left(\mathrm{I}_{n} + \Delta_{2}\right), \quad \|\Delta_{2}\|_{\mathrm{op}} = o_{d, \mathbb{P}}(1);
    \end{aligned}
    \end{equation}
    with
    \begin{align*}
    \kappa_{1} = \Theta_{d}(1), \quad \kappa_{2} = \Theta_{d}\left( d^{-(\ell_{\gamma}+1)} \right);
\end{align*}
Moreover, denote $\tilde{N}_{\ell_{\gamma}} = \sum_{k=0}^{\ell_{\gamma}}\tilde{N}(d, k)$, then we have
    \begin{equation}\label{eq k ge l in lemma_general_3}
         \left\| \frac{\tilde{\Psi}_{\leq \ell}^{\top} \tilde{\Psi}_{\leq \ell} }{n} - \mathrm{I}_{\tilde{N}_{\ell_{\gamma}}} \right\|_{\mathrm{op}} = o_{d, \mathbb{P}}(1).
    \end{equation}
\end{lemma}
\begin{proof}
While Theorem 6 in \cite{mei2022generalization} is stated under the setting of random feature kernels, its proofs (mainly relying on Propositions 3 and 4 in \cite{mei2022generalization}) do not rely on any property specific to this class. Consequently, the results remain valid for general kernels, and we adopt them in this broader setting.

\noindent {\bf Proof of (\ref{eq k ge l in lemma_general_1}) and (\ref{eq k ge l in lemma_general_3}). }
We first state the following conditions (i.e., Assumption 6 (a), Assumption 6 (b), and Assumption 7 (a) in \cite{mei2022generalization} with the following substitution of notations: $u(d) \to \sum_{k=0}^{2\ell_{\gamma} +1}\tilde{N}(d, k)$, $N(d) \to n$, $M(d) \to \sum_{k=0}^{\ell_{\gamma}}\tilde{N}(d, k)$, and $\{\lambda_{d, j}^2\} \to \{\tilde{\lambda}_{k, j}\}$):
\begin{itemize}
    \item For any integer $q \geq 1$, there exists $C=C(q, \ell)$, such that for any $k \leq 2\ell + 1$ and any $j \leq \tilde{N}(d, k)$, we have
\begin{equation}\label{eqn_condi_general_1}
 \|\tilde{\psi}_{k, j}\|_{L^{2q}} \leq C(q, \ell) \|\tilde{\psi}_{k, j}\|_{L^{2}}.
    \end{equation}

    \item There exists $\delta_1>0$, such that
 \begin{equation}\label{eqn_condi_general_2}
    n^{2+\delta_1} \leq \frac{\left[\sum_{k \geq 2\ell_{\gamma}+2} \sum_{j \leq \tilde{N}(d, k)} \tilde{\lambda}_{k, j} \right]^2}{\sum_{k \geq 2\ell_{\gamma}+2} \sum_{j \leq \tilde{N}(d, k)} \tilde{\lambda}_{k, j}^2 }.
    \end{equation}

    \item There exists $\delta_2>0$, such that
\begin{equation}\label{eqn_condi_general_3}
    n^{1+\delta_2} \leq \frac{\sum_{k \geq \ell_{\gamma}+1} \sum_{j \leq \tilde{N}(d, k)} \tilde{\lambda}_{k, j} }{\tilde{\lambda}_{\ell_{\gamma}+1, 1}}.
    \end{equation}

\end{itemize}

Now let's verify these conditions. Since $\|\tilde{\psi}_{k, j}\|_{L^{2}}=1$, Assumption 9 implies (\ref{eqn_condi_general_1}). 
Assumption~8 implies that $\sum_{j \leq \tilde{N}(d, k)} \tilde{\lambda}_{k, j} =\Theta(1)$ for any $k\leq 2\ell_{\gamma}+2$. 
Together with Assumption 7, it further implies that there exist positive constants $\delta_1$ and $\delta_2$ such that
\begin{align*}
    \frac{\left[\sum_{k \geq 2\ell_{\gamma}+2} \sum_{j \leq \tilde{N}(d, k)} \tilde{\lambda}_{k, j} \right]^2}{\sum_{k \geq 2\ell_{\gamma}+2} \sum_{j \leq \tilde{N}(d, k)} \tilde{\lambda}_{k, j}^2 } \geq \frac{\left[\sum_{j \leq \tilde{N}(d, 2\ell_{\gamma}+2)} \tilde{\lambda}_{2\ell_{\gamma}+2, j} \right]^2}{\tilde{\lambda}_{2\ell_{\gamma}+2, 1}\sum_{k = 0}^{\infty} \sum_{j \leq \tilde{N}(d, k)} \tilde{\lambda}_{k, j}} = \Omega_d(d^{2(\ell_{\gamma}+1)}) \gg n^{2+\delta_1},
\end{align*}
% \dhnote{the sum $\sum_{k = 0}^{\infty} \sum_{j \leq \tilde{N}(d, k)} \tilde{\lambda}_{k, j}$ needs to be bounded as $O_d(1)$; being finite is not enough. }
% {\green [I have modified the assumption $\sum_{k = 0}^{\infty} \sum_{j \leq \tilde{N}(d, k)} \tilde{\lambda}_{k, j} \leq \tilde{\mathfrak{C}}$ in Section 4.]}
and
\begin{align*}
    \frac{\sum_{k \geq \ell_{\gamma}+1} \sum_{j \leq \tilde{N}(d, k)} \tilde{\lambda}_{k, j} }{\tilde{\lambda}_{\ell_{\gamma}+1, 1}} \geq \frac{\sum_{j \leq \tilde{N}(d, \ell_{\gamma}+1)} \tilde{\lambda}_{\ell_{\gamma}+1, j} }{\tilde{\lambda}_{\ell_{\gamma}+1, 1}}
    = \Omega_d(d^{\ell_{\gamma}+1}) \gg n^{1+\delta_2}.
\end{align*}
Therefore, Theorem 6 in \cite{mei2022generalization} ensures that both (\ref{eq k ge l in lemma_general_1}) and (\ref{eq k ge l in lemma_general_3}) hold.

\noindent {\bf Proof of (\ref{eq k ge l in lemma_general_2}). } 
Assumption 7 and 8 imply that there exist $\delta_3, \delta_4>0$, such that
\begin{align*}
    \frac{\left[\sum_{k \geq 2\ell_{\gamma}+2} \sum_{j \leq \tilde{N}(d, k)} \tilde{\lambda}_{k, j}^2 \right]^2}{\sum_{k \geq 2\ell_{\gamma}+2} \sum_{j \leq \tilde{N}(d, k)} \tilde{\lambda}_{k, j}^4 } 
    & \geq 
    \frac{\tilde{\lambda}_{2\ell_{\gamma}+2, \tilde{N}(d, 2\ell_{\gamma}+2)}^2\left[\sum_{j \leq \tilde{N}(d, 2\ell_{\gamma}+2)} \tilde{\lambda}_{2\ell_{\gamma}+2, j} \right]^2}{\tilde{\lambda}_{2\ell_{\gamma}+2, 1}^3\sum_{k = 0}^{\infty} \sum_{j \leq \tilde{N}(d, k)} \tilde{\lambda}_{k, j}}
   \\
   & = 
    \Omega_d(d^{2(\ell_{\gamma}+1)})\\
   & \gg n^{2+\delta_3},
\end{align*}
and
\begin{align*}
    \frac{\sum_{k \geq \ell_{\gamma}+1} \sum_{j \leq \tilde{N}(d, k)} \tilde{\lambda}_{k, j}^2 }{\tilde{\lambda}_{\ell_{\gamma}+1, 1}^2} \geq \frac{\tilde{\lambda}_{\ell_{\gamma}+1, \tilde{N}(d, \ell_{\gamma}+1)} \sum_{j \leq \tilde{N}(d, \ell_{\gamma}+1)} \tilde{\lambda}_{\ell_{\gamma}+1, j} }{\tilde{\lambda}_{\ell_{\gamma}+1, 1}^2}
    = \Omega_d(d^{\ell_{\gamma}+1}) \gg n^{1+\delta_4}.
\end{align*}
Similar as above, from Theorem 6 in \cite{mei2022generalization} with the substitution 
$$
\tilde{\mathscr{K}}(x,x^\prime) := \sum_{k=0}^{\infty} \sum_{j=1}^{\tilde{N}(d, k)} \tilde{\lambda}_{k, j}  \tilde{\psi}_{k, j}(x) \tilde{\psi}_{k, j}\left(x^\prime\right) \to \sum_{k=0}^{\infty} \sum_{j=1}^{\tilde{N}(d, k)} \tilde{\lambda}_{k, j}^2  \tilde{\psi}_{k, j}(x) \tilde{\psi}_{k, j}\left(x^\prime\right),
$$
we have that (\ref{eq k ge l in lemma_general_2}) holds.
\end{proof}

The next lemma generalizes the results in Lemma \ref{lemma ge l l2}. We can obtain its proof by modified the corresponding proof for Lemma S8 and S10 in \cite{zhang2024phase}, and Lemma \ref{lemma ge l l2}, and hence we omit the proof.

\begin{lemma}\label{lemma ge l l2_general}
We have
\begin{itemize}
    \item[(i)] $\left\| \tilde{\Psi}_{>\tilde{\ell}} \theta_{>\tilde{\ell}} \right\|_{2}^{2} = O_{d, \mathbb{P}}\left( n  \| \theta_{>\tilde{\ell}} \|_{2}^{2}\right)$.

    \item[(ii)] $\| \theta_{>\tilde{\ell}} \|_{2}^{2}  = \Theta_d \left(  d^{-(\tilde{\ell}+1)s} \right)$.

    \item[(iii)] $\left\| \tilde{\Sigma}_{\leq \tilde{\ell}}^{-1}  \theta_{\leq \tilde{\ell}} \right\|_{2}^{2}  = \Theta_{d} \left(  d^{(2-\min\{s, 2\})\tilde{\ell} } \right)$.
\end{itemize}
\end{lemma}






\subsection{Proof of Theorem 4.2}
\begin{proof}
    Adopt the notations in (\ref{eqn_order_in_bias_and_variance}). We divide the proof into four steps.

{\bf Step (I). } We first show that results in Appendix \ref{append_kernel_app_sphere} can be generalized to the cases where kernels are defined on general domains. 

Recall that $\gamma > \ell_{\gamma}$. We have:
\begin{itemize}
    \item Lemma \ref{lemma psi psi top} and Lemma \ref{lemma psi top psi} can be replaced by Lemma \ref{lemma psi psi top_general}.

    \item Since the proof of Lemma \ref{lemma lambda max min} relies on Lemma \ref{lemma psi psi top} and Lemma \ref{lemma psi top psi}, it can be established in a similar manner by replacing them with Lemma \ref{lemma psi psi top_general}.

    \item Lemma \ref{lemma ge l l2} can be replaced by Lemma \ref{lemma ge l l2_general}.

    \item Lemma \ref{lemma trans in V2} still holds with each symbol augmented with a tilde.

    
\end{itemize}



{\bf Step (II). } We next show that the results for the variance term in Appendix \ref{append_variance} can also be extended to the case where kernels are defined on general domains.

Since the proof of Theorem \ref{thm_variance_spe} depends on the results established in Appendix \ref{append_kernel_app_sphere}, Step (I) directly implies that
\begin{equation*}
    \begin{aligned}
      \mathrm{var}\left( \hat{f}_{\lambda} \right) =
        \Theta_{d, \mathbb{P}}\left(
        d^{v_1} + d^{v_2}
\right).
    \end{aligned}
\end{equation*}


{\bf Step (III). } We then bound the bias term by modifying the proof in Appendix \ref{append_bias}.

From (\ref{eq derivation of bias_spe}) we have
\begin{equation*}
\begin{aligned}
    \mathrm{bias}^{2}(\hat{f}_{\lambda})  =  &~
    \|B_1+B_2\|_{2}^{2} + \|B_3-B_4-B_5\|_{2}^{2}.
\end{aligned} 
\end{equation*}

\begin{itemize}
    \item In the proof of Theorem \ref{thm_bias_spe}, since results for $B_2, B_3, B_4$ and $B_5$ depend on the results established in Appendix \ref{append_kernel_app_sphere}, Step (I) directly implies that (see also (\ref{eqn_learning_curve_pre_results_1})):
    \begin{equation*}
    \begin{aligned}
                \|B_2\|_{2}^{2} =&~ 
            O_{d, \mathbb{P}}\left(  d^{b_1}\right)\\
            \|B_3-B_5\|_{2}^{2} 
        =&~
        \Theta_{d, \mathbb{P}} \left( d^{b_1}\right)\\
        \|B_4\|_{2}^{2} =&~ o_{d, \mathbb{P}} \left(d^{b_2(u, 1)}+d^{b_3}\right).
    \end{aligned}
\end{equation*}

    \item We then bound the term $\|B_1\|_2^2$. Notice that for kernel ridge regression, we have 
    \begin{equation}\label{eqn_inter_general_1}
        \varphi_{\lambda, \tilde{K}}:=n^{-1}\reg^{\krr}(\tilde{K}/n) = (\tilde{K} + n\lambda \mathrm{I}_n)^{-1} = \tilde{M}^{-1}.
    \end{equation}
    Moreover, from Step (I), the generalization of Lemma \ref{lemma trans in V2} and Lemma \ref{lemma lambda max min} (v) imply that
    \begin{align}
            \tilde{\Sigma}_{\leq \tilde{\ell}} - \tilde{\Sigma}_{\leq \tilde{\ell}} \tilde{\Psi}^{\top}_{\leq \tilde{\ell}} \tilde{M}^{-1} \tilde{\Psi}_{\leq \tilde{\ell}} \tilde{\Sigma}_{\leq \tilde{\ell}}
        =
            \left( \tilde{\Sigma}_{\leq \tilde{\ell}}^{-1} + \tilde{\Psi}_{\leq \tilde{\ell}}^{\top} \tilde{M}_{>\tilde{\ell}}^{-1} \tilde{\Psi}_{\leq \tilde{\ell}}  \right)^{-1}, \label{eqn_inter_general_2} \\[6pt]
         \lambda_{\text{max}} \left(  \left( \tilde{\Sigma}_{\leq \tilde{\ell}}^{-1} + \tilde{\Psi}_{\leq \tilde{\ell}}^{\top} \tilde{M}_{>\tilde{\ell}}^{-1} \tilde{\Psi}_{\leq \tilde{\ell}}  \right)^{-1} \right)
        = \Theta_{d, \mathbb{P}}\left(\lambda + \frac{1}{n}\right),
        \label{eqn_inter_general_3} \\[6pt]
        \lambda_{\text{min}} \left(  \left( \tilde{\Sigma}_{\leq \tilde{\ell}}^{-1} + \tilde{\Psi}_{\leq \tilde{\ell}}^{\top} \tilde{M}_{>\tilde{\ell}}^{-1} \tilde{\Psi}_{\leq \tilde{\ell}}  \right)^{-1} \right) = \Theta_{d, \mathbb{P}}\left(\lambda + \frac{1}{n}\right). \label{eqn_inter_general_4} 
    \end{align}
    Therefore, for $B_1$, we have
    \begin{equation*}
        \begin{aligned}
            \|B_1\|_{2}^{2} =&~ \left\| \left( \mathrm{I}_{\tilde{N}_{\tilde{\ell}}} - \tilde{\Sigma}_{\leq \tilde{\ell}} \tilde{\Psi}^{\top}_{\leq \tilde{\ell}} \varphi_{\lambda, \tilde{K}} \tilde{\Psi}_{\leq \tilde{\ell}}\right)    \theta_{\leq \tilde{\ell}} \right\|_{2}^{2}\\
           (\because (\ref{eqn_inter_general_1})) \quad {=}&~ \left\| \left( \tilde{\Sigma}_{\leq \tilde{\ell}} - \tilde{\Sigma}_{\leq \tilde{\ell}} \tilde{\Psi}^{\top}_{\leq \tilde{\ell}} \tilde{M}^{-1} \tilde{\Psi}_{\leq \tilde{\ell}} \tilde{\Sigma}_{\leq \tilde{\ell}} \right) \tilde{\Sigma}_{\leq \tilde{\ell}}^{-1}   \theta_{\leq \tilde{\ell}} \right\|_{2}^{2}\\
           (\because (\ref{eqn_inter_general_2})) \quad {=}&~ \left\| \left( \tilde{\Sigma}_{\leq \tilde{\ell}}^{-1} + \tilde{\Psi}_{\leq \tilde{\ell}}^{\top} \tilde{M}_{>\tilde{\ell}}^{-1} \tilde{\Psi}_{\leq \tilde{\ell}}  \right)^{-1}
           \tilde{\Sigma}_{\leq \tilde{\ell}}^{-1}   \theta_{\leq \tilde{\ell}} \right\|_{2}^{2}\\
           (\because (\ref{eqn_inter_general_3}) \text{ and } (\ref{eqn_inter_general_4})) \quad {=}&~
           \Theta_{d, \mathbb{P}}(d^{-2u} + d^{-2\gamma}) \left\| 
           \tilde{\Sigma}_{\leq \tilde{\ell}}^{-1}   \theta_{\leq \tilde{\ell}} \right\|_{2}^{2}\\
           (\because \text{ Lemma \ref{lemma ge l l2_general}}) \quad {=}&~ \Theta_{d, \mathbb{P}}\left((d^{-2u} + d^{-2\gamma}) d^{(2-\min\{s, 2\})\tilde{\ell} }\right)\\
           {=}&~
           \Theta_{d, \mathbb{P}}\left(d^{b_2(u, 1)}+d^{b_3}\right).
        \end{aligned}
    \end{equation*}
\end{itemize}

{\bf Step (IV). } Finally we combine the estimates from Step (II) (variance) and Step (III) (bias decomposition) to obtain the learning curve.

From Lemma \ref{prop_e_2} we have $b_3 \leq v_2$. Therefore: 
\begin{itemize}
    \item Upper bound:
            \begin{align*}
&~ E_{x, \epsilon} \left[ \left(\hat{f}_{\lambda}(x) - f_{\star}(x) \right)^{2} \right]
\leq  \mathrm{var}\left( \hat{f}_{\lambda} \right) + \|B_1\|_{2}^{2} + \|B_2\|_{2}^{2} + \|B_3 - B_5\|_{2}^{2} + \|B_4\|_{2}^{2}\\
= &~
O_{d, \mathbb{P}} \left( d^{v_1} + d^{v_2}\right) + O_{d, \mathbb{P}}\left( d^{b_1} + d^{b_2(u, 1)}  + d^{b_3} \right) =  O_{d, \mathbb{P}} \left( d^{v_1} + d^{v_2} + d^{b_1} + d^{b_2(u, 1)}\right).
\end{align*}

    \item For lower bound, we analyze three possible cases:
    \begin{itemize}
    
    \item Suppose $b_3 < b_2(u, 1)$ and $b_1 < b_2(u, 1)$. Then
\begin{align*}
  &~  E_{x, \epsilon} \left[ \left(\hat{f}_{\lambda}(x) - f_{\star}(x) \right)^{2} \right]
\geq 
\Omega_{d, \mathbb{P}} \left( d^{v_1} + d^{v_2}\right) + \|B_1\|_{2}^{2} - \|B_2\|_{2}^{2} \\
=&~  \Omega_{d, \mathbb{P}}\left( d^{v_1} + d^{v_2} + d^{b_2(u, 1)}\right) = \Omega_{d, \mathbb{P}}\left( d^{v_1} + d^{v_2} + d^{b_1} + d^{b_2(u, 1)}\right).
\end{align*}

    \item Suppose $b_3 < b_2(u, 1)$ and $b_1 \geq b_2(u, 1)$. Then
    \begin{align*}
  &~  E_{x, \epsilon} \left[ \left(\hat{f}_{\lambda}(x) - f_{\star}(x) \right)^{2} \right]
\geq 
\Omega_{d, \mathbb{P}} \left( d^{v_1} + d^{v_2}\right) + \|B_3 - B_5\|_{2}^{2} - \|B_4\|_{2}^{2} \\
=&~  \Omega_{d, \mathbb{P}}\left( d^{v_1} + d^{v_2} + d^{b_1}\right) = \Omega_{d, \mathbb{P}}\left( d^{v_1} + d^{v_2} + d^{b_1} + d^{b_2(u, 1)}\right).
\end{align*}

    \item Suppose $b_3 \geq b_2(u, 1)$. Then we further have $\max\{b_2(u, 1), b_3\} \leq v_2$, and hence
    \begin{align*}
  &~  E_{x, \epsilon} \left[ \left(\hat{f}_{\lambda}(x) - f_{\star}(x) \right)^{2} \right]
\geq 
\Omega_{d, \mathbb{P}} \left( d^{v_1} + d^{v_2}\right) + \|B_3 - B_5\|_{2}^{2} - \|B_4\|_{2}^{2} \\
=&~  \Omega_{d, \mathbb{P}}\left( d^{v_1} + d^{v_2} + d^{b_1}\right) = \Omega_{d, \mathbb{P}}\left( d^{v_1} + d^{v_2} + d^{b_1} + d^{b_2(u, 1)}\right).
\end{align*}
\end{itemize}
\end{itemize}
In sum, we have
\begin{align*}
   E_{x, \epsilon} \left[ \left(\hat{f}_{\lambda}(x) - f_{\star}(x) \right)^{2} \right]
= 
\Theta_{d, \mathbb{P}}\left( d^{v_1} + d^{v_2} + d^{b_1} + d^{b_2(u, 1)}\right),
\end{align*}
and we obtain the desired results.
\end{proof}





\section{Proof of Other Results in the Main Text}\label{appendix_proof_other_results}




\subsection{Proof of Theorem 3.3}

\begin{proof}
\noindent {\bf Proof of (i). } Suppose $0  \leq s \leq \Gamma(\gamma)$. 
% By the definition of $\Gamma(\cdot)$ in Eq.~(16), it follows that $\gamma$ cannot be an integer. 
Recall that we have $p := \lfloor \gamma/(s+1) \rfloor \leq \ell_{\gamma} := \lfloor \gamma \rfloor$. 
The proof is split into the following three steps:


\noindent {\bf Step I. } We first consider the case where $\gamma > 0$ is an integer. Then (16) and $0 \leq s \leq \Gamma(\gamma)$ imply that $s=0$. 
Corollary \ref{coroll_learn_curve_case_1} implies that
$$
E_{x, \epsilon} \left[ \left(\hat{f}_{\lambda}(x) - f_{\star}(x) \right)^{2} \right]
=
\Theta_{d, \mathbb{P}}(1)
= \Theta_{d, \mathbb{P}} \left( 
d^{p-\gamma} + 
    d^{-(p+1)s}
\right),
$$
where the last equation follows from 
$p \leq \gamma$ and $s=0$. 




\noindent {\bf Step II. }  
We then consider the case where $\gamma > 0$ is not an integer and show $p=\ell_{\gamma}$ under the following three cases:
\begin{itemize}
    \item When $\gamma \leq 0.5$, we have $p=\ell_{\gamma}=0$;

    \item When $\gamma \in (\ell_{\gamma}, \ell_{\gamma}+0.5]$ for $\ell_{\gamma}=1, 2, \cdots$, we have
    $$
    \ell_{\gamma} \geq p = \lfloor \gamma/(s+1) \rfloor \geq \lfloor \gamma/(\Gamma(\gamma)+1) \rfloor = \lfloor \ell_{\gamma} \rfloor = \ell_{\gamma};
    $$

    \item When $\gamma \in (\ell_{\gamma}+0.5, \ell_{\gamma}+1)$, we have
    $$
    \ell_{\gamma} \geq p \geq \lfloor \gamma/(\Gamma(\gamma)+1) \rfloor = 
     \left\lfloor \frac{\gamma (\ell_{\gamma}+1)}{2(\ell_{\gamma}+1)-\gamma} \right\rfloor
    \geq
    \left\lfloor \frac{(\ell_{\gamma}+0.5) (\ell_{\gamma}+1)}{2(\ell_{\gamma}+1)-(\ell_{\gamma}+1)} \right\rfloor
    = \ell_{\gamma}.
    $$

\end{itemize}
Furthermore, Proposition \ref{prop_e_1} implies that $p \leq u^{\prime}(\tau) < \gamma$. Therefore, for any $u \in [u^{\prime}(\tau), \infty]$, we have 
\begin{equation}\label{eqn_benign_overfit_1}
    \tilde{\ell} := \min\{\ell_{\gamma}, \lfloor u \rfloor\} = \min\{\ell_{\gamma}, p\} = \ell_{\gamma} = p.
\end{equation}


\noindent {\bf Step III. } Now we are ready to determine the convergence rate of the excess risk when $\gamma > 0$ is not an integer.
We first recall the definitions of certain notations, e.g., $b_i = b_i(u)$ and $v_i = v_i(u)$ for $i=1,2$, in Appendix \ref{append_learning_curve}.

From Theorem 3.1, we have
\begin{equation*}
\begin{aligned}
    E_{x, \epsilon} \left[ \left(\hat{f}_{\lambda}(x) - f_{\star}(x) \right)^{2} \right]
= \Theta_{d, \mathbb{P}} \left( d^{v_1} + d^{v_2} + d^{b_1} + \mathbf{1}\{\tau<\infty\} d^{b_2(u, \tau)}\right);
\end{aligned}
\end{equation*}
Furthermore, Proposition \ref{prop_e_3} implies that, for any $u \in [u^{\prime}(\tau), \infty]$, we have
\begin{equation}\label{eqn_benign_overfit_2}
\begin{aligned}
    E_{x, \epsilon} \left[ \left(\hat{f}_{\lambda}(x) - f_{\star}(x) \right)^{2} \right]
= \Theta_{d, \mathbb{P}} \left( d^{v_1} + d^{v_2} + d^{b_1}\right).
\end{aligned}
\end{equation}









We divide the remaining proof in two parts:
\begin{itemize}
\item When $u \geq \gamma$, we have
    \begin{equation*}
    \begin{aligned}
    &~v_1 =  \gamma - \ell_{\gamma} - 1,\\
    &~v_2 =   \ell_{\gamma} - \gamma,\\
    &~b_1 =  -(\ell_{\gamma} + 1)  s.
    \end{aligned}
\end{equation*}

Therefore, (\ref{eqn_benign_overfit_2}) implies 
\begin{equation}
\begin{aligned}
    E_{x, \epsilon} \left[ \left(\hat{f}_{\lambda}(x) - f_{\star}(x) \right)^{2} \right]
= \Theta_{d, \mathbb{P}} \left( d^{\ell_{\gamma} - \gamma} + d^{\gamma - \ell_{\gamma} - 1} + d^{-(\ell_{\gamma} + 1)  s}\right).
\end{aligned}
\end{equation}
Notice that the above rate is the same as the rate given in (13) of \cite{zhang2024phase}. Following the same discussion under Corollary 1 in \cite{zhang2024phase}, we have 
\begin{equation}\label{eqn_zhang_coroll_1_prove_part_1}
\begin{aligned}
    E_{x, \epsilon} \left[ \left(\hat{f}_{\lambda}(x) - f_{\star}(x) \right)^{2} \right]
= \Theta_{d, \mathbb{P}} \left( d^{\ell_{\gamma} - \gamma} + d^{\gamma - \ell_{\gamma} - 1} + d^{-(\ell_{\gamma} + 1)  s}\right) = \Theta_{d, \mathbb{P}} \left( 
d^{p-\gamma} + 
    d^{-(p+1)s}
\right).
\end{aligned}
\end{equation}
(For completeness, we provide an alternative proof of \eqref{eqn_zhang_coroll_1_prove_part_1} in Remark \ref{remark_zhang_coroll_1_prove_part_1}.)



\iffalse
{\color{red}
If $v_2<v_1$, then $0.5< \gamma-\ell_{\gamma}<1$. 
By $s\leq \Gamma(\gamma)$ and (16), we have 
$$
b_1-v_1=(\ell_{\gamma}+1)(1-s)-\gamma \geq \gamma - \gamma=0. 
$$
Therefore, either $v_2\geq v_1$ or $b_1\geq v_1$ holds, which implies 
}
\dhcom{
Same as the discussion under Corollary 1 in \cite{zhang2024phase}, we have
}{I cannot find this discussion so I just write the proof above}
{\color{red}
\begin{equation*}
\begin{aligned}
    E_{x, \epsilon} \left[ \left(\hat{f}_{\lambda}(x) - f_{\star}(x) \right)^{2} \right]
= \Theta_{d, \mathbb{P}} \left( d^{v_1} + d^{v_2} + d^{b_1}\right) = \Theta_{d, \mathbb{P}} \left( 
d^{p-\gamma} + 
    d^{-(p+1)s}
\right).
\end{aligned}
\end{equation*}
}
\fi 

\item When $u^{\prime}(\tau) \leq u < \gamma$, from (\ref{eqn_benign_overfit_1}) we have
    \begin{equation*}
    \begin{aligned}
    &~v_1 =  -\gamma - p - 1 +2u,\\
    &~v_2 =   p - \gamma,\\
    &~b_1 =  -(p + 1)  s.
    \end{aligned}
\end{equation*}
\begin{itemize}

    \item When $\gamma \in (\ell_{\gamma}, \ell_{\gamma}+0.5]$, from (\ref{eqn_benign_overfit_1}) we have $u<\ell_{\gamma}+0.5=p+0.5$, hence
    $$
    v_2 > v_1.
    $$

    \item When $\gamma \in (\ell_{\gamma}+0.5, \ell_{\gamma}+1)$, we have
    $$
    b_1 - v_1 = (p+1)(1-s) + \gamma - 2u \overset{(16)}{\geq} 2(\gamma - u)>0.
    $$

\end{itemize}
Therefore, for any $u \in [u^{\prime}(\tau), \infty]$, from (\ref{eqn_benign_overfit_2}) we have
\begin{equation*}
\begin{aligned}
    E_{x, \epsilon} \left[ \left(\hat{f}_{\lambda}(x) - f_{\star}(x) \right)^{2} \right]
= \Theta_{d, \mathbb{P}} \left( d^{v_2} + d^{b_1}\right) = \Theta_{d, \mathbb{P}} \left( 
d^{p-\gamma} + 
    d^{-(p+1)s}
\right).
\end{aligned}
\end{equation*}
\end{itemize}
%Since the minimax rate is \(\frac{d^{p}}{n} + d^{-(p+1)s}\) (see, e.g., Proposition \ref{prop_pinsker_constant}), we obtained the desired results.





\noindent {\bf Proof of (ii). } Suppose $s > \Gamma(\gamma)$. 
When $u=\gamma$, we have $\tilde{\ell}=\ell_{\gamma}$, hence Theorem 3.1 implies
\begin{equation*}
    \begin{aligned}
        E_{x, \epsilon} \left[ \left(\hat{f}_{d^{\gamma}}(x) - f_{\star}(x) \right)^{2} \right]
        =&~
        \Omega_{d, \mathbb{P}} \left( d^{\ell_{\gamma} - \gamma} + d^{\gamma - \ell_{\gamma} - 1} + d^{-(\ell_{\gamma} + 1)  s}\right).
    \end{aligned}
\end{equation*}
Notice that the above rate is the same as the one given in (13) of \cite{zhang2024phase}. Therefore, following the same discussion under Corollary 1 in \cite{zhang2024phase}, we have 
\begin{equation}\label{eqn_zhang_coroll_1_prove_part_2}
    \begin{aligned}
        E_{x, \epsilon} \left[ \left(\hat{f}_{d^{\gamma}}(x) - f_{\star}(x) \right)^{2} \right]
        = \Omega_{d, \mathbb{P}} \left( d^{\ell_{\gamma} - \gamma} + d^{\gamma - \ell_{\gamma} - 1} + d^{-(\ell_{\gamma} + 1)  s}\right) \gg d^{p-\gamma} + 
    d^{-(p+1)s}
 \quad \text{ in prob.}
    \end{aligned}
\end{equation}

(For completeness, we provide an alternative proof of \eqref{eqn_zhang_coroll_1_prove_part_2} in Remark \ref{remark_zhang_coroll_1_prove_part_2}.)


% {\color{red}
% Similar to the discussion below Corollary 1 in \cite{zhang2024phase}, we can show that 
% $$
% E_{x, \epsilon} \left[ \left(\hat{f}_{d^{\gamma}}(x) - f_{\star}(x) \right)^{2} \right] \overset{\text{Theorem 3.1}}{=}
% \Omega_{d, \mathbb{P}} \left( 
% d^{\gamma - \ell_{\gamma} - 1}
% + d^{\ell_{\gamma} - \gamma}
% + d^{-(\ell_{\gamma} + 1)  s}
% \right)
%  \gg  
% d^{p-\gamma} + 
%     d^{-(p+1)s}
%  \quad \text{ in prob.}
% $$
% }


Recall that Proposition \ref{prop_e_1} implies $\gamma > u^{\prime}$.
Notice that, for any non-integer $u$, Theorem 3.1 implies that the convergence rate of the excess risk $E_{x, \epsilon} \left[ \left(\hat{f}_{\lambda}(x) - f_{\star}(x) \right)^{2} \right]$ is continuous with respect to $u$, hence there exists $\tilde{u} \in (u^{\prime}, \gamma)$, such that for $\tilde{\lambda}=d^{-\tilde{u}}$, we have
\begin{equation*}
\begin{aligned}
    E_{x, \epsilon} \left[ \left(\hat{f}_{\tilde{\lambda}}(x) - f_{\star}(x) \right)^{2} \right]
 \gg  
d^{p-\gamma} + 
    d^{-(p+1)s} \quad \text{ in prob.}
\end{aligned}
\end{equation*}
\end{proof}

\vspace{10pt}

\begin{remark}\label{remark_zhang_coroll_1_prove_part_1}
    
    Here we provide an alternative proof for (\ref{eqn_zhang_coroll_1_prove_part_1}).
    Notice that
    \begin{align*}
        \gamma - \ell_{\gamma} - 1 
        \left\{\begin{matrix}
%
% \overset{0< \gamma - \ell_{\gamma} \leq 0.5}{\leq} 
\leq
[1-(\gamma - \ell_{\gamma})]-1 = \ell_{\gamma} - \gamma & \text{ if } \gamma \in (\ell_{\gamma}, \ell_{\gamma}+0.5] ~\text{ and }~ \ell_{\gamma} \geq 0, \\
\vspace{3pt}
= -(\ell_{\gamma}+1)\Gamma(\gamma) \leq -(\ell_{\gamma}+1)s & \text{ if } \gamma \in (\ell_{\gamma}+0.5, \ell_{\gamma}+1),
\end{matrix}\right.
    \end{align*}
    where the first case follows from $0< \gamma - \ell_{\gamma} \leq 0.5$ and the second case follows from the definition of $\Gamma(\cdot)$ and the assumption that $s\geq \Gamma(\gamma)$. 
Recall that in {\bf Step II} we have shown that $p = \ell_{\gamma}$, therefore we have
$$
\Theta_{d, \mathbb{P}} \left( d^{\ell_{\gamma} - \gamma} + d^{\gamma - \ell_{\gamma} - 1} + d^{-(\ell_{\gamma} + 1)  s}\right) 
= 
\Theta_{d, \mathbb{P}} \left( d^{\ell_{\gamma} - \gamma} + d^{-(\ell_{\gamma} + 1)  s}\right)
= \Theta_{d, \mathbb{P}} \left( 
d^{p-\gamma} + 
    d^{-(p+1)s}
\right).
$$
\end{remark}

\vspace{10pt}


\begin{remark}\label{remark_zhang_coroll_1_prove_part_2}
    Here we provide an alternative proof for (\ref{eqn_zhang_coroll_1_prove_part_2}). We consider four different cases given by (16).

\noindent {\bf Case 1: $\gamma \in (0, 0.5]$. } In this case we have $\Gamma(\gamma)=\infty$, hence it is impossible to let $s>\Gamma(\gamma)$.

\noindent {\bf Case 2: $\gamma \in  (\ell_{\gamma}, \ell_{\gamma}+0.5] ~\text{ and }~ \ell_{\gamma} \geq 1$. } In this case we have $\gamma / (s+1) < \gamma / (\Gamma(\gamma)+1)=\ell_{\gamma}$, and hence $p<\ell_{\gamma}$. Therefore, we have

$$
\left\{
\begin{aligned}
\ell_{\gamma} - \gamma
&\ge \gamma - \ell_{\gamma} - 1,
&&\\[4pt]
\ell_{\gamma} - \gamma
&= -\ell_{\gamma}\Gamma(\gamma)
   > -(\ell_{\gamma}+1)\Gamma(\gamma)
   > -(\ell_{\gamma}+1)s,\\[4pt]
\ell_{\gamma} - \gamma
&> p - \gamma,
\end{aligned}\right.
$$
where the first inequality follows from $0 < \gamma - \ell_{\gamma} \le \tfrac{1}{2}$. 
% \begin{align*}
%         \ell_{\gamma} - \gamma 
% \left\{\begin{matrix}
% \overset{0< \gamma - \ell_{\gamma} \leq 0.5}{\geq}  \gamma - \ell_{\gamma} - 1\\
% \vspace{3pt}
% =-\ell_{\gamma}\Gamma(\gamma)> -(\ell_{\gamma}+1)\Gamma(\gamma)>-(\ell_{\gamma}+1)s\\
% \vspace{3pt}
% > p - \gamma.
% \end{matrix}\right.
%     \end{align*}
Therefore, we have 
$$\Omega_{d, \mathbb{P}} \left( 
d^{\gamma - \ell_{\gamma} - 1}
+ d^{\ell_{\gamma} - \gamma}
+ d^{-(\ell_{\gamma} + 1)  s}
\right)
=
\Omega_{d, \mathbb{P}} \left( 
d^{\ell_{\gamma} - \gamma}
\right). 
$$

Furthermore, we have
$$
\left\{
\begin{aligned}
-(p+1)s 
&< -\ell_{\gamma}\Gamma(\gamma) = \ell_{\gamma} - \gamma, 
&& \text{if } p = \ell_{\gamma} - 1,\\
-(p+1)s 
&< p+1-\gamma < \ell_{\gamma} - \gamma,
&& \text{if } p < \ell_{\gamma} - 1, 
\end{aligned}
\right.
$$
where the first case uses $s>\Gamma(\gamma)$ and (16), and the second case uses $\gamma < (p+1)(s+1)$ by the definition of $p$. 

% \begin{align*}
%         -(p+1)s 
% \left\{\begin{matrix}
% < -\ell_{\gamma}\Gamma(\gamma) = \ell_{\gamma} - \gamma, & \text{ if } p=\ell_{\gamma} - 1\\
% \vspace{3pt}
% \overset{\gamma<(p+1)(s+1)}{<} p+1-\gamma < \ell_{\gamma} - \gamma, & \text{ if } p<\ell_{\gamma} - 1.
% \end{matrix}\right.
%     \end{align*}

Together with $p<\ell_\gamma$, we have
$$
\Omega_{d, \mathbb{P}} \left( 
d^{\ell_{\gamma} - \gamma}
\right)
\gg  
d^{p-\gamma} + 
    d^{-(p+1)s}
 \quad \text{ in prob.}
$$

\noindent {\bf Case 3: $\gamma \in (\ell_{\gamma}+0.5, \ell_{\gamma}+1)$. } In this case we have

$$
\left\{
\begin{aligned}
\gamma - \ell_{\gamma} - 1
&> \ell_{\gamma} - \gamma,
\qquad (\text{since } \gamma - \ell_{\gamma} > 0.5)\\[4pt]
\gamma - \ell_{\gamma} - 1
&= -(\ell_{\gamma}+1)\Gamma(\gamma)
> -(\ell_{\gamma}+1)s.
\end{aligned}
\right.
$$

% \begin{align*}
%         \gamma - \ell_{\gamma} - 1 
% \left\{\begin{matrix}
% \overset{\gamma - \ell_{\gamma} > 0.5}{>}   \ell_{\gamma} - \gamma\\
% \vspace{3pt}
% =-(\ell_{\gamma}+1)\Gamma(\gamma)> -(\ell_{\gamma}+1)s.
% \end{matrix}\right.
%     \end{align*}
On one hand, when $p=\ell_{\gamma}$, we have
$$
\Omega_{d, \mathbb{P}} \left( 
d^{\gamma - \ell_{\gamma} - 1}
+ d^{\ell_{\gamma} - \gamma}
+ d^{-(\ell_{\gamma} + 1)  s}
\right)
\gg  
d^{\ell_{\gamma} - \gamma}
+ d^{-(\ell_{\gamma} + 1)  s}
=
d^{p-\gamma} + 
    d^{-(p+1)s}
 \quad \text{ in prob.}
$$
On the other hand, when $p<\ell_{\gamma}$, we have
$$
\left\{\begin{aligned}
& \ell_{\gamma} - \gamma > p - \gamma, \\
\vspace{3pt}
 &\ell_{\gamma} - \gamma \geq p + 1 - \gamma > -(p+1)s,
\end{aligned}\right.
$$
where the last inequality follows from $\gamma<(p+1)(s+1)$. 
Therefore, we have
$$
\Omega_{d, \mathbb{P}} \left( 
d^{\gamma - \ell_{\gamma} - 1}
+ d^{\ell_{\gamma} - \gamma}
+ d^{-(\ell_{\gamma} + 1)  s}
\right)
\gg  
d^{\ell_{\gamma} - \gamma}
\gg
d^{p-\gamma} + 
    d^{-(p+1)s}
 \quad \text{ in prob.}
$$



\noindent {\bf Case 4: $\gamma =1, 2, \cdots$. } In this case we have $s>\Gamma(\gamma)=0$. Since $\gamma / (s+1) < \gamma = \ell_{\gamma}$, we further have $p<\ell_{\gamma}$. Therefore, $d^{\ell_{\gamma} - \gamma}$ is of order $\Theta_d(1)$ while $d^{p-\gamma} + d^{-(p+1)s}$ is of order $o_d(1)$. Hence we have
$$
\Omega_{d, \mathbb{P}} \left( 
d^{\gamma - \ell_{\gamma} - 1}
+ d^{\ell_{\gamma} - \gamma}
+ d^{-(\ell_{\gamma} + 1)  s}
\right)
=  
\Omega_{d, \mathbb{P}}(1)
\gg
d^{p-\gamma} + 
    d^{-(p+1)s}
 \quad \text{ in prob.}
$$

\end{remark}










\subsection{Estimator of sequence model in (17)}
\begin{proposition}\label{prop_sequence_estimators}
    Consider the sequence model in (17). 
    Suppose $\lambda_j$'s are known and consider the parametrization $f_j = \lambda_j^{1/2} \beta_j$, $j=1, 2, \cdots$.
    Then, applying ridge regression with regularization parameter $\lambda$ or gradient flow with early-stopping time $t = \lambda^{-1}$ yields estimators of the form 
    $$
    \hat{f}_j^{\lambda} = \lambda_j \reg(\lambda_j) z_j,\quad j=1, 2, \cdots,
    $$
    where $\reg$ is the regularization function $\reg^{\krr}$ or $\reg^{\gf}$ defined in Section 2.3.1.
\end{proposition}
\begin{proof}
    \noindent {\bf Ridge regression. } Considering the loss $L(\hat{\beta})=\frac{1}{2}\sum_{j=1}^{\infty} (z_j - \lambda_j^{1/2} \hat{\beta}_j)^2 + \frac{\lambda}{2} \sum_{j=1}^{\infty} \hat{\beta}_j^2$. Minimizing the above loss with respect to each $\hat{\beta}_j$ yields the closed-form solution $\hat{\beta}_j = \lambda_j^{1/2} (\lambda_j + \lambda)^{-1} z_j = \lambda_j^{1/2} \reg^{\krr}(\lambda_j) z_j$. Consequently, the estimator for $\theta_j$ is given by:
    $$
        \hat{f}_j^{\lambda} = \lambda_j \reg^{\krr}(\lambda_j) z_j.
    $$
    \noindent {\bf Gradient flow. }
    Considering the loss $L(\hat{\beta})=\frac{1}{2}\sum_{j=1}^{\infty} (z_j - \lambda_j^{1/2} \hat{\beta}_j)^2$. 
    The gradient flow dynamics, initialized at $\hat{\beta}_j(0) = 0$, are governed by the differential equation:
\[
\frac{\mathsf d}{\mathsf d t} \hat{\beta}_j(t) = -\frac{\partial L}{\partial \hat{\beta}_j} = \lambda_j^{1/2} (z_j - \lambda_j^{1/2} \hat{\beta}_j(t)),
\]
and the solution to this equation is  $\hat{\beta}_j = \lambda_j^{-1/2} (1+e^{-\lambda_j t}) z_j= \lambda_j^{1/2} \reg^{\gf}(\lambda_j) z_j$. Therefore, the corresponding estimator for $\theta_j$ at time $t = \lambda^{-1}$ is:
    $$
        \hat{f}_j^{\lambda} = \lambda_j \reg^{\gf}(\lambda_j) z_j,
    $$
    and we obtain the desired results.
\end{proof}


\subsection{Proof of Proposition 3.4}
\begin{proof}
Denote
$$
\mathcal{N}_{2, \varphi}(\lambda) = \sum_{j =1}^\infty \left[ \lambda_j \reg(\lambda_j) \right]^2
\quad \text{ and } \quad
\mathcal{M}_{2, \varphi}(\lambda) = \sum\limits_{j=1}^{\infty} \left( \rem(\lambda_j) f_{j}\right)^{2}.
$$
Then, when $\lambda > 0$, Lemma D.14 in \cite{lu2024saturation} implies that
$$
\mathcal{N}_{2, \varphi}(\lambda) = \Theta_{d} \left( 
d^{2u - \ell_{\lambda} - 1}
+ d^{\ell_{\lambda}}
\right)
\quad \text{ and } \quad
\mathcal{M}_{2, \varphi}(\lambda) = \Theta_{d} \left( 
d^{-(\ell_{\lambda} + 1)  s}
+ \mathbf{1}\{\tau<\infty\} d^{-2  \tau  u + (2  \tau - \tilde{s})  \ell_{\lambda}}
\right).
$$
Therefore, (19) implies
$$
E_{\xi}\left[\sum_{j=1}^{\infty}(\hat{f}_j^{\lambda} - f_j)^2\right]
= 
\Theta_{d} \left( 
d^{2u - \gamma - \ell_{\lambda} - 1}
+ d^{\ell_{\lambda} - \gamma}
+ d^{-(\ell_{\lambda} + 1)  s}
+ \mathbf{1}\{\tau<\infty\} d^{-2  \tau  u + (2  \tau - \tilde{s})  \ell_{\lambda}}
\right).
$$
\iffalse
Then, the proof for the case where $\lambda = \Omega_d(n^{-1})$ can be proved by following the same arguments in the proof of Lemma D.14 in \cite{lu2024saturation}. 
\fi
\iffalse
Now let's consider the case where $\lambda=0$.
    We have
    \begin{align*}
        \sum_{j =1}^\infty \left[ \lambda_j \varphi_{0}(\lambda_j) \right]^2 \overset{(\ref{eq:Filter_case1})}{\geq} \sum_{j =1}^\infty C_1^2 = +\infty,
    \end{align*}
    hence for any finite $n$, we have
    $$
    E_{\xi}\left[\sum_{j=1}^{\infty}(\hat{f}_j^{\lambda} - f_j)^2\right] \geq \frac{\sigma^2}{n}\sum_{j =1}^\infty \left[ \lambda_j \reg(\lambda_j) \right]^2 = +\infty.
    $$
    \fi
\end{proof}


\subsection{Examples of Assumption 8}\label{appendix_kernel_satisfy_assump_eigenval_general}

In this subsection, we first define the Product Laguerre kernels. Then, we will verify that Assumption 8 holds for the kernels mentioned below it.


\subsubsection{Definition of Product Laguerre kernels on Gamma measures}\label{appendix_Laguerre_polynomials}

In this subsection, we will define a type of kernels based on Laguerre polynomials and Gamma measures, and show that it satisfies Assumption 8.

Fix positive constants $\{\alpha_i >0\}_{i=1}^{\infty}$ and $r > 0$, such that we have that $\sup_{i \geq 1} \alpha_i < \infty$.
For dimensions $d \geq 2r$, let $\mathcal{X} = [0, \infty)^d$ be the positive orthant of $\mathbb{R}^d$, and let $\rho_{\mathcal{X}}$ be the product measure where each coordinate $x_{(i)}$ follows an independent Gamma distribution $\Gamma(\alpha_i + 1, 1)$, $i \leq d$. 
Denote $\alpha = (\alpha_1, \cdots, \alpha_d)^{\top}$.


For any $c>0$, let $I_{c}(z):=\sum_{m=0}^{\infty} [m!\Gamma(m+c+1)]^{-1}(z/2)^{2m+c}$ be the modified Bessel function of the first kind. Define
\begin{equation*}
    \begin{aligned}
        \mathscr{K}_{i}(z, z^{\prime}) =
\frac{\Gamma(\alpha_i+1)}{\left(1- r /d\right)}
\left(\frac{r z z^{\prime}}{d}\right)^{-\alpha_i/2}
\exp\left\{-\frac{ \frac{r}{d}(z + z^{\prime})}{1- r / d}\right\}
I_{\alpha_i}\left(\frac{2\sqrt{ \frac{r}{d}z z^{\prime}}}{1- r /d}\right), \quad z, z^{\prime} \in [0, \infty), \quad i \leq d.
    \end{aligned}
\end{equation*}
Finally, let's define a mapping from $\mathcal{X} \times \mathcal{X}$ to $\mathbb{R}$:
\begin{equation}\label{eqn_def_product_Laguerre_kernels}
    \tilde{\mathscr{K}}_{\alpha, r}(x,x^\prime) := \prod_{i=1}^{d} \mathscr{K}_{i}(x_{(i)}, x_{(i)}^{\prime}), \quad x, x^{\prime} \in [0, \infty)^{d}.
\end{equation}


\subsubsection{Verifying Assumption 8}

In this part, we will show that Assumption 8 holds for several types of common kernels in large dimensions, that is, for these kernels, there exists an absolute constant $\tilde{\mathfrak{C}}$ such that $\sum_{k=0}^{\infty} \sum_{j=1}^{\tilde{N}(d, k)} \tilde{\lambda}_{k, j}\leq \tilde{\mathfrak{C}}$; moreover, for any $k  = 0, 1, \cdots, 2\ell_{\gamma} +2$, we assume that
\begin{equation}\label{eqn_appendix_eigen_gap_condition}
    \tilde{N}(d, k) = \Theta_d(d^{k}) \quad \text{ and } \quad 
        \tilde{\lambda}_{k, j} = \Theta_d(d^{-k}), \quad \text{ for } 1\leq j \leq \tilde{N}(d, k).
\end{equation}
    


\noindent {\bf Inner product kernels defined on $\mathbb{S}^{d-1}$ with uniform distribution. } 
Assumption 2 assumes that $\sum\nolimits_{i=0}^{\infty} a_i \leq a_{-1} < \infty$, hence $\sup_{x \in \mathcal X} \mathscr{K}(x, x) 
    = \sup_{x \in \mathcal X} \Phi(\left\langle x, x \right\rangle) 
    \leq \sum\nolimits_{i=0}^{\infty} a_i \leq a_{-1}$. Therefore, the integral operator defined in (3) is bounded:
    $$
        \sum_{k=0}^{\infty} \sum_{j=1}^{\tilde{N}(d, k)} \tilde{\lambda}_{k, j}  = \text{Tr}(T) :=\int \mathscr{K}(x, x)  ~\mathsf{d} \rho_{\calX}(x) \leq a_{-1}.
    $$
    Moreover, since Proposition 2.1 implies (\ref{eqn_appendix_eigen_gap_condition}), we know that Assumption 8 holds for inner product kernels.


    
    

\noindent {\bf Random feature kernels defined on the hypercube $\{-1, 1\}^{d}$ with uniform distribution. } 
From Appendix D.2 and E.2 in \cite{mei2022generalization}, eigen-decomposition for the random feature kernels defined on the hypercube is similar to inner-product kernels on the sphere. Especially, the addition formula holds for both of them (see, e.g., Appendix E.2.2 in \cite{mei2022generalization}).
    This means that we can write $\mathscr{K}(x, x^\prime) = \Phi(\left\langle x, x^\prime \right\rangle)$, since
    $$
        \mathscr{K}(x, x^\prime) = \sum_{k=0}^{\infty} \mu_{k} \sum_{j=1}^{\tilde{N}(d, k)}  Y_{k,j}(x) Y_{k,j}(x^\prime) \overset{\text{Addition formula}}{=} \sum_{k=0}^{\infty} \mu_{k} \tilde{N}(d, k) Q_k^{(d)}(<x, x^\prime>),
    $$
    where $\tilde{N}(d, k)$ is the multiplicity of $\mu_{k}$ defined in \cite{mei2022generalization}, and $Q_k^{(d)}(\cdot)$ is the $k$-th degree  hyper-cubic Gegenbauer defined as Appendix E.2.2 in \cite{mei2022generalization}. Therefore, if we assume Assumption 2 again,
    then the proof for the random feature kernels defined on the hypercube follows from the same proof as for the inner-product kernels on the sphere, and thus Assumption 8 holds for random feature kernels.



\noindent {\bf Gaussian kernels defined on $\mathbb{R}^d$ with standard Gaussian distribution. } 
Notice that from the definition of Gaussian kernels, we have $\sup_{x \in \mathcal X} \mathscr{K}(x, x) \leq \tilde{\mathfrak{C}}$. Hence
    $$
        \sum_{k=0}^{\infty} \sum_{j=1}^{\tilde{N}(d, k)} \tilde{\lambda}_{k, j}  = \text{Tr}(T) :=\int \mathscr{K}(x, x)  ~\mathsf{d} \rho_{\calX}(x) \leq \tilde{\mathfrak{C}}.
    $$
    Moreover, from Section 4.3 in \cite{rasmussen2003gaussian}, we know that the eigenspaces of Gaussian kernels defined on $\mathbb{R}^d$ with standard Gaussian distribution have $\Theta(d^k)$ eigenvalues of order $\Theta(d^{-k})$.
    Therefore, Assumption 8 holds for Gaussian kernels.

\noindent {\bf Product Laguerre kernels defined in (\ref{eqn_def_product_Laguerre_kernels}). }
The following lemma shows that $\tilde{\mathscr{K}}_{\alpha, r}$ defined in (\ref{eqn_def_product_Laguerre_kernels}) is a kernel satisfies Assumption 8.

\begin{lemma}
    Let $\tilde{\mathscr{K}}_{\alpha, r}$ be defined as (\ref{eqn_def_product_Laguerre_kernels}). Then it is a continuous and positive-definite kernel on $\mathcal{X} \times \mathcal{X}$ with respect to $\rho_{\mathcal{X}}$, where $\mathcal{X} = [0, \infty)^d$, and $\rho_{\mathcal{X}}$ is the product measure with each coordinate $x_{(i)}$ following an independent Gamma distribution $\Gamma(\alpha_i + 1, 1)$. Moreover, it admits a Mercer decomposition given in (22), with eigenvalues and eigenfunctions satisfying Assumption 8 with $\tilde{\mathfrak{C}} = 4^r$.
\end{lemma}
\begin{proof}
It is obvious that $\tilde{\mathscr{K}}_{\alpha, r}$ is a continuous and positive-definite kernel. 
\iffalse
Notice that for any $z>0$ and any $c>0$, we have $I_{c}(z) \leq (z/2)^{c} [\Gamma(c+1)]^{-1} e^{z}$ (see, e.g., Equation (6.25) in \cite{luke1972inequalities}).
% https://math.stackexchange.com/questions/415834/bounding-the-modified-bessel-function-of-the-first-kind
Hence for any $i \leq d$ and any $z, z^{\prime}>0$, we have
\begin{align*}
    |K_{i}(z, z^{\prime})| \leq &~
 \frac{\Gamma(\alpha_i+1)}{\left(1- r /d\right)}
\left(\frac{r z z^{\prime}}{d}\right)^{-\alpha_i/2}
\exp\left\{-\frac{ \frac{r}{d}(z + z^{\prime})}{1- r / d}\right\}
\left(\frac{\sqrt{ \frac{r}{d}z z^{\prime}}}{1- r /d}\right)^{\alpha_i} [\Gamma(\alpha_i+1)]^{-1} e^{\frac{2\sqrt{ \frac{r}{d}z z^{\prime}}}{1- r /d}}\\
= &~
\frac{ \Gamma(\alpha_i+1)}{\Gamma(\alpha_i+1)} \cdot \left(1- r /d\right)^{-\alpha_i-1} \cdot \exp\left\{-\frac{ \frac{r}{d}(z + z^{\prime}) - 2\sqrt{ \frac{r}{d}z z^{\prime}}}{1- r / d}\right\}\\
\leq &~ 1 \cdot \left(1- r /d\right)^{-\sup_{i \geq 1}\alpha_i -1} \cdot 1.
\end{align*}
where the last inequality is implied by that $\sup_{i \geq 1} \alpha_i < \infty$.
Therefore, for any $x, x^{\prime} \in [0, \infty)^{d}$, we have
\begin{align*}
    |\tilde{\mathscr{K}}_{\alpha, r}(x,x^\prime)| \leq \left(1- r /d\right)^{-(d /r) \cdot r(\sup_{i \geq 1}\alpha_i + 1)} \overset{d \geq 2r}{\leq}
    4^{r(\sup_{i \geq 1}\alpha_i + 1)},
\end{align*}
hence $\tilde{\mathscr{K}}_{\alpha, r}$ is bounded. 
\fi
We then show that
$$
\tilde{\mathscr{K}}_{\alpha, r}(x,x^\prime) = \sum_{k=0}^{\infty} \sum_{j=1}^{\tilde{N}(d, k)} \tilde{\lambda}_{k, j}  \tilde{\psi}_{k, j}(x) \tilde{\psi}_{k, j}\left(x^\prime\right),
$$
and provide explicit expressions of $\tilde{\lambda}_{k, j}$'s and $\tilde{\psi}_{k, j}$'s. 
For any $c>0$, let $\{L_{k}^{(c)}(\cdot)\}_{k = 0}^{\infty}$ be the generalized Laguerre polynomials (see, e.g., Chapter 6.2 in \cite{andrews1999special})
% https://en.wikipedia.org/wiki/Laguerre_polynomials
and $\tilde{L}_{k}^{(c)}:=\sqrt{k! \Gamma(c+1) / \Gamma(k+c+1)}L_{k}^{c}(\cdot)$, $k \geq 0$, then we have (see, e.g., Eq.~(6.2.3) in \cite{andrews1999special}):
% https://en.wikipedia.org/wiki/Laguerre_polynomials
\begin{align*}
    &~\int_{0}^{\infty} \tilde{L}_{k}^{(c)}(x) \tilde{L}_{m}^{(c)}(x) \frac{1}{\Gamma(c+1)} x^{c}  e^{-x} ~ \mathrm{d} x\\
    =&~
    \sqrt{\frac{k! m!}{\Gamma(k+c+1)\Gamma(m+c+1)}}\int_{0}^{\infty} L_{k}^{(c)}(x) L_{m}^{(c)}(x) x^{c}  e^{-x} ~ \mathrm{d} x\\
    =&~
    \sqrt{\frac{k! m!}{\Gamma(k+c+1)\Gamma(m+c+1)}} \frac{\Gamma(k+c+1)}{k!} \mathbf{1}\{k=m\} = \mathbf{1}\{k=m\}.
\end{align*}
Furthermore, it is known that the generalized Laguerre polynomials satisfy the Hardy–Hille formula (see, e.g., Eq. (2.13) in \cite{chatterjea1966operational}), that is, 
% https://en.wikipedia.org/wiki/Laguerre_polynomials
\begin{equation}\label{eqn_Hardy_Hille_Laguerre_poly}
\begin{aligned}
    \mathscr{K}_{i}(z, z^{\prime}) =&~  \Gamma(\alpha_i+1) \cdot \sum_{k=0}^{\infty} \frac{k! (r/d)^{k}}{\Gamma(k+\alpha_i+1)} L_{k}^{(\alpha_i)}(z) L_{k}^{(\alpha_i)}(z^{\prime})\\
    =&~ \sum_{k=0}^{\infty}  (r/d)^{k} \tilde{L}_{k}^{(\alpha_i)}(z) \tilde{L}_{k}^{(\alpha_i)}(z^{\prime}).
\end{aligned}
\end{equation}
Combining (\ref{eqn_def_product_Laguerre_kernels}) and (\ref{eqn_Hardy_Hille_Laguerre_poly}), we have
\begin{equation*}
    \begin{aligned}
        \tilde{\mathscr{K}}_{\alpha, r}(x,x^\prime) =&~ \sum_{k_1=0}^{\infty} \cdots \sum_{k_d=0}^{\infty} (r/d)^{k_1 + \cdots + k_d} \prod_{i=1}^{d} \tilde{L}_{k_i}^{(\alpha_i)}(x_{(i)}) \tilde{L}_{k_i}^{(\alpha_i)}(x_{(i)}^{\prime})\\
        =&~
        \sum_{k=0}^{\infty} (r/d)^{k} \sum_{\substack{k_1, \cdots, k_d \geq 0 \\ \sum_{i \leq d} k_{i} = k}} \left[\prod_{i=1}^{d} \tilde{L}_{k_i}^{(\alpha_i)}(x_{(i)}) \right]
        \left[\prod_{i=1}^{d} \tilde{L}_{k_i}^{(\alpha_i)}(x_{(i)}^{\prime}) \right],
    \end{aligned}
\end{equation*}
implying that $\tilde{\lambda}_{k, j} = (r/d)^{k}$ and $\tilde{N}(d, k) = C_{k+d-1}^{d-1}$. Therefore, we have
    $$
    \sum_{k=0}^{\infty} \sum_{j=1}^{\tilde{N}(d, k)} \tilde{\lambda}_{k, j} = \sum_{k=0}^{\infty} (r/d)^{k} C_{k+d-1}^{d-1} \overset{\text{Generating function identity}}{=} (1-r/d)^{-(d-1+1)} \overset{d \geq 2r}{\leq} 4^r;
    $$

    % Generating function identity: https://en.wikipedia.org/wiki/Generating_function

    Moreover, for any $k  = 0, 1, \cdots, 2\ell_{\gamma} +2$, we have 
    $$
        \tilde{\lambda}_{k, j} = (r/d)^{k}= \Theta_d(d^{-k}), \quad j \leq \tilde{N}(d, k) = C_{k+d-1}^{d-1} = \Theta_d(d^{k}).
    $$
\end{proof}











\subsection{Assumption 9 may hold even when the assumption in \cite{misiakiewicz2024non} fails}\label{appendix_counter_mis}

We present an example that Assumption 9 holds while the assumption in \cite{misiakiewicz2024non} does not hold. 

Consider $\calX=[0, 1]$. We have that $\{e_j(x):=e^{2\pi i j x}\}_{j \in \mathbb{Z}}$ form an orthonormal basis of $L^2([0, 1])$, and that
$$
\sup_{j \in \mathbb{Z}} \|e_j\|_{L^{q}} = \left(\int_0^{1} |e_j(x)|^q ~ \mathsf{d} x \right)^{\frac{1}{q}} = 1, \quad q \geq 2,
$$
hence Assumption 9 holds.

On the contrary, define $N=\lfloor (4C)^{2(\ell_{\gamma}+1)} \pi^4 / 8 \rfloor +1$ and $f_{\star}(x):= N^{-1/2} \sum_{j=1}^{N} e_j(x)$. Then $\left\|f_{\star}\right\|_{L^2}^2 ={\sum_{j=1}^{N} N^{-1}} = 1$. However, for $q=4$, we have
\begin{align*}
    \left\|f_{\star}\right\|_{L^4}^4 
    = &~
    \frac{1}{N^2} \int_0^{1} \left|\frac{e^{2\pi i (N+1)x} - e^{2\pi i x}}{e^{2\pi i x} - 1}\right|^4 ~ \mathsf{d} x = 
    \frac{1}{N^2} \int_0^{1} \left|
    \frac{e^{\pi i (N+2)x}(e^{\pi i Nx} - e^{-\pi i Nx})}{e^{\pi i x}(e^{\pi i x} - e^{-\pi i x})}
    \right|^4 ~ \mathsf{d} x\\
    = &~
    \frac{1}{N^2} \int_0^{1} \left|
    e^{\pi i (N+1)x}
    \right|^4\left|
    \frac{\sin(\pi N x)}{\sin(\pi x)}
    \right|^4 ~ \mathsf{d} x = \frac{1}{N^2} \int_0^{1} \left|
    \frac{\sin(\pi N x)}{\sin(\pi x)}
    \right|^4 ~ \mathsf{d} x\\
    \geq &~
    \frac{1}{N^2} \int_0^{(2N)^{-1}} \left|
    \frac{\sin(\pi N x)}{\sin(\pi x)}
    \right|^4 ~ \mathsf{d} x\\
    \overset{\text{(A)}}{\geq} &~ 
    \frac{1}{N^2} \int_0^{(2N)^{-1}} \left|
    \frac{\frac{2}{\pi}\pi N x}{\pi x}
    \right|^4 ~ \mathsf{d} x = \frac{8}{\pi^4} N > \left[(C q)^{(\ell_{\gamma}+1) / 2}\left\|f_{\star}\right\|_{L^2}\right]^4,
\end{align*}
 where the inequality (A) holds since for any $\theta \in [0, \pi /2]$, we have $2\theta / \pi \leq \sin(\theta) \leq \theta$.
 This implies that the assumption in \cite{misiakiewicz2024non} does not hold. 



\section{Properties of $\reg$}\label{appendix_filter_property}

\subsection{Additional requirements on $\reg$}\label{appendix_addtion_require_of_reg}

We list two additional requirements on $\reg$ that are borrowed from \cite{lu2024saturation}.

\begin{itemize}
    \item If $\tau<\infty$, then there exist positive constants $\mathfrak{C}_{7}$ and $\mathfrak{C}_{8}$ only depending on $\tau$, such that we have
    \begin{align*}
    \label{eq:Filter_Rem_finite_case2}
        (z/\lambda)^{2\tau} \rem^2(z)  \geq \mathfrak{C}_{7},
        \quad 
        &\forall \lambda \in (0,1), z > \lambda\\
        (z/\lambda)^{2\tau} \rem^2(z)  \leq \mathfrak{C}_{8} z  \reg(z),
        \quad 
        &\forall \lambda \in (0,1), z \leq \lambda.
    \end{align*}

    \item Let
  \begin{align*}
    D_{\lambda} &= \left\{ z \in \bbC : \Re z \in [-\lambda/2,\kappa^2], ~ \abs{\Im z} \leq \Re z + \lambda/2 \right\} \\
    & \quad \cup \left\{ z \in \bbC : \abs{z - \kappa^2} \leq \kappa^2 + \lambda/2,~ \Re z \geq \kappa^2  \right\};
  \end{align*}
  Then $\reg(z)$ can be extended to be an analytic function on some domain containing $D_\lambda$
  and the following conditions hold for all $\lambda \in (0,1)$:
  \begin{enumerate}
    \item[(C1)] $\left|(z+\lambda)\reg(z)\right| \leq \tilde{E}$ for all $z \in D_\lambda$;
    \item[(C2)] $\left| (z+\lambda)\rem(z)\right| \leq \tilde{F} \lambda$ for all  $z \in D_\lambda$;
  \end{enumerate}
  where  $\tilde{E}, \tilde{F}$ are positive constants.
\end{itemize}





\subsection{An identity of $\reg$ as a matrix function}

The following identity will be used in the proof of our main theorem.

\begin{proposition}\label{prop_matrix_prop_1_of_filter}
    Suppose $A, B$ are semi-positive definite matrices, $A+B$ is a positive definite matrix. Then we have
    \begin{align*}
        \mathrm{I}-A\reg(A+B) =&~ A(A+B)^{-1}\rem(A+B) + B(A+B)^{-1}.
    \end{align*}
\end{proposition}
\begin{proof}
Denote $M:=A+B$. Recall that $\rem(M) = \mathrm{I} - M \reg(M)$. Then
\begin{align*}
    &\quad~ A M^{-1}\rem(M) + B M^{-1} \\
    &= A M^{-1}\big(\mathrm{I} - M \reg(M)\big) + B M^{-1} \\
    &= A M^{-1} - A M^{-1} M \reg(M) + B M^{-1} \\
    &= (A+B) M^{-1} - A \reg(M) \\
    &= \mathrm{I} - A \reg(M).
\end{align*}
\end{proof}


\subsection{Verification of Assumption 5 for Common Spectral Algorithms}\label{appendix_verify_assump_filter}

In this subsection, we demonstrate that several commonly used spectral algorithms, as listed in Section 2.3.1, satisfy Assumption 5.

Given $0 \leq u < \gamma <1$, we have $n^{-1}=o_d(\lambda)$. By Proposition \ref{lemma psi psi top}, there exist constants $\mathfrak{C}_5, \mathfrak{C}_6 >0$, such that we have $\mathfrak{C}_5 \mathrm{I}_n \leq K_{>0} \leq \mathfrak{C}_6 \mathrm{I}_n$. Recalling the decomposition that $K = K_{\leq 0} + K_{>0} = K_{0} + K_{>0}$, we now verify the assumption for each algorithm. 

\begin{itemize}
    \item {\bf KRR.} We have
    \begin{equation*}
    \begin{aligned}
        \rem(K/n) = \lambda (K_{0}/n + K_{> 0}/n + \lambda \mathrm{I}_n)^{-1} 
        =\Omega_{d, \mathbb{P}}(1)\lambda (K_{0}/n  + \lambda \mathrm{I}_n)^{-1}
        = \Omega_{d, \mathbb{P}}(1)\rem(K_{0}/n).
    \end{aligned}
\end{equation*}


    \item {\bf Kernel gradient flow.} We have
    \begin{equation*}
    \begin{aligned}
        \rem(K/n) = e^{- (K_{0}/n + K_{> 0}/n)/\lambda} 
        = \Omega_{d, \mathbb{P}}(1)\rem(K_{0}/n).
    \end{aligned}
\end{equation*}

\item {\bf Iterated ridge regression.} We have
\begin{equation*}
    \begin{aligned}
        \rem(K/n) = \lambda^{q} (K_{0}/n + K_{> 0}/n + \lambda \mathrm{I}_n)^{-q} 
        =\Omega_{d, \mathbb{P}}(1)\lambda^{q} (K_{0}/n  + \lambda \mathrm{I}_n)^{-q}
        = \Omega_{d, \mathbb{P}}(1)\rem(K_{0}/n).
    \end{aligned}
\end{equation*}

\item {\bf Kernel gradient descent.} Denote $\Delta = -\eta K_{> 0}/n$, $S_0=\mathrm{I}_n-\eta K_{0}/n$, and $S_1=S_0+\Delta$. Therefore, when $\eta$ is sufficiently small (depending on $\mathfrak{C}_5, \mathfrak{C}_6$), we have $\lambda_{\min}(S_0)=\Omega_{d, \mathbb{P}}(1)$ and $\Delta=o_{d, \mathbb{P}}(1)$, hence from Taylor's expansion, we have
\begin{equation*}
    \begin{aligned}
        \ln(S_1) \geq \ln(S_0) + \int_{0}^{1} (S_0 + t\Delta)^{-1} \Delta (S_0 + t\Delta)^{-1} ~\mathrm{d}t - O_{d, \mathbb{P}}(\|\Delta\|_{2}^{2})\mathrm{I}_n.
    \end{aligned}
\end{equation*}
Recall that $\rem^{\mr{GD}}(z) = (1-\eta z)^{1/(\eta \lambda)}$ 
and notice that $\frac{1}{\eta \lambda}\|\Delta\|_2^2 =  \eta \cdot (n^2\lambda)^{-1} \|K_{>0}\|_2^2 \to 0$ in probability, since we have $n^2\lambda = d^{2\gamma-u} \to \infty$ by the assumption that $u<\gamma$. 
Hence,
\begin{equation*}
    \begin{aligned}
        \rem(K/n) = &~ e^{\frac{1}{\eta \lambda}\ln(S_1)}\\
        \geq &~ 
        \Omega_{d, \mathbb{P}}(1) 
        e^{\frac{1}{\eta \lambda}\ln(S_0)}
        \exp\left\{-\frac{1}{n\lambda}
        \int_{0}^{1} (S_0 + t\Delta)^{-1} K_{>0} (S_0 + t\Delta)^{-1} ~\mathrm{d}t
        \right\}\\
        \geq &~
        \Omega_{d, \mathbb{P}}(1) 
        e^{\frac{1}{\eta \lambda}\ln(S_0)}.
    \end{aligned}
\end{equation*}

\end{itemize}



\iffalse

\noindent {\bf KRR.} 
We have
\begin{equation*}
    \begin{aligned}
        &~\mathrm{I}-A\reg(A+B) =
        (B+\lambda \mathrm{I})(A+B+\lambda \mathrm{I})^{-1}\\
        =&~
        \lambda A(A+B)^{-1}(A+B+\lambda \mathrm{I})^{-1} + \lambda B(A+B)^{-1}(A+B+\lambda \mathrm{I})^{-1} + B(A+B+\lambda \mathrm{I})^{-1}\\
        =&~
        A(A+B)^{-1}\rem(A+B) + B(A+B)^{-1}\lambda \mathrm{I}(A+B+\lambda \mathrm{I})^{-1}+
        B(A+B)^{-1}(A+B)(A+B+\lambda \mathrm{I})^{-1}\\
        =&~A(A+B)^{-1}\rem(A+B) + B(A+B)^{-1}
    \end{aligned}
\end{equation*}
If $0 \leq u < \gamma <1$, then we have $n^{-1}=o_d(\lambda)$. From Proposition \ref{lemma psi psi top}, there exist constants $\mathfrak{C}_5, \mathfrak{C}_6 >0$, such that we have $\mathfrak{C}_5 \mathrm{I}_n \leq K_{>0} \leq \mathfrak{C}_6 \mathrm{I}_n$. Hence,
\begin{equation*}
    \begin{aligned}
        \rem(K/n) = \lambda (K_{0}/n + K_{> 0}/n + \lambda \mathrm{I}_n)^{-1} 
        =\Omega_{d, \mathbb{P}}(1)\lambda (K_{0}/n  + \lambda \mathrm{I}_n)^{-1}
        = \Omega_{d, \mathbb{P}}(1)\rem(K_{0}/n).
    \end{aligned}
\end{equation*}

\noindent {\bf Kernel gradient flow.} 
We have
\begin{equation*}
    \begin{aligned}
        &~\mathrm{I}-A\reg(A+B) =
        \mathrm{I} - A(A+B)^{-1}\left[\mathrm{I} - e^{-(A+B)/\lambda}\right]\\
        =&~A(A+B)^{-1}\rem(A+B) + B(A+B)^{-1}
    \end{aligned}
\end{equation*}
If $0 \leq u < \gamma <1$, then we have $n^{-1}=o_d(\lambda)$. From Proposition \ref{lemma psi psi top}, there exist constants $\mathfrak{C}_5, \mathfrak{C}_6 >0$, such that we have $\mathfrak{C}_5 \mathrm{I}_n \leq K_{>0} \leq \mathfrak{C}_6 \mathrm{I}_n$. Hence,
\begin{equation*}
    \begin{aligned}
        \rem(K/n) = e^{- (K_{0}/n + K_{> 0}/n)/\lambda} 
        = \Omega_{d, \mathbb{P}}(1)\rem(K_{0}/n).
    \end{aligned}
\end{equation*}



\noindent {\bf Iterated ridge regression.} We have
\begin{equation*}
    \begin{aligned}
        &~\mathrm{I}-A\reg(A+B) =
        \mathrm{I} - A(A+B)^{-1}\left[\mathrm{I} - \lambda^{q} (A+B+\lambda \mathrm{I})^{-q}\right]\\
        =&~A(A+B)^{-1}\rem(A+B) + B(A+B)^{-1}
    \end{aligned}
\end{equation*}

If $0 \leq u < \gamma <1$, then we have $n^{-1}=o_d(\lambda)$. From Proposition \ref{lemma psi psi top}, there exist constants $\mathfrak{C}_5, \mathfrak{C}_6 >0$, such that we have $\mathfrak{C}_5 \mathrm{I}_n \leq K_{>0} \leq \mathfrak{C}_6 \mathrm{I}_n$. Hence,
\begin{equation*}
    \begin{aligned}
        \rem(K/n) = \lambda^{q} (K_{0}/n + K_{> 0}/n + \lambda \mathrm{I}_n)^{-q} 
        =\Omega_{d, \mathbb{P}}(1)\lambda^{q} (K_{0}/n  + \lambda \mathrm{I}_n)^{-q}
        = \Omega_{d, \mathbb{P}}(1)\rem(K_{0}/n).
    \end{aligned}
\end{equation*}


\noindent {\bf Kernel gradient descent.} We have
\begin{equation*}
    \begin{aligned}
        &~\mathrm{I}-A\reg(A+B) =
        \mathrm{I} - A(A+B)^{-1}\left[\mathrm{I} - (\mathrm{I}-\eta(A+B))^{t}\right]\\
        =&~A(A+B)^{-1}\rem(A+B) + B(A+B)^{-1}
    \end{aligned}
\end{equation*}


If $0 \leq u < \gamma <1$, then we have $n^{-1}=o_d(\lambda)$. From Proposition \ref{lemma psi psi top}, there exist constants $\mathfrak{C}_5, \mathfrak{C}_6 >0$, such that we have $\mathfrak{C}_5 \mathrm{I}_n \leq K_{>0} \leq \mathfrak{C}_6 \mathrm{I}_n$. Denote $\Delta = -\eta K_{> 0}/n$, $S_0=\mathrm{I}_n-\eta K_{0}/n$, and $S_1=S_0+\Delta$. Therefore, when $\eta$ is sufficiently small (depending on $\mathfrak{C}_5, \mathfrak{C}_6$), we have $\lambda_{\min}(S_0)=\Omega_{d, \mathbb{P}}(1)$ and $\Delta=o_{d, \mathbb{P}}(1)$, hence from Taylor's expansion, we have
\begin{equation*}
    \begin{aligned}
        \ln(S_1) \geq \ln(S_0) + \int_{0}^{1} (S_0 + t\Delta)^{-1} \Delta (S_0 + t\Delta)^{-1} ~\mathrm{d}t - O_{d, \mathbb{P}}(\|\Delta\|_{2}^{2})\mathrm{I}_n.
    \end{aligned}
\end{equation*}
Recall that $\rem^{\mr{GD}}(z) = (1-\eta z)^{1/(\eta \lambda)}$ {\green and notice that $\frac{1}{\eta \lambda}\|\Delta\|_2^2 =  \eta \cdot (n^2\lambda)^{-1} \|K_{>0}\|_2^2 \to 0$ in probability.} Hence,
\begin{equation*}
    \begin{aligned}
        \rem(K/n) = &~ e^{\frac{1}{\eta \lambda}\ln(S_1)}\\
        \geq &~ 
        \Omega_{d, \mathbb{P}}(1) 
        e^{\frac{1}{\eta \lambda}\ln(S_0)}
        \exp\left\{-\frac{1}{n\lambda}
        \int_{0}^{1} (S_0 + t\Delta)^{-1} K_{>0} (S_0 + t\Delta)^{-1} ~\mathrm{d}t
        \right\}\\
        \geq &~
        \Omega_{d, \mathbb{P}}(1) 
        e^{\frac{1}{\eta \lambda}\ln(S_0)}.
    \end{aligned}
\end{equation*}
\fi




% \dhnote{
% It requires $\frac{1}{\eta \lambda}\|\Delta\|_2^2 \to 0$. 
% Do we need $\eta \lambda = \Theta(1)$?
% }

% {\green [$\eta$ is a constant. I have added a proof showing that $\frac{1}{\eta \lambda}\|\Delta\|_2^2 \to 0$ in above.]}






\section{Auxiliary Lemmas}\label{appendix_auxiliary}


The following proposition is known as Weyl's inequality.

\begin{proposition}[Weyl's inequality]\label{prop_weyl_ine}
Let $A_1$ and $A_2$ be symmetric matrices on $\mathbb{R}^{n \times n}$, with their eigenvalues ordered in descending order. Then, for any indices $i$ and $j$, we have:
\[
\lambda_{i+j-1}(A_1 + A_2) \leq \lambda_i(A_1) + \lambda_j(A_2) \leq \lambda_{i+j-n}(A_1 + A_2).
\]
\end{proposition}






The following proposition is a restatement of von Neumann's trace inequality; see also \cite{marshall1979inequalities}.




\begin{proposition}\label{prop_ri}
    For Hermitian $n \times n$ positive semi-definite complex matrices $A$ and $B$, where now the eigenvalues are sorted decreasingly ($a_1 \geq a_2 \geq \cdots \geq a_n \geq 0$ and $b_1 \geq b_2 \geq \cdots \geq b_n \geq 0$, respectively), we have
$$
\mathrm{tr}\left( AB \right) \geq \sum_{i=1}^{n} a_i b_{n-i+1}.
$$
\end{proposition}






The following proposition is an extension of the symmetric matrix Bernstein inequality.

\begin{proposition}[Restate Exercise 6.10 in \cite{wainwright2019high}]\label{prop_non_sym_matrix_bernstein}
Let $\{A_i\}_{i=1}^n$ be a collection of $n$ independent, zero-mean matrices in $\mathbb{R}^{d_1 \times d_2}$, such that $\|A_i\|_2 \leq b$ almost surely. 
Let 
\[
\sigma^2 := \max\left\{
\left\|\frac{1}{n}\sum_{i=1}^n \mathbb{E}[A_i A_i^\top]\right\|_2,
\left\|\frac{1}{n}\sum_{i=1}^n \mathbb{E}[A_i^\top A_i]\right\|_2
\right\}.
\]
Then, for any $\delta > 0$, the following tail bound holds:
\[
\mathbb{P}\left(\left\|\sum_{i=1}^n A_i\right\|_2 \geq \delta\right) \leq 2(d_1 + d_2) \exp\left(-\frac{\delta^2}{2(n\sigma^2 + b\delta)}\right).
\]
\end{proposition}


The following proposition restates the Sherman-Morrison-Woodbury formula.


\begin{proposition}[Sherman-Morrison-Woodbury formula]\label{prop_smw_formula}
    Suppose $k < n$, $A \in \mathbb{R}^{n \times n}$ is an invertible matrix, and $Z \in \mathbb{R}^{n \times k}$ is such that $ZZ^{\top} + A$ is invertible. Then we have
    \[
(ZZ^{\top} + A)^{-1} = A^{-1} - A^{-1} Z(\mathrm{I}_{k} + Z^{\top} A^{-1} Z)^{-1} Z^{\top} A^{-1}.
\]
\end{proposition}

















\section{Further Experiments and Supporting Proofs}\label{appendix_experiments}

\subsection{Additional Experiments}\label{appendix_addtional_experiments}

Figures \ref{experiment_1_settings_3} and \ref{experiment_1_settings_4} report results of the remaining Type 1 experiments with optimal $C_{\lambda}$ value.


Figure \ref{experiment_1_settings_1_all}, \ref{experiment_1_settings_2_all}, \ref{experiment_1_settings_3_all}, and \ref{experiment_1_settings_4_all} report results of all Type 1 experiments with different $C_{\lambda}$ values.


\subsection{Verification of the regression function in (23)}\label{append_proof_true_function_in_experiments}

In this subsection, we show that $f_{\star}$ defined as (23) satisfies Assumption 3 and Assumption 6. Without loss of generality, we assume that $d \geq \mathfrak{C}$ such that results in Proposition 2.1 hold.

In the following proof, we will frequently use the addition formula (see, e.g., Proposition 1.18 in \cite{gallier2009notes}, page 56 in \cite{ghorbani2021linearized}, or (39) in \cite{lu2024pinsker}):
\begin{equation}\label{eqn_addition_formula}
    P_{k, d}(\langle x, x^{\prime}\rangle) = \frac{1}{N(d, k)} \sum_{j=1}^{N(d, k)} \psi_{k,j}(x) \psi_{k,j}(x^{\prime}),\quad k=0, 1, \cdots
\end{equation}
where $P_{k, d}(\cdot)$ is the $k$-th normalized Gegenbauer polynomial (such that $P_{k, d}(1)=1$) with parameter $(d-2)/2$, and we have
\begin{equation}\label{eqn_former_gegenbauer}
    \begin{aligned}
    P_{0, d}(t)=1, \quad P_{1, d}(t)=t, \quad P_{2, d}(t)=\frac{dt^2-1}{d-1}.
    %P_{3, d}(t)=\frac{t[(d+2)t^2-3]}{d-1}, \quad P_{4, d}(t)=\frac{(d+2)(d+4)t^4-6(d+2)t^2+3}{(d-1)(d+1)}.
    \end{aligned}
\end{equation}



\noindent {\bf Case $s=1$. } In this case we have
\begin{equation*}
    \begin{aligned}
        f_{\star}(x) = &~ \sum_{i=1}^{3} \mathscr{K}(\xi_{i},x)
        \overset{(4)}{=}
        \sum_{j = 1}^{\infty} \left[\sum_{i=1}^{3} \lambda_j \phi_{j}(\xi_{i}) \right]\phi_{j}(x)\\
        (\because \text{Proposition 2.1} ) \quad = &~
        \sum_{k=0}^{\ell_{\gamma}+1} \sum_{j=1}^{N(d,k)} \left[\sum_{i=1}^{3} \mu_{k} \psi_{k,j}(\xi_{i})\right]\psi_{k,j}(x) + \sum_{j =N_{\ell_{\gamma}+1}+1}^{\infty} \left[\sum_{i=1}^{3} \lambda_j \phi_{j}(\xi_{i}) \right] \phi_{j}(x).
    \end{aligned}
\end{equation*}
We first note that 
\begin{align*}
 \|f\|_{[\calH]}^2 =&~ \sum_{j=1}^{\infty} \lambda_j^{-1} \left[\sum_{i=1}^{3} \lambda_j \phi_{j}(\xi_{i}) \right]^2 \leq 3\sum_{i=1}^{3}\sum_{j=1}^{\infty} \lambda_j \phi_{j}(\xi_{i}) \phi_{j}(\xi_{i}) = 3\sum_{i=1}^{3} \mathscr{K}(\xi_{i}, \xi_{i}) \leq 9K_{\max},
\end{align*}
thus $f_{\star}$ satisfies Assumption 3. 

Since we have set $\gamma<2$ and $d>10$ in our experiments, we have $\ell_{\gamma} \leq 1$ and $3-6/(d-1)>2$. Therefore, for any $m \leq \ell_{\gamma}+1 \leq 2$, we have

    % \mu_{m}^{-1} \sum\limits_{j=1}^{N(d, m)}  \theta_{m,j}^{2} =&~ \mu_{m}^{-1} \sum\limits_{j=1}^{N(d, m)} \left[\sum_{i=1}^{3} \mu_{m} \psi_{m,j}(\xi_{i})\right]^2 \geq \mu_{m}^{-1} \sum\limits_{j=1}^{N(d, m)} \mu_{m}^2 \psi_{m,j}(\xi_{1})^2
    % \dhnote{\text{Why is this step true? The cross terms such as $\psi_{m,j}(\xi_{1})\psi_{m,j}(\xi_{2})$ may be negative. }}\\ 
    % \overset{(\ref{eqn_addition_formula})}{=} &~ \mu_{m} N(d, m) \overset{\text{Proposition 2.1}}{=} a_m (1+o_d(1));\\
\begin{equation}\label{eqn_verify_assump_6_for_experiments}
    \begin{aligned}
    &~ \mu_{m}^{-1} \sum\limits_{j=1}^{N(d, m)}  \theta_{m,j}^{2}\\
    =&~ \mu_{m}^{-1} \sum\limits_{j=1}^{N(d, m)} \left[\sum_{i=1}^{3} \mu_{m} \psi_{m,j}(\xi_{i})\right]^2\\
    =&~ \mu_{m} \left[\sum_{i=1}^{3} \sum\limits_{j=1}^{N(d, m)} \psi_{m,j}^2(\xi_{i}) + \sum_{i \neq i^{\prime}}^{3} \sum\limits_{j=1}^{N(d, m)} \psi_{m,j}(\xi_{i}) \psi_{m,j}(\xi_{i^{\prime}}) \right]\\
    \overset{(\ref{eqn_addition_formula})}{=} &~ \mu_{m} N(d, m) \left[ \sum_{i=1}^{3} 
    P_{k, d}(\langle \xi_{i}, \xi_{i}\rangle) + \sum_{i \neq i^{\prime}}^{3} P_{k, d}(\langle \xi_{i}, \xi_{i^{\prime}}\rangle)\right]\\
    \overset{(\ref{eqn_former_gegenbauer})}{=}&~
    \left\{\begin{matrix}
\mu_{0} N(d, 0) [3+6] = 9 \mu_{0} N(d, 0), &  m=0;\\
\mu_{1} N(d, 1) [3 + \sum_{i \neq i^{\prime}}^{3} \langle \xi_{i}, \xi_{i^{\prime}}\rangle] \overset{(24)}{\geq} 0.1 \mu_{1} N(d, 1), &  m=1;\\
\mu_{2} N(d, 2) [3 + \sum_{i \neq i^{\prime}}^{3} \frac{d\langle \xi_{i}, \xi_{i^{\prime}}\rangle^2-1}{d-1}] \geq \mu_{2} N(d, 2) [3 - 6/(d-1)] > 2\mu_{2} N(d, 2), &  m=2;\\
\end{matrix}\right.
\end{aligned}
\end{equation}
Therefore, from Proposition 2.1 we have $\mu_{m}^{-1} \sum_{j=1}^{N(d, m)}  \theta_{m,j}^{2} = \Omega_d(1)$, $m=0, 1, 2$. 
Furthermore, 
\begin{align*}
    \sum\limits_{k=0}^{\tilde{\ell}} \sum\limits_{j =1}^{N(d, k)}  \theta_{k, j}^{2} \geq &~ \theta_{0, 1}^{2} = \left[\sum_{i=1}^{3} \mu_{0} \psi_{0,1}(\xi_{i})\right]^2 = 9\mu_0^2 =\Theta_d(1),
\end{align*}
where the last equation follows from Proposition~2.1. 
We conclude that $f_{\star}$ satisfies Assumption 6.



\noindent {\bf Case $s \neq 1$ and $\mathscr{K}=\mathscr{K}^{\mathrm{rbf}}$. } In this case we have
\begin{equation*}
    \begin{aligned}
        f_{\star}(x) = &~ \sum_{i=1}^{3} \sum_{k=0}^{2} d^{k(1-s)/2} P_{k, d}(\langle\xi_{i}, x\rangle)
        \overset{(\ref{eqn_addition_formula})}{=}
        \sum_{k=0}^{2} \sum_{j=1}^{N(d, k)} \left[\sum_{i=1}^{3} \frac{d^{k(1-s)/2}}{N(d, k)} \psi_{k,j}(\xi_{i}) \right]  \psi_{k,j}(x) 
    \end{aligned}
\end{equation*}
By Proposition 2.1, we have
$$
\begin{aligned}
 \|f\|_{[\calH]^{s}}^2 =&~ \sum_{k=0}^{2} \sum_{j=1}^{N(d, k)} \mu_k^{-s} \left[\sum_{i=1}^{3} \frac{d^{k(1-s)/2}}{N(d, k)} \psi_{k,j}(\xi_{i}) \right]^2 \overset{(\ref{eqn_addition_formula})}{\leq}
        3 \sum_{k=0}^{2} \frac{\mu_k^{-s}d^{k(1-s)}}{N(d, k)} = O_d(1),
\end{aligned}
$$
which implies that $f_{\star}$ satisfies Assumption 3. 



Since we have set $\gamma<2$ and $d>10$ in our experiments, we have $\ell_{\gamma} \leq 1$ and $3-6/(d-1)>2$. 
Therefore, for any $m \leq \ell_{\gamma}+1 \leq 2$, similar to (\ref{eqn_verify_assump_6_for_experiments}), we have
\begin{equation*}
    \begin{aligned}
        \mu_{m}^{-s} \sum\limits_{j=1}^{N(d, m)}  \theta_{m,j}^{2} =&~ \mu_{m}^{-s} \sum\limits_{j=1}^{N(d, m)}  \left[\sum_{i=1}^{3} \frac{d^{m(1-s)/2}}{N(d, m)} \psi_{m,j}(\xi_{i}) \right]^2\\
        =&~  \frac{\mu_{m}^{-s} d^{m(1-s)}}{N(d, m)^2} \left[\sum_{i=1}^{3} \sum\limits_{j=1}^{N(d, m)} \psi_{m,j}^2(\xi_{i}) + \sum_{i \neq i^{\prime}}^{3} \sum\limits_{j=1}^{N(d, m)} \psi_{m,j}(\xi_{i}) \psi_{m,j}(\xi_{i^{\prime}}) \right]\\
        =&~ \Omega_d\left(
        \frac{\mu_{m}^{-s} d^{m(1-s)}}{N(d, m)}
        \right). 
    \end{aligned}
\end{equation*}
 By Proposition 2.1, the right-hand side of the last display is $\Omega_d(1)$. Furthermore, 
$$
\begin{aligned}
\sum\limits_{k=0}^{\tilde{\ell}} \sum\limits_{j =1}^{N(d, k)}  \theta_{k, j}^{2} \geq &~ \theta_{0, 1}^{2} = \left[\sum_{i=1}^{3} \frac{d^{0(1-s)/2}}{N(d, 0)} \psi_{0,1}(\xi_{i}) \right]^2 = \frac{9}{N(d, 0)^2} = \Theta_d(1),
\end{aligned}
$$
where the last equation follows from Proposition 2.1. 
Therefore, $f_{\star}$ satisfies Assumption 6.


% \begin{equation*}
%     \begin{aligned}
%         \mu_{m}^{-s} \sum\limits_{j=1}^{N(d, m)}  \theta_{m,j}^{2} =&~ \mu_{m}^{-s} \sum\limits_{j=1}^{N(d, m)}  \left[\sum_{i=1}^{3} \frac{d^{m(1-s)/2}}{N(d, m)} \psi_{m,j}(\xi_{i}) \right]^2 \geq  \mu_{m}^{-s} \sum\limits_{j=1}^{N(d, m)}  \left[\frac{d^{m(1-s)/2}}{N(d, m)} \psi_{m,j}(\xi_{1}) \right]^2 \dhnote{\text{Similar to the last comment; the cross terms can be negative.} } \\
%         \overset{(\ref{eqn_addition_formula})}{=} &~ \frac{d^{m(1-s)}}{\mu_{m}^{s} N(d, m)} \overset{\text{Proposition 2.1}}{=} \Omega_d(1);\\
        
%      \end{aligned}
% \end{equation*}



\noindent {\bf Case $s \neq 1$ and $\mathscr{K}=\mathscr{K}^{\mathrm{ntk}}$. } 
In our NTK experiments, we always set $\ell_{\gamma} \leq 1$. Moreover, for the NTK kernel we have $a_{j}>0$ for $j \in \{0,1\}\cup \{2,4,6,\cdots\}$. Therefore, the argument follows the same steps as in the RBF kernel case.



\clearpage


\begin{figure}[!htbp]
\centering
\subfigure{\includegraphics[width=0.9\columnwidth]{figures_experiments_1/plot_setting_3_best_C_NTK.pdf}}


\vspace{-15pt}

\subfigure{\includegraphics[width=0.9\columnwidth]{figures_experiments_1/plot_setting_3_best_C_RBF.pdf}}



\caption{
Type 1 experiments with $(\gamma, s, u) = (1.0, 1.0, 2.0)$. The setup and analysis are the same as in Figure 3. Results for additional $C_{\lambda}$ values are shown in Figure \ref{experiment_1_settings_3_all}.
}
\label{experiment_1_settings_3}
\end{figure}

\begin{figure}[!htbp]
\centering
\subfigure{\includegraphics[width=0.9\columnwidth]{figures_experiments_1/plot_setting_4_best_C_NTK.pdf}}

\vspace{-15pt}

\subfigure{\includegraphics[width=0.9\columnwidth]{figures_experiments_1/plot_setting_4_best_C_RBF.pdf}}


\caption{
Type 1 experiments with $(\gamma, s, u) = (1.2, 2.0, 1.0)$. The setup and analysis are the same as in Figure 3. Results for additional $C_{\lambda}$ values are shown in Figure \ref{experiment_1_settings_4_all}.
}
\label{experiment_1_settings_4}
\end{figure}

\begin{figure}
\centering
\subfigure{\includegraphics[width=0.9\columnwidth]{figures_experiments_1/plot_setting_1_all_C_NTK.pdf}}


\vspace{-15pt}

\subfigure{\includegraphics[width=0.9\columnwidth]{figures_experiments_1/plot_setting_1_all_C_RBF.pdf}}



\caption{Complete results for Type 1 experiments with $(\gamma, s, u) = (1.5, 1.5, 0.5)$.}
\label{experiment_1_settings_1_all}
\end{figure}

\begin{figure}
\centering
\subfigure{\includegraphics[width=0.9\columnwidth]{figures_experiments_1/plot_setting_2_all_C_NTK.pdf}}


\vspace{-15pt}

\subfigure{\includegraphics[width=0.9\columnwidth]{figures_experiments_1/plot_setting_2_all_C_RBF.pdf}}



\caption{Complete results for Type 1 experiments with $(\gamma, s, u) = (0.8, 1.0, 2.0)$.}
\label{experiment_1_settings_2_all}
\end{figure}

\begin{figure}
\centering
\subfigure{\includegraphics[width=0.9\columnwidth]{figures_experiments_1/plot_setting_3_all_C_NTK.pdf}}


\vspace{-15pt}

\subfigure{\includegraphics[width=0.9\columnwidth]{figures_experiments_1/plot_setting_3_all_C_RBF.pdf}}



\caption{Complete results for Type 1 experiments with $(\gamma, s, u) = (1.0, 1.0, 2.0)$.}
\label{experiment_1_settings_3_all}
\end{figure}

\begin{figure}
\centering
\subfigure{\includegraphics[width=0.9\columnwidth]{figures_experiments_1/plot_setting_4_all_C_NTK.pdf}}


\vspace{-15pt}

\subfigure{\includegraphics[width=0.9\columnwidth]{figures_experiments_1/plot_setting_4_all_C_RBF.pdf}}



\caption{Complete results for Type 1 experiments with $(\gamma, s, u) = (1.2, 2.0, 1.0)$.}
\label{experiment_1_settings_4_all}
\end{figure}


\clearpage







\end{appendix}



\end{document}
